\documentclass[11pt]{article}

\usepackage[T1]{fontenc}
\usepackage[utf8]{inputenc}
\usepackage[english]{babel}
\usepackage[margin=1in]{geometry}
\usepackage{lmodern}
\usepackage{microtype}
\usepackage{mathtools}
\usepackage{amssymb}
\usepackage{amsthm}
\usepackage{graphicx}
\usepackage{tikz}
\usetikzlibrary{arrows.meta,positioning}
\usepackage{booktabs}
\usepackage{longtable}
\usepackage{array}
\usepackage{tabularx}
\usepackage{pdflscape}
\usepackage{caption}
\usepackage{enumitem}
\usepackage{fancyhdr}
\usepackage[numbers,sort&compress]{natbib}
\usepackage{orcidlink}
\usepackage{titling}
\usepackage{doi}
\usepackage{xurl}
\usepackage{hyperref}

\hypersetup{
  pdftitle={Six Birds Foundations II: Admissibility, Meta-Theory, and the Exact-Six Program},
  pdfauthor={Ioannis Tsiokos},
  hidelinks
}
\urlstyle{same}
\captionsetup{
  font=small,
  labelfont=bf,
  labelsep=period,
  margin=1em
}
\captionsetup[table]{skip=6pt,position=top}
\captionsetup[figure]{skip=6pt}
\setlist[enumerate]{leftmargin=*, itemsep=2pt, topsep=4pt}
\setlist[itemize]{leftmargin=*, itemsep=2pt, topsep=4pt}
\setlength{\droptitle}{-3.2em}
\pretitle{\begin{center}\LARGE\bfseries}
\posttitle{\par\end{center}\vspace{0.5em}}
\preauthor{\begin{center}\normalsize}
\postauthor{\par\end{center}\vspace{-0.4em}}
\predate{\begin{center}\small}
\postdate{\par\end{center}\vspace{0.6em}}
\setlength{\emergencystretch}{2em}
\clubpenalty=10000
\widowpenalty=10000
\displaywidowpenalty=10000
\interfootnotelinepenalty=10000
\allowdisplaybreaks[2]
\setlength{\bibsep}{2pt plus 0.3ex}

\pagestyle{fancy}
\fancyhf{}
\fancyhead[L]{\small\itshape Six Birds Foundations II}
\fancyhead[R]{\small\thepage}
\renewcommand{\headrulewidth}{0.4pt}
\renewcommand{\footrulewidth}{0pt}

\fancypagestyle{firstpage}{\fancyhf{}
  \fancyfoot[C]{\scriptsize Automorph Inc., Wilmington, DE, USA \quad \texttt{ioannis@automorph.io} \quad \textcopyright\ 2026 \quad CC-BY 4.0 \quad DOI: \href{https://doi.org/10.5281/zenodo.23096828}{10.5281/zenodo.23096828} \quad Preprint --- not peer reviewed.}
  \renewcommand{\headrulewidth}{0pt}
  \renewcommand{\footrulewidth}{0pt}
}

\usepackage{titlesec}
\titlespacing*{\section}{0pt}{2.8ex plus 1.0ex minus 0.4ex}{1.6ex plus 0.4ex}
\titlespacing*{\subsection}{0pt}{2.2ex plus 0.8ex minus 0.3ex}{1.1ex plus 0.3ex}

\newcommand{\classid}[1]{\texttt{\small\hyphenchar\font=`\_ #1}}

\newtheorem{theorem}{Theorem}[section]
\newtheorem{proposition}[theorem]{Proposition}
\newtheorem{lemma}[theorem]{Lemma}
\newtheorem{corollary}[theorem]{Corollary}
\theoremstyle{definition}
\newtheorem{definition}[theorem]{Definition}
\newtheorem{remark}[theorem]{Remark}

\newcommand{\R}{\mathbb{R}}
\newcommand{\N}{\mathbb{N}}

\newcommand{\FATCD}{\ensuremath{\mathsf{FATCD}}}
\newcommand{\Raw}[1]{\ensuremath{\operatorname{Raw}\!\left(#1\right)}}
\newcommand{\Audit}[1]{\ensuremath{\operatorname{Audit}\!\left(#1\right)}}
\newcommand{\Closed}[1]{\ensuremath{\operatorname{Closed}\!\left(#1\right)}}
\newcommand{\Dec}[1]{\ensuremath{\operatorname{Dec}\!\left(#1\right)}}
\newcommand{\Residual}[1]{\ensuremath{\operatorname{Residual}\!\left(#1\right)}}
\newcommand{\Chan}[2]{\ensuremath{\operatorname{Chan}_{#1}\!\left(#2\right)}}
\newcommand{\Status}[2]{\ensuremath{\operatorname{Status}_{#1}\!\left(#2\right)}}
\newcommand{\ScopedExactSix}[1]{\ensuremath{\operatorname{ScopedExactSix}\!\left(#1\right)}}

\newcommand{\activeproj}{\textit{active\_\allowbreak projection}}
\newcommand{\belowthr}{\textit{below\_\allowbreak threshold}}
\newcommand{\collapsedbc}{\textit{collapsed\_\allowbreak boundary\_\allowbreak case}}
\newcommand{\absentrec}{\textit{absent\_\allowbreak with\_\allowbreak record}}
\newcommand{\outsidescope}{\textit{inapplicable\_\allowbreak outside\_\allowbreak scope}}

\newcommand{\Pone}{\ensuremath{\mathrm{P1}}}
\newcommand{\Ptwo}{\ensuremath{\mathrm{P2}}}
\newcommand{\Pthree}{\ensuremath{\mathrm{P3}}}
\newcommand{\Pfour}{\ensuremath{\mathrm{P4}}}
\newcommand{\Pfive}{\ensuremath{\mathrm{P5}}}
\newcommand{\Psix}{\ensuremath{\mathrm{P6}}}
 
\begin{document}

\title{Six Birds Foundations II:\\[4pt]
Admissibility, Meta-Theory, and the Exact-Six Program}
\author{Ioannis Tsiokos\,\orcidlink{0009-0009-7659-5964}}
\date{Version 3: October 2, 2026\\[2pt] \footnotesize (Version 1: April 20, 2026)}

\maketitle
\thispagestyle{firstpage}

\begin{abstract}
\noindent
Six Birds Foundations I~\citep{Tsiokos2026SixBirdsFoundations}
described how a layer of description closes over what it tracks, and
identified six roles that such closure uses: descent (whether an update
respects a quotient), representability of a macro-level specification,
enriched route and coherence mismatch, staged order and refinement, packaging into
quotients, and audit. This paper gives the six roles precise meanings
and asks whether exactly six are needed. For each role it fixes what
counts as evidence for it, the level at which it acts, and where it
meets the other roles, and it states the bookkeeping rules a
description must follow before any role claim counts.

The question is posed on a precisely delimited domain, \FATCD{}, of finite
audited typed closure descriptions: finite collections of typed records in
which every claim carries an honest audit. Without further hypotheses the paper
proves three things. Each of the six roles occurs in some description of the
domain. Every description has a decomposition record with one channel per role,
each channel in one of five statuses (active, collapsed, below threshold,
outside scope, absent), and every feature is accounted for by an active
channel, an inactive channel-status entry, or one of eleven residual labels;
the last two, \emph{unresolved candidate} and \emph{recorded seventh-role
screen}, apply only when none of the first nine places the feature. The six channels are part of the record's definition, so the
count is of channels, not of active roles, and the record need not be unique.
Probability, causality and agency, the most natural candidates for a seventh
role, are placed without a seventh role in the forms the paper specifies.

The upper bound is conditional. Its recognition hypothesis is assumed, not
derived, and asks, without reference to any decomposition record, that every
feature be content of one of the six roles, an annotation, data that asserts
nothing, a checked restatement of such content, or content whose check is not
finite over the description (placed as outside-domain structure); the
upper-bound theorem shows that each of these ways yields one of the first nine
labels. Under the hypothesis a feature that no channel accounts for receives
one of labels (iii)--(ix), and no content record, annotation or leftover record
is licensed the last two. The hypothesis is a genuine restriction and depends
on presentation. A standalone formula stating that finitely many recorded
weights are nonnegative and sum to one fails it, while the same weights
recorded as the single row of a transition kernel, with a certificate that the
row is a probability distribution, satisfy it. On the other hand, any closed formula
of the paper's finite check language that does not mention the six witness
predicates can be recorded, provided the completed description lies in \FATCD{}, with its truth value, as a one-variable
specification over two records of the same kind, which the hypothesis accepts;
a false formula is recorded through its checked negation, and the rest of the
description must still satisfy the hypothesis. The paper does not claim six roles
for all of emergence, a full algebra of the six roles, their independence in
every setting, or an empirical realization.
Appendix~\ref{app:mechanization} describes the accompanying Lean
development and what it covers.
\end{abstract}
 \section{Introduction}\label{sec:introduction}


Six Birds Foundations I~\citep{Tsiokos2026SixBirdsFoundations} studied
how a layer of description closes over what it
tracks~\citep{Anderson1972}. It worked with finite theory packages
\(\mathcal{T} = (Z, f, \Sigma_f, E, \mathcal{A})\): a finite space \(Z\)
of micro-descriptions, a lens \(f\) recording what the layer can see,
the definable content \(\Sigma_f\) that the lens induces, a completion
operator \(E\) whose fixed points are the layer's
theories~\citep{DaveyPriestley2002}, and an audit functional
\(\mathcal{A}\). From two assumptions, that constructions compose and
that the layer sees its states only through a limited interface, it
derived six primitive closure roles \Pone{}, \ldots, \Psix{}.

This paper asks whether the list of six is complete. Could a
description need a seventh role, for probability, say, or for causation
or agency? Asked of all possible descriptions, the question is too
loose to answer, so the paper makes it precise. It gives each role an
exact meaning, states the rules a description must follow before a
role claim counts, and fixes a domain of finite descriptions on which
the question can be settled. On that domain it proves that each role
occurs, classifies the stated forms of probability, causality and
agency, and gives every description a record with exactly six role
channels. Under an explicit recognition hypothesis, every residual feature is placed by
one of nine criteria (labels (i)--(ix), listed in
Section~\ref{subsec:objects-at-a-glance}): membership in one or two active role clusters,
role content not yet absorbed, a checked re-presentation or level translation, an annotation,
bookkeeping or assertion-free data, or content outside the covered
scope. Residual features are of two kinds: every content record and
structural annotation, even one an active role uses (a \emph{kind residual}),
and every record a decomposition record leaves unassigned (a \emph{record
residual}). Under the hypothesis, therefore, no residual feature of either kind is left as
an unresolved seventh-role candidate. The exact-six statement is the
last step of this chain, not its whole content.

\subsection{The three parts}\label{subsec:three-pillars}

The paper has three parts, each supplying structure the other two lack.
\begin{itemize}[topsep=2pt,itemsep=2pt]
  \item \emph{Role architecture}
    (Section~\ref{sec:primitive-role-architecture}). It fixes what each
    role means; the three levels at which a role can be present
    (behavior, verification, structure); where each pair of roles meets
    without one collapsing into the other; and how data moves between
    levels.
  \item \emph{Admissibility}
    (Section~\ref{sec:admissibility-meta-theory}). It states the
    bookkeeping discipline: a role claim counts only when its
    hypotheses, changes of level, translations and visibility
    assumptions are explicitly recorded. Under this discipline the roles
    stay distinct as types (typed non-collapse), with no claim that they
    are independent generators.
  \item \emph{The exact-six chain}
    (Sections~\ref{sec:fatcd-domain}--\ref{sec:scoped-exact-six}). It
    defines the domain \FATCD{} of finite audited typed closure
    descriptions and proves lower bounds (each role occurs), an
    upper-bound classification under the recognition hypothesis, with
    separate closures for the recognized forms of probability, causality
    and agency, and a decomposition theorem (six channels, each with a
    status), a checklist of misreadings that these results rule out, and
    the terminal theorem (Section~\ref{sec:scoped-exact-six}).
\end{itemize}

\subsection{The terminal theorem, informally}\label{subsec:terminal-informal}

In plain terms: every finite audited description can be sorted into six role
channels, each marked active or given one of four reasons for being inactive,
together with a set of leftover records; under a stated hypothesis every
leftover is of a listed kind (a re-presentation or level shift of placed
content, role content not placed in an active channel, an artifact of what an
instrument can see, an empirical-bridge annotation, structure outside the
domain, or bookkeeping and data without an assertion of its own), so no
leftover calls for a seventh role.

A finite audited typed closure description \(D\) is a finitely
presented collection of typed raw records and annotations subject to
finiteness, closure and audit conditions (every claim it carries is referred
to by an honest audit annotation, which need not certify it). The role
vocabulary \Pone{}--\Psix{} enters \(D\) not as a constraint on
admissibility but as a family of derived projections \(\pi_i(D)\),
each available when threshold, witness-schema, collapse-case, level,
nonclaim and audit data are
attached to raw features. Theorems~\ref{thm:decomposition-existence} and~\ref{thm:six-channel-exactness} state that
every \(D \in \FATCD{}\) has a well-formed decomposition record
\(\Dec{D}\) with exactly the six role channels \Pone{}, \ldots, \Psix{},
each in exactly one of five statuses (active, or one of four reasons for
inactivity), and that each \(r\in\Residual{D}\) receives one of eleven labels.
Under an explicit recognition hypothesis
(Definition~\ref{def:recognition-hypothesis}), every such label lies in
(iii)--(ix), and every kind residual and every record residual is licensed only
labels among (i)--(ix), so none is left as an unresolved candidate for a
seventh role. The claim is about role channels, not active projections, and
the decomposition is existential, not unique.

\paragraph{Where the content lies.} The parts of the terminal theorem have
different sources. That a decomposition record has six channels is part
of its definition (Definition~\ref{def:decomposition-record}), which is
indexed by the six roles. That every description has a well-formed record
in which each channel carries one status and every raw feature is
accounted for is proved by construction, for every description and
without the recognition hypothesis
(Theorem~\ref{thm:decomposition-existence} and
Theorem~\ref{thm:six-channel-exactness}(i)--(ii)). That each role occurs somewhere is
proved by explicit witnesses (Theorem~\ref{thm:six-role-lower-bounds}).
That the recognized forms of probability, causality and agency need no
seventh role is proved case by case
(Propositions~\ref{prop:probability-closure}--\ref{prop:agency-closure}).
For a description satisfying the recognition hypothesis, no kind residual
and no record residual is licensed label (x) or (xi)
(Theorem~\ref{thm:upper-bound-classification}).
Theorem~\ref{thm:scoped-exact-six} states these parts together, under
that hypothesis. The recognition hypothesis lists five ways a feature can be
recognized; each is a version, independent of the decomposition record, of
label criteria among (i)--(ix), and the proof of
Theorem~\ref{thm:upper-bound-classification} checks the correspondence. That
check uses no property of the particular witness rows or seed clauses
(Remark~\ref{rem:classification-scope}); the number six enters through
Definition~\ref{def:decomposition-record}, and the particular predicates
through Theorem~\ref{thm:six-role-lower-bounds} and
Propositions~\ref{prop:probability-closure}--\ref{prop:agency-closure}.

\paragraph{What the hypothesis adds.} Criterion (xi) presupposes (x). Hence a description meets the condition of Definition~\ref{def:scoped-exact-six-question} exactly when some well-formed decomposition record licenses label (x) for none of its kind residuals and none of its record residuals. The recognition hypothesis is a sufficient condition for this that refers to no decomposition record: Theorem~\ref{thm:upper-bound-classification} proves sufficiency, and Remark~\ref{rem:classification-scope} shows it is not necessary. Recognized features are therefore placed by (i)--(ix) whatever decomposition
record is chosen.

\paragraph{How restrictive the hypothesis is.} The exclusion of a seventh role rests
on the hypothesis, which some descriptions in the domain fail
(Remark~\ref{rem:recognition-open}). That remark also calibrates the
hypothesis: the finite normalization description \(D_{\mathrm{norm}}\) fails it;
the same weights presented as the row of a transition kernel satisfy it; and a closed check-language formula \(\varphi\) without atoms \(W_j\), in a description with two distinct declared records of one kind, can be re-presented as a covered \Ptwo{} seed recording its truth value in the extended description, provided that description stays in \FATCD{}; a false formula is represented by its checked negation. In the completed description channel \Ptwo{} is active and channel \Pfive{} is not \(\absentrec\). When wrapping preserves the truth value of a true closed assertion, the seed records the same checked content, so the hypothesis constrains only presentation, and the
original record can be dropped when nothing else depends on it (the removal
conditions of Remark~\ref{rem:recognition-open}); if it must be kept, it is
classified on its own and may still fail the hypothesis. Label (xi) is realized in \FATCD{} by the
description \(D_{\mathrm{scr}}\) of the same remark.

\subsection{What the paper does and does not claim}\label{subsec:intro-nonclaims}

The terminal theorem is scoped exact-six, not unrestricted exact-six.
It makes none of the following stronger claims:
\begin{itemize}[topsep=2pt,itemsep=1pt]
  \item unrestricted exact-six, or any claim about all emergence
    phenomena;
  \item that every \(D \in \FATCD{}\) activates all six roles or
    carries six active nontrivial witnesses;
  \item that the decomposition/exhaustion record \(\Dec{D}\) is unique
    or canonical;
  \item a full six-primitives algebra or global primitive
    independence;
  \item empirical realization;
  \item internal instrument-family completeness;
  \item any statement about broader domains without further work.
\end{itemize}
These are scope boundaries, collected in
Section~\ref{sec:limitations-nonclaims} and enforced locally by the
theorem-chain results that would otherwise overread them.

\subsection{The objects at a glance}\label{subsec:objects-at-a-glance}

The paper works with a small set of objects, defined formally in
Sections~\ref{sec:fatcd-domain} and~\ref{sec:decomposition-exhaustion};
this preview fixes their roles so that the earlier sections can be read with
them in mind. A \emph{description} \(D\) is a finite audited typed closure
description (Definition~\ref{def:fatcd}). Its raw content \(\Raw{D}\) is a
finite collection of typed raw records, such as objects, maps, relations,
paths, transitions, traces and checks (Definition~\ref{def:raw-record}),
together with annotations of level, threshold, lift, forgetting, visibility
and empirical bridge (Definition~\ref{def:annotations}). Its audit part
\(\Audit{D}\subseteq\Raw{D}\) holds the bookkeeping record of every claim it
carries (Definition~\ref{def:audited}). A \emph{role projection} \(\pi_i(D)\)
reads a cluster of raw features as the role \(P_i\), and is admissible only
when threshold, witness-schema, collapse-case, level, nonclaim and audit
data are attached to the cluster
(Definitions~\ref{def:threshold-attachment} and~\ref{def:role-projection}). A
\emph{decomposition record} \(\Dec{D}\) gives each of the six role channels
exactly one of five statuses (Definition~\ref{def:channel-status}): an
\activeproj{}, when an admissible projection exists; \collapsedbc{}, when
otherwise a candidate cluster offered as the witness fails its witness
schema; \belowthr{}, when otherwise a covered seed is present (so no
admissible projection and no failed offer exists; the name refers to incomplete
witness or attachment data, not necessarily a numerical threshold); \outsidescope{}, when otherwise
role content is declared outside the covered scope; and \absentrec{}, with the
absence recorded, in the remaining case. The
raw features that no channel accounts for form the set \(\Residual{D}\) of
\emph{record residuals}, each classified under the scheme of
Definition~\ref{def:residual-scheme}; that definition calls every content
record and structural annotation a \emph{residual feature} (\emph{kind
residual}), including records used by an active cluster. The
scheme has eleven labels: (i)--(ii) membership in, or sharing between, active
clusters; (iii)--(iv) presentation variants and level shifts; (v) role-aligned or covered-seed
content not absorbed by an active projection, content checked in a failed offer, or an activation-threshold refinement; (vi)--(vii) visibility and
empirical-bridge annotations; (viii) content outside the covered scope; (ix)
bookkeeping, annotation and assertion-free data; (x) an unresolved candidate,
placed by none of (i)--(ix); (xi) a recorded seventh-role screen.

\paragraph{Working vocabulary.} Several words carry a technical meaning
throughout. The entries below anticipate definitions of Sections~\ref{sec:fatcd-domain}--\ref{sec:decomposition-exhaustion} and are meant to be consulted when a term is first met.
\begin{itemize}[topsep=2pt,itemsep=1pt]
  \item \emph{Honest bookkeeping}: every claim a description carries, and
    every change of level, translation, lift, forgetting map, visibility
    assumption and empirical bridge that the claim uses, has an explicit
    record carried in \(\Audit{D}\) (Definitions~\ref{def:annotations} and~\ref{prop:honest-bookkeeping}).
  \item \emph{Threshold annotation}: a declared record attached to a
    cluster of raw features before the cluster is read as a role; the
    reading also needs a witness schema, the collapse cases, a level
    annotation, a nonclaim annotation and an audit record
    (Definition~\ref{def:role-projection}).
  \item \emph{Witness schema}: a statement of what data in the cluster
    would support that reading.
  \item \emph{Instrument}: the declared means under which content is
    visible or suppressed (Definition~\ref{prop:visibility}).
  \item \emph{Covered scope}: the records not satisfying the outside-scope
    criterion of Definition~\ref{def:residual-scheme}. A record lies outside it when an assertion quantifies over a
    sort with no recorded finite carrier or uses a predicate whose extension
    \(D\) does not record, the assertion is true in some expansion of
    \(D\) and false in another, and its other assertions are of the same kind or
    serve a role witness (exact criterion: paragraph \emph{Terms used in the
    label criteria} after Definition~\ref{def:residual-scheme}). Such a
    record receives label (viii) and is certified for no role. Whether a
    record lies outside depends on presentation: a finite assertion written
    as a formula over recorded field values lies inside the covered scope,
    while the same assertion written with a declared predicate name whose
    extension \(D\) does not record lies outside
    (Remark~\ref{rem:recognition-open}). Not every finite, recorded feature
    inside the covered scope meets the recognition criteria: the normalization
    predicate of Remark~\ref{rem:recognition-open} does not.
  \item \emph{Covered}: in \emph{covered seed} the word names seed clauses
    (1)--(6) of Definition~\ref{def:channel-status}; in \emph{covered scope}
    it refers to the outside-scope criterion of Definition~\ref{def:residual-scheme}, under which a record with an assertion that is not finitely checkable lies inside the covered scope when another of its assertion fields is finitely checkable and serves no witness; in label (ix), \emph{already-covered}
    bookkeeping is an annotation whose non-exempt fields are data or
    references. The three uses are unrelated.
  \item \emph{Closure}: several senses occur. The \emph{closure content} of a
    feature (Definition~\ref{def:recognition-hypothesis}) is the subdescription
    reached by its references; a \emph{closure operation} in criterion (xi) is an
    extensive, monotone, idempotent map on a finite order; \emph{reference-closed}
    means \(\Closed{D}\). A \emph{closure role} is one
    of the six roles, in the sense of Foundations~I; the \emph{closure
    condition} \(\Closed{D}\) requires the references and bookkeeping
    specified in Definition~\ref{def:closed}; and the \emph{closures} in
    Section~\ref{sec:upper-bound-classification} classify the stated
    recognized forms of probability, causality and agency without the
    recognition hypothesis.
  \item \emph{Residual feature}: in
    Definition~\ref{def:residual-scheme}, every content record and every
    structural annotation (a \emph{kind residual}: the raw grammar fixes no
    role reading for any of them, so members of active clusters are included); in a
    decomposition record (Definition~\ref{def:decomposition-record}),
    \(\Residual{D}\) contains the raw records and annotations that no channel
    accounts for (the \emph{record residuals}); Section~\ref{sec:decomposition-exhaustion}
    compares the two sets. The
    upper-bound classification and the recognition hypothesis cover both
    classes; the decomposition theorems apply them to \(\Residual{D}\).
  \item \emph{Admissible}: the meaning depends on the object qualified. An
    admissible claim satisfies Definition~\ref{prop:honest-bookkeeping};
    admissible forgetting and lifts, visibility records and empirical-bridge
    claims satisfy Definitions~\ref{prop:forgetting-lifts},
    \ref{prop:visibility} and~\ref{prop:empirical-bridge}; an admissible
    description is a member of \FATCD{} (Definition~\ref{def:fatcd}); an
    admissible projection is a derived role projection
    (Definition~\ref{def:role-projection}), whose cluster satisfies its
    witness schema \(W_i\), with the attachment data of
    Definition~\ref{def:threshold-attachment}; and an admissible role witness carries the common witness data of
    Definition~\ref{def:common-witness-data}. In
    Section~\ref{sec:integrated-stress}, a misreading is called excluded when it
    is not licensed by the results of
    Sections~\ref{sec:fatcd-domain}--\ref{sec:decomposition-exhaustion}.
  \item \emph{Closed unconditionally}: said of probability, causality and
    agency, it means closed without the recognition hypothesis, in the
    recognized forms of
    Propositions~\ref{prop:probability-closure}--\ref{prop:agency-closure}.
    Other finite, audited, in-domain forms remain subject to the
    recognition hypothesis; content outside the covered scope is
    classified as an outside-domain structure.
  \item \emph{Recognition hypothesis}
    (Definition~\ref{def:recognition-hypothesis}): every residual feature is
    recognized in one of five ways: (R1) it is role content: it contributes
    a required field to an aligned cluster (that is, removing the record from
    the cluster falsifies its witness predicate; the phrase concerns
    membership, not a particular field) or carries a covered seed, and
    each of its assertion fields is required for a witness at an aligned
    cluster containing it, an input to that witness's check, or a named
    field of its seed; or it is a member of an honestly failed offer whose
    assertion fields are inputs to the recorded check; (R2) it is a
    structural annotation carrying only the fields its declared operation
    requires and no independent assertion, or a bookkeeping annotation whose
    non-exempt fields are data or references; (R3) it is role-neutral data
    without an assertion of its own, or a predicate record supplying no witness field whose only assertion field is a formula with free variables whose truth value varies under assignments respecting the declared sorts of the free variables, using only predicates whose extensions \(D\) records and quantifying only over sorts whose carriers \(D\) records, referenced by no claim annotation and by certificates only through \(s(t)\) (this branch constrains presentation, not content; Remark~\ref{rem:recognition-open}); (R4) it is a checked re-presentation or
    level translation of such content; (R5) an assertion of it uses an unrecorded predicate or uncarried sort and is true in some expansion of \(D\) and false in another, its other assertions being of the same kind or serving a witness (an outside-domain structure).
    The upper bound, that no seventh role is needed, is proved under this
    hypothesis. It is a genuine restriction: some descriptions in the
    domain do not satisfy it. Recognition constrains presentation, not content:
    Remark~\ref{rem:recognition-open} shows how a closed check-language
    formula without atoms \(W_j\) can be recorded, with its truth value, as a
    covered \Ptwo{} seed; recognition alone does not show that the feature
    concerns representability.
\end{itemize}

\paragraph{Symbols used before their definition.} Sections~\ref{sec:primitive-role-architecture}
and~\ref{sec:admissibility-meta-theory} use the following symbols, which are
defined later.
\begin{center}
\small
\begin{tabularx}{\linewidth}{@{}l>{\raggedright\arraybackslash}X>{\raggedright\arraybackslash}p{0.30\linewidth}@{}}
\toprule
Symbol & Meaning & Defined in \\
\midrule
\FATCD{} & the domain of finite, closed, audited descriptions & Definition~\ref{def:fatcd} \\
\(\Raw{D}\) & all declared records of \(D\), content records (``raw records'') and annotations alike; its elements are the raw features & Definition~\ref{def:annotations} \\
\(\mathrm{Ann}_\Theta\) & a threshold annotation & Definition~\ref{def:annotations} \\
\(\Audit{D}\) & the honest audit annotations of \(D\); claim, promotion, translation and similar annotations they reference are \emph{carried in} \(\Audit{D}\) & Definition~\ref{def:annotations} \\
\(\pi_i(D)\), \(W_i\) & a derived role projection and its witness predicate & \(W_i\): \emph{Witness types} table, Section~\ref{sec:primitive-role-architecture} (paragraph \emph{Witness predicates}); \(\pi_i\): Definition~\ref{def:role-projection} \\
\(\Status{i}{D}\) & the status of channel \(i\), one of five values & Definition~\ref{def:channel-status} \\
\(\Dec{D}\), \(\Residual{D}\) & a decomposition record and its record residuals & Definition~\ref{def:decomposition-record} \\
kind residual & every content record and structural annotation, including members of active clusters & Definition~\ref{def:residual-scheme} \\
\(\operatorname{Classify}\) & the scheme assigning a label to each residual feature & Definition~\ref{def:residual-scheme} \\
covered seed & the core record of a role witness, which may be present without the rest of that witness: an update map with its source relation; a one-variable macro-specification whose recorded evaluations take both truth values, with a realization relation targeting it; a route non-reducibility certificate; a law-bearing stage, order, transition or refinement record; a quotient relation or packaging map; an event, provenance or trace record, or a trace check & Definition~\ref{def:channel-status} \\
\(\mathrm{Audit}_0(D)\) & the honest audits whose certificates and conclusion contain no atom \(W_j\) & Definition~\ref{def:audited} \\
declared use & the typed field of a map, partial-map or relation record, with value update, packaging, quotient, realization or other & Definition~\ref{def:raw-record} \\
assertion field & a field of role predicate, comparison, invariant, law or claim & Definition~\ref{def:raw-record} \\
honest audit & an audit annotation meeting the honesty condition & Definition~\ref{def:audited} \\
offer, failed offer & a claim naming \(P_i\) and \(C\) with an honest audit recording \(W_i(C)\) (for a failed offer, recording \(W_i(C)=\mathrm{false}\)) & Definitions~\ref{def:audited} and~\ref{def:channel-status} \\
well-typed interpretation & an assignment of values to the datum fields of \(D\) keeping its records, kinds, declared uses, distribution-valued markers, field roles, references and formula- and term-valued fields fixed, unrecorded predicate extensions and unlisted-sort carriers varying arbitrarily; required certificates must match their prescribed conditions in all of them & Definition~\ref{def:audited}(c1) \\
\bottomrule
\end{tabularx}
\end{center}

\begin{figure}[t]
\centering
\begin{tikzpicture}[>={Stealth[length=2mm]},
  box/.style={rectangle,rounded corners,draw,align=center,font=\footnotesize,minimum height=8mm,text width=2.6cm}]
\node[box] (s2) at (0,0) {\S\ref{sec:primitive-role-architecture}\\ roles and levels};
\node[box] (s3) at (3.6,0) {\S\ref{sec:admissibility-meta-theory}\\ admissibility};
\node[box] (s4) at (7.2,0) {\S\ref{sec:fatcd-domain}\\ the domain \FATCD{}};
\node[box] (s5) at (3.6,-1.6) {\S\ref{sec:lower-bounds}\\ lower bounds};
\node[box] (s6) at (7.2,-1.6) {\S\ref{sec:upper-bound-classification}\\ upper-bound classification};
\node[box] (s7) at (5.4,-3.2) {\S\ref{sec:decomposition-exhaustion}\\ decomposition};
\node[box] (s8) at (1.8,-3.2) {\S\ref{sec:integrated-stress}\\ misreading checks};
\node[box] (s9) at (1.8,-4.8) {\S\ref{sec:scoped-exact-six}\\ scoped exact six};
\draw[->] (s2) -- (s3);
\draw[->] (s3) -- (s4);
\draw[->] (s4) -- (s5);
\draw[->] (s4) -- (s6);
\draw[->] (s5) -- (s7);
\draw[->] (s6) -- node[right,font=\scriptsize] {Thm.~\ref{thm:six-channel-exactness}(iii)} (s7);
\draw[->] (s7) -- (s9);
\draw[->] (s5) -- (s9);
\draw[->,dotted] (s7) -- (s8);
\draw[->,dotted] (s8) -- node[right,font=\scriptsize] {consistency check} (s9);
\draw[->,dashed] (s4) to[bend right=25] node[above,font=\scriptsize] {definitions} (s3);
\draw[->,dashed] (s4) to[bend right=35] (s2);
\draw[->,dashed] (s7) to[bend left=20] node[below left,font=\scriptsize] {channel statuses} (s2);
\draw[->,dashed] (s9) to[bend right=30] node[right,font=\scriptsize] {inventory of \(D_i\)} (s7);
\end{tikzpicture}
\caption{Dependencies between the sections that build the terminal theorem.
Solid arrows show the order of the argument; dashed arrows mark forward uses:
Sections~\ref{sec:primitive-role-architecture}
and~\ref{sec:admissibility-meta-theory} use definitions from
Section~\ref{sec:fatcd-domain}, and Section~\ref{sec:primitive-role-architecture}
uses the channel statuses of Section~\ref{sec:decomposition-exhaustion}.
Propositions~\ref{prop:channel-vs-activation} and~\ref{prop:non-uniqueness} also
use the record-by-record inventory of \(D_1,\dots,D_6\) in the proof of the last
clause of Theorem~\ref{thm:scoped-exact-six}, which uses only
Theorem~\ref{thm:six-role-lower-bounds} and the definitions; this forward use
creates no cycle. The
solid arrow from Section~\ref{sec:admissibility-meta-theory} to
Section~\ref{sec:fatcd-domain} marks that Definition~\ref{def:role-projection}
uses Definitions~\ref{prop:level-profile} and~\ref{prop:forgetting-lifts}, and
label criteria (vi) and (vii) of Definition~\ref{def:residual-scheme} use
Definitions~\ref{prop:visibility} and~\ref{prop:empirical-bridge}. Among the
definitions carried by the dashed arrow back to
Section~\ref{sec:admissibility-meta-theory}, Definition~\ref{def:audited}
requires every claim of a description to be referred to by an honest audit,
which Definition~\ref{prop:honest-bookkeeping} further requires to certify the claim's recorded statement or, for an offer, record the value of \(W_i(C)\); the dotted arrow from Section~\ref{sec:integrated-stress} to
Section~\ref{sec:scoped-exact-six} marks a consistency check that the proof of
the terminal theorem does not use. The solid arrow from Section~\ref{sec:lower-bounds} to
Section~\ref{sec:decomposition-exhaustion} marks Propositions~\ref{prop:channel-vs-activation}
and~\ref{prop:non-uniqueness}, which reuse the descriptions \(D_i\); the existence argument in the proof of
Theorem~\ref{thm:decomposition-existence} does not use Section~\ref{sec:lower-bounds};
only its closing non-vacuity remark cites Theorem~\ref{thm:six-role-lower-bounds}.}\label{fig:dependencies}
\end{figure}

\subsection{Organization}\label{subsec:organization}

Section~\ref{sec:primitive-role-architecture} fixes the role meanings,
the three levels, the boundaries between roles, and the selected-lift
atlas. Section~\ref{sec:admissibility-meta-theory} states the
admissibility rules. Section~\ref{sec:fatcd-domain} defines \FATCD{} and
poses the exact-six question over it. Section~\ref{sec:lower-bounds}
shows that each role occurs. Section~\ref{sec:upper-bound-classification}
classifies the recognized forms of probability, causality and agency
without the recognition hypothesis, and proves the general upper bound
under it.
Section~\ref{sec:decomposition-exhaustion} builds the six-channel
decomposition record and shows that it exists but need not be unique.
Section~\ref{sec:integrated-stress} checks the chain against misreadings,
and Section~\ref{sec:scoped-exact-six} states and proves the terminal
theorem. Section~\ref{sec:discussion} discusses the result and places it
beside earlier Six Birds work, and Sections~\ref{sec:limitations-nonclaims}
and~\ref{sec:future-work} give its limits and open directions. Appendix~\ref{app:provenance} records development provenance and
evidence hygiene, and Appendix~\ref{app:mechanization} documents the
Lean mechanization, with a note on the Lean counterpart of each result. Appendix~\ref{app:version-history} records the version history and is not needed for reading the paper. The Lean development checks the typed objects and rejects
malformed cases; the proofs of the numbered results are the ones in the body.
Figure~\ref{fig:dependencies} shows how the sections depend on one another.
Sections~\ref{sec:primitive-role-architecture}
and~\ref{sec:admissibility-meta-theory} use the domain \FATCD{} and the
channel statuses before they are formally defined. The record, witness and audit definitions refer to one another. A reader
checking proofs can use the following route, consulting forward references
as needed:
\begin{enumerate}[label=(\arabic*),topsep=2pt,itemsep=1pt]
  \item Section~\ref{subsec:objects-at-a-glance};
  \item Definitions~\ref{def:raw-record} and~\ref{def:annotations};
  \item the paragraphs of Section~\ref{sec:primitive-role-architecture} from
    \emph{Witness types} through \emph{Route certificates}, with the witness
    table and the full Definition~\ref{def:audited} at hand, using its clause
    (f) to order the evaluation of honesty and witness predicates;
  \item Definitions~\ref{def:finite}--\ref{def:fatcd}, then
    Definitions~\ref{prop:honest-bookkeeping}, \ref{prop:level-profile}
    and~\ref{prop:forgetting-lifts};
  \item the rest of Section~\ref{sec:fatcd-domain} and
    Definition~\ref{def:channel-status};
  \item the rest of Section~\ref{sec:primitive-role-architecture};
  \item Definitions~\ref{prop:visibility} and~\ref{prop:empirical-bridge} and
    Proposition~\ref{prop:nonclaim-closure};
  \item Sections~\ref{sec:lower-bounds}--\ref{sec:decomposition-exhaustion}
    and~\ref{sec:scoped-exact-six}.
\end{enumerate}
Propositions~\ref{prop:channel-vs-activation} and~\ref{prop:non-uniqueness}
also use the record inventory in the proof of
Theorem~\ref{thm:scoped-exact-six} (p.~\pageref{thm:scoped-exact-six}). The other definitions and results of
Section~\ref{sec:admissibility-meta-theory} are not cited in
Sections~\ref{sec:fatcd-domain}--\ref{sec:scoped-exact-six}, and
Section~\ref{sec:integrated-stress} is a separate check against misreadings;
the proof of Theorem~\ref{thm:scoped-exact-six} does not use it.

\paragraph{Example descriptions.} The running example \(D_{\mathrm{ex}}\)
(Sections~\ref{sec:primitive-role-architecture}, \ref{sec:fatcd-domain}
and~\ref{sec:decomposition-exhaustion}); the lower-bound descriptions
\(D_1,\dots,D_6\) (Propositions~\ref{prop:p1-witness}--\ref{prop:p6-witness});
\(D_0\) and the normalization description \(D_{\mathrm{norm}}\)
(Remark~\ref{rem:fatcd-minimality}, analysed in
Remark~\ref{rem:recognition-open}, paragraphs \emph{The counterexample} and \emph{Verification}), the latter failing the recognition
hypothesis; \(D_{\mathrm{scr}}\), which realizes label (xi)
(Remark~\ref{rem:recognition-open}, paragraph \emph{A description passing the screen}); and \(D_{\mathrm c}\), \(D_{\mathrm o}\),
\(D_T\) and \(D_{\mathrm{all}}\), which realize the five channel statuses
(proof of Proposition~\ref{prop:channel-vs-activation}).
 \section{Primitive role architecture}\label{sec:primitive-role-architecture}


This section fixes the meanings of the six typed closure roles
\Pone{}, \ldots, \Psix{}, the behavioral, verification, and
structural-categorical level pattern they share, the local boundary
distinctions between them, and the selected-lift atlas vocabulary.
These meanings are fixed for the remainder of the paper and govern
every subsequent definition, lower bound, and upper-bound
classification. Sections~\ref{sec:primitive-role-architecture}
and~\ref{sec:admissibility-meta-theory} use only the following from later
sections: a description is a finite set of typed records and annotations
(Definitions~\ref{def:raw-record} and~\ref{def:annotations}); a claim counts
when an honest audit certifies it (Definition~\ref{def:audited}); an
annotation is attached to a cluster when its target field lists exactly the
cluster's members (Definition~\ref{def:threshold-attachment}); and a channel's
status is the first of five cases that applies
(Definition~\ref{def:channel-status}). The paragraph \emph{Symbols used before
their definition} (Section~\ref{subsec:objects-at-a-glance}) lists the rest. A reader checking
proofs may read the paragraphs \emph{Witness types} to \emph{Route
certificates} below, then Section~\ref{sec:fatcd-domain} and
Definition~\ref{def:channel-status}, before the rest of this section.
The proofs of Sections~\ref{sec:fatcd-domain}--\ref{sec:decomposition-exhaustion}
and~\ref{sec:scoped-exact-six} use from this section only the role names of
Definition~\ref{def:primitive-roles}, the \emph{Witness types} table with the
check-language, certificate and law conventions before it and the paragraph
\emph{Route certificates} and the prescribed verification conditions after it,
Remark~\ref{rem:repaired-alignments}, the level values of
Definition~\ref{def:level-trichotomy} and Table~\ref{tab:boundary-matrix};
Proposition~\ref{prop:boundary-distinctions} is used by
Proposition~\ref{prop:typed-non-collapse}, and the other remarks and
Definition~\ref{def:selected-lift-atlas} fix vocabulary.

\subsection{Role meanings}\label{subsec:role-meanings}

\begin{definition}[Primitive closure roles]\label{def:primitive-roles}
The six primitive closure roles are
\begin{itemize}[topsep=2pt,itemsep=1pt]
  \item \Pone{} --- update-vs-quotient compatibility / descent
    obstruction~\citep{GrothendieckSGA1,JanelidzeTholen1994};
  \item \Ptwo{} --- macro-specification
    representability~\citep{MacLane1998}; when a satisfying realizer exists,
    the witness also requires a comparison with the classes of the kernel of
    a declared packaging map, whose declared domain contains the declared
    candidate class, certifying that some class of that kernel meets both
    the realizer set and its complement in the candidate class;
  \item \Pthree{} --- enriched route/coherence
    mismatch~\citep{BaaderNipkow1998,Nakahara2003}; here ``enriched'' and
    ``coherence'' name only the recorded rewrite rules (and, in certificate
    (b), the move weights), not enriched-category or coherence-theorem
    structure; \Pthree{} is not a directionality certificate;
  \item \Pfour{} --- stage/order/transition/refinement
    structure~\citep{Milner1989,BackVonWright1998};
  \item \Pfive{} --- quotient packaging /
    quotienting~\citep{KemenySnell1960};
  \item \Psix{} --- audit/bookkeeping/provenance/replay/checkability~\citep{GreenKarvounarakisTannen2007,Merkle1987}.
\end{itemize}
The operative content of role \(P_i\) is the witness predicate \(W_i\) fixed by
the \emph{Witness types} table below (the table in the paragraph \emph{Witness predicates},
p.~\pageref{tab:witness-types}, read with the paragraphs that paragraph lists); its record kinds and declared uses are those
of Definition~\ref{def:raw-record}, and \(\Audit{D}\) and \(\mathrm{Audit}_0(D)\)
are those of Definition~\ref{def:audited}. The phrases above name the roles.
\end{definition}

\paragraph{Reading the roles.} In plain terms, and with the lower-bound
witnesses of Section~\ref{sec:lower-bounds} as instances:
\begin{itemize}[topsep=2pt,itemsep=1pt]
  \item \Pone{} asks whether an update respects a grouping of states. For
    an update \(F\) and a grouping \(q\), \(F\) \emph{descends} when
    \(q(x)=q(y)\) implies \(q(Fx)=q(Fy)\); two grouped states of the declared carrier obstruct descent when their defined images lie in different target classes, or when exactly one lies in the domain of a partial \(F\). A \(W_1\) witness of descent (alternative (a)) further requires the domain of \(F\) to be saturated under \(q\) (every \(q\)-class meeting it lies inside it; elements in no listed class are allowed) and some class to contain two of its elements. Trivial descent is not excluded: if some class of \(q\) has two elements, source and target are both the relation of \(q\), that relation lists every class of \(q\), and the update's declared domain is the domain of \(q\), then the identity update satisfies (a), as does every constant update.
  \item \Ptwo{} asks whether a macro-level specification, a predicate
    \(s\), has a realizer through a declared realization relation. In short, \(W_2\)
    holds in two cases: no candidate in a declared class \(K\), the full
    declared domain of a packaging or update map, satisfies \(s\), although \(s\)
    is true at a declared record of the same kind outside \(K\); or some
    candidates satisfy \(s\) and some class of a declared packaging kernel whose
    domain contains the candidate class contains both a satisfying and a
    non-satisfying candidate. The
    specification \(s\) is, relative to a declared finite nonempty candidate
    class, \emph{realizable} when some candidate satisfies it and
    \emph{unrealizable} when none does, and \emph{representable over \(E\)} when
    its realizer set is a union of the sets \([x]_E\cap K\) for the candidate
    class \(K\), that is, when no class of \(E\) meets both the realizer set and
    its complement in \(K\). The empty realizer set of an
    unrealizable \(s\) is such a union; unrealizability is counted
    as a separate witness, not as a failure of representability over \(E\). The role is named for the question it settles; its witnesses are the two answers the packaging does not already supply: no candidate realizes \(s\) although \(s\) is satisfiable outside the candidate class, or some packaging class splits the realizers. A \(W_2\) assertion of
    unrealizability requires the recorded finite absence certificate, and a
    satisfying realizer yields a \(W_2\) witness only when the comparison
    certifies that the realizer set is not a union
    of \(E\)-classes. If it is such a union, as for the pullback of a
    predicate on packaged states, the comparison fails and there is no \(W_2\)
    witness; the pullback alone does not establish a \(W_5\) witness.
    ``Macro'' refers to the declared status of \(s\), not to definability over
    the quotient: in the satisfying branch the witness certifies that
    realizability of \(s\) is not already decided by the packaging, so the
    \Ptwo{} content is not \Pfive{} content. For example, with the states \(0,1,2,3\) and the packaging \(q\) of the
    running example, the specification \(s(r)\equiv(r=0\lor r=2)\) has realizer
    set \(\{0,2\}\); the \(q\)-class \(\{0,1\}\) meets both \(\{0,2\}\) and its
    complement, so realizability of \(s\) is not decided by \(q\)
    (Proposition~\ref{prop:p2-witness}).
  \item \Pthree{} compares two routes with the same start and end. A route
    grammar has generating moves and rewrite rules, each rule relating two
    words with the same source and target; two routes from \(u\)
    to \(v\) mismatch when they are distinct in the category presented by the
    moves and rules (a monoid when all moves are endomorphisms of one object), that is, when no chain of rule applications, in either
    direction and in any context, connects them,
    and a certificate of this, of one of the two types in the \(W_3\) row, is a
    \Pthree{} witness. For example, let
    \(a,b:u\to u\) be generating moves; if a certificate shows that no
    chain of rule applications, in either direction and in any context,
    connects the routes \(ab\) and
    \(ba\), they give a \Pthree{} mismatch. With the single rule
    \(aa\to a\), the monoid \(\{e,\bar a,\bar b\}\), with \(\varphi(a)=\bar a\) and
    \(\varphi(b)=\bar b\), in which \(e\) is the identity and \(xy=x\) for
    \(x\ne e\), records the first move: it identifies \(aa\) with \(a\) and
    separates \(ab\) from \(ba\).
  \item \Pfour{} is ordered structure: stages \(i_0<i_1<i_2\), transitions
    between them, and refinements with stated laws. The \(W_4\) row requires no order: one transition record with a certified law suffices, including the identity transition \(T(z)=z\) on a one-point set (\(D_T\), proof of Proposition~\ref{prop:channel-vs-activation}).
  \item \Pfive{} packages distinct states into one object, for example
    \(x\ne y\) with \(\mathrm{pack}(x)=\mathrm{pack}(y)\), together with
    evidence that the resulting quotient is valid.
  \item \Psix{} is checkable record-keeping: events, a trace, their
    provenance, and a replay or check of the trace.
\end{itemize}

\paragraph{A running example.} Throughout the paper we follow one small
description. Its states are \(0,1,2,3\), its update \(F\) sends
\(0\mapsto2\), \(1\mapsto0\), \(2\mapsto3\) and \(3\mapsto1\), and its lens
\(q\) groups the states into the classes \(A=\{0,1\}\) and \(B=\{2,3\}\); the
grouping is recorded as a relation record \(E\) of declared use quotient whose
classes are \(A\) and \(B\), so \(E=\ker q\). It
also records a trace \(0\to2\to3\) of the update, together with a replay and a
check of that trace. The lens packages the states into a quotient object
\(\{A,B\}\) with packaging map \(q\); this is \Pfive{} content. Descent of
\(F\) through that quotient fails: \(q(0)=q(1)=A\), while \(q(F0)=q(2)=B\) and
\(q(F1)=q(0)=A\). This split pair is \Pone{} content, and the example shows the
permitted interaction of the row \Pone{}/\Pfive{} of
Table~\ref{tab:boundary-matrix}: \Pone{} uses the \Pfive{} packaging when it
checks descent, and neither is the other. The recorded trace, with its replay
and check, is \Psix{} content. The description also records \(F\) as a
transition record with a certified law and an index record naming \(A\) and
\(B\). Its full record list is the table in the paragraph \emph{The running example as a description} (Section~\ref{sec:fatcd-domain}), and its decomposition record is the paragraph \emph{The running example, decomposed} (Section~\ref{sec:decomposition-exhaustion}); both use these two records.

\subsection*{Witness predicates and the check language}\addcontentsline{toc}{subsection}{Witness predicates and the check language}\label{subsec:witness-predicates}
\paragraph{Witness types.} Here a description is a finite set of typed
records (Definition~\ref{def:raw-record}); map and relation records declare a
use (update, packaging, quotient, realization or other); \(\Audit{D}\) is the
set of honest audit annotations (Definition~\ref{def:audited}); and a cluster
is a finite set of records. The table below fixes, for each role, when a
finite cluster \(C\) of records \emph{supplies the witness type} of \(P_i\),
that is, when \(W_i(C)\) holds; this table, with the paragraphs listed under \emph{Witness predicates}, is the definition of \(W_i\) that Definition~\ref{def:role-projection} uses. ``Recorded'' means
given by finite typed fields of, or references to, declared records of the
description. For \(W_i(C)\), witness fields must be carried by members of
\(C\). A reference to a declared record outside \(C\) establishes its identity
and type but does not import its fields as witness data. A certificate
carried by a member of \(C\) is evaluated over the recorded values of \(D\). A formula \emph{reads the values of} a record \(d\) when \(d\)'s identifier occurs in it as the record argument of a field-value term, as the record whose graph, table or relation an application or relation atom uses, in \(t\in\mathrm{dom}\,d\) or \(t<_d t'\), as the record whose field a list operation or a quantifier bound uses, or as the predicate record of an abbreviation \(d(t)\); it also reads \(d\) when a variable or an indexed list entry occurring as the record argument of a field-value term can denote \(d\) under the formula's quantifier bounds and recorded lists (for a variable bounded by all declared records, or by those of one kind, every such record). This is a syntactic condition, met even when the truth value does not depend on those values. A
record whose values a certificate reads is an input to it and, when the row names that
record, must belong to \(C\). A row names the records it lists as
constituents, together with the declared endpoint or index records of a
\(W_4\) law-bearing record; the map, partial-map or transition record of a \(W_6\) replay certificate is named by the row and belongs to \(C\); other records a certificate reads, such as that map's
declared domain, need not belong to
\(C\). The \Psix{} audit
entry is checked separately in \(\mathrm{Audit}_0(D)\subseteq\Audit{D}\) and need not be a content member
of \(C\).

\paragraph{Check language.} Checks are written in a fixed language: in summary, terms are record identifiers, rational constants, recorded field values, list lengths, entries and concatenations, applications of recorded maps and tables, finite sums and products, and word products \(\varphi^{m,e}(w)\); atoms are equality, order, \(t<_o t'\), \(t\in\mathrm{dom}\,F\), list membership and subword occurrence, relation atoms and \(W_i(C)\) (defined at the end of the paragraph \emph{Sorts}); formulas close atoms under \(\neg,\wedge,\vee,\to\) and bounded quantifiers. In detail: terms are record identifiers,
rational constants, recorded field values and applications of recorded finite
maps, that is, of the graphs listed by map, partial-map and transition records, combined by finite sums and products; a recorded field value is written \(f(d)\), with \(f\) the name of a datum or reference field and \(d\) a record-sorted term, its record argument (for instance \(v(w_k)\), \(A(E)\)), and \(f(d)_k\) is entry \(k\) of a list-valued field; the record argument may be an identifier, a variable or an indexed list entry, provided the denoted record carries that field, and a field name is not a record identifier, so \(f(d)\) is never read as an application \(F(t)\) of a map, partial-map, transition or table record; formulas are built from equality
and order by \(\neg,\wedge,\vee,\to\) (with \(\leftrightarrow\) as an abbreviation,
expanded before polarity is counted) and by quantifiers bounded by declared finite
records or by the finite set of all declared records or of the declared records of one kind. An object record listing no members is not a declared finite record in this sense: a quantifier over it is a sort quantifier of the assertion language, and no certificate contains one.
\par\emph{Lists and maps.} For every recorded finite list (a word in generating moves, the entry list
of a trace, a list held in a datum field, or the identifier list of a
reference field, such as the records an event or provenance record references
or the members an object record lists) the language also has list
length, indexing (\(\tau_k\) for the entry of a trace \(\tau\) at position \(k\), positions counted from \(0\)),
membership, concatenation and contiguous-subword occurrence, with their
standard finite-list meanings; the rewrite rules recorded in \(D\), taken in a
fixed order of their rule fields, form a finite list \(R\) whose \(k\)-th entry is a
pair \((\lambda_k,\mu_k)\), and \(\lambda_k\), \(\mu_k\) and \(\lambda_{k,p}\) (the entry of
\(\lambda_k\) at position \(p\)) are terms; an entry of a recorded finite list that is itself a
finite list is again a recorded finite list, so indexing and the other list
operations apply to it at every finite depth (\(L_{k,p}\), \(L_{k,p,q}\), as for
\(\lambda_{k,p}\)); for a recorded finite list \(w\) and a term
\(u(x)\), \(\sum_{k<|w|}u(w_k)\) is a term; and for a finite table field \(m\)
recording a binary operation on a recorded list \(N\) of rational constants,
an element \(e\) of \(N\) and a recorded finite map \(\varphi\) into \(N\), the
term \(\varphi^{m,e}(w)\) is the left-to-right \(m\)-product of
\(\varphi(w_0),\dots,\varphi(w_{|w|-1})\), with \(\varphi^{m,e}(())=e\). In route
certificate (a) the extension of \(\varphi\) to words is \(\varphi^{m,e}\), and in
(b) the weight of a word \(w\) is \(\sum_{k<|w|}\mathrm{wt}(w_k)\); quantifiers over rules, positions and list
entries are bounded by declared finite lists. Applications of recorded finite
maps include the graphs of refinement records and finite tables recorded in a
field; for a datum field \(L\) holding a finite list of finite tables on a
declared finite record \(X\), \(L_k(t)\) is the value at \(t\) of the \(k\)-th
listed table, so that closure of \(L\) under composition is
\(\forall k,l<|L|\,\exists j<|L|\,\forall x\in X\,(L_j(x)=L_k(L_l(x)))\) and ``\(c\)
differs from every listed map'' is
\(\forall k<|L|\,\exists x\in X\,(c(x)\neq L_k(x))\). \par\emph{Domains, rows and orders.} For a map, partial-map or transition record \(F\) the language has the atom
\(t\in\mathrm{dom}\,F\), true exactly when \(t\) is defined and its value
occurs as the first coordinate of a pair in the graph of \(F\) (an
interpretation assigns a map or transition record a graph defined at every
member of its declared domain), and an atomic formula containing an undefined
map or table application is false; for a transition record \(t\) with
codomain \(\mathrm{Dist}(S)\), \(t(x)(s)\) is the entry of the row \(t(x)\)
at the position of \(s\) in the member list of \(S\); for a declared finite record \(X\) and a term \(u(x)\),
\(\sum_{x\in X}u(x)\) is a term; for each order record \(o\) it has the atom
\(t<_o t'\).
\par\emph{Relations and predicates.} For a recorded finite relation record \(R\)
of arity \(k\), the language also has atoms \(R(t_1,\dots,t_k)\), true exactly
when the tuple of values of \(t_1,\dots,t_k\) is among the tuples \(R\) lists
(for a relation of declared use quotient, when the values lie in one listed
class); \(x\sim y\) denotes such an atom. For a predicate record \(s\) given by
an explicit check-language formula with one free variable, \(s(t)\) in a
certificate abbreviates that formula with \(t\) substituted for the variable; the abbreviation is unfolded at each recomputation using the formula then recorded in \(s\), so a certificate containing \(s(t)\) reads the formula field of \(s\) and every record the unfolded formula reads.

\paragraph{Certificates, assertion language and laws.} A \emph{certificate} is a record of a closed formula \(\varphi\) of
this language, the inputs it names, and a truth value \(b\); it is correct
when \(\varphi\) evaluates to \(b\) over the recorded values of the
description, and \emph{recomputation} is this evaluation. For an atom \(W_j(C)\) the evaluation is classical: its value is fixed by its row over the description, but the free-form conjuncts allowed in \(W_3\)(b) need only express their conditions in every well-typed interpretation, so recomputing a certificate that contains \(W_3\) atoms is not in general an effective procedure; ``finite'' in the outside-scope criterion refers only to recorded extensions and carriers. Predicate,
invariant and claim fields hold formulas of the \emph{assertion language}, and
a comparison or law field holds such a formula or a certificate, whose formula
is in the check language. The assertion language is the check language extended by the predicate
symbols declared by predicate records of the description, each with a declared
arity, and by quantifiers over sorts declared by object records that list no
members. Given an assignment to any free variables, such a formula has a truth value over \(D\) when every predicate extension it
uses and every quantified sort's carrier are recorded
(Section~\ref{sec:fatcd-domain}); otherwise it is evaluated in \emph{expansions} of \(D\): an expansion keeps every recorded value, extension and carrier of \(D\), assigns to each sort declared by an object record listing no members an arbitrary set (possibly empty), and assigns to each declared predicate whose extension \(D\) does not record an arbitrary relation of its declared arity on the carriers of its argument sorts (\(\mathbb Q\) for a rational-sorted argument). Certificates,
the rows of the table below and the seed clauses use the check language only;
in particular the formula of a macro-specification in the \(W_2\) row and in
seed clause (2) of Definition~\ref{def:channel-status} is a check-language formula containing no atomic \(W_j\)
formula.
\par\emph{Sorts.} Record identifiers and rational constants are distinct sorts. The declared
types of fields, list entries and map or table arguments determine their
sorts. A numeral or rational constant in record-sorted position (a record-typed argument of a
recorded map, an argument of \(W_j\), a term substituted for a record-sorted
variable, as in \(s(t)\), or compared with a record-valued term)
denotes the record, with \(=\) as identity of records and \(<\) between record-sorted terms as the atom \(t<_o t'\) of the order
record \(o\) the formula names (in an order law, the order record carrying it; in a refinement law for
\(\rho:I\to I'\), \(<\) and \(<'\) are the atoms of the declared order records on
\(I\) and on \(I'\)); a numeral in rational-sorted position (an argument of
\(+\) or \(\cdot\), a rational-typed argument of a map or table, or compared
with a value field) denotes the rational constant. In particular, monoid-table
elements and arguments are rational-sorted, while the lists supplied to
\(W_j\) contain record identifiers; equality compares terms of the same sort. The language also
contains, for each \(i\le 6\) and each finite set \(C\) of declared record
identifiers (written as a list of its members), an atomic formula \(W_i(C)\); its value is determined by the \(P_i\)
row of the table below, with the \(\Audit{D}\) conjunct of the \(W_6\) row
evaluated as in Definition~\ref{def:audited}. This is not circular: every certificate a row, law or seed clause reads contains no atom \(W_j\) (paragraph \emph{Failure and placement}), and the \(\Audit{D}\) conjunct of \(W_6\) ranges over \(\mathrm{Audit}_0(D)\), whose certificates and conclusions contain no such atoms (Definition~\ref{def:audited}(f)); so the rows are evaluated before any \(W_j\) atom is. Record kinds, declared uses and
field roles are recorded field values. These tag values, and elements of declared finite sets whose declared type is neither record-valued nor rational-valued, form a third sort, compared only by equality; record identifiers and rational values retain their respective sorts. Recorded field values are the values of datum and reference fields together with these; formula-valued fields and the formula of a certificate field are not terms, and a term-valued field enters only through its recorded term.
\par\emph{What a required certificate must express.} Whenever a row of
the table below, a law, or a covered seed requires a certificate of a
property, the certificate's formula must express that property: over the named inputs
it must be equivalent, in every well-typed interpretation of
Definition~\ref{def:audited}, to the verification condition prescribed for
that certificate by its row, law or seed clause. In addition, its formula must be
obtained from that condition, with the named inputs substituted, using only the
following steps: omitting conjuncts true in every well-typed interpretation (as
described below), expanding a quantifier or sum bounded by a reference-field
list into the finite conjunction, disjunction or sum over its members, and
replacing an application of a term-valued field by its recorded term. No other
conjunct may be added, even one true in every well-typed interpretation. \par\emph{Exceptions.} Two rows adjust this rule where they are displayed: the \(W_6\) certificate has one of two fixed forms, and for \(W_3\)(b) only the displayed conjuncts and the step checks are fixed in form, the other conjuncts needing only to express their conditions. \par\emph{Further conventions.} A well-typed interpretation is an interpretation in the sense of Definition~\ref{def:audited}(c1); in particular it
keeps the records of \(D\), their kinds, declared uses, field roles, references
and formula- and term-valued fields, and lets every other datum field, such as
a map graph or a multiplication table, take any value of its declared type.
The prescribed conditions are displayed after the table: for \(W_6\) the exact
forms, and for the other rows and for laws the formulas in the list beginning
``For the other rows and for laws''. The certificate's truth value must be
true, and it must be correct; an evaluation certificate \((s(t),b)\) recorded
under the \(W_2\) row or seed clause (2) records a value, not a required
property, and may have \(b=\mathrm{false}\). \par\emph{Named elements.} When a row, law or seed clause prescribes that some element, pair or class has a property, the prescribed
verification condition is that property for the element, pair or class the
certificate names among its inputs, as in alternative (b) of the \(W_1\) row, except where the displayed condition states the existence by a quantifier (the second conjunct of \(W_1\)(a), the coverage conjunct \(\forall z\in Q\,\exists u\in X\,(q(u)=z)\) of \(W_5\), the last conjunct of the order law, the formula \(\mathrm{Step}\) in the step checks of \(W_3\)(b), and the conditions of criterion (xi)(e)). \par\emph{Omitted conjuncts.} A conjunct that is true in every well-typed interpretation, such as \(x\in X\) for a declared domain \(X\) or \(k_0\in K\) for a declared candidate class \(K\) (reference-field lists being fixed by Definition~\ref{def:audited}(c1)), may be omitted without affecting equivalence. The \(W_1\)(b) certificates of \(D_{\mathrm{ex}}\) and \(D_1\), the \(W_2\) certificates of \(D_2\), and the comparison and realization-relation certificates of Remark~\ref{rem:recognition-open} omit exactly these conjuncts. The certificate \((\neg s(c),\mathrm{true})\) of the paragraph \emph{One record suffices for recognition} instead expands \(\forall k\in K\) over \(K=\{c\}\). These conditions must imply the
asserted property; for \(W_3\) they are the monoid or rewriting checks, not an
equivalence to route non-reducibility itself. \par\emph{Failure and placement.} A certificate that fails recomputation supplies no
evidence for a witness, law or seed clause requiring that certificate; the
record's other fields are assessed independently. A certificate used to
establish a witness-table row, law or covered seed has a formula containing no
atomic \(W_j(C)\); such atoms may instead be certified in an audit after the
witness predicates have been determined. A certificate may be
carried in a comparison or check record, or in the fields of the relevant
relation, map, order, transition or refinement record.
\par\emph{Laws.} A \emph{law} is an
order, transition or refinement record (a binary relation record whose declared source and target are one object record \(I\) listing its members carries an order when it carries a certificate whose formula is the order law displayed after the table on the declared indices \(I\), with \(<\) read as its relation atom; it then counts as an order record in the \(W_4\) row, in seed
clause (4) of Definition~\ref{def:channel-status} and for the atom
\(t<_o t'\), read as its relation atom) together with a certificate of a stated property of it: that an order relation on the
declared indices is irreflexive and transitive and relates at least one pair, that a transition map is
defined on its declared domain (for a transition record this holds in every well-typed interpretation, so that conjunct may be omitted) with values in its declared codomain, where the
codomain is a declared record listing finitely many members, or the set
\(\mathrm{Dist}(X)\) of probability distributions on such a listed set \(X\), recorded by a codomain
reference to \(X\) together with a datum field marking the codomain as
distribution-valued (this marker, like a declared use, is kept fixed by every
well-typed interpretation of Definition~\ref{def:audited}(c1)),
membership in \(\mathrm{Dist}(X)\) being checked by nonnegativity and sum one,
or that a
refinement map between declared orders sends its declared source into its declared target and preserves the order.

\paragraph{Witness predicates.}\label{par:witness-predicates} A role witness does not by
itself supply the separate threshold, level, collapse-case and audit
attachments of Definition~\ref{def:threshold-attachment}. Operationally, \(W_i(C)\) holds exactly when \(C\) records the data and relationships required by row \(P_i\) and its defining paragraphs: named constituent records belong to \(C\); certificates of required properties are carried by members of \(C\), recorded true, recompute true over \(D\), and satisfy the prescribed form requirements; for \(W_2\), recorded evaluations of \(s\), including false evaluation certificates, are correct; and the \(W_6\) audit entry satisfies that row's reference and conclusion requirements in \(\mathrm{Audit}_0(D)\). The predicates \(W_i\) are
defined by the table together with the paragraphs \emph{Check language} and \emph{Certificates, assertion language and laws} before it (the latter fixes the laws of the \(W_4\) row), the paragraph \emph{Route certificates} (for \(W_3\)), the paragraph \emph{What a required
certificate must express} with its continuations \emph{Exceptions} and \emph{Failure and placement}, the verification conditions displayed after the
table, the membership convention of the paragraph \emph{Witness types}, and, for the \(\Audit{D}\) conjunct of \(W_6\), Definition~\ref{def:audited}(f); Proposition~\ref{prop:alignment-rules}
derives from it which candidate features can support each role.
\begin{center}
\small
\begin{tabularx}{\linewidth}{@{}l>{\raggedright\arraybackslash}X@{}}
\toprule
Role & \(W_i(C)\) holds exactly when \(C\) records \\
\midrule
\Pone{} & a map or partial-map record \(F\) whose declared use is \emph{update} (Definition~\ref{def:raw-record}), source and target relation records of declared use quotient (the two may be the same record; they are the relation records that the row's certificate names as its source relation \(\sim\) and target relation \(\sim'\)), and either (a) a finite certificate that the domain of \(F\) is saturated under the source relation, that some class of the source relation contains at least two elements of that domain, and that every related pair in that domain has related images, or (b) a comparison or check record of \(C\) naming a pair \((x,y)\) of members of the declared domain of \(F\) and carrying a certificate that \(x\) and \(y\) are related by the source relation and that either exactly one of them lies in the domain of \(F\) or \(F(x)\), \(F(y)\) are defined, each lies in a class of the target relation, and they are not related by the target relation; in either case the certificate identifies its inputs and the update \\
\Ptwo{} & (1)~a macro-specification predicate \(s\) whose formula has exactly one free variable \(r\), occurring in it and ranging over records of the candidates' kind; (2)~a relation record \(\mathrm{Real}(r,s)\) of declared use realization and, in fields of that realization relation or of comparison or check records of \(C\), the evaluation of the formula of \(s\) at each member of the declared finite, nonempty class \(K\) of candidate realizers, every evaluation of \(s\) recorded in \(C\) being correct (equal to the value of the formula of \(s\) at that record in \(D\)); (3)~a certificate either exhibiting a satisfying realizer or checking that none in that class exists.\newline \emph{No-realizer branch}: also (a) an evaluation, recorded in a field of the realization relation or of a comparison or check record of \(C\), at which the formula of \(s\) is true at a declared same-kind record outside that class and (b) a candidate class that is the full declared domain of a map record of \(C\) of declared use packaging or update.\newline \emph{Satisfying branch}: also a comparison with the classes of the kernel \(E\) of a map record of \(C\) whose declared use is \emph{packaging} and whose declared domain contains the declared candidate class, certifying that some class of \(E\) meets both the realizer set and its complement in the candidate class, that is, that the realizer set is not a union of the sets \([x]_E\cap K\), so that in the satisfying branch realizability is not reducible to quotient-class equality \\
\Pthree{} & two defined (paragraph \emph{Route certificates}, \emph{Definedness}) path or derivation records with the same declared source and target, written as words in the generating moves (map or partial-map records; since definedness reads their graphs, they belong to \(C\)), the finite rewrite rules \(\lambda\to\mu\) recorded with those routes or in referenced records that are members of \(C\), these being all rewrite rules recorded in \(D\), with the images \(\varphi(m)\) in (a) and the weights in (b) given for every move occurring in the two routes or in those rules, and a comparison certificate that the two routes are not related by those rules, in one of the two forms of the paragraph \emph{Route certificates} below \\
\Pfour{} & at least one stage, order, transition or refinement record with declared endpoints or indices and a law for it in the sense above; an index label without such a law does not suffice \\
\Pfive{} & distinct declared object or state records \(x,y\) (the micro-elements), a relation record \(E\) whose declared use is \emph{quotient}, a map record \(q\) whose declared use is \emph{packaging} to a declared quotient-object record \(Q\) (named by the row, so it belongs to \(C\)), and a check record of \(C\) carrying a certificate of quotient validity: \(E\) is the declared equivalence, \(q\) identifies \(x\) and \(y\), and the fibers of \(q\) agree with \(E\), with the quotient object covered by \(q\); a \Pone{} split-pair certificate alone does not suffice \\
\Psix{} & a declared event \(e\), a trace record \(\tau\), a record tied to \(e\) (one of the two references the other) that is \(\tau\) itself or a provenance record also tied to \(\tau\), a replay or check record tied to \(\tau\) carrying a certificate, in the sense above, whose formula states a replay property of the trace (each successive trace entry is the image of the preceding one under a declared map, partial-map or transition record whose codomain is not of the form \(\mathrm{Dist}(S)\)) or a provenance property of it (each trace entry is a record referenced by a provenance record \(p\) of \(C\) tied to \(\tau\), and the trace begins at a record the event references), and an entry \(a\in\mathrm{Audit}_0(D)\) that refers to a claim annotation, to the event and trace records of \(C\) and to that replay or check record, and whose conclusion is the formula of that record's certificate \\
\bottomrule
\end{tabularx}
\captionof*{table}{Witness types table. Witness predicates \(W_1,\dots,W_6\): for a fixed description \(D\), \(W_i(C)\) holds exactly when the finite cluster \(C\) records the data in row \(P_i\); the \(W_3\) row also refers to all rewrite rules recorded in \(D\), and the \(W_6\) row to \(\mathrm{Audit}_0(D)\). Every certificate a row requires must also meet the form requirements displayed after the table and in the paragraph \emph{What a required certificate must express}.}\phantomsection\label{tab:witness-types}
\end{center}

\paragraph{Route certificates.} (a)~\emph{Monoid certificate}: a finite monoid \(N\) given by its multiplication table (its elements recorded as rational constants in a list field and its multiplication as a finite table field) and the images \(\varphi(m)\in N\) of the moves, such that the extension \(\varphi\) of these images to words satisfies \(\varphi(\lambda)=\varphi(\mu)\) for every rule and \(\varphi(\rho_1)\neq\varphi(\rho_2)\), together with the checks that \(e\in N\), that \(N\) is closed under the recorded multiplication, and that \(N\) is associative with identity \(e\).
(b)~\emph{Rewriting certificate}: a certificate that, for a recorded rational weight at least \(1\) of each move (finitely many such weights generate a discrete set of word weights, so strict decrease terminates) extended additively to words, each rule strictly decreases the weight. It also records the list of all critical pairs (every overlap of left-hand sides, as in the prescribed condition below, including proper overlaps of a left-hand side with itself or with another, and inclusion overlaps of left-hand sides; non-equivalence in the free monoid on the moves implies non-equivalence under rule applications restricted to composable routes), with a check that the list is complete. For each critical pair and for each of \(\rho_1,\rho_2\) it records a finite rewrite sequence, each step checked to replace one occurrence of some \(\lambda\) by its \(\mu\), words being recorded as finite lists of moves; the two sequences of each critical pair end in one word, and those from \(\rho_1,\rho_2\) end in distinct words containing no left-hand side. By the critical-pair lemma the rewrite system on all words is locally confluent, and by Newman's lemma confluent. The distinct normal forms therefore exclude any sequence of applications of the recorded rules, in either direction, relating the two routes, even when intermediate words need not be composable.
\emph{Definedness}: a path or derivation record is \emph{defined} when it lists the successive records it passes through and each move's map or partial-map record is recorded as defined at the preceding record, with value the next one.
A record is \emph{tied to} another when one of the two references the other.
The recorded rule list may be empty; a \(W_3\) witness then certifies only that
two parallel words are distinct, as for two distinct moves \(a,b:u\to v\), and
in \(D_3\) the rule \(aa\to a\), applied in either direction, never identifies \(ab\) with \(ba\), since it changes only the lengths of runs of \(a\). Because the \(W_3\) row and both certificates range over all rewrite rules recorded in \(D\), a \(W_3\) witness is in general not preserved when records carrying further rewrite rules are adjoined; Proposition~\ref{prop:channel-vs-activation} therefore forms disjoint unions in which \(D_3\) alone records rules.

For \(W_6\), the certificate formula is exactly \(\bigwedge_{k+1<|\tau|}\tau_{k+1}=F(\tau_k)\), for \(|\tau|\ge2\) and a map, partial-map or transition record \(F\) whose codomain is not of the form \(\mathrm{Dist}(S)\), or exactly \(\bigwedge_{k<|\tau|}\tau_k\in P\wedge\tau_0\in R_e\), with \(P\) and \(R_e\) the concatenations of the reference-field identifier lists of \(p\) and of \(e\), expanded over the entry list of \(\tau\); it contains no additional assertion. The row admits degenerate instances: a trace \((s,s)\)
replayed by an identity map, and, since the provenance branch imposes no
length condition, a one-entry trace whose entry the event and provenance record
reference. In the provenance branch, the provenance record \(p\) and every
record whose fields the check reads belong to \(C\).

For the other rows and for laws, the prescribed verification condition is the
following formula over the named inputs (\(\sim,\sim'\) the source and target
relations, \(X\) the declared domain, and \(|{\sim}|\), \(|E|\) the lists of
members of the listed classes of the relation records \(\sim\), \(E\); \(\forall u\in|E|\,\chi\) abbreviates \(\forall u\,(E(u,u)\to\chi)\), \(u\) ranging over all declared records).
\begin{itemize}[topsep=2pt,itemsep=1pt]
  \item \(W_1\)(a):
\(\forall x\in X\,\forall y\in|{\sim}|\,((x\sim y\wedge x\in\mathrm{dom}F)\to y\in\mathrm{dom}F)
\wedge\exists x,y\in X\,(x\in\mathrm{dom}F\wedge y\in\mathrm{dom}F\wedge x\ne y\wedge x\sim y)
\wedge\forall x,y\in X\,((x\in\mathrm{dom}F\wedge y\in\mathrm{dom}F\wedge x\sim y)\to F(x)\sim'F(y))\);
  \item \(W_1\)(b), for the named pair:
\(x\in X\wedge y\in X\wedge x\sim y\wedge((x\in\mathrm{dom}F\leftrightarrow y\notin\mathrm{dom}F)\vee
(x\in\mathrm{dom}F\wedge y\in\mathrm{dom}F\wedge F(x)\sim'F(x)\wedge F(y)\sim'F(y)\wedge\neg F(x)\sim'F(y)))\);
  \item \(W_2\), no-realizer branch: \(\forall k\in K\,\neg s(k)\) for the
non-realizer certificate, together with a separately recorded true evaluation
\(s(t)\) for a named same-kind record \(t\notin K\); satisfying branch:
\(k_0\in K\wedge s(k_0)\) for the named \(k_0\) and, for the comparison,
\(x\in K\wedge y\in K\wedge q(x)=q(y)\wedge s(x)\wedge\neg s(y)\) for the named
\(x,y\);
  \item \(W_5\): \(\forall u,v\in X\,(E(u,v)
\leftrightarrow q(u)=q(v))\wedge x\ne y\wedge q(x)=q(y)\wedge\forall z\in Q\,\exists
u\in X\,(q(u)=z)\wedge\forall u\in|E|\,(u\in X)\wedge\forall u\in X\,(q(u)\in Q)\);
  \item order law, on declared indices \(I\): \(\forall i\in I\,\neg(i<i)\wedge\forall
i,j,k\in I\,((i<j\wedge j<k)\to i<k)\wedge\exists i,j\in I\,(i<j)\);
  \item transition law, with declared source
\(X\) and codomain \(Y\): \(\forall x\in X\,(x\in\mathrm{dom}\,t)\wedge\forall x\in X\,(t(x)\in
Y)\), where, when \(Y=\mathrm{Dist}(S)\) for a declared finite carrier \(S\),
membership means \(\forall s\in S\,t(x)(s)\ge0\) and
\(\sum_{s\in S}t(x)(s)=1\) (a transition law imposes no nontriviality; the
identity transition \(T\) of \(D_T\) in the proof of
Proposition~\ref{prop:channel-vs-activation} satisfies it);
  \item refinement law, for a declared map
\(\rho:I\to I'\): \(\forall i\in I\,(\rho(i)\in I')\wedge\forall i,j\in I\,(i<j\to\rho(i)<'\rho(j))\) (the
order-preservation conjunct is vacuous when the order on \(I\) relates no pair;
the range conjunct remains required; like the transition law it imposes no
nontriviality). \end{itemize}
The \(W_3\) certificates are those of the paragraph \emph{Route
certificates}. For \(W_3\)(a), with \(R\) the recorded rule list, \(M\) a list field of the comparison record, and \(w_1,w_2\) the recorded words of
\(\rho_1,\rho_2\), the prescribed condition is
\(e\in N\wedge\forall x,y\in N\,(m(x,y)\in N)\wedge\forall x,y,z\in N\,(m(m(x,y),z)=m(x,m(y,z)))\wedge\forall x\in N\,(m(e,x)=x\wedge m(x,e)=x)\wedge\forall a\in M\,(\varphi(a)\in N)\wedge\forall k<|R|\,(\varphi^{m,e}(\lambda_k)=\varphi^{m,e}(\mu_k))\wedge\neg\,\varphi^{m,e}(w_1)=\varphi^{m,e}(w_2)\wedge\forall k<|R|\,\forall p<|\lambda_k|\,(\lambda_{k,p}\in M)\wedge\forall k<|R|\,\forall p<|\mu_k|\,(\mu_{k,p}\in M)\wedge\forall p<|w_1|\,(w_{1,p}\in M)\wedge\forall p<|w_2|\,(w_{2,p}\in M)\).
For \(W_3\)(b) it is the conjunction of the positivity of the weight of every move occurring in \(R\), \(w_1\) or \(w_2\), written \(\forall k<|R|\,\forall p<|\lambda_k|\,(\mathrm{wt}(\lambda_{k,p})\ge1)\wedge\forall k<|R|\,\forall p<|\mu_k|\,(\mathrm{wt}(\mu_{k,p})\ge1)\wedge\forall p<|w_1|\,(\mathrm{wt}(w_{1,p})\ge1)\wedge\forall p<|w_2|\,(\mathrm{wt}(w_{2,p})\ge1)\),
\(\forall k<|R|\,(\mathrm{wt}(\lambda_k)>\mathrm{wt}(\mu_k))\), completeness and
correctness of the critical-pair list (for all rule indices \(k,l<|R|\), including \(k=l\), and every overlap of \(\lambda_k\) with \(\lambda_l\), whether an inclusion overlap (an occurrence of \(\lambda_k\) as a contiguous subword of \(\lambda_l\) for \(k\ne l\), including one at position \(0\) even if \(\lambda_k=\lambda_l\); for \(k=l\) there is no inclusion overlap) found by contiguous-subword occurrence or a proper overlap (a nonempty proper suffix of \(\lambda_k\) equal to a prefix of \(\lambda_l\)) found by indexing, the list has an entry consisting
of the overlap word and the two words obtained from it by the two overlapping
rule applications, and every entry has this form), the step check of each
recorded rewrite sequence, the condition that the two sequences of each entry
start at its two words and end in one word, and that the sequences for
\(\rho_1,\rho_2\) start at \(w_1,w_2\), and irreducibility and
distinctness of the normal forms of \(w_1,w_2\), each written with the list
operations of the check language. The step check of a recorded sequence \(L\) is
\(\forall j<|L|\,(j+1<|L|\to\mathrm{Step}(L_j,L_{j+1}))\), where
\(\mathrm{Step}(u,v)\) abbreviates
\(\exists k<|R|\,\exists p<|u|\,(p+|\lambda_k|\le|u|\wedge|v|+|\lambda_k|=|u|+|\mu_k|\wedge\forall q<|u|\,(q<p\to v_q=u_q)\wedge\forall q<|\lambda_k|\,(u_{p+q}=\lambda_{k,q})\wedge\forall q<|\mu_k|\,(v_{p+q}=\mu_{k,q})\wedge\forall q<|u|\,(p+|\lambda_k|+q<|u|\to v_{p+|\mu_k|+q}=u_{p+|\lambda_k|+q}))\).
For \(W_3\)(b) the exact-form requirement of the paragraph \emph{What a required certificate must express} applies to the displayed conjuncts and to the step checks; each remaining conjunct (critical-pair completeness and correctness, the endpoint conditions, irreducibility and distinctness) may be any check-language formula that expresses its condition in every well-typed interpretation. Equivalence in every well-typed interpretation is not decidable in general (datum lists may change length across interpretations and the language has \(+\) and \(\cdot\) on list lengths, so solvability of Diophantine equations over \(\mathbb N\) reduces to it), and neither are \(W_3\), the entailment requirement (c2) or criterion (viii); the paper uses these predicates classically and claims no decision procedure.

In the no-realizer branch of the \(W_2\) row, condition (a) supplies a true
evaluation of \(s\) outside the candidate class, so that the recorded
evaluations of \(s\) take both truth values, as seed clause (2) of
Definition~\ref{def:channel-status} requires; condition (b) fixes the candidate
class as the full declared domain of a packaging or update map, so that it is
not chosen ad hoc. These conditions restrict the recorded evaluations of \(s\).
Given the map required by (b) and the evaluations and certificates of the row, \(s(r)\equiv(r=t)\), for any
declared same-kind record \(t\notin K\), satisfies the no-realizer
requirements on \(s\). In the satisfying branch, given a packaging map \(q\) of the cluster whose declared domain contains \(K\) and the evaluations and certificates of the row, \(s(r)\equiv(r=x)\)
suffices whenever \(x\in K\) and its \(q\)-class contains a
second member of \(K\). A \(W_2\) witness thus certifies the recorded form of a
specification relative to the declared class
(Remark~\ref{rem:classification-scope}).

The three levels of Definition~\ref{def:level-trichotomy} separate what
happens (behavior), the certificate that it happens (verification), and
the structure that explains it, such as a quotient object with its
universal property (structure). The running example above shows \Pone{}, \Pfive{} and \Psix{}
in one small description.

\begin{remark}[Excluded alignments]\label{rem:repaired-alignments}
Three alignments that Definition~\ref{def:primitive-roles} and the
\emph{Witness types} table rule out
are used throughout later classification: generic transition data is
not \Pthree{}; generic path or derivation data is not \Pfive{}; and
generic semantic or interpretation data is not \Pone{}. These
constraints determine which feature clusters can later be interpreted
as active projections of each role. Precisely: a cluster without two defined
path or derivation records, or without a comparison certificate of the
\(W_3\) row, fails \(W_3\); a cluster without a relation record of declared use quotient,
or without a map record of declared use packaging, fails \(W_5\); and a cluster without
a map or partial-map record of declared use update fails \(W_1\). Each is read
off its row of the \emph{Witness types} table, so these alignments add no
hypothesis beyond that table.
\end{remark}

\begin{remark}[Drive face of \Psix{}]\label{rem:p6-drive}
Foundations I uses a specialization \(\mathrm{P6}_{\mathrm{drive}}\) of
\Psix{} for arguments about affinities and detailed
balance~\citep{Seifert2012}. This paper assigns it no witness type and no
result uses it; the point retained here is that a \Pthree{}
route/coherence-mismatch witness does not by itself establish directionality
(Definition~\ref{def:primitive-roles}).
\end{remark}

No proof in this paper uses the drive face; the terms of the remark are as
follows. For a Markov chain with stationary law \(\pi\), \emph{detailed balance}
means \(\pi_x P_{xy}=\pi_y P_{yx}\) for every pair of states. For a
cycle of states with positive stationary mass and positive transition
probabilities in both directions,
its \emph{affinity} is the logarithm of the product of the forward
probabilities divided by the product of the reverse ones; a nonzero
affinity witnesses a failure of detailed balance. The remark records that a \Pthree{} mismatch alone does
not establish a direction, the distinction also recorded in
Definition~\ref{def:primitive-roles}.

\subsection{Behavioral, verification, and structural-categorical levels}\label{subsec:levels}

Each role admits a common three-level pattern.

\begin{definition}[Level trichotomy]\label{def:level-trichotomy}
For each role \Pone{}, \ldots, \Psix{}, we distinguish
\begin{itemize}[topsep=2pt,itemsep=1pt]
  \item the \emph{behavioral level}: status, occurrence, or observed
    behavior of the role in a description;
  \item the \emph{verification level}: proof, witness, certificate, or
    checkable record supporting a behavioral
    claim~\citep{MartinLof1984};
  \item the \emph{structural-categorical level}: reusable structure
    supporting the primitive role.
\end{itemize}
Forgetting proceeds from structural-categorical to verification to
behavioral and is in general lossy. Lifting proceeds from behavioral to
verification to structural-categorical only after adding explicitly
recorded data, and selected lifts are not unique. Formally, a level is one of
the three values recorded in the level field of a level annotation; the
proofs of Sections~\ref{sec:lower-bounds}--\ref{sec:scoped-exact-six} require
that such a value be recorded and use only whether two recorded values coincide
(for the level crossings of Definition~\ref{def:role-projection}), not which
value is recorded.
\end{definition}

\begin{figure}[htbp]
\centering
\begin{tikzpicture}[>={Stealth[length=2mm]},
  lv/.style={rectangle,rounded corners,draw,minimum width=4.4cm,minimum height=7mm,font=\small}]
\node[lv] (st) at (0,2.6) {structural-categorical level};
\node[lv] (ve) at (0,1.3) {verification level};
\node[lv] (be) at (0,0) {behavioral level};
\draw[->] ([xshift=-6mm]st.south) -- node[left,font=\footnotesize] {forget (lossy)} ([xshift=-6mm]ve.north);
\draw[->] ([xshift=-6mm]ve.south) -- node[left,font=\footnotesize] {forget (lossy)} ([xshift=-6mm]be.north);
\draw[->,dashed] ([xshift=6mm]be.north) -- node[right,font=\footnotesize] {selected lift, with added data} ([xshift=6mm]ve.south);
\draw[->,dashed] ([xshift=6mm]ve.north) -- node[right,font=\footnotesize] {selected lift, with added data} ([xshift=6mm]st.south);
\end{tikzpicture}
\caption{The three levels of Definition~\ref{def:level-trichotomy}. Forgetting
moves down and loses data; lifting moves up only with recorded added data, and
the lift is a selection, not in general unique.}\label{fig:levels}
\end{figure}

\begin{remark}[Level content per primitive]\label{rem:per-primitive-levels}
The level trichotomy instantiates as follows.
\begin{itemize}[topsep=2pt,itemsep=1pt]
  \item \Pone{}: descent success or descent obstruction status; descent
    proof or obstruction witness; quotient data equipped with a descended
    map, or an explicit no-descended-map obstruction record.
  \item \Ptwo{}: certified unrealizability or failure of representability
    over the declared packaging kernel; the certificate of that status;
    macro-specification layer together with a realization relation.
  \item \Pthree{}: witness behavior of route disagreement; generated route
    terms and non-reducibility certificates; categorical route grammar or a
    reduced categorical universe.
  \item \Pfour{}: stage status, transition occurrence, or refinement
    behavior; order/preorder, transition, or refinement certificates;
    stage/refinement architecture.
  \item \Pfive{}: packaging or identification behavior; equivalence and
    quotient-validity certificate; quotient object with a packaging map
    and universal-property or tower/coarsening data.
  \item \Psix{}: audit-relevant event occurrence or record--event
    association; checkable audit record; audit schema with provenance,
    replay, and check relations.
\end{itemize}
The meta-theoretic properties of this trichotomy --- lossy forgetting,
selected non-unique lifts, and the absence of lower-to-higher determinacy
without added data --- are formalized in
Section~\ref{sec:admissibility-meta-theory}.
\end{remark}

\subsection{Local boundary distinctions}\label{subsec:boundary-matrix}

With role meanings and level pattern fixed, we turn to how the six
roles remain distinct. Table~\ref{tab:boundary-matrix} records the
pairwise local rules. These are local distinctions under scoped
bookkeeping, not a global primitive-independence theorem: two roles
may interact legitimately in a single description, but neither
collapses into the other. The rules below use the domain \FATCD{} and the
channel statuses before they are defined; a reader who prefers definitions
first can read Sections~\ref{subsec:raw-grammar}--\ref{subsec:threshold-attachment}
and Definition~\ref{def:channel-status} first.

\begingroup\footnotesize
\setlength{\LTleft}{0pt}\setlength{\LTright}{0pt}
\renewcommand{\arraystretch}{1.15}
\begin{longtable}{@{}>{\raggedright\arraybackslash}p{1.2cm}>{\raggedright\arraybackslash}p{4.3cm}>{\raggedright\arraybackslash}p{3.0cm}>{\raggedright\arraybackslash}p{3.0cm}>{\raggedright\arraybackslash}p{2.9cm}@{}}
\caption{Local boundary distinctions. Each row records the distinction
rule between two roles, a permitted interaction, the collapse that
is forbidden under scoped bookkeeping, and the bare partner datum \(X_b\)
on which that collapse would credit a role.}\label{tab:boundary-matrix}\\
\toprule
pair & distinction rule & permitted interaction & forbidden collapse & bare datum \(X_b\) \\
\midrule
\endfirsthead
\toprule
pair & distinction rule & permitted interaction & forbidden collapse & bare datum \(X_b\) \\
\midrule
\endhead
\bottomrule
\endfoot
\Pone{}/\Pfive{} & \Pfive{} supplies quotient packaging; \Pone{} tests selected update compatibility with quotient packaging & \Pone{} uses \Pfive{} when checking descent or obstruction & packaging alone is not \Pone{} & the \Pfive{} packaging map record alone \\
\Pone{}/\Pfour{} & \Pfour{} supplies stage/refinement structure; \Pone{} may vary by stage & stage-indexed descent or obstruction & stage index alone is not \Pone{} & the \Pfour{} stage-index record alone \\
\Pone{}/\Ptwo{} & \Pone{} is update compatibility; \Ptwo{} is macro-specification representability & representability may be studied over quotient structures affected by descent & representability failure is not automatically descent obstruction & the \Ptwo{} non-representability certificate alone \\
\Pone{}/\Pthree{} & \Pone{} concerns update maps, possibly partial, whose undefined values leave a route undefined; \Pthree{} compares already-defined routes & \Pone{} may block a route before \Pthree{} comparison & route undefinedness is not route mismatch & the partial update map and the undefined route alone \\
\Pone{}/\Psix{} & \Psix{} audits a \Pone{} event; \Pone{} is the event content & audited descent success or obstruction & audit record is not descent obstruction & the \Psix{} audit record alone \\
\Ptwo{}/\Pfive{} & \Pfive{} is quotient packaging; \Ptwo{} requires a macro-specification, a realization relation and a finite existence or nonexistence certificate, plus a comparison with quotient classes when a realizer exists & a specification representable over \(E\) supplies no \Ptwo{} witness in the satisfying branch; a separately witnessed \Ptwo{} projection may accompany \Pfive{} packaging & packaging alone is not \Ptwo{} & the \Pfive{} packaging map record alone \\
\Ptwo{}/\Pfour{} & \Pfour{} supplies stages; \Ptwo{} representability may vary by stage & stage-indexed representability & stage variation alone is not \Ptwo{} & the \Pfour{} stage records alone \\
\Ptwo{}/\Pthree{} & \Ptwo{} concerns representability of specifications; \Pthree{} concerns route disagreement & route grammar may be built over represented macro structure & non-representability is not route mismatch & the \Ptwo{} non-representability record alone \\
\Ptwo{}/\Psix{} & \Psix{} audits representability claims; \Ptwo{} is the representability content & audited representability or non-representability & audit trail is not representability & the \Psix{} audit-trail records alone \\
\Pthree{}/\Pfour{} & \Pfour{} supplies stage/regime structure; \Pthree{} route mismatch may be stage-indexed & route systems vary over valid stages & stage-indexed route mismatch is not \Pfour{} itself & the \Pthree{} route-disagreement certificate alone \\
\Pthree{}/\Pfive{} & \Pfive{} supplies quotient objects; \Pthree{} compares defined routes over objects & routes over quotient objects & quotient object is not route mismatch & the \Pfive{} quotient-object record alone \\
\Pthree{}/\Psix{} & \Psix{} audits route comparisons; \Pthree{} is the route mismatch content & audited route mismatch & audit record is not route mismatch & the \Psix{} audit record alone \\
\Pfour{}/\Pfive{} & \Pfour{} supplies stage/refinement structure; \Pfive{} supplies packaging at a stage & stage-indexed packaging & packaging is not stage structure & the \Pfive{} packaging map record alone \\
\Pfour{}/\Psix{} & \Pfour{} supplies staged structure; \Psix{} audits staged histories or transitions & audited stage transitions & audit record is not staging & the \Psix{} audit record of staged histories alone \\
\Pfive{}/\Psix{} & \Pfive{} supplies packaging; \Psix{} audits packaging events & audited packaging & audit record is not packaging map & the \Psix{} audit record of packaging events alone \\
\end{longtable}
\endgroup
 
\begin{proposition}[Boundary distinctions under scoped bookkeeping]\label{prop:boundary-distinctions}
In each row of Table~\ref{tab:boundary-matrix}, isolate the bare partner
datum \(X_b\) listed in its last column, presented as a boundary record in
the sense of Definition~\ref{def:annotations} (carrying exactly the fields of that
datum, the typing fields of its kind, the datum and reference fields its certificates read, and reference fields to other members of its cluster). That datum alone lacks a record that the credited role's witness row requires a witness cluster to contain: \(W_i(C)\) is false for a cluster \(C\) formed only from that
datum, without importing fields of referenced records or adding
independently supplied credited-role records. Consequently a claim that the credited \(P_i\)
projection follows from that datum, without the further data required by
\(W_i\), is inadmissible (in the sense of Section~\ref{sec:admissibility-meta-theory}) under the alignments of
Remark~\ref{rem:repaired-alignments} together with the honest bookkeeping
required by a \FATCD{}. A cluster carrying independently supplied witnesses
for both roles may activate both channels; this is the permitted interaction,
not their identification.
\end{proposition}

The proof uses the domain \FATCD{}, the truth of honest audit conclusions and the certification condition for claims (Definitions~\ref{def:fatcd}, \ref{def:audited}(c3) and~\ref{prop:honest-bookkeeping}); the row-by-row verdicts after it use the offers of Definition~\ref{def:audited}(d) and are contrasted there with the statuses of Definition~\ref{def:channel-status}; a first reading can skip it. The idea of
the proof is short. Each forbidden collapse would credit one
role with data that only the other role's witness type supplies. The
borrowed content does not supply the credited role's witness, so the
attempted identification receives an inactive local verdict and is
inadmissible; independent data may still activate the description's
channel. The proof checks this
row by row; the three rows worked in full show the two shapes the
argument takes. The local verdicts \emph{no offer} and \emph{failed offer} below describe
two failures of a local attempt: respectively, no credited-role witness is
offered, or an offered cluster fails its witness check. They are not channel
statuses \(\Status{i}{D}\), which Definition~\ref{def:channel-status} assigns
from the whole description. In the general argument, \(P_i\) is the role the collapse
would credit and \(P_j\) the role whose content it borrows; in the
enumeration of the rows, each row is written \(P_i/P_j\) in the order of
Table~\ref{tab:boundary-matrix}, and each clause names the credited role,
which may be either member of the pair.


\begin{proof}
In short, each credited row of the witness table requires a record that the
row's bare datum lacks: an update map or a source relation record of declared use quotient for \(W_1\) (rows P1/P5, P1/P4, P1/P2, P1/P6); a
macro-specification with a realization relation for \(W_2\) (rows P2/P5,
P2/P4, P2/P6); two defined path or derivation records for \(W_3\) (rows
P1/P3, P2/P3, P3/P5, P3/P6); a law-bearing order, transition or refinement
record for \(W_4\) (rows P3/P4, P4/P5, P4/P6); and a quotient relation with a
packaging map for \(W_5\) (row P5/P6). Hence \(W_i(C)=\mathrm{false}\). A claim that the credited \(P_i\) projection
follows from the datum has statement \(W_i(C)\). By
Definition~\ref{def:audited}(c3), an outright conclusion of an honest audit that
uses no unrecorded predicate symbol and quantifies only over recorded carriers
is true in \(D\), and \(W_i(C)\) is such a formula; so no honest audit certifies
\(W_i(C)\), the certification condition of
Definition~\ref{prop:honest-bookkeeping} fails, and the claim is inadmissible.
This proves the proposition.
\end{proof}

\paragraph{Row-by-row local verdicts.} The local attempt verdicts below
are illustrative and are not used by later results. The enumeration and
table also identify each row's credited role and missing witness field,
used in Proposition~\ref{prop:typed-non-collapse}. A reader may continue
at Remark~\ref{rem:not-global-independence}
(p.~\pageref{rem:not-global-independence}). The verdict is not derived
from the row: in every row it is
failed offer if the attempt is recorded as an offer and no offer otherwise,
and only \(W_i(C)=\mathrm{false}\) is proved. No later result uses these
verdicts. Each row of Table~\ref{tab:boundary-matrix} names an
ordered pair \((P_i,P_j)\), a permitted interaction, and a forbidden
collapse. The forbidden collapse always has the same logical form: it would
license treating witness or structural data carried under the meaning of one
role as if it were, by itself, witness or structural data for the other role
--- that is, it would identify the two roles' admissible data types. The following
enumeration records, row by row, the credited role and the missing witness
field that the proof of Proposition~\ref{prop:boundary-distinctions}
summarizes, and maps each row's forbidden collapse to a common schema. The
enumeration discharges
all fifteen rows of Table~\ref{tab:boundary-matrix}, mapping each row's
forbidden collapse to the schema, with the role-specific data types read off
Definition~\ref{def:primitive-roles} and the per-primitive level content of
Remark~\ref{rem:per-primitive-levels}. Three rows are then worked in full as
illustrations.

\emph{Schema.} Fix a row and isolate the particular partner-role datum named
in its forbidden-collapse entry. Compare that datum with the fields required
by the credited role's witness row. If a required field is absent, the bare
datum cannot certify \(W_i\); hence the implication from that datum alone to
an active \(P_i\) projection is invalid. An offered cluster containing only the
bare datum may honestly record a failed \(W_i\) check. Adding the missing
\(P_i\) data may make \(W_i\) true and may activate both channels, without
identifying their meanings; this is the permitted interaction of the row.

\emph{Application to the rows.} Every entry in the forbidden-collapse column
has one of two shapes. In the \emph{carried form}, ``\(X\) alone is not
\(P_i\)'' or ``\(X\) is not \(P_i\)'', the datum \(X\) is content named by the
distinction-rule column as the meaning of the partner role; the credited-role
field absent from the bare datum is listed below. In the
\emph{non-implication form}, ``\(P_j\)-status \(S\) is not automatically
\(P_i\)-status \(S'\)'', the partner's status or certificate is not the
credited role's required certificate. These are checks of the stated bare
datum, not assertions that every \(P_j\)-bearing cluster fails \(W_i\). The
fifteen clauses record the missing fields.

\emph{Enumeration of all fifteen rows.} We list each row as
\(P_i/P_j\), name the credited channel, the borrowed content, and the
distinct witness type the credited channel requires but does not receive, so
that the schema applies with no residue. Throughout, the local attempt verdict of a clause describes the attempted
identification; it is not the channel status \(\Status{i}{D}\) that
Definition~\ref{def:channel-status} assigns from a whole description; a
description may carry such an attempt and, independently, an active witness
for the same role. In every clause the attempt may be recorded as an offer, with check result
\(W_i(C)=\mathrm{false}\) (verdict failed offer), or not recorded (verdict no
offer); the clauses below illustrate the first case in (3), (4), (8), (10),
(11) and (13) and the second in the others.
These stipulations concern the listed attempts, not the status
\(\Status{i}{D}\) of a full description; the two verdicts describe the
attempts as listed, and the statement uses only
\(W_i(C)=\mathrm{false}\), which holds in every clause.
\begin{enumerate}[topsep=2pt,itemsep=2pt,label=\textup{(\arabic*)}]
  \item \(\Pone{}/\Pfive{}\), ``packaging alone is not \Pone{}'': credits
    \Pone{} on the \Pfive{} packaging map record alone; the update or partial-map record and the
    certificate or related-pair witness required by \(W_1\) are absent; the displayed local attempt verdict for \Pone{}
    is no offer.
  \item \(\Pone{}/\Pfour{}\), ``stage index alone is not \Pone{}'': credits
    \Pone{} on \Pfour{} stage-index content; no descent or obstruction witness
    is thereby present; the displayed local attempt verdict for \Pone{} is no offer.
  \item \(\Pone{}/\Ptwo{}\), ``representability failure is not automatically
    descent obstruction'': credits \Pone{} on a \Ptwo{} non-representability
    status; a non-representability certificate is not a descended-map
    obstruction record and does not instantiate the \Pone{} obstruction
    witness; the displayed local attempt verdict for \Pone{} is failed offer.
  \item \(\Pone{}/\Pthree{}\), ``route undefinedness is not route mismatch'':
    credits \Pthree{} on the undefinedness of a route through a partial update
    map, where the map is undefined at a value the route passes through; a
    non-reducibility certificate between \emph{already-defined} routes is
    required, and an undefined route supplies no such certificate, so the candidate cluster does
    not instantiate the \Pthree{} witness schema; recorded as an offer, the
    attempt has local verdict failed offer.
  \item \(\Pone{}/\Psix{}\), ``audit record is not descent obstruction'':
    credits \Pone{} on a \Psix{} audit record; an audit record is not a
    descended-map or obstruction witness; the displayed local attempt verdict for \Pone{} is
    no offer.
  \item \(\Ptwo{}/\Pfive{}\), ``packaging alone is not \Ptwo{}'':
    credits \Ptwo{} on \Pfive{} packaging content alone. The packaging map
    supplies neither the realization relation nor its finite existence
    or nonexistence certificate; in the satisfying-realizer branch it also
    supplies no required comparison. The local attempt verdict for \Ptwo{}
    is no offer. A separately witnessed \Ptwo{} projection may still
    accompany the packaging.
  \item \(\Ptwo{}/\Pfour{}\), ``stage variation alone is not \Ptwo{}'':
    credits \Ptwo{} on \Pfour{} stage content; a representability certificate
    with realization relation is required, and stage variation supplies none, so
    the local attempt verdict for \Ptwo{} is no offer.
  \item \(\Ptwo{}/\Pthree{}\), ``non-representability is not route mismatch'':
    credits \Pthree{} on a \Ptwo{} non-representability status; a
    non-representability record is not a non-reducibility certificate between
    defined routes and does not instantiate the \Pthree{} witness schema, so the
    local attempt verdict for \Pthree{} is failed offer.
  \item \(\Ptwo{}/\Psix{}\), ``audit trail is not representability'': credits
    \Ptwo{} on a \Psix{} audit trail; an audit record is not a representability
    certificate; the displayed local attempt verdict for \Ptwo{} is no offer.
  \item \(\Pthree{}/\Pfour{}\), ``stage-indexed route mismatch is not \Pfour{}
    itself'': credits \Pfour{} on the \Pthree{} route-disagreement certificate
    alone, without the stage law that indexes it; a non-reducibility
    certificate carries no order, transition, or
    refinement structure and does not instantiate the \Pfour{} witness schema;
    if the certificate is carried in the fields of an order, transition or
    refinement record, that record still carries no law, since its only
    certificate expresses non-reducibility of two routes, so \(W_4\) fails and
    the local attempt verdict for \Pfour{} is failed offer.
  \item \(\Pthree{}/\Pfive{}\), ``quotient object is not route mismatch'':
    credits \Pthree{} on a \Pfive{} quotient object; a quotient object is not a
    non-reducibility certificate between defined routes and does not instantiate
    the \Pthree{} witness schema; the displayed local attempt verdict for \Pthree{} is
    failed offer.
  \item \(\Pthree{}/\Psix{}\), ``audit record is not route mismatch'': credits
    \Pthree{} on a \Psix{} audit record; an audit record is not a
    non-reducibility certificate; the displayed local attempt verdict for \Pthree{} is
    no offer.
  \item \(\Pfour{}/\Pfive{}\), ``packaging is not stage structure'': credits
    \Pfour{} on the \Pfive{} packaging map alone; a packaging map is not an
    order/transition/refinement certificate and does not instantiate the
    \Pfour{} witness schema; the displayed local attempt verdict for \Pfour{} is failed offer.
  \item \(\Pfour{}/\Psix{}\), ``audit record is not staging'': credits \Pfour{}
    on a \Psix{} audit record of staged histories; an audit record is not an
    order/transition/refinement certificate; the displayed local attempt verdict for \Pfour{} is
    no offer.
  \item \(\Pfive{}/\Psix{}\), ``audit record is not packaging map'': credits
    \Pfive{} on a \Psix{} audit record of packaging events; an audit record is
    not a quotient object with a packaging map; the displayed local attempt verdict for \Pfive{} is
    no offer.
\end{enumerate}
Clauses (1)--(15) exhaust the rows of Table~\ref{tab:boundary-matrix}: each
names the credited channel, the borrowed partner-role content, and the
distinct witness type the credited channel requires and does not receive, and
each clause shows that the named bare datum lacks a field required by the
credited witness; the displayed local verdict describes an absent or failed
offer, not a status of the full description.

We now work three of these rows in full; they illustrate the carried and
non-implication shapes.

\emph{Row \(\Pone{}/\Pfive{}\) (packaging alone is not \(\Pone{}\)), clause
(1).} By Definition~\ref{def:primitive-roles}, \Pfive{} is quotient packaging
and \Pone{} is update-vs-quotient compatibility, i.e.\ descent. By the
\(W_5\) row of the witness table, a \Pfive{}-witness is a quotient object
with a packaging map and quotient-validity evidence, whereas a
\Pone{}-witness (the \(W_1\) row) is an update map or partial map with source
and target relation data and either its compatibility certificate or a related
pair witnessing failed domain saturation or separated defined images. Possession of the packaging
map is precisely the data that says \emph{what} the quotient is; it says nothing
about \emph{whether} a selected update descends through it. In terms of the
\emph{Witness types} table, \(W_1\) requires a map record whose declared use is
\emph{update}, and the \Pfive{} cluster of the packaging datum contains only
the packaging map, whose declared use is \emph{packaging}; so \(W_1\) fails on
that cluster. The collapse
``packaging alone is \Pone{}'' would credit the \Pone{} channel on the strength
of the packaging map alone, but no descended-map witness and no obstruction
record is thereby present, so by the schema the verdict on this attempted
identification is no offer; an independently witnessed \Pone{} channel may
still be active. The descent question has not been answered by the packaging
map. The
permitted interaction --- \Pone{} using \Pfive{} when checking descent or
obstruction --- is untouched, since there \Pone{} consumes the quotient
packaging as the object \emph{over which} a descended map or obstruction is then
exhibited, and the descent witness remains the \Pone{} channel's own content.

\emph{Row \(\Pthree{}/\Pfour{}\) (stage-indexed route mismatch is not
\(\Pfour{}\) itself), clause (10).} By Definition~\ref{def:primitive-roles},
\Pfour{} is stage/order/transition/refinement structure and \Pthree{} is
enriched route/coherence mismatch; by Remark~\ref{rem:repaired-alignments}
generic transition data is \emph{not} \Pthree{} and, symmetrically, a
route-mismatch witness is not \Pfour{}. A \Pfour{}-witness is an
order/transition/refinement certificate or a stage architecture; a
\Pthree{}-witness is a non-reducibility certificate exhibiting disagreement
between already-defined routes. When route mismatch is indexed by stages, the
stage index is \Pfour{}-content and the route disagreement at each index is
\Pthree{}-content; these are recorded as data of different roles. The collapse
``stage-indexed route mismatch is \Pfour{}'' would credit the \Pfour{} channel
on the strength of the route-disagreement certificate considered without its
stage-law record, but a
non-reducibility certificate is not an order/transition/refinement certificate
--- it carries no stage, order, or refinement structure of its own --- so by the
schema the local attempt verdict for \Pfour{} is failed offer: the candidate
cluster does not instantiate the required \Pfour{} witness schema. The permitted
interaction --- route systems varying over valid stages --- is untouched, since
there the \Pfour{} stage structure is supplied independently and the \Pthree{}
mismatch is merely indexed by it, each channel keeping its own witness.

\emph{Row \(\Ptwo{}/\Pfive{}\), clause (6).} \Ptwo{} concerns a
macro-specification and its realization relation; \Pfive{} concerns
quotient packaging. In the satisfying-realizer branch of \(W_2\), the recorded
comparison must show that some class of the kernel of the declared packaging
map meets both the realizer set and its complement in the candidate class. In the no-realizer branch, recorded satisfaction data, a
finite certificate that no candidate satisfies the specification, and a
recorded true evaluation at a declared same-kind record outside the class,
and a candidate class that is the full declared domain of a declared
packaging or update map, suffice; no quotient comparison is required. A packaging map alone
supplies neither branch's realization relation and certificate. Thus the
proposed credit of \Ptwo{} from \Pfive{} packaging alone makes no \(W_2\)
offer; the local attempt verdict for \Ptwo{} is no offer. A
\Ptwo{} witness may still be carried alongside \Pfive{} quotient data when its
own required records and attachments are present.

The enumeration (1)--(15) covers every row of
Table~\ref{tab:boundary-matrix}, and the three worked rows confirm both shapes
in detail. The argument is local to each listed pair under the covered scope and
the audit condition of Definition~\ref{def:audited}; it establishes pairwise
non-collapse, not global independence of the six roles (cf.\
Remark~\ref{rem:not-global-independence}).

The clauses of the proof follow the rows of Table~\ref{tab:boundary-matrix};
the credited role of each, and the verdict its clause illustrates, are as follows; in every row the verdict is \emph{failed offer} if the attempt is recorded as an offer and \emph{no offer} otherwise, and only \(W_i(C)=\mathrm{false}\) is proved:
\begin{center}
\footnotesize
\begin{tabular}{@{}llll@{\qquad}llll@{}}
\toprule
& pair & credited & local verdict & & pair & credited & local verdict \\
\midrule
(1) & \Pone{}/\Pfive{} & \Pone{} & no offer & (9) & \Ptwo{}/\Psix{} & \Ptwo{} & no offer \\
(2) & \Pone{}/\Pfour{} & \Pone{} & no offer & (10) & \Pthree{}/\Pfour{} & \Pfour{} & failed offer \\
(3) & \Pone{}/\Ptwo{} & \Pone{} & failed offer & (11) & \Pthree{}/\Pfive{} & \Pthree{} & failed offer \\
(4) & \Pone{}/\Pthree{} & \Pthree{} & failed offer & (12) & \Pthree{}/\Psix{} & \Pthree{} & no offer \\
(5) & \Pone{}/\Psix{} & \Pone{} & no offer & (13) & \Pfour{}/\Pfive{} & \Pfour{} & failed offer \\
(6) & \Ptwo{}/\Pfive{} & \Ptwo{} & no offer & (14) & \Pfour{}/\Psix{} & \Pfour{} & no offer \\
(7) & \Ptwo{}/\Pfour{} & \Ptwo{} & no offer & (15) & \Pfive{}/\Psix{} & \Pfive{} & no offer \\
(8) & \Ptwo{}/\Pthree{} & \Pthree{} & failed offer & & & & \\
\bottomrule
\end{tabular}
\end{center}

\begin{remark}[Not a global independence theorem]\label{rem:not-global-independence}
The boundary matrix asserts local distinctions, not global
independence. Primitives interact legitimately in a single
description: \Pone{} descent is audited by \Psix{}; \Pthree{} route
terms may be indexed by \Pfour{} stages; in the satisfying branch a
\Ptwo{} witness is carried over \Pfive{} packaging data.
No claim is made that the six roles are independent generators of a
global algebra. The typed-non-collapse counterpart of this observation
appears in Section~\ref{sec:admissibility-meta-theory}.
\end{remark}

\subsection{Selected-lift atlas and global nonclaims}\label{subsec:lift-atlas}

The remaining architectural data are the selected lifts between levels
and the nonclaims that accompany them. Upgrading from behavioral or
verification content to the structural-categorical level requires
explicitly recorded added data. The resulting lifts are selected, not
unique, and the atlas carries structural content that is lost under
forgetting.

\begin{definition}[Selected-lift atlas]\label{def:selected-lift-atlas}
The \emph{selected-lift atlas} for the six roles records, for each
\Pone{}, \ldots, \Psix{}:
\begin{itemize}[topsep=2pt,itemsep=1pt]
  \item \Pone{}: proof or witness together with quotient structure, carrying
    either a descended map or an explicit no-map obstruction;
  \item \Ptwo{}: certificates of unrealizability or of failure of
    representability over the declared packaging kernel, together with the macro-specification layer and realization relation;
  \item \Pthree{}: the selected lift data,
    carrying generated route terms, non-reducibility certificates, and
    categorical route grammar or a reduced categorical universe;
  \item \Pfour{}: validity certificates together with structural
    stage/refinement architecture;
  \item \Pfive{}: quotient-validity certificates together with the
    structural quotient data, a packaging map, and universal-property or
    tower/coarsening data;
  \item \Psix{}: record-event or check evidence together with the audit
    schema and its provenance and replay/check relations.
\end{itemize}
Each lift is selected, and not unique in general; where uniqueness holds, as
for the descended map \(\bar F([x])=[F(x)]\) of a descending update, it is
recorded with its proof. Forgetting from a higher level
loses structural content that is not recoverable without re-adding data,
unless, as for the descended map, the higher datum is determined and this is
recorded with its proof. Formally, the selected-lift atlas of a description
\(D\) consists of those lift annotations of \(D\) that satisfy
Definition~\ref{prop:forgetting-lifts}, are attached to a cluster \(C\) such that
some witness-schema annotation targeting exactly \(C\) names \(P_i\), and carry
the corresponding added data listed above; such a lift is indexed by each such
\(P_i\). Other admissible lifts need
not belong to this role-indexed atlas.
\end{definition}

\begin{remark}[Global nonclaims for the role architecture]\label{rem:atlas-global-nonclaims}
The role architecture preserves the following nonclaims, all carried
forward into the admissibility regime of
Section~\ref{sec:admissibility-meta-theory} and the theorem chain of
Sections~\ref{sec:lower-bounds}--\ref{sec:decomposition-exhaustion}:
\begin{itemize}[topsep=2pt,itemsep=1pt]
  \item no full six-primitives algebra is claimed;
  \item no global primitive-independence theorem is claimed;
  \item no empirical realization is claimed;
  \item no uniqueness of primitive presentations is claimed;
  \item no uniqueness of selected lifts is claimed;
  \item no lower-level-to-higher-level determinacy is claimed;
  \item no instrument-free formulation is claimed;
  \item direct cross-level comparison remains inadmissible without
    translation.
\end{itemize}
\end{remark} \section{Admissibility and meta-theory}\label{sec:admissibility-meta-theory}


This section states when a claim about a description counts. In
Definitions~\ref{prop:honest-bookkeeping}--\ref{prop:empirical-bridge} and
Propositions~\ref{prop:typed-non-collapse} and~\ref{prop:nonclaim-closure}, a
claim is a claim annotation of some \(D\in\FATCD{}\), and ``inadmissible''
means that condition (a) or (b) of Definition~\ref{prop:honest-bookkeeping}
fails, or that a record required by Definition~\ref{prop:level-profile},
\ref{prop:forgetting-lifts}, \ref{prop:visibility} or~\ref{prop:empirical-bridge}
is missing; in particular a claim other than an offer that no honest audit annotation of
\(D\) certifies is inadmissible; an offer satisfies condition (a) of Definition~\ref{prop:honest-bookkeeping} when an honest audit
records the value of \(W_i(C)\), even when that value is false. The central
rule is honest bookkeeping (Definition~\ref{prop:honest-bookkeeping}):
every claim, and every change of level, translation, visibility
assumption or empirical bridge it relies on, must be explicitly
recorded. Together with the role meanings and boundary distinctions of
Section~\ref{sec:primitive-role-architecture}, that rule governs the
claims developed below.
\begin{itemize}[topsep=2pt,itemsep=1pt]
  \item The six roles are kept apart as types (typed non-collapse); the
    paper does not claim that they are independent generators.
  \item A claim moves between levels only through recorded forgetting,
    which loses data, and recorded selected lifts, which add it.
  \item A role counts as active only with the required threshold,
    witness, level, collapse, nonclaim and audit data.
  \item An instrument decides what is visible and what is suppressed.
  \item An empirical claim enters only through a recorded bridge.
\end{itemize}
The section uses the domain \FATCD{}, the audit part \(\Audit{D}\) and
threshold attachment, which are defined in
Section~\ref{sec:fatcd-domain}; a reader checking proofs can follow the reading order at the end of
Section~\ref{subsec:organization}, keeping Definition~\ref{def:audited} at
hand while reading Definition~\ref{prop:honest-bookkeeping}. Several claim-level rules
below use records and annotations defined in
Section~\ref{sec:fatcd-domain}. The main correspondences are:
\begin{itemize}[topsep=2pt,itemsep=1pt]
  \item Definition~\ref{prop:honest-bookkeeping} and
    Definition~\ref{def:audited}: Definition~\ref{def:audited} requires every
    claim of a description to be referred to by an honest audit;
    Definition~\ref{prop:honest-bookkeeping} further requires that audit to
    certify the claim's recorded statement (for an offer, the value of
    \(W_i(C)\));
  \item Remark~\ref{prop:activation-thresholds} and
    Definition~\ref{def:threshold-attachment}: when a role reading of a
    cluster is licensed;
  \item Definition~\ref{prop:forgetting-lifts} and the lift and forgetting
    annotations of Definition~\ref{def:annotations}: how a claim changes
    level;
  \item Definition~\ref{prop:visibility} and the visibility annotations:
    what an instrument shows and suppresses;
  \item Definition~\ref{prop:empirical-bridge} and the empirical-bridge
    annotations: how an empirical claim enters.
\end{itemize} The rules keep the global
nonclaims of Remark~\ref{rem:atlas-global-nonclaims} in force. Of this section, only
Definitions~\ref{prop:honest-bookkeeping}, \ref{prop:level-profile},
\ref{prop:forgetting-lifts}, \ref{prop:visibility}
and~\ref{prop:empirical-bridge} and Proposition~\ref{prop:nonclaim-closure}
are cited in Sections~\ref{sec:fatcd-domain}--\ref{sec:scoped-exact-six}, the
last only in Section~\ref{sec:integrated-stress}. The other results of this section
fix how claims are stated; several return in
Sections~\ref{sec:discussion}--\ref{sec:future-work}.

\subsection{Honest bookkeeping as an admissibility regime}\label{subsec:honest-bookkeeping}

The audit condition of Definition~\ref{def:audited} records enough
data for each claim, promotion, translation, lift, forgetting,
visibility, and empirical-bridge assertion to be interrogated on its
own terms. The following definition stipulates that these record
requirements jointly constitute the admissibility regime, in a spirit
analogous to the axiomatic specification discipline of Hoare
logic~\citep{Hoare1969}.

\begin{definition}[Admissible claim under honest bookkeeping]\label{prop:honest-bookkeeping}
Honest bookkeeping is the admissibility regime for the claim annotations
carried by a description \(D \in \FATCD{}\). A claim annotation of \(D\) is
\emph{admissible} if and only if (a) some honest audit annotation of
\(D\) refers to it and certifies the formula recorded as its statement (for an offer, \(W_i(C)\) or
\(\neg W_i(C)\) according to the recorded result, Definition~\ref{def:audited}),
and (b) every assertion below has an explicit record carried in \(\Audit{D}\)
in the sense of Definition~\ref{def:annotations}:
\begin{itemize}[topsep=2pt,itemsep=1pt]
  \item the claim itself, with its statement, hypotheses, and scope;
  \item every promotion, with its source level, target level, added data,
    losses, and preserved nonclaims;
  \item every translation, with its source context, target context,
    preserved data, suppressed data, and claim strength;
  \item every selected lift, marked selected and non-unique unless uniqueness for its specified source data and target class is proved and recorded;
  \item every forgetting map, with its record of losses, empty only when
    losslessness is proved and recorded;
  \item every visibility claim, with both its visible and its suppressed
    data;
  \item every empirical-bridge assertion, with its substrate, observable,
    instrument, level map, threshold-witness relation, and
    suppressed-structure nonclaims;
  \item every cross-level claim, with a level annotation targeting it and
    translation annotations for each pair of levels recorded on the records it
    names, as in Definition~\ref{prop:level-profile}.
\end{itemize}
Accordingly, unrecorded claims, implicit promotions, silent translations,
asserted-but-unselected lifts, lossless forgetting without proof,
visibility overreads, and bare empirical-realization assertions are all
inadmissible.
\end{definition}

\begin{remark}[Why the regime is necessary]\label{rem:why-honest-bookkeeping}
The bookkeeping requirement is a structural check, not a stylistic
one. Without explicit records, passive observation can be misread as
\Psix{}, mere packaging as \Pone{} descent, and route failure as
\Pthree{} mismatch. The regime prevents those misreadings from
entering the theorem chain.
\end{remark}


\subsection{Typed non-collapse}\label{subsec:typed-non-collapse}

The regime does not include flat primitive independence --- the claim
that \Pone{}, \ldots, \Psix{} are independent generators of a global
algebra. Its meta-theoretic counterpart here is typed non-collapse.
Proposition~\ref{prop:boundary-distinctions} excludes each forbidden collapse
of Table~\ref{tab:boundary-matrix}; Proposition~\ref{prop:typed-non-collapse}
states the same exclusion as a rule of the admissibility regime of
Definition~\ref{prop:honest-bookkeeping} and adds that each row's permitted
interaction is not thereby excluded. Each row is checked in one direction, for the role it credits (listed in the
table after the proof of Proposition~\ref{prop:boundary-distinctions}); the
row label names the pair without order.

\begin{proposition}[Typed non-collapse]\label{prop:typed-non-collapse}
Under honest bookkeeping (Definition~\ref{prop:honest-bookkeeping}) and
the alignments of Remark~\ref{rem:repaired-alignments}, consider each
forbidden-collapse row of Table~\ref{tab:boundary-matrix}. Write \(P_c\)
for the role that row would credit and \(X_b\) for the bare partner datum
it borrows, as listed in the last column of Table~\ref{tab:boundary-matrix}
and presented as a boundary record in the sense of
Definition~\ref{def:annotations}.
For every \(D\in\FATCD{}\), no honest audit annotation of \(D\)
certifies \(W_c(C)\) for a cluster \(C\) formed only from \(X_b\); the ``permitted interaction'' of the same row is
not excluded by this proposition. More generally, for each \(k\le6\), no
honest audit annotation of \(D\) certifies \(W_k(C)\) for a cluster \(C\)
lacking any record that the \(W_k\) row requires a witness cluster to contain (among them, for \(k=1\): a map or partial-map record of declared use update, or a relation record of declared use quotient; \(k=2\): a macro-specification predicate record or a realization relation; \(k=3\): two defined path or derivation records; \(k=4\): a law-bearing stage, order, transition or refinement record; \(k=5\): a relation record of declared use quotient, a packaging map record or a check record; \(k=6\): an event record, a trace record or a replay or check record); in particular each row's exclusion holds with its two
roles exchanged whenever the exchanged datum lacks that record.
\end{proposition}

\begin{proof}
Fix a forbidden-collapse row and its credited role \(P_c\) and borrowed datum \(X_b\). The row records a ``forbidden collapse'': an
identification asserting that the content of its two roles is the same
closure role, so that one
projection may be read off the other without further data. To be
admissible under Definition~\ref{prop:honest-bookkeeping}, such an
identification would have to enter the theorem chain as a claim with a
recorded statement and scope. The role meanings are distinct, and
Proposition~\ref{prop:boundary-distinctions} checks, for each
forbidden-collapse row, a field required by the credited witness that the
named partner datum alone does not supply. An audit certifying the credited
projection solely from that datum would therefore certify a conclusion
unsupported by its recorded check and would not be honest under
Definition~\ref{def:audited}: by that definition the check-language conclusion
\(W_c(C)\) of an honest audit would be true in \(D\), whereas \(W_c(C)\) is false by
Proposition~\ref{prop:boundary-distinctions}. Independently supplied data may satisfy both
witness conditions; the proposition excludes the claimed automatic
identification, not that joint occurrence. Hence the collapse is
inadmissible. More generally, if \(C\) lacks a record required by the \(W_k\) row, that row gives \(W_k(C)=\mathrm{false}\). Definition~\ref{def:audited}(c3) therefore excludes an honest audit certifying \(W_k(C)\); this also proves the stated exchanged-role exclusions.

It remains to check that ruling out the collapse does not also rule out
the row's ``permitted interaction''. The permitted interaction does not
identify \(P_i\) with \(P_j\); it records a legitimate joint occurrence
of the two distinct roles in a single description --- for instance,
\Pone{} descent audited by \Psix{}, or \Pthree{} route terms indexed by
\Pfour{} stages, as in Remark~\ref{rem:not-global-independence}. Such
an interaction carries its own honest-bookkeeping record for each role
separately and asserts no identification. Nothing in the argument above
touches it: the inadmissibility derived for the collapse rests on the unsupported
inference from the bare partner datum to the credited witness, which the
permitted interaction never makes. Therefore the permitted interaction remains admissible
whenever its own bookkeeping conditions are met.
\end{proof}

\begin{remark}[Typed non-collapse versus global independence]\label{rem:typed-non-collapse-vs-independence}
Typed non-collapse is not a global primitive-independence statement. It does not assert that the six roles are logically or
categorically independent in a global algebra; it asserts only that
the pairwise identifications listed as forbidden collapses in
Table~\ref{tab:boundary-matrix} are inadmissible under honest
bookkeeping, leaving the corresponding typed interactions available.
\end{remark}


\subsection{Level-profile architecture}\label{subsec:level-profile-architecture}

Typed non-collapse fixes which identifications are inadmissible; cross-level
claims also require recorded levels and translations. The level trichotomy of
Definition~\ref{def:level-trichotomy} therefore supplies the profile for such
claims.

\begin{definition}[Level profile of an admissible claim]\label{prop:level-profile}
A claim annotation is \emph{cross-level} when records it names are listed in the target fields of level annotations recording different levels; it is a \emph{direct comparison} when
its statement contains an equality or order atom, or an implication, whose two sides read records listed under different recorded levels.
Its \emph{level profile} is a level annotation targeting the claim, together
with, for each such pair of levels, a translation annotation carried in
\(\Audit{D}\) whose source and target contexts are those levels.
Every admissible cross-level claim about a description \(D \in \FATCD{}\) carries a
\emph{level profile}: it locates the claim on the behavioral,
verification, or structural-categorical level, and each cross-level
assertion is accompanied by a translation record in the sense of
Definition~\ref{prop:honest-bookkeeping}. Direct comparison of claims at
different levels without a translation record is inadmissible.
\end{definition}

\subsection{Forgetting and selected lifts}\label{subsec:forgetting-lifts}

Level transitions are controlled by forgetting and selected lifts.

\begin{definition}[Admissible forgetting and selected lifts]\label{prop:forgetting-lifts}
Forgetting from a higher to a lower level (structural-categorical to
verification, verification to behavioral, or structural-categorical directly to
behavioral) is admissible if and only if it carries a record of
its losses carried in \(\Audit{D}\); the loss record may be empty only when
losslessness is proved and recorded. Lifting from a lower level to a higher level is
admissible only when accompanied by
\begin{itemize}[topsep=2pt,itemsep=1pt]
  \item the added data required to support the higher-level statement;
  \item the mark that the lift is selected and non-unique unless uniqueness for its specified source data and target class is proved and recorded;
  \item explicit preservation of lower-level invariants and explicit
    losses where the lift is partial.
\end{itemize}
Lossless forgetting and unique lifts are inadmissible unless the losslessness
or uniqueness is itself proved and recorded.
\end{definition}

\begin{remark}[Lower-to-higher non-determinacy]\label{rem:lower-to-higher-non-determinacy}
Definition~\ref{prop:forgetting-lifts} formalizes the non-determinacy of
lower-level data: given behavioral or verification content alone, the admissibility regime
forces no choice of structural-categorical content, and any asserted choice
must be marked as a selected lift. Where the lower-level data do determine the
higher-level datum, as the descended map of a descending update is determined,
that determinacy is admissible with its recorded proof.
\end{remark}

\subsection{Activation thresholds}\label{subsec:activation-thresholds}

The regime must also state when a projected role is active. The
attachment protocol of Definition~\ref{def:threshold-attachment} is
lifted to a regime-level principle.

\begin{remark}[Activation thresholds]\label{prop:activation-thresholds}
At the regime level, the threshold-attachment discipline of
Definition~\ref{def:threshold-attachment} reads as follows. An
admissible active derived role projection \(\pi_i(D)\) exists only when
\(D\) carries the threshold, witness, level, collapse, nonclaim and audit data
required by Definition~\ref{def:threshold-attachment}. A channel with no
admissible projection receives \(\collapsedbc\) if a failed offer is
recorded, and otherwise \(\belowthr\) if a covered seed is present
(Definition~\ref{def:channel-status}); these are statuses of channels, not of
clusters, and a channel with an admissible projection is \(\activeproj\)
whatever other clusters do. Arbitrary enrichment, adding structure
solely to force a role projection, is inadmissible unless the additions
are themselves honestly bookkept as raw features or annotations. This
restates the attachment protocol at regime level: the sentences before the
enrichment sentence restate Definitions~\ref{def:threshold-attachment}
and~\ref{def:channel-status}, and the enrichment sentence
follows from Definition~\ref{def:fatcd}, since every added record of a member of
\FATCD{} is subject to Definitions~\ref{def:finite}--\ref{def:audited}.
\end{remark}

\subsection{Instrument-relative visibility and suppression}\label{subsec:visibility-suppression}

Threshold-governed activation still leaves open what an instrument makes
visible or suppresses. Admissibility of visibility claims is therefore
instrument-relative.

\begin{definition}[Admissible visibility record]\label{prop:visibility}
Every admissible visibility claim carries an explicit \emph{visibility record}
comprising its visible content, its suppressed content, the instruments
under which the visibility and suppression are asserted, and the
resulting instrument-relative claim strength. Overreading a suppressed
claim as visible under a different instrument is inadmissible. Partial
visibility, in which a claim is visible only up to some coarsening or
translation, is admitted when the visibility record itself makes the
partiality explicit. An \emph{instrument} is a declared record of the
description named in an instrument-visibility annotation's instruments field
or as the measuring instrument of an empirical-bridge claim carried in
\(\Audit{D}\); the \emph{claim strength} is a value of a declared finite ordered
set recorded in that annotation; a visibility claim \emph{asserts
instrument-free truth} when it asserts visible content without naming an
instrument. The \emph{instrument family} of a description is the set of its
instrument records.
\end{definition}

\begin{remark}[No instrument-free truth]\label{rem:no-instrument-free}
Definition~\ref{prop:visibility} does not claim that instruments are
free of distortion or that any single instrument is canonical. It asserts only that admissible visibility claims carry an instrument-relative
visibility record. Instrument completeness is deferred to
Section~\ref{sec:future-work}.
\end{remark}

\subsection{Empirical-bridge admissibility}\label{subsec:empirical-bridge}

Any claim crossing from the closure-algebraic side of the framework
to a substrate-level observable must pass through the
empirical-bridge discipline, related to but narrower than
the empirical-adequacy view of theories~\citep{vanFraassen1980}.

\begin{definition}[Admissible empirical-bridge claim]\label{prop:empirical-bridge}
An empirical-bridge claim about \(D \in \FATCD{}\) is admissible only when
it has, carried in \(\Audit{D}\),
\begin{itemize}[topsep=2pt,itemsep=1pt]
  \item the substrate and observable to which the bridge relates;
  \item the instrument under which the observable is measured;
  \item the level map between the closure-algebraic content and the
    substrate-level content;
  \item a threshold-witness relation recording when the bridge is
    crossed;
  \item the nonclaims suppressed by the bridge record, in particular any
    structural content of \(D\) that is not tracked by the substrate or
    instrument.
\end{itemize}
Empirical realization --- the claim that \(D\) is itself a substrate
observation --- is stronger than empirical-bridge admissibility and is
not asserted by any empirical-bridge claim.
\end{definition}


\subsection{Nonclaim closure across the admissibility regime}\label{subsec:nonclaim-closure}

The rules of this section, from honest bookkeeping
(Definition~\ref{prop:honest-bookkeeping}) to admissible empirical-bridge
claims (Definition~\ref{prop:empirical-bridge}), keep the global
nonclaims of Remark~\ref{rem:atlas-global-nonclaims} in force. The next
proposition lists the overclaims they block.

\begin{proposition}[Admissibility-regime nonclaim closure]\label{prop:nonclaim-closure}
Under the admissibility regime, each of the following claims is
inadmissible; the second kind of claim in the first item cannot occur, since
no member of \FATCD{} carries it:
\begin{itemize}[topsep=2pt,itemsep=1pt]
  \item a claim identifying two of the six roles along a forbidden collapse
    of Table~\ref{tab:boundary-matrix}, or a claim annotation whose statement applies
    a recorded finite map with domain the role names \(P_1,\dots,P_6\) (or
    tuples of them), or an independence predicate on them, whose referenced
    records are only boundary and witness records, and whose statement is not
    represented by their fields in the sense of
    Definition~\ref{def:annotations} (by Definition~\ref{def:audited}(e));
  \item an unconditional empirical-realization claim, that \(D\) is itself a
    substrate observation, where \(D\) records no extension of that
    predicate (criterion after Definition~\ref{def:residual-scheme}), whether
    or not an empirical-bridge record is present;
  \item a lift asserted to be unique without a recorded proof of uniqueness;
  \item a claim that lower-level content determines higher-level content,
    without a recorded proof of that determinacy;
  \item a visibility claim asserting instrument-free truth;
  \item a direct cross-level comparison without a translation record;
  \item a forgetting asserted to be lossless without a recorded proof of
    losslessness;
  \item an empirical-bridge claim lacking any of: a threshold-witness
    relation, a level map, the instrument under which the observable is
    measured, or suppressed-structure nonclaims.
\end{itemize}
\end{proposition}

\begin{proof}
Each requirement of the regime blocks a specific class of overclaims;
we check that it blocks the listed claims.

Honest bookkeeping (Definition~\ref{prop:honest-bookkeeping}) declares
inadmissible every unrecorded claim, implicit promotion, silent
translation, asserted-but-unselected lift, lossless forgetting without
proof, visibility overread, and bare empirical-realization assertion;
these are requirements of Definition~\ref{prop:honest-bookkeeping},
beyond the description-level audit condition of Definition~\ref{def:audited}.

Typed non-collapse (Proposition~\ref{prop:typed-non-collapse}) blocks the
forbidden identifications of Table~\ref{tab:boundary-matrix}. A claim
annotation asserting a full algebra (a map applied to role names) or a global
independence theorem (an independence predicate on role names) must represent
that substantive assertion in fields of its referenced records
(Definition~\ref{def:annotations}). Boundary and \(W_i\)-witness records carry
exactly the fields of their declared checks, and those checks supply neither a
map on role names nor a global-independence certificate. An assertion not represented by the fields of the cited boundary and witness
records requires an additional declared content record. A claim citing only
those records and failing the representation condition cannot be carried by a member of
\FATCD{}. This does not prove that no full
algebra or global independence theorem can exist. Both broader statements remain nonclaims of
Remark~\ref{rem:atlas-global-nonclaims}.

Level-profile architecture (Definition~\ref{prop:level-profile}) blocks
direct cross-level comparison without a translation record, which is
exactly the non-claim that cross-level comparison remains inadmissible
without translation.

Admissible forgetting and selected lifts
(Definition~\ref{prop:forgetting-lifts}) block lossless forgetting,
unique lifts, and lower-to-higher determinacy: forgetting is lossy
unless losslessness is proved and recorded, no lift is admitted as
unique unless uniqueness is proved and recorded, and, by
Remark~\ref{rem:lower-to-higher-non-determinacy}, no structural-categorical
content is asserted to be forced by behavioral or verification content without
a recorded proof. These are
the non-claims that forgetting remains lossy, that no unique-lift
theorem is claimed, and that no lower-to-higher determinacy is claimed.

Activation thresholds (Remark~\ref{prop:activation-thresholds}) block
arbitrary enrichment and threshold-free nontriviality claims, since an
active projection is admissible only with the threshold, witness, level,
collapse, nonclaim and audit data of Definition~\ref{def:threshold-attachment}.

Instrument-relative visibility (Definition~\ref{prop:visibility}) blocks
instrument-free visibility claims and visible-to-suppressed overreads; it does
not impose instrument indexing on every claim kind.

Empirical-bridge admissibility (Definition~\ref{prop:empirical-bridge})
blocks empirical-realization overclaims and any bridge claim lacking a
threshold-witness relation, a level map, the instrument under which the
observable is measured, or suppressed-structure nonclaims. An
empirical-realization claim asserts a predicate, being a substrate
observation, whose extension \(D\) does not record; its truth does not follow
from any finite check of \(D\), so by Definition~\ref{def:audited} no honest
audit certifies it outright, and it is inadmissible whether or not a bridge
record is present. An unrecorded realization predicate may occur as a
hypothesis of an audited conditional consequence; no audit certifies
realization outright or as a consequent. This yields the non-claim that no empirical realization
is claimed and that empirical-bridge claims carry the required data.

Collecting these residues gives each of the listed non-claims. Those that
restate a global nonclaim of Remark~\ref{rem:atlas-global-nonclaims} are
thereby carried forward into the regime; the two further items --- that
forgetting remains lossy absent a recorded proof of losslessness, and that
empirical-bridge claims carry their threshold-witness relation, level map,
instrument and suppressed-structure nonclaims --- are additional disciplines the regime imposes
at this level.
\end{proof}
 \section{The FATCD domain and the exact-six question}\label{sec:fatcd-domain}


This section defines the class of descriptions over which the
exact-six theorem is stated. A \emph{finite audited typed closure
description}, abbreviated \FATCD{}, is a finitely presented collection
of typed raw records and annotations that is finite, closed under its
own references, and audited in the sense of Definition~\ref{def:audited}; the setting is that of finite
model theory~\citep{Libkin2004}. Roles are not part of the definition.
A role enters only when a threshold annotation is attached to a cluster
of raw features, and the role projections \(\pi_1(D),\dots,\pi_6(D)\)
are readings of the raw material, not conditions for membership in the
domain. Every content record and structural annotation is a
residual feature in the sense of Definition~\ref{def:residual-scheme} and
is licensed some of the eleven labels of Definition~\ref{def:residual-scheme}; the unresolved label (x) is licensed exactly when none of labels (i)--(ix) is; the features
that no channel of a chosen decomposition record accounts for are its record
residuals. So a seventh role is not ruled out by definition. The section closes with
non-circularity and non-overbreadth schemas and formulates the
exact-six question precisely over \FATCD{}: whether every
\(D \in \FATCD{}\) admits a decomposition/exhaustion record with
exactly six role channels relative to which no kind residual and no record residual is licensed
label (x) or (xi) of Definition~\ref{def:residual-scheme}. A six-channel record
always exists (Theorem~\ref{thm:decomposition-existence}); the question is
decided by the residual condition. The theorem chain answering that question
appears in
Sections~\ref{sec:lower-bounds}--\ref{sec:decomposition-exhaustion}.

The domain is framed in the finite theory-package setting of
Foundations~I~\citep{Tsiokos2026SixBirdsFoundations}, where one works
with tuples \(\mathcal{T} = (Z, f, \Sigma_f, E, \mathcal{A})\)
carrying a lens, a completion operator, and an audit functional. \FATCD{} does not
identify a description with such a tuple; it extracts from that
setting the structural requirements relevant to the scoped exact-six
question: finite typed raw data, closure under references, sources,
targets, endpoints, and annotations, and the audit condition of
Definition~\ref{def:audited}, on which the regime of
Section~\ref{sec:admissibility-meta-theory} builds. Optimization,
objective, and richer semantic material are not excluded at the
outset; they belong in the domain when finitely represented and
honestly recorded as raw or annotated data.

\subsection{Raw grammar}\label{subsec:raw-grammar}

We begin with the neutral typed raw grammar from which every element
of \FATCD{} is built. The alphabet is fixed in advance of the role
vocabulary and follows standard typed-record
conventions~\citep{CardelliMitchell1991}.

\paragraph{How to read this section.} A description is a finite collection of
records. Each record is either a raw record of one of the kinds of
Definition~\ref{def:raw-record} or an annotation of one of the six structural or seven bookkeeping kinds
of Definition~\ref{def:annotations}; it carries finitely many typed fields
and may refer to other records of the same description. The conditions
of Section~\ref{subsec:conditions} then require three things: everything
is finite (Definition~\ref{def:finite}); every reference, endpoint and
annotation target is a record of the description, and every claim it
carries has a bookkeeping record (Definition~\ref{def:closed}); and those
bookkeeping records are complete (Definition~\ref{def:audited}). A
\emph{claim} is a claim annotation (Definition~\ref{def:audited}(a)), an assertion the description makes about its own
records, such as that a cluster is read as a role; a field with field role claim on another record is not a claim. Its bookkeeping record
is an audit annotation in \(\Audit{D}\) that refers to the claim, lists
the certificates it uses and records its conclusion
(Definition~\ref{def:audited}(a)); the claim annotation records the
statement, hypotheses and scope. For a claim admissible under
Definition~\ref{prop:honest-bookkeeping}, the promotions, translations,
lifts, forgetting maps, visibility and empirical bridges it uses have further
annotations carried in \(\Audit{D}\). Appendix~\ref{app:mechanization}
describes the typed Lean mirror of these definitions.

In set terms, a description can be written as a finite set \(R\) of records
with a kind map from \(R\) to the kinds of Definitions~\ref{def:raw-record}
and~\ref{def:annotations}, finitely many typed fields on each record, a
reference relation on \(R\) (an annotation refers to its target), a finite
set of claims, and a map sending each claim to its bookkeeping record in
\(\Audit{D}\subseteq R\). Definition~\ref{def:finite} makes all of these
finite, Definition~\ref{def:closed} requires every reference to land in
\(R\) and every claim to have a bookkeeping record, and
Definition~\ref{def:audited} lists the entries each bookkeeping record must
carry. Here \(R\) is \(\Raw{D}\). A \emph{typed field} is a triple \((n,\rho,v)\): a
name \(n\), a field role \(\rho\) (datum, reference, predicate, comparison,
invariant, law or claim; Definition~\ref{def:raw-record}), and a value \(v\) of
the type fixed by the record kind and \(\rho\). For a datum field this is a
rational, an element of a declared finite set, or a finite list or table of
such values and record identifiers, such as the graph of a map record; for a
reference field an identifier or a finite list of identifiers (the member
list of an object record, the candidate class of a realization relation, the
entry list of a trace, the list of successive records and the word of a path
or derivation record, and
every declared domain, codomain, source, target, endpoint and index field are
reference fields); for a term- or
formula-valued field a check- or assertion-language expression (a term-valued field, such as the row
\(v(w_1),\dots,v(w_n)\) of the transition record of Remark~\ref{rem:recognition-open}, has
field role datum; a formula-valued field has an assertion role); and for a
certificate field a triple \((\varphi,\text{inputs},b)\). A field may carry a
value-type tag. Deleting a field removes one triple; changing its value
replaces \(v\) by another value of the same type. A \emph{declared record} of \(D\) is an
element of \(R\), and a \emph{raw feature} is any declared record. A
\emph{content record} is a declared record of a kind of
Definition~\ref{def:raw-record}, and an \emph{annotation} one of a kind of
Definition~\ref{def:annotations}. Thus \(\Raw{D}\) contains annotations, while
``raw record'' means a content record.

\begin{definition}[Raw record]\label{def:raw-record}
A \emph{raw record} is a finite typed syntactic or relational feature of a
description. The admissible record kinds are
\begin{itemize}[topsep=2pt,itemsep=1pt]
  \item \emph{object} and \emph{state} records;
  \item \emph{relation} records, \emph{map} records, and \emph{partial-map}
    records;
  \item \emph{predicate} records and \emph{comparison} records;
  \item \emph{path} or \emph{derivation} records;
  \item \emph{index}, \emph{order}, \emph{transition}, and \emph{refinement}
    records; a \emph{stage} is an index record together with an order
    record ranging over it; in the \(W_4\) row and seed clause (4) of Definition~\ref{def:channel-status} a stage
    counts through its order record;
  \item \emph{trace}, \emph{event}, and \emph{provenance} records;
  \item \emph{replay} and \emph{check} records.
\end{itemize}
The rest of the definition fixes, in order: declared uses of map and relation records (including realization relations), the sort declarations of predicate records, rule fields, field roles, and value-type tags.
A map or partial-map record declares its use in a typed field: \emph{update},
for a map from a declared state set to itself, \emph{packaging}, for a map to
a declared quotient object (surjectivity is certified by the \(W_5\) check, not assumed), or \emph{other}; a relation record likewise declares its use
as \emph{quotient}, for an equivalence relation presented as the classes of a
quotient, or \emph{other}. The graph of a map or transition record lists exactly
one value at each member of its declared domain; only a partial-map record may
leave members of its declared domain without a value. The listed classes of a relation record of declared use quotient are nonempty and pairwise disjoint, so its relation atom is an equivalence relation on the members of its classes. These conditions are part of the
record's type, so the recorded values of \(D\) form a well-typed interpretation
in the sense of Definition~\ref{def:audited}(c1). A relation record may instead declare its use as
\emph{realization}: it declares a finite nonempty candidate class \(K\) of
declared records and, as its target, a predicate record \(s\) declared as a
macro-specification whose formula has exactly one free variable, which occurs in the formula and
ranges over records of the kind of the members of \(K\), and whose formula is a formula of the check language, so that it has a truth value over \(D\) at every member of \(K\); the relation lists
exactly those \(r\in K\) at which that formula evaluates to true in \(D\) (a condition on the recorded values of \(D\), required for membership in \FATCD{}; unlike the graph and class conditions above, it is not imposed on other well-typed interpretations of Definition~\ref{def:audited}(c1)); it may also record the value of that formula at further declared records
of the same kind. A macro-specification's formula may be closed; the
one-free-variable condition is imposed only by realization relations, seed
clause (2) and the \(W_2\) row. A predicate record declares in a
typed field whether it is a \emph{macro-specification}. It also declares, in a datum field, for each free variable of its formula and, for a predicate record without a formula field, for each argument place of its declared arity, the record kind over whose declared records that variable or argument ranges, the rational sort when it is applied to rational-valued terms, or a sort declared by an object record listing no members; for a macro-specification targeted by a realization relation this is the common kind of the members of \(K\), which must all be of one kind. A path, derivation or comparison
record may carry \emph{rule fields}: datum fields whose declared type, fixed by
the record kind, is a pair \((\lambda,\mu)\) of finite lists of identifiers of map
or partial-map records; the rewrite rules recorded in \(D\) are the values of
all rule fields of records of \(D\). Each typed field of a
record declares one of the field roles \emph{datum}, \emph{reference},
\emph{predicate}, \emph{comparison}, \emph{invariant}, \emph{law} or
\emph{claim}; a field recording a checked equality, isomorphism or level
translation is a comparison field. A field whose value is a proposition about
recorded values (a formula, constraint or asserted property) declares one of
the assertion roles predicate, comparison, invariant, law or claim; datum and
reference fields record values and identities only. ``Predicate field'', ``comparison field''
and the like refer to these declared roles. A field holding a certificate has
an assertion role: law when it is the law field of an order, transition or
refinement record, or of a relation record carrying an order in the sense of
the paragraph \emph{Laws} (Section~\ref{sec:primitive-role-architecture}), and
comparison otherwise; the certificate of an equality,
isomorphism or level translation is part of the comparison field it checks. A typed field may also carry a
declared value-type tag from a finite list recorded in \(D\) (for instance
probability, weight, kernel, confidence, error, intervention, counterfactual,
causal dependence, goal, objective, plan, policy or choice); tags make no
assertion and enter no witness or residual-label criterion;
Section~\ref{sec:upper-bound-classification} uses them to delimit its threat
classes. A raw record is not, by itself, a
\Pone{}--\Psix{} role.
\end{definition}

\begin{definition}[Annotations]\label{def:annotations}
In addition to raw records, a description admits \emph{annotations} of six
declared structural kinds:
\begin{itemize}[topsep=2pt,itemsep=1pt]
  \item level annotations \(\mathrm{Ann}_L\);
  \item threshold annotations \(\mathrm{Ann}_\Theta\);
  \item lift annotations \(\mathrm{Ann}_{\mathrm{Lift}}\);
  \item forgetting annotations \(\mathrm{Ann}_{\mathrm{Forget}}\);
  \item instrument-visibility annotations \(\mathrm{Ann}_{\mathrm{Vis}}\);
  \item empirical-bridge annotations \(\mathrm{Ann}_{\mathrm{Bridge}}\).
\end{itemize}
A description also admits \emph{bookkeeping annotations} of kinds audit,
claim, witness-schema, collapse-case, promotion, translation and nonclaim. A
\emph{promotion} annotation records a claim's move from a lower to a higher
level of Definition~\ref{prop:level-profile}, with the fields of
Definition~\ref{def:audited}(g); the selected higher-level datum the move adds
is recorded by a lift annotation.
They belong to \(\Raw{D}\), have finite typed fields and declared targets,
and are subject to the same reference-closure condition as the other
annotations. \(\Audit{D}\) is the collection of its honest audit
annotations (as in Definition~\ref{def:audited}(e)); a record is \emph{carried in} \(\Audit{D}\) when it is such an
annotation or a claim, promotion, translation, collapse-case,
witness-schema, nonclaim, lift, forgetting, visibility or bridge annotation
referenced by one.
\par\emph{Claim annotations and the representation condition.} A claim annotation records a statement, its referenced
records, its hypotheses and its scope; the scope field is a reference field
listing the records and clusters the claim concerns, and makes no assertion. It may additionally carry a
\emph{witness-cluster field}: a typed reference field listing exactly
the members of a cluster proposed for a named role. In every
construction in this paper, a claim described as naming \(P_i\)
and a cluster \(C\) for a witness check or derived projection
lists \(C\) in this field. Merely naming a content record does
not supply a witness-cluster field. A claim annotation is bookkeeping about
the declared content records it names. It satisfies the
\emph{representation condition} when every substantive assertion in its
statement, that is, any predicate, comparison, invariant or law content of the
statement other than naming its referenced records, hypotheses and scope, is
represented by fields of those records, and an assertion
introducing independent content requires an additional declared content
record. The audit identifies the fields and checks the asserted consequence.
Precisely, an assertion of the statement is \emph{represented by fields} of the
named records when it is the formula of one of their predicate, comparison,
invariant, law or certificate fields, a conjunction of such formulas, or a
formula logically entailed by a specified finite list of signed certificate
formulas carried in those fields (true in every well-typed interpretation, under every assignment of truth values to the atomic \(W_j\) formulas, in which those signed formulas are true; this is the entailment of Definition~\ref{def:audited}(c2) with the \(W_j\) atoms uninterpreted). For this representation test,
\((\varphi,\mathrm{true})\) supplies \(\varphi\) and
\((\varphi,\mathrm{false})\) supplies \(\neg\varphi\), and atomic \(W_j\)
formulas are treated as propositional atoms. Certificate correctness,
completeness and audit honesty are checked separately under
Definition~\ref{def:audited}. The statement \(W_i(C)\) or \(\neg W_i(C)\) of an
offer, and the statement \(W_i(C)\) of a role claim, is represented by the
members of \(C\) listed in its witness-cluster field, the \(\Audit{D}\) conjunct
of the \(W_6\) row being handled as in the next sentence.
For a \(W_6(C)\) claim, its event, trace and check content is represented by
the named content records; its separate \(\Audit{D}\)-entry conjunct is
evidenced by the honest event-or-trace audit entry required by the \(W_6\)
row.
\par\emph{Citation roles.} Two citation roles, (1) and (2), are fixed as follows; (3) is a condition common to both.
\begin{enumerate}[label=(\arabic*),topsep=2pt,itemsep=1pt]
  \item A \emph{\(W_i\)-witness record} is a content record of a cluster offered or used as a \(P_i\) witness that meets condition (3).
  \item A \emph{boundary record} is a record meeting condition (3) of a cluster on which a row of the boundary matrix (Table~\ref{tab:boundary-matrix}) is checked: a content record or, for the rows whose bare datum is a \Psix{} audit record, that audit annotation; a bare datum consisting of several records is presented as the set of their boundary records.
  \item A \(W_i\)-witness record carries exactly the fields specified by its declared witness check together with the datum and reference fields that check uses and reference fields to other members of the same cluster; a boundary record carries exactly the fields constituting the bare datum of its row of Table~\ref{tab:boundary-matrix}, the typing fields of its kind (kind and declared use), the datum and reference fields its certificates read, and reference fields to other members of the same cluster. A record carrying further fields is not a \(W_i\)-witness or boundary record; this is a condition on which records a claim may cite as such, not a membership condition of \FATCD{}.
\end{enumerate}
\noindent Independent operations on the six primitives, axioms for those operations, and global-independence certificates require separate content records.
\par\emph{Witness-schema annotations.} When a witness-schema annotation names
\(P_i\), it carries a target field listing the cluster to be checked. A
witness-schema annotation either names
\(P_i\) and the corresponding row of the \emph{Witness types} table of
Section~\ref{sec:primitive-role-architecture}, or names a candidate role
outside \(P_1,\ldots,P_6\) and
a check-language formula schema \(\psi(C)\) evaluated on declared clusters. A \emph{formula schema} is a check-language formula without atomic \(W_j\) formulas, with a cluster parameter \(C\) occurring only in finite disjunctions \(\bigvee_{d\in C}\chi(d)\) and conjunctions \(\bigwedge_{d\in C}\chi(d)\), \(d\) record-sorted; its \emph{instance} \(\psi(C')\) at a finite set \(C'\) of declared identifiers expands each over the members of \(C'\) (the empty disjunction being false and the empty conjunction true). We write \(t\in C\) for \(\bigvee_{d\in C}t=d\).
The latter annotation supplies no \(W_i\) witness or derived projection. An
assertion that \(W_i(C)\) holds additionally requires the finite evidence
specified in its table row.
These annotation kinds belong to the raw grammar; the finiteness and closure
conditions on how they attach to declared records are imposed in the next
subsection. For a description \(D\), we write \(\Raw{D}\) for the finite
collection of raw records and annotations of \(D\).
\end{definition}

\begin{remark}[Independence of the feature alphabet]\label{rem:alphabet-independence}
The record and annotation kinds in Definitions~\ref{def:raw-record} and
\ref{def:annotations} are syntactic or relational entities. No kind is named
\Pone{}, \ldots, \Psix{}. Consequently, membership in \FATCD{} does not
require any derived role projection, does not require an exact-six
decomposition/exhaustion record, and does not exclude candidate seventh-role
features by definition. An optional witness-schema annotation may name a
role, but no content kind is itself a role, and membership requires no
witness schema or derived projection.
\end{remark}

\subsection{Finiteness, closure, and audit conditions}\label{subsec:conditions}

Three further restrictions turn the raw grammar into a domain of
admissible descriptions: finiteness, closure, and audit.

\begin{definition}[Finite description]\label{def:finite}
A description \(D\) is \emph{finite} when
\begin{itemize}[topsep=2pt,itemsep=1pt]
  \item \(\Raw{D}\) contains finitely many records and annotations;
  \item every record or annotation carries finite typed fields;
  \item every relation, map, path, trace, or annotation is finitely
    represented;
  \item every reference in a record points to a declared record of
    \(D\);
  \item in particular, semantic, probabilistic, topological, geometric,
    causal, optimization, or resource structure enters \(D\) only through
    such finitely represented records, by the preceding conditions.
\end{itemize}
\end{definition}

\begin{definition}[Closed description]\label{def:closed}
A finite description \(D\) satisfies the \emph{closure condition} \(\Closed{D}\)
when
\begin{itemize}[topsep=2pt,itemsep=1pt]
  \item every source and target of every relation or map record is a declared
    record of \(D\);
  \item every endpoint of every path or derivation record is a declared record
    of \(D\);
  \item every index, transition, order, or refinement reference is a declared
    record of \(D\);
  \item every trace, event, provenance, replay, or check record references
    declared records of \(D\);
  \item every level, threshold, lift, forgetting, visibility, or
    empirical-bridge annotation references declared records of \(D\);
  \item every claim annotation carried by \(D\) is referred to by an audit annotation
    of \(D\) (its bookkeeping record).
\end{itemize}
\end{definition}

\begin{definition}[Audited description]\label{def:audited}
Summary, made precise in (a)--(g) below. An audit is honest exactly when its
references resolve (b1); its certificates recompute correctly (b2); it lists
every check its claim requires (b3); predicates with unrecorded extensions
occur in its conclusion only negatively (b4); its conclusion follows from the
signed certificates in every well-typed interpretation (c1)--(c2); and
unverified hypotheses occur only as antecedents of a conditional conclusion,
every hypothesis used in an outright conclusion being recorded and checked
true, and the conclusion uses no field it lists as non-certified (c4). These
conditions are evaluated in the order fixed by (f). Consequently (c3), every conclusion, outright or conditional, using only recorded predicates and carriers is true in
\(D\).
\par \emph{(a) Claims and audits.} A \emph{claim} is a claim annotation with a statement, referenced records,
hypotheses and scope. An audit annotation for a claim records references to
the claim and to the records it names, a finite list of certificates
(Section~\ref{sec:primitive-role-architecture}) with their recomputed truth
values, and a conclusion field holding a closed formula of the assertion language;
it \emph{certifies} exactly the formula in its conclusion field. \par\emph{(b) Honesty.} It is
\emph{honest} exactly when (b1)--(b4) below, the entailment requirement (c2)
and the requirements of (c4) all hold, evaluated in the order fixed by (f);
(c1) defines the interpretations used in (c2), and (c3) is a consequence:
(b1) its references resolve in \(D\); (b2) every listed
certificate is correct; (b3) the list contains every check the claim requires,
namely: (b3-i) for a claim naming \(P_i\) and listing a cluster \(C\) in a
witness-cluster field (the claim of an offer in the sense of (d)), every certificate that the \(W_i\) row prescribes and a member of \(C\) carries (for \(W_2\), including every recorded evaluation certificate \((s(t),b)\)), each listed as \((\varphi,b')\) with \(b'\) its recomputed truth value, together with the certificate \((W_i(C),b)\) recording
the result, or, when \(C\) lacks a record that every alternative of the row requires, that certificate
alone, with \(b=\mathrm{false}\); (b3-ii) for \(W_6(C)\), the content checks of the
\(W_6\) row; (b3-iii) for an assertion represented by fields of content records, the
certificates in those fields; and (b4) the conclusion field contains a closed formula of the assertion language in which every atom whose
predicate symbol has an unrecorded extension occurs only negatively (within an
odd number of implication antecedents and negations). Here \(D\) records the extension of a declared predicate \(p\) when at least
one complete presentation of \(p\) is supplied, either by a relation record
that references the declaring record of \(p\) in a field declaring it the
extension of \(p\), its listed tuples being by definition the true tuples of
\(p\), provided the argument sorts of \(p\) have recorded finite carriers, or
by its declaring predicate record giving an explicit formula of the fixed check
language over recorded inputs and declared finite ranges (arguments of
rational sort being then admitted and evaluated by that formula), all such presentations agreeing on
every tuple; the criterion is repeated, with the definition of recorded
carriers, in the paragraph \emph{Terms used in the label criteria} after
Definition~\ref{def:residual-scheme}.
\par\emph{(c) Premises.} (c1)~\emph{Interpretations.} An \emph{interpretation} (or \emph{well-typed interpretation}) of \(D\) does three things. First, it keeps the records of \(D\), their kinds, declared uses, the markers declaring a transition codomain distribution-valued, field roles and references; so reference-field carrier lists and declared map and transition domains are fixed across interpretations, while quantifier bounds supplied by datum fields may vary. Second, it keeps every formula-valued field (predicate, comparison, invariant, law, claim, certificate and conclusion fields) and every term-valued field syntactically fixed. A term-valued field is a field explicitly recorded as a symbolic check-language term or a finite list of such terms, including a transition row recorded as \(v(w_1),\dots,v(w_n)\); these expressions are evaluated using the interpretation's datum values, componentwise for lists; ordinary numeric datum fields are not term-valued merely because their values are rational constants. Third, it assigns each remaining datum field a value of its declared type: the type of a map, partial-map or transition graph, or of a table field, fixes the sorts of its values but not their membership in a declared codomain or list, and a row into \(\mathrm{Dist}(S)\) has exactly one coordinate for each member of the fixed reference-field list \(S\), with arbitrary rational entries. Unrecorded predicate extensions and unlisted-sort carriers vary arbitrarily, and \(W_j(C)\) is evaluated by its row, the \(\Audit{D}\) conjunct of the \(W_6\) row ranging over \(\mathrm{Audit}_0(D)\) of \(D\) itself.
(c2)~\emph{Entailment.} For a correct certificate
\((\varphi,b)\), an audit may use \(\varphi\) as a premise if \(b\) is true
and \(\neg\varphi\) if \(b\) is false. Honesty requires the signed premises to entail the conclusion in every
interpretation; the entailment requirement is therefore a validity condition over these interpretations; certificate correctness is checked separately by recomputation. When an audit conclusion is one of its signed premises, entailment holds by syntactic identity, but reference closure, certificate correctness, completeness and the remaining honesty requirements must also hold. (c3)~\emph{Truth of conclusions.} In particular, every conclusion of an honest audit, outright or conditional, that uses no unrecorded predicate
symbol and quantifies only over recorded carriers is true in \(D\): the recorded
datum values of \(D\), with any choice of unrecorded extensions and carriers,
form an interpretation in which every correct signed premise is true, so the
conclusion is true there, and its value does not depend on extensions it does
not use. (c4)~\emph{Hypotheses and non-certification.} An unverified stated hypothesis may occur only as
an antecedent of a conditional conclusion; an outright conclusion requires
every hypothesis used to be recorded and checked true. A failed check may thus
be honestly recorded without supporting a positive witness claim. An audit
annotation may carry a \emph{non-certification field}, a reference field
listing fields of the
records it names whose asserted property it does not certify; it is honest
only if its conclusion does not use those fields.
\par\emph{(d) Offers.} In particular, an \emph{offer} of a cluster \(C\) as a \(P_i\)
witness is a claim annotation naming \(P_i\) and listing \(C\) in a
witness-cluster field, together with an honest
audit annotation recording the check of \(W_i(C)\) and its result; an offer
requires no threshold, level or collapse-case annotation; its statement is
\(W_i(C)\) when the recorded result is true and \(\neg W_i(C)\) when it is false,
and its audit concludes that statement from the certificate \((W_i(C),b)\). A
witness-schema annotation records the predicate to be checked, not a claim
that every referenced cluster satisfies it.
\par\emph{(e) Audited descriptions.} A description \(D\) is
\emph{audited} when every audit annotation it carries is honest and every
claim it carries is referred to by at least one such annotation and
satisfies the representation condition of Definition~\ref{def:annotations}; the honest
audit annotations form \(\Audit{D}\subseteq\Raw{D}\).
\par\emph{(f) Order of evaluation.} Honesty is determined in
two stages. Let \(\mathrm{Audit}_0(D)\) be the honest audit annotations whose
certificates and conclusion contain no atomic \(W_j\) formula; their honesty
does not depend on any \(W_j\). The \(\Audit{D}\) conjunct of the \(W_6\) row
ranges over \(\mathrm{Audit}_0(D)\), so an audit of \(W_6(C)\) cannot itself
discharge it. Every atom \(W_j(C')\) is then determined, since
the rows of \(W_1,\dots,W_5\) refer to no audit annotation and so do not depend
on which audits are honest, and \(W_6\) depends on audits only through
\(\mathrm{Audit}_0(D)\); this determines the
honesty of the remaining audit annotations.
\par\emph{(g) Required fields.} Promotions,
translations, lifts, forgetting, visibility and bridge assertions must also
carry the following fields, with their stated losses and limits preserved:
\begin{itemize}[topsep=2pt,itemsep=1pt]
  \item every promotion records its source level, target level, added data,
    losses, and nonclaims;
  \item every translation records its source context, target context,
    preserved data, suppressed data, and claim strength;
  \item every lift is marked as selected and non-unique unless uniqueness for its specified source data and target class is proved and recorded;
  \item every forgetting map records its losses;
  \item every visibility claim records both visible and suppressed data;
  \item every empirical-bridge claim records substrate, observable,
    instrument, level map, threshold witness relation, and
    suppressed-structure nonclaims.
\end{itemize}
\end{definition}

\subsection{Finite audited typed closure descriptions}\label{subsec:fatcd}

Combining the grammar with the three admissibility conditions yields
the domain over which the exact-six question is posed.

\begin{definition}[Finite audited typed closure description]\label{def:fatcd}
A \emph{finite audited typed closure description} is a finitely presented,
closed, audited collection of typed raw records and annotations: every claim
it carries is referred to by an honest audit annotation, and its promotions,
translations, lifts, forgetting, visibility and empirical-bridge assertions
carry the fields listed in Definition~\ref{def:audited}. Membership does not
require every claim to be certified: a member may carry a claim whose honest
audit does not certify its statement, making that claim inadmissible under
Definition~\ref{prop:honest-bookkeeping}. Precisely, it is a
description \(D\) whose raw content \(\Raw{D}\) is drawn from the kinds
declared in Definitions~\ref{def:raw-record} and~\ref{def:annotations}, such
that \(D\) is finite in the sense of Definition~\ref{def:finite}, satisfies
the closure condition \(\Closed{D}\) of Definition~\ref{def:closed}, is
audited in the sense of Definition~\ref{def:audited}, and each realization relation of \(D\) lists exactly the members of its candidate class at which its target's formula is true in \(D\) (Definition~\ref{def:raw-record}). We write \FATCD{} for
the class of all such descriptions, and \(D \in \FATCD{}\) for membership.
\end{definition}

\begin{remark}[No minimum richness]\label{rem:fatcd-minimality}
\FATCD{} imposes no lower bound on description size. A description \(D_0\)
consisting of one declared object record and no claims is finite, closed and
audited. Its record has no witness, predicate, comparison, invariant, law or
claim field, so its record satisfies (R3) of
Definition~\ref{def:recognition-hypothesis} and \(D_0\) satisfies the
recognition hypothesis. Richer descriptions carrying
transitions, paths, derivations or annotations of every kind may also belong
to \FATCD{}; whether they satisfy the recognition hypothesis is checked record
by record. The domain is defined by structural discipline, not by the presence
of any particular feature kind. A second small member, used in
Proposition~\ref{prop:non-circularity}, is the following description \(D_{\mathrm{norm}}\), analysed in
Remark~\ref{rem:recognition-open}: state records \(w_1,\dots,w_n\) (\(n\ge1\))
with nonnegative rational value fields \(v(w_k)\) summing to one (for instance
\(n=2\) and \(v(w_1)=v(w_2)=1/2\)), a predicate record \(N\) (unrelated to the monoid \(N\) of route certificate (a)) declared a
macro-specification whose formula is \(\bigwedge_k v(w_k)\ge0\wedge\sum_k v(w_k)=1\), a claim annotation
naming \(N\) whose statement is the formula of \(N\), and an honest audit
annotation listing the certificate \((\bigwedge_k v(w_k)\ge0\wedge\sum_k v(w_k)=1,\mathrm{true})\) and
certifying that formula.
\end{remark}

\subsection{Threshold attachment}\label{subsec:threshold-attachment}

Membership in \FATCD{} is fixed; the next step is to specify how
role-theoretic claims may attach to the raw data of an admissible
description. Role projections in \FATCD{} enter only through a
disciplined attachment of threshold and witness data to raw features.
The condition below is the structural precondition for the role
vocabulary introduced in the next subsection.

\begin{definition}[Threshold attachment]\label{def:threshold-attachment}
Let \(C\subseteq\Raw{D}\) be a finite cluster of raw features of \(D\), possibly
a single feature; attachment of an annotation to \(C\) is defined below. A derived role projection on \(C\) is
admissible only when all of the following are present in \(\Raw{D}\):
\begin{itemize}[topsep=2pt,itemsep=1pt]
  \item the raw features of \(C\) themselves;
  \item a threshold annotation \(\mathrm{Ann}_\Theta\) attached to \(C\);
  \item a witness schema, in the sense of
    Definition~\ref{def:role-projection}, for the attempted projection, recorded by a witness-schema annotation whose target field lists exactly the members of \(C\);
  \item a collapse-case annotation targeting exactly \(C\), a record of the collapse
    cases that would invalidate the witness; like a threshold label, its
    content is recorded and is not checked by the results of this paper;
  \item a level annotation \(\mathrm{Ann}_L\) attached to \(C\);
  \item a nonclaim annotation targeting exactly \(C\);
  \item an audit record in \(\Audit{D}\) honestly bookkeeping a claim annotation
    that names the attempted role and lists \(C\) in its witness-cluster field.
\end{itemize}
These are necessary conditions; the complete condition is the admissibility
clause of Definition~\ref{def:role-projection}.
An annotation is \emph{attached to} \(C\) when its declared target field
lists exactly the members of \(C\). A threshold annotation may carry a \emph{threshold predicate} on recorded
fields of \(C\); when it does, the projection on \(C\) is attempted only if
\(\Audit{D}\) records that predicate checked true, that is, only if some audit annotation in \(\Audit{D}\) lists a correct certificate \((\theta,\mathrm{true})\) whose formula is the threshold predicate \(\theta\), a closed check-language formula. A threshold annotation may
also carry a label field naming the attempted reading, and a level annotation
carries a level field. Label and level fields are operation data in the sense
of (R2) of Definition~\ref{def:recognition-hypothesis} and of criterion (ix-b),
and make no assertion. A threshold predicate is operation data under (R2) and
(ix-b) only when it reads the value of at least one record, every record whose values it reads is a member of \(C\), and each such record satisfies (R1); otherwise
it is an independent assertion.
\end{definition}

\begin{remark}[Pre-projection discipline]\label{rem:pre-projection}
The items of Definition~\ref{def:threshold-attachment} form a structural
precondition, not a process narrative. What matters is that each item is
present explicitly in \(\Raw{D}\) before a role projection is asserted, so
that no threshold, witness, collapse case, or level is supplied tacitly by
the act of projection itself. The lower-bound examples use threshold
annotations as recorded tags; Definition~\ref{def:threshold-attachment} does
not require a numerical cut-off. By the ordered rule of
Definition~\ref{def:channel-status}, \belowthr{} means that a covered seed is
present but neither an admissible active projection nor a failed offer takes
precedence. The missing data may be witness fields or attachments.
\end{remark}

\subsection{Derived role projections}\label{subsec:derived-role-projections}

The role vocabulary \Pone{}--\Psix{} now enters \FATCD{}, not as a
constraint on membership but as a derived interpretation of raw
features conditioned on the attachment data of
\S\ref{subsec:threshold-attachment}.

\begin{definition}[Derived role projection]\label{def:role-projection}
A \emph{derived role projection} for a description \(D \in \FATCD{}\) at
role index \(i \in \{1,\ldots,6\}\) is a tuple \(\pi_i(D,C)\), written \(\pi_i(D)\) when the cluster \(C\) is fixed
by context (a description may carry admissible projections on several clusters
for the same \(i\)), specified as follows, interpreting raw features of \(D\) as the role \(\mathrm{P}_i\), supported
by threshold data, witness schema, level assignment, and honest bookkeeping.
Concretely, \(\pi_i(D)\) is specified by
\begin{itemize}[topsep=2pt,itemsep=1pt]
  \item the proposed role name, one of \Pone{}, \Ptwo{}, \Pthree{},
    \Pfour{}, \Pfive{}, \Psix{};
  \item a raw feature cluster \(C \subseteq \Raw{D}\) on which the role is
    interpreted;
  \item the threshold data \(\Theta_i\) attached to \(C\) in the sense of
    Definition~\ref{def:threshold-attachment};
  \item a witness schema \(W_i\): the predicate on pairs \((D,C)\) of a finite
    description and a finite \(C\subseteq\Raw{D}\), not necessarily with
    \(D\in\FATCD{}\), given by the \(P_i\) row of the
    \emph{Witness types} table in Section~\ref{sec:primitive-role-architecture}
    (depending on annotations only through the \(\Audit{D}\) conjunct of the
    \(W_6\) row, which is evaluated over \(\mathrm{Audit}_0(D)\), not over \(C\)),
    fixed by a witness-schema annotation \(w_i\in\Raw{D}\) naming \(P_i\) and
    that row, which the tuple also records; we write \(W_i(C)\) when \(D\) is
    fixed;
  \item a level assignment \(\mathrm{Level}_i\);
  \item an evidence or certificate type;
  \item a record of collapse cases avoided or accepted;
  \item loss and lift annotations for any level crossings involved;
  \item the explicit nonclaims preserved by the projection.
\end{itemize}
The supporting data for \(\pi_i(D)\) are carried by
\(\Raw{D}\) and certified by \(\Audit{D}\) in the sense of
Definition~\ref{def:audited}. Where the cluster matters we write
\(\pi_i(D,C)\). The projection \(\pi_i(D,C)\) makes a \emph{level crossing} when the level
recorded in its level annotation differs from the level recorded in some other
level annotation whose target field lists a member of \(C\). The projection \(\pi_i(D,C)\) is \emph{admissible} if and only if
\(W_i(C)\) holds with its check recorded by an honest audit annotation; the
threshold, witness-schema, collapse-case and level annotations target exactly
\(C\); every threshold predicate \(\theta\) they carry is recorded in \(\Audit{D}\) as
checked true in the sense of Definition~\ref{def:threshold-attachment}; a claim naming \(P_i\) and \(C\) is admissible in the sense of
Definition~\ref{prop:honest-bookkeeping}, some honest audit annotation
referring to it and certifying \(W_i(C)\); a nonclaim annotation targets exactly \(C\); if the
projection makes a level crossing, it carries the translation and the lift or
forgetting annotation required for that crossing, with the audited losses or
selected-lift data required by Definitions~\ref{prop:level-profile}
and~\ref{prop:forgetting-lifts}; and every lift or forgetting annotation whose
targets meet \(C\) satisfies Definition~\ref{prop:forgetting-lifts}. The
evidence or certificate is then an applicable alternative of the \(W_i\) row,
since \(W_i(C)\) holds, and no lift or forgetting annotation is required when
no crossing occurs. An admissible derived role projection is called
\emph{active}; a channel is active when its status is \activeproj{}
(Definition~\ref{def:channel-status}).
\end{definition}

\begin{remark}[Fixed role meanings]\label{rem:fixed-role-meanings}
The role names \Pone{}--\Psix{} are defined by their meanings in
Section~\ref{sec:primitive-role-architecture}: \Pone{} concerns
update-vs-quotient compatibility and descent obstruction; \Ptwo{} concerns
macro-specification representability; \Pthree{} concerns enriched
route and coherence mismatch, not a directionality certificate; \Pfour{}
concerns stage, order, transition, and refinement structure; \Pfive{}
concerns quotient packaging; and
\Psix{} concerns audit, bookkeeping, provenance, replay, and checkability.
The witness schema \(W_i\) in Definition~\ref{def:role-projection} must be
aligned with these meanings; a cluster whose raw features do not support the
meaning of the proposed role contributes no \(\pi_i(D)\).
\end{remark}

\paragraph{What a witness type is.} Proposition~\ref{prop:alignment-rules}
lists, for each role, the data a candidate cluster must carry (tabulated
under \emph{Witness types} in Section~\ref{sec:primitive-role-architecture}): for \Pone{},
for instance, a map or partial-map record of declared use update, source and target relation records of declared use quotient, and a certificate of alternative (a) or (b) of the \(W_1\) row.
A cluster or feature \emph{supplies the witness type} of \(P_i\) when it
carries the data listed there for \(P_i\). The witness schema \(W_i\), recorded in
\(\Raw{D}\) under Definition~\ref{def:threshold-attachment}, specifies what in
the candidate cluster would support the projection.

\begin{remark}[Projection and membership are independent]\label{rem:projection-vs-membership}
Whether a description \(D\) supports a given derived role projection
\(\pi_i(D)\) is not part of the admissibility conditions of
Definitions~\ref{def:finite}--\ref{def:audited}. A description
\(D \in \FATCD{}\) may carry no derived role projection, projections
for some roles, or projections for all six; its membership in
\FATCD{} is unaffected. The decomposition/exhaustion record of
Section~\ref{sec:decomposition-exhaustion} accordingly assigns a
status to every role channel of \(D\), only some of which need be
active.
\end{remark}

\subsection{Residual classification scheme}\label{subsec:residual-classification}

Some raw features remain role-theoretically unsettled at the
raw-grammar level. \FATCD{} provides an explicit scheme for
classifying them.

The scheme is stated in four pieces: the notions it borrows from later
sections (next paragraph); the list of labels (Definition~\ref{def:residual-scheme});
the criterion for each label (the table after it, part of the definition); and
the terms those criteria use (the paragraph after the table).
Section~\ref{sec:upper-bound-classification} restates the criteria in
record-independent form as (R1)--(R5).

\paragraph{Forward notions.} The label criteria below use six notions
defined later; we state them here, and the cited definitions are the official
source. A cluster is \emph{aligned} with \(P_i\) when \(W_i\) holds of it, and
the required fields are defined after the label table. A \emph{covered seed}
for \(P_i\) (Definition~\ref{def:channel-status}, p.~\pageref{def:channel-status}, which also defines failed offers and outside-scope declarations) is a record of one of the
kinds (1)--(6) of that definition, with the named seed fields listed there.
Briefly: (1) an update map with its source relation; (2) a one-free-variable
macro-specification taking both truth values, with a realization relation;
(3) a comparison certificate for two routes with the same endpoints; (4) a
law-bearing stage, order, transition or refinement record; (5) a quotient
relation or packaging map; (6) an event, provenance or trace record, or a
replay or check record certifying a property of its trace. \(A(\Dec{D})\) and \(S(\Dec{D})\) are the records of
the active projections and of the inactive-status entries of a decomposition
record (Definition~\ref{def:decomposition-record}), fixed before its residual
features are labelled. An \emph{(R4) chain} from \(r\) is a finite sequence of
distinct declared records \(r=r_0,\ldots,r_n\), \(n\ge1\), in which, for each
\(j<n\), \(r_j\) records a checked equality or isomorphism of presentation data (a bijection between specified datum fields that are not term-valued, or member lists of object records, of \(r_j\) and \(r_{j+1}\), with a certificate that each value equals its image; a level-translation step carries such a bijection and certificate as well; no structural isomorphism is meant) with,
or a checked level translation to, \(r_{j+1}\) and, apart from one comparison field recording the equality, isomorphism or level translation with \(r_{j+1}\), and the datum and reference fields it compares, has no predicate, comparison, invariant, law or claim field (Definition~\ref{def:recognition-hypothesis}).
Two further notions come from Definition~\ref{def:recognition-hypothesis}. A
record satisfies \emph{(R1)} when it contributes a required field to a cluster
aligned with some \(P_i\) or carries a covered seed, and each of its predicate,
comparison, invariant, law or claim fields is required for, or an input to the
check of, \(W_{i_\varphi}\) at some cluster \(C_\varphi\ni r\) aligned with
\(P_{i_\varphi}\), with the index and cluster chosen separately for each field
\(\varphi\), or is a named field of a seed clause it carries (so \(r\) carries no assertion beyond these fields; content beyond them must be recorded in a separate declared content record, which is classified on its own); or when it is a member of the cluster of an honestly
recorded failed offer and each such field is an input to the recorded check.
The \emph{exempt fields of (R2)} are the statement and hypothesis fields of a
claim annotation, the certificate field of an audit annotation and its conclusion field when its conclusion is true in every well-typed interpretation or is represented, in the sense of Definition~\ref{def:annotations}, by fields of the content records named by the claim it audits,
and the formula schema of a witness-schema annotation; a conclusion field that is not exempt is a formula-valued field, so an audit annotation carrying one does not have only datum and reference fields besides its exempt fields.

\begin{definition}[Residual feature; residual classification scheme]\label{def:residual-scheme}
Let \(D \in \FATCD{}\). A \emph{residual feature}, or \emph{kind residual}, of \(D\) is
any content record of Definition~\ref{def:raw-record} or any of the six
structural annotations of Definition~\ref{def:annotations} in \(\Raw{D}\). The
seven bookkeeping annotations are excluded from this kind-residual set; they
may still be record residuals of a particular decomposition. Raw-grammar kind
alone assigns none of these features a role; ``residual'' means residual to
the raw kinds, not left over by a decomposition, and every content record is a
kind residual, including the constituents of active clusters. The records left
over by a particular decomposition record form the set \(\Residual{D}\) of
Section~\ref{sec:decomposition-exhaustion}. The \emph{residual classification scheme}
specifies, relative to the channel entries of a decomposition record (its six
statuses and the clusters, offers, seeds and declarations they name, which fix
\(A(\Dec{D})\), \(S(\Dec{D})\) and \(\Residual{D}\) before any label is chosen),
which labels are licensed for each residual feature; a record chooses one licensed label
\(\operatorname{Classify}(r)\) for each \(r\in\Residual{D}\)
(Definition~\ref{def:decomposition-record}), drawn from the list
\begin{enumerate}[label=(\roman*),topsep=2pt,itemsep=1pt]
  \item a declared constituent of the cluster that the decomposition record
    names for an active \Pone{}--\Psix{} projection \(\pi_i(D)\), whose
    cluster satisfies \(W_i\);
  \item shared constituent of two such projections;
  \item presentation variant of an already classified feature;
  \item level shift of an already classified feature;
  \item unabsorbed role-aligned content, covered-seed content with no
    certified witness \(W_i\), content checked in an honestly failed offer, or activation-threshold refinement;
  \item instrument-visibility artifact;
  \item empirical-bridge annotation;
  \item outside-domain structure;
  \item already-covered bookkeeping or annotation data, or role-neutral data
    without an assertion of its own;
  \item unresolved candidate seventh-role threat: an in-scope feature placed by none of (i)--(ix), which need be neither true nor a proposed role (Remark~\ref{rem:recognition-open});
  \item recorded seventh-role screen.
\end{enumerate}
The same criteria apply to every member of \(\Residual{D}\) for a chosen
decomposition record, including a bookkeeping annotation not assigned to a
channel. They are evaluated in the same way for every bookkeeping annotation of \(D\), assigned to a channel or not; Theorems~\ref{thm:upper-bound-classification} and~\ref{thm:scoped-exact-six} use them so. A claim annotation's bookkeeping fields may receive (ix); its substantive
assertion is classified through the content records representing it,
including any additional record required for independent content. For a feature with no claim about it, ``honestly audited'' in
criterion (x) means that every claim concerning it, if any, has an honest
audit; it does not require an audit annotation for an unclaimed raw record.
A decomposition record fixes one such assignment; where the label criteria
license more than one label for \(r\), different records may choose
differently (Proposition~\ref{prop:non-uniqueness}). A label applies to \(r\) under the criterion given for it in the table
after this definition, which is part of the definition. In particular,
(x) applies when \(r\) has finite, honestly audited content in the covered scope and
no criterion (i)--(ix) holds for it, so that it remains an unresolved
candidate; and (xi) applies under the criterion in the last row of the table. In words, (xi)
holds for an unresolved feature \(r\) when a recorded formula \(\psi\), attached
to a named candidate role outside \(P_1,\dots,P_6\), holds at some cluster
containing \(r\) at which no \(W_i\) holds, fails at some other cluster and at
every cluster at which some \(W_i\) holds, and \(r\) records a closure operator \(c\) on
a declared finite order lying outside the monoid generated by \(\mathrm{id}_X\) and the restrictions to \(X\) of the graphs of those map, partial-map, transition and refinement records of
\(D\) that are defined at every member of \(X\) with values in \(X\) (a map recorded only in a field of a record of another kind is not counted). There the
separation certificate is the one whose formula is the finite conjunction over
all \(C'\subseteq\Raw{D}\) of
\((W_1(C')\vee\dots\vee W_6(C'))\rightarrow\neg\psi(C')\), with \(\psi(C')\) the
instance of the schema at the list \(C'\); this separation certificate
establishes no witness row, law or seed, so its formula may contain \(W\)
atoms, although \(\psi\) itself contains none, and ``\(\psi\) false at some other cluster'' is the certificate
\(\neg\psi(C_0)\) for a named \(C_0\ne C\). Criterion (xi) is a recorded screen, not the identification of a new closure
role.
\par\emph{Example (illustrative; not part of the definition).} For \(X=\{0,1\}\) with a declared order record for \(0<1\) and a certificate of its order law, a relation record of declared use other may
carry a finite table field \(c\) with \(c(0)=1\), \(c(1)=1\), and a datum field
listing \((\mathrm{id}_X)\), with correct certificates, recorded true, of extensivity,
monotonicity, \(c^2=c\), closure of this list under composition, and
\(c\ne\mathrm{id}_X\). If no listed generating-record kind records a total map
on \(X\), its generated monoid is \(\{\mathrm{id}_X\}\). These records establish only the closure-operation part of (xi). The full
criterion additionally requires (x), a witness-schema annotation, and the
certified \(\psi\)-separation specified above. For that full criterion, add a check record carrying the stated recomputable
certificates, and a claim annotation referencing both \(r\) and that check
record whose statement is exactly the certified screen conditions. An honest
audit of the claim cites those records and certifies that statement. This ends
the example.
\par The scheme is declarative. At the domain-definition level, \FATCD{} records
the available labels but neither proves that every residual feature lands in
categories (i)--(ix) nor proves that categories (x) and (xi) are empty. Those
claims are deferred to Section~\ref{sec:upper-bound-classification}.
The terms aligned, contributes a required field, required for \(W_i\) at \(C\),
input to the check and outside the covered scope are defined in the paragraph
\emph{Terms used in the label criteria} after the table below.
\end{definition}

\paragraph{When each label applies.} Labels are assigned relative to the
channel entries of a fixed decomposition record \(\Dec{D}\)
(Definition~\ref{def:decomposition-record}); well-formedness of \(\Dec{D}\)
then requires each chosen label to be licensed in this sense. Labels (i)--(v) use
the named channel entries where their criteria require them, together with the
records and clusters specified in those criteria. The residual scheme allows the
following labels for a residual feature \(r\) in the sense of this
definition; the criteria are given in the table below, and
Theorem~\ref{thm:upper-bound-classification} uses them to exclude (x) and
(xi) under the recognition hypothesis. Labels~(i) and~(ii) do not occur in the
set \(\Residual{D}\) of a decomposition record
(Section~\ref{sec:decomposition-exhaustion}).
\begingroup\small
\begin{longtable}{@{}l>{\raggedright\arraybackslash}p{0.88\linewidth}@{}}
\toprule
Label & The residual feature \(r\) \\
\midrule
\endhead
(i) & belongs to the cluster that the decomposition record names for an active projection satisfying \(W_i\), and every predicate, comparison, invariant, law or claim field of \(r\) is required for \(W_i\) at the cluster \(C\) of the clause, or is an input to the check of \(W_i\) at \(C\), or is a named field of a seed clause that \(r\) carries (so \(r\) carries no assertion beyond these fields; content beyond them must be recorded in a separate declared content record, which is classified on its own) \\
(ii) & there are indices \(i\neq j\) and clusters \(C_i,C_j\) named by the decomposition record for active projections such that \(r\) contributes a required field to both, and every predicate, comparison, invariant, law or claim field of \(r\) is required for, or an input to the check of, \(W_i\) at \(C_i\) or \(W_j\) at \(C_j\); the relation record \(E\) of the running example, which lies in the \Pone{} and \Pfive{} clusters, is an example, as is the map record for \(F\), which lies in \(C_1\) and \(C_6\) (the clusters of the paragraph \emph{The running example as a description} below) \\
(iii) & re-presents a classified feature: \(r=r_0\) begins a finite chain \(r_0,\ldots,r_n\) as in (R4), with \(n\geq1\), distinct records and only checked equality or isomorphism steps, whose terminal \(r_n\) lies in \(A(\Dec{D})\cup S(\Dec{D})\) or satisfies a criterion among (i), (ii) and (v)--(ix) \\
(iv) & restates a classified feature at a different level: \(r=r_0\) begins a finite chain \(r_0,\ldots,r_n\) as in (R4), with \(n\geq1\), distinct records and at least one checked level-translation step, whose terminal \(r_n\) lies in \(A(\Dec{D})\cup S(\Dec{D})\) or satisfies a criterion among (i), (ii) and (v)--(ix) \\
(v) & one of: (v-a) \(r\) contributes a required field to a six-role-aligned cluster, or itself carries a covered seed for one of the six roles, but criterion (i) does not hold for \(r\); this includes a cluster missing an attachment and a seed lacking a complete witness; and every predicate, comparison, invariant, law or claim field \(\varphi\) of \(r\) is, for some index \(i_\varphi\) and some finite cluster \(C_\varphi\ni r\) aligned with \(P_{i_\varphi}\), required for \(W_{i_\varphi}\) at \(C_\varphi\) or an input to its check, or is named in a seed clause that \(r\) carries (so \(r\) carries no assertion beyond these fields; content beyond them must be recorded in a separate declared content record, which is classified on its own);\newline (v-b) \(r\) is a threshold annotation targeting a cluster named for an active projection, which carries a threshold predicate and no field other than its typing, target and label fields and that predicate, in the sense of Definition~\ref{def:threshold-attachment}, and the threshold predicate reads at least one record, every record whose values it reads is listed in the annotation's target field, and each such record satisfies (R1);\newline (v-c) \(r\) is a member of the cluster of an honestly recorded failed offer for some \(P_i\), each of whose predicate, comparison, invariant, law or claim fields is an input to the recorded check of \(W_i\) on that cluster \\
(vi) & is an instrument-visibility annotation whose fields other than its typing and target fields are its visible data, suppressed data, instruments and claim strength (Definition~\ref{prop:visibility}), each formula-valued one of which reads at least one record, reads only records listed in its target field, and reads only records satisfying (R1) \\
(vii) & is an empirical-bridge annotation identifying an admissible bridge claim whose associated entry in \(\Audit{D}\) carries the five items of Definition~\ref{prop:empirical-bridge}, and whose own fields other than its typing and target fields only identify that bridge claim and those items \\
(viii) & lies outside the covered scope: a field of \(r\) asserts a property whose check is not finite over the records of \(D\): it quantifies over a declared domain for which \(D\) supplies no finite checking range, or uses a declared predicate whose extension \(D\) does not record, and its truth value is not fixed by the recorded values of \(D\) (it is true in some expansion of \(D\) and false in another, in the sense of the paragraph \emph{Certificates, assertion language and laws} of Section~\ref{sec:primitive-role-architecture}; for a formula with free variables: for some assignment \(\alpha\) of its free variables and two expansions of \(D\) whose carriers contain the values \(\alpha\) gives to variables of unlisted sorts, it is true under \(\alpha\) in one and false under \(\alpha\) in the other). In addition, every other predicate, comparison, invariant, law or claim field of \(r\) either has this property or is required for \(W_i\) at a cluster \(C\) aligned with \(P_i\), an input to the check of \(W_i\) at such a cluster, or named in a seed clause that \(r\) carries \\
(ix) & is one of: (ix-a) a bookkeeping annotation whose fields other than the exempt fields of (R2) are datum or reference fields;\newline (ix-b) a structural annotation whose fields contain only its declared operation, its targets, and the finite typing, checking, preserved or suppressed data, losses and nonclaims required for that operation, with no independent assertion, when (v), (vi) and (vii) do not apply. Here a formula-valued field of a structural annotation (a threshold predicate, a preserved or lost invariant) is part of its declared operation only when it reads at least one record, every record whose values it reads is listed in the annotation's target field, and each such record satisfies (R1); otherwise it is an independent assertion and (ix-b) does not apply;\newline (ix-c) role-neutral data without an assertion of its own: it supplies no witness field and has no predicate, comparison, invariant, law or claim field, such as an input used in checking a required or seed field, or a declared object or state that no role datum uses; also a predicate record supplying no witness field whose only assertion field is a formula with at least one free variable whose truth value is not the same under all assignments respecting the declared sorts of its free variables, using only predicates whose extensions \(D\) records, and quantifying only over sorts whose carriers \(D\) records, each free variable ranging over a record kind or the rational sort, referenced by no claim annotation, and referenced by certificates only through the abbreviation \(s(t)\) of the paragraph \emph{Check language} (a definition, which has no truth value over \(D\)) \\
(x) & has finite, honestly audited content in the covered scope, and no criterion (i)--(ix) holds for it; for \(D\in\FATCD{}\) the first two conditions hold for every record (Definitions~\ref{def:finite} and~\ref{def:audited}(e)), so (x) applies exactly when \(r\) lies in the covered scope and no criterion (i)--(ix) holds \\
(xi) & satisfies (x), and: (a) \(D\) carries a witness-schema annotation naming a role label outside \(P_1,\dots,P_6\) with a formula schema \(\psi(C)\) (Definition~\ref{def:annotations}) each of whose instances \(\psi(C')\), at the identifier list of a cluster \(C'\), is a closed formula of the fixed check language containing no atomic \(W_j\) formula;\newline (b) for some cluster \(C\ni r\), an honest audit records \(\psi(C)\) true, \(W_1(C),\dots,W_6(C)\) all false, and \(\psi\) false at some other cluster of \(D\);\newline (c) the same audit records \(\psi(C')\) false for every cluster \(C'\subseteq\Raw{D}\) at which some \(W_i(C')\) holds, by the separation certificate specified in the body of this definition;\newline (d) \(r\) records a closure operation \(c\) on a declared finite order \((X,\le)\), with \(\le\) the reflexive closure of the relation of a declared order record whose law is certified, that is, a total map \(c:X\to X\) recorded in a field of \(r\) with correct certificates, recorded true, whose formulas are obtained, by the steps permitted in the paragraph \emph{What a required certificate must express} of Section~\ref{sec:primitive-role-architecture}, from the prescribed conditions \(\forall x\in X\,(x<_o c(x)\vee x=c(x))\), \(\forall x,y\in X\,((x<_o y\vee x=y)\to(c(x)<_o c(y)\vee c(x)=c(y)))\) and \(\forall x\in X\,(c(c(x))=c(x))\), \(o\) being that order record;\newline (e) a datum field of \(r\) lists finitely many total maps \(X\to X\), among them \(\mathrm{id}_X\) and, for every map, partial-map, transition or refinement record of \(D\) whose graph is defined at every member of \(X\) with values in \(X\), the restriction of that graph to \(X\), and correct certificates, recorded true, whose formulas are obtained in the same way from the prescribed conditions \(\forall k,l<|L|\,\exists j<|L|\,\forall x\in X\,(L_j(x)=L_k(L_l(x)))\) and \(\forall k<|L|\,\exists x\in X\,(c(x)\neq L_k(x))\), \(L\) being that datum field, stating that the list is closed under composition and that \(c\) differs from every listed map; the list then contains the monoid these maps generate, so \(c\) lies outside it \\
\bottomrule
\end{longtable}
\endgroup

Criterion (viii) states the same condition as the paragraph \emph{Covered scope} below and as (R5); criterion (ix-c) states the same condition as (R3); and criterion (v-a) is the first branch of (R1) together with the failure of (i).

\paragraph{Terms used in the label criteria.}
\emph{Declared operations.} The \emph{declared operation} of a structural annotation of
Definition~\ref{def:annotations} is the operation of its kind (a level
assignment, a threshold, a lift, a forgetting, an instrument restriction or an
empirical bridge) together with the fields that the definition of that kind
prescribes for it; an \emph{independent assertion} of such an annotation is a
predicate, comparison, invariant, law or claim field other than those
prescribed fields, or a formula-valued prescribed field that is not part of
its declared operation under the condition of criterion (ix-b).
\par\smallskip\noindent\emph{Alignment and required members.} A cluster \(C\subseteq\Raw{D}\) is \emph{aligned} with \(P_i\), or
\emph{six-role-aligned}, when \(W_i(C)\) holds; it may lack the attachment
data of Definition~\ref{def:threshold-attachment}. A record \(r\in C\) is a \emph{required member} of such a cluster, and is said
to \emph{contribute a required field} to it, when
\(W_i(C\setminus\{r\})\) fails (the condition concerns membership, so a record
without fields, such as a micro-element of the \(W_5\) row, can contribute a
required field), and \(r\) \emph{supplies a witness field}
when it contributes a required field to some aligned cluster or carries a
covered seed. \par\noindent\emph{Required fields and inputs.} For \(r\in C\) and a field \(\varphi\) of \(r\), let \(D[r\setminus\varphi]\) be
the description with \(\varphi\) deleted from \(r\), all else unchanged, and
\(C[r\setminus\varphi]\) the corresponding cluster; the rows of the witness table
define \(W_i\) on every such pair, whether or not it lies in \FATCD{}. The field
\(\varphi\) is \emph{required for \(W_i\) at \(C\)} when \(W_i(C)\) holds and
\(W_i(D[r\setminus\varphi],C[r\setminus\varphi])\) fails, and is an \emph{input to
the check} of \(W_i\) at \(C\) when changing only its value (for a datum
field its recorded datum; for a formula-valued or certificate field its
recorded formula together with, for a certificate, its named inputs; the
recorded truth value of a certificate is not read by recomputation) can change the
recomputed result of some certificate \(\kappa\) carried by a member of \(C\)
whose recorded formula, before the change, expresses, in the sense of the paragraph \emph{What a required certificate must
express} (Section~\ref{sec:primitive-role-architecture}), the
conditions that the \(W_i\) row prescribes
for that certificate (for
\(W_6\), a property of a trace in \(C\) tied to an event in \(C\); for \(W_3\), a
non-reducibility certificate for two routes in \(C\); for \(W_4\), a law of an
order, transition or refinement record in \(C\); \(\kappa\) may be \(\varphi\)
itself only when \(\varphi\)'s recorded formula already expresses such a
condition). For a failed offer on \(C\), a
field \(\varphi\) of \(r\in C\) is an \emph{input to the recorded check} when the
offer's audit lists a certificate carried by a member of \(C\), other than
\((W_i(C),\mathrm{false})\), whose formula expresses a condition prescribed by
the \(W_i\) row and whose recomputed result changes when only the value of
\(\varphi\) changes; when the audit lists \((W_i(C),\mathrm{false})\) alone, no
field is such an input. This field-level notion differs from the record-level
phrase of the paragraph \emph{Witness types}, which governs only cluster
membership; a reference field is never an input to the check. For a covered seed, the required fields are the
fields named in its seed clause of Definition~\ref{def:channel-status}. \par\noindent\emph{Covered scope.} A record \(r\) of \(D\) lies \emph{outside the covered scope} when a field of \(r\) asserts a property whose check is not finite over the records of \(D\): it quantifies over a declared domain for which \(D\) supplies no finite checking range, or uses a declared predicate whose extension \(D\) does not record, and its truth value is not fixed by the recorded values of \(D\) (it is true in some expansion of \(D\) and false in another, in the sense of the paragraph \emph{Certificates, assertion language and laws} of Section~\ref{sec:primitive-role-architecture}; for a formula with free variables: for some assignment \(\alpha\) of its free variables and two expansions of \(D\) whose carriers contain the values \(\alpha\) gives to variables of unlisted sorts, it is true under \(\alpha\) in one and false under \(\alpha\) in the other). In addition, every other predicate, comparison, invariant, law or claim field of \(r\) either has this property or is required for \(W_i\) at a cluster \(C\) aligned with \(P_i\), an input to the check of \(W_i\) at such a cluster, or named in a seed clause that \(r\) carries; the \emph{covered scope} of \(D\) consists of its other records. \par\noindent\emph{Recorded extensions and carriers.} Here \(D\) \emph{records the extension} of a declared predicate \(p\) when at least one complete presentation of \(p\) is supplied, either by a relation record that references the declaring record of \(p\) in a field declaring it the extension of \(p\), its listed tuples being by definition the true tuples of \(p\), provided the argument sorts of \(p\) have recorded finite carriers, or by its declaring predicate record (possibly \(r\) itself) giving an explicit formula of the fixed check language of Section~\ref{sec:primitive-role-architecture} over recorded inputs and declared finite ranges (arguments of rational sort being then admitted and evaluated by that formula); all such presentations of \(p\) must agree on every tuple, and otherwise \(D\) does not record its extension. This depends only on \(D\) and \(p\).\par\noindent\emph{Finite checking ranges.} \(D\) supplies a finite
checking range for, and records the carrier of, a declared domain \(\sigma\)
when \(\sigma\) is declared by an object record listing its members or is the
set of declared records of one kind; a sort declared by an object record
listing no members has neither. Both conditions are decided by inspecting the fields of \(r\) and the declarations of \(D\). The record and its references remain finite and declared. \par\noindent\emph{Outside-scope declarations; examples.} An outside-scope declaration (Definition~\ref{def:channel-status}) additionally requires the honest audit of its claim to list that field in its non-certification field; the outside-scope criterion itself imposes no condition on audits. A finite, audited feature is not outside scope merely because it supplies
none of the six role witnesses. Label~(viii) applies to a record outside the
covered scope: the record itself is finite and belongs to \(D\), but the check
its assertion needs is not finite over the records of \(D\). Examples are an audited metaphysical-causation assertion whose declared
predicate has no recorded extension (Proposition~\ref{prop:causality-closure})
and a proposed intentional or teleological predicate whose extension is
unrecorded (Proposition~\ref{prop:agency-closure}), in each case
provided the asserted property is true in some expansion of \(D\)
and false in another.

\begin{remark}[Optimization and objective data are not excluded]\label{rem:optimization-not-excluded}
Raw features expressing optimization, objective, or selection criteria
pass through the scheme of Definition~\ref{def:residual-scheme}; they
are not dismissed by definition. Under the fixed role meanings
(Remark~\ref{rem:fixed-role-meanings}), typical classifications
include: a predicate or specification record possibly supporting a
\Ptwo{} projection; a map or transition-selection record possibly
supporting a \Pone{} or \Pfour{} projection; a route-comparison or
route-choice record possibly supporting a \Pthree{} projection; an
audit or check record possibly supporting a \Psix{} projection; an
empirical-bridge annotation; or an outside-domain structure when an assertion of the
feature uses an unrecorded predicate or uncarried sort and is true in some expansion of \(D\) and false in another, and every other assertion field meets the additional conditions of criterion (viii). A feature that
fits none of these receives whichever of criteria (iii)--(ix) it meets (for
instance (ix-c) for a record with only datum fields, such as a value tagged
objective), and category~(x) only when none holds.
\end{remark}

\paragraph{The running example as a description.} The example of
Section~\ref{sec:primitive-role-architecture} is a description
\(D_{\mathrm{ex}}\in\FATCD{}\) whose raw records are listed below, with the
references the text of the example fixes and the projection that uses each
record. Every reference points to a declared record, and the audit part
carries a bookkeeping record for each claim, so \(D_{\mathrm{ex}}\) is finite,
closed and audited. Threshold, witness, level, collapse and audit annotations
are attached to the three clusters \(C_1\), \(C_5\) and \(C_6\) described after
the table.
\begin{center}
\footnotesize
\begin{tabularx}{\linewidth}{@{}>{\raggedright\arraybackslash}p{0.3\linewidth}>{\raggedright\arraybackslash}p{0.22\linewidth}>{\raggedright\arraybackslash}X@{}}
\toprule
Record & Relates to & Used by \\
\midrule
four state records \(0,1,2,3\) & & \\
state-set object record \(S_0\) listing \(0,1,2,3\) & state records & declared domain and codomain of \(F\) and of its transition record, domain of \(q\); \Pone{} and \Pfive{} projections \\
object records \(A\), \(B\) and a quotient-object record \(\{A,B\}\) listing them & & \Pfive{} projection \\
map record for \(F\), declared use \emph{update} & \(S_0\) & \Pone{} projection \\
relation record \(E\), of declared use quotient, whose datum fields list the classes \(\{0,1\}\) and \(\{2,3\}\) under class-name fields \(A\) and \(B\) & & \Pone{} and \Pfive{} projections \\
map record for \(q\), declared use \emph{packaging}, declared domain \(S_0\), codomain \(\{A,B\}\) & \(S_0\), state records, quotient-object record & \Pfive{} projection \\
quotient-validity check record for \(q\) and \(E\), checking the two fibers, their equality with the \(E\)-classes (so \(E=\ker q\)), and coverage of \(\{A,B\}\), and that every member of a class of \(E\) lies in \(S_0\) & state, object, \(q\)-map and relation records & \Pfive{} projection \\
comparison record for the split pair \((0,1)\) & map and relation records & \Pone{} projection \\
trace record \(0\to2\to3\) & & \Psix{} projection \\
a replay record carrying only reference fields to the trace and to the map record for \(F\), and a check record carrying reference fields to the trace and to the map record for \(F\) together with the certificate \(\tau_1=F(\tau_0)\wedge\tau_2=F(\tau_1)\), where \(\tau=(0,2,3)\) & the trace and the map record for \(F\) & \Psix{} projection \\
event and provenance records & the trace & \Psix{} projection \\
transition record for \(F\), whose graph lists the same four pairs as the map record for \(F\), carrying in its law field the certificate \((\forall x\in S_0\,(t(x)\in S_0),\mathrm{true})\), \(t\) being this record (the transition law of Section~\ref{sec:primitive-role-architecture} with its definedness conjunct omitted), where \(S_0=\{0,1,2,3\}\) (no separate check record) & map record for \(F\), \(S_0\) and state records & \\
index record \(\mathrm{idx}\) with two datum fields named \(A\) and \(B\) (field names, distinct from the identifiers of the object records \(A\) and \(B\), as are the class-name fields of \(E\)) holding the lists \((0,1)\) and \((2,3)\), with one comparison field recording the bijection \(\beta\) from these two fields to the class fields \(A\) and \(B\) of \(E\), together with its certificate that each recorded value equals that of its \(\beta\)-image, written entrywise with indexing (\(A(\cdot)_k\) denoting entry \(k\) of the field \(A\)): \(\bigwedge_{k<2}(A(\mathrm{idx})_k=A(E)_k\wedge B(\mathrm{idx})_k=B(E)_k)\) & relation record \(E\) & \\
\bottomrule
\end{tabularx}
\end{center}
The split-pair comparison carries the certificate
\((E(0,1)\wedge((0\in\mathrm{dom}F\leftrightarrow1\notin\mathrm{dom}F)\vee(0\in\mathrm{dom}F\wedge1\in\mathrm{dom}F\wedge E(F(0),F(0))\wedge E(F(1),F(1))\wedge\neg E(F(0),F(1)))),\mathrm{true})\); \(F(0)\) and \(F(1)\) are
defined because \(F\) is a total map record on \(S_0\), so this is alternative
(b) of the \(W_1\) row. The quotient-validity check
carries the certificate, recomputed true,
\(\forall u,v\in S_0\,(E(u,v)\leftrightarrow q(u)=q(v))\wedge 0\ne1\wedge q(0)=q(1)\wedge\forall z\in\{A,B\}\,\exists u\in S_0\,(q(u)=z)\wedge\forall u\in|E|\,(u\in S_0)\wedge\forall u\in S_0\,(q(u)\in\{A,B\})\),
the \(W_5\) formula with \(X=S_0\), \(Q=\{A,B\}\), \(x=0\), \(y=1\) (the numerals in \(0\ne1\) replace the record-sorted variables \(x,y\) and denote records). Both certificates name their declared inputs.
No threshold annotation targets a cluster containing the transition record. The \Pone{}
cluster \(C_1\) contains \(F\), \(E\), \(S_0\), the split-pair comparison and their referenced
state records. The \Pfive{} cluster \(C_5\) contains \(S_0\), the four state records, the object
records \(A\) and \(B\), the quotient-object record \(\{A,B\}\), the map record for \(q\), the relation
\(E=\ker q\) and the finite quotient-validity check. The \Psix{} cluster \(C_6\)
contains the event, trace, provenance, replay and check records and the map
record for \(F\). For each
\(i\in\{1,5,6\}\), \(D_{\mathrm{ex}}\) carries a threshold, witness-schema,
level, collapse-case, claim, honest audit and nonclaim annotation assigned to
the displayed active projection, the threshold, witness-schema,
collapse-case, level and nonclaim annotations each targeting exactly
\(C_i\), each of the three level annotations recording the verification level,
so that no projection makes a level crossing (Definition~\ref{def:role-projection}); for \(i=6\) these annotations also include a claim annotation
\(c_\tau\) referencing the event, trace and check records, whose statement is the formula of the check record's certificate (so \(c_\tau\) satisfies the representation condition of Definition~\ref{def:annotations}), with an
honest audit \(a_\tau\) referencing \(c_\tau\), the event, trace, \(F\) and
check records, and certifying precisely those equations from the check
record, both assigned to the active \Psix{} projection, which discharges the
\(\Audit{D}\) conjunct of \(W_6(C_6)\) certified by the \Psix{} audit. It
carries no lift or forgetting annotation and no other annotation. Each threshold annotation carries only its typing, target and label, each level annotation only its typing, target and level, and each claim, collapse-case and nonclaim annotation only datum and reference fields besides its statement and hypothesis fields. For each \(i\in\{1,5,6\}\), the role claim has statement \(W_i(C_i)\).
Its audit references that claim and its named records, lists the
row's content certificates, recomputed true, together with
\((W_i(C_i),\mathrm{true})\), and concludes \(W_i(C_i)\).
For \(i=6\), the latter certificate is recomputed after
\(a_\tau\) is established honest in \(\mathrm{Audit}_0(D_{\mathrm{ex}})\). The
\Pfour{} status names the law-bearing transition record as its seed. The index record references \(E\); its equality-check field records that its
two lists equal, entry by entry, the class lists of \(E\). Recomputing that equality
verifies its presentation-variant label. The table and
these annotations exhaust \(\Raw{D_{\mathrm{ex}}}\). Consequently, for the decomposition record that assigns all these annotations to their projections and names the transition record as the \Pfour{} seed, \(A(\Dec{D_{\mathrm{ex}}})\cup S(\Dec{D_{\mathrm{ex}}})\) contains every feature except the index record, and \(\Residual{D_{\mathrm{ex}}}\) consists of the index record alone, classified as a presentation variant of the relation record; a well-formed record leaving \(c_\tau\) and \(a_\tau\) unassigned has them as further record residuals, each licensed (ix-a).

\subsection{Non-circularity and non-overbreadth}\label{subsec:non-circularity-non-overbreadth}

Two proposition schemas record the definition-theoretic facts used by
the later theorem chain: the domain does not presuppose the answer to
the exact-six question, and it is structurally restrictive enough
for that question to remain substantive.

\begin{proposition}[Non-circularity schema for \FATCD{}]\label{prop:non-circularity}
Membership in \FATCD{} is not defined by requiring a
decomposition/exhaustion record or by requiring any derived role
projection to exist. Concretely, the description \(D_0\) of
Remark~\ref{rem:fatcd-minimality} lies in \FATCD{} and has no derived role
projection, and the description \(D_{\mathrm{norm}}\) of
Remark~\ref{rem:fatcd-minimality} (analysed in
Remark~\ref{rem:recognition-open}) lies in \FATCD{}, and its raw record \(N\)
is licensed only label (x) in every well-formed decomposition record.
\end{proposition}

\begin{proof}
We inspect the clauses that govern membership, namely those of
Definitions~\ref{def:raw-record}--\ref{def:fatcd}, and the residual
classification scheme of Definition~\ref{def:residual-scheme}.
\begin{enumerate}[label=(\alph*),topsep=2pt,itemsep=1pt]
  \item The raw grammar permits a witness-schema annotation to name \(P_i\),
    and honesty checks any role claim a description actually makes. Neither
    such an annotation nor a role claim is required for membership in
    \FATCD{}. Its membership conditions require finite fields, closed
    references, and honest audits for claims present; they require no
    derived \(\pi_i\), role-channel assignment, or decomposition record.
    Thus optional role-naming data do not presuppose a role projection.
  \item No membership clause excludes a candidate seventh-role feature by
    definition. A candidate seventh-role feature is, at the raw-grammar
    level, a record or annotation of one of the declared kinds; nothing in
    Definitions~\ref{def:finite}--\ref{def:audited} filters a feature out on
    the ground that it might resist absorption into a
    \Pone{}--\Psix{} projection. Such a feature is therefore admissible
    whenever its containing description satisfies all membership conditions of \FATCD{}, and is carried forward to
    the residual classification scheme rather than barred at the door.
  \item The residual classification scheme
    (Definition~\ref{def:residual-scheme}) is declarative and carries the
    categories (x) and (xi)---an unresolved candidate seventh-role threat,
    and a recorded seventh-role screen---as explicit failure
    options. At the domain-definition level the scheme records the available
    labels but does not assert that those two categories are empty, so it
    cannot smuggle the exact-six answer into the definition of the domain.
\end{enumerate}
Since none of (a)--(c) ties admissibility to the existence of a role
projection or a six-channel record, membership in \FATCD{} is independent of
the exact-six question, which is what was claimed. For the concrete clause,
\(D_0\) carries no annotation, so no witness-schema annotation exists and no
tuple \(\pi_i(D_0)\) can be formed (Definition~\ref{def:role-projection}); \(D_{\mathrm{norm}}\) lies in \FATCD{}
and its record \(N\) receives label (x) in every well-formed decomposition record: \(N\)
has a predicate field containing a closed formula and is referenced by a claim annotation, so both branches of (ix-c) fail; \(D_{\mathrm{norm}}\) has no map,
relation, path, order, transition, refinement, event, provenance, trace,
replay or check record, so no cluster satisfies any \(W_i\), no record is a
covered seed, every channel is \(\absentrec\), and (i), (ii) and (v) fail; \(N\)
has a predicate field and records no checked equality, isomorphism or level
translation, so it begins no (R4) chain and (iii), (iv) fail; \(N\) is not an
annotation, so (vi), (vii), (ix-a) and (ix-b) fail; \(N\) gives its predicate
by an explicit formula over recorded values, so (viii) fails; and
\(D_{\mathrm{norm}}\) has no witness-schema annotation naming a candidate
outside \(P_1,\dots,P_6\), so (xi) fails
(Remark~\ref{rem:recognition-open}).
\end{proof}

\begin{proposition}[Non-overbreadth schema for \FATCD{}]\label{prop:non-overbreadth}
\FATCD{} is structurally restricted and does not admit arbitrary
syntactic material: every record of a description in \FATCD{} is finite,
typed and reference-closed, every claim it carries is referred to by an honest
audit annotation,
unclaimed records may occur but support no projection without the attachment
data of Definition~\ref{def:threshold-attachment}, and mentioning a role name without the supporting threshold,
witness, level, collapse, nonclaim and bookkeeping data does not by itself produce a
derived role projection. In particular, if no witness-schema annotation of
\(D\) names \(P_i\), then no \(\pi_i(D,C)\) is admissible.
\end{proposition}

\begin{proof}
Let \(D \in \FATCD{}\). We inspect the admissibility conditions of
Definitions~\ref{def:finite}--\ref{def:threshold-attachment} in turn.
\begin{enumerate}[label=(\roman*),topsep=2pt,itemsep=1pt]
  \item By Definition~\ref{def:finite}, \(\Raw{D}\) contains finitely many
    records and annotations, each record or annotation carries finite typed
    fields, and no infinite semantic, probabilistic, topological, geometric,
    causal, optimization, or resource structure is admitted unless finitely
    represented and bookkept as a raw feature. Thus unbounded or
    unrepresented material is excluded at the level of finiteness.
  \item By Definition~\ref{def:closed}, every source, target, endpoint,
    index, transition, order, refinement, trace, event, provenance, replay,
    check, and annotation reference is bound to a declared record of \(D\).
    Hence no dangling or externally defined reference can occur, and the
    syntactic material of \(D\) is self-contained.
  \item By Definition~\ref{def:audited}(e), every claim carried by \(D\)
    is referred to by an honest audit annotation. Promotions, translations,
    lifts, forgetting, visibility and empirical-bridge assertions carry the
    fields required by Definition~\ref{def:audited}(g). Thus a claim lacking
    an honest audit, or one of these annotations lacking its required
    fields, cannot occur in a member of \FATCD{}; an assertion field of an
    unclaimed content record needs no audit.
  \item By Definition~\ref{def:threshold-attachment}, a role projection
    cannot be asserted on a cluster \(C\) until the raw features of \(C\), a
    threshold annotation, a witness schema, a record of collapse cases, a
    level annotation, a nonclaim annotation targeting exactly \(C\), and an
    audit record are all present in \(\Raw{D}\).
    Hence a role name supplied without this attachment data does not, by the
    act of naming alone, produce a projection.
\end{enumerate}
Combining (i)--(iii), every record of \(D\) is finite, typed and
reference-closed, and every claim of \(D\) is honestly audited; a record about
which no claim is made needs no audit (Definition~\ref{def:audited}), and by
(iv) it supports no projection unless the attachment data are present. By (iv), mentioning a role name without the supporting
threshold, witness, level, collapse, nonclaim and bookkeeping data does not by itself
produce a derived role projection. This establishes both halves of the
claim. For the last clause, admissibility of \(\pi_i(D,C)\) requires a
witness-schema annotation naming \(P_i\) and targeting exactly \(C\)
(Definition~\ref{def:role-projection}).
\end{proof}

\subsection{The scoped exact-six question over \FATCD{}}\label{subsec:scoped-exact-six-question}

We can now state the scoped exact-six question precisely.

\begin{definition}[Scoped exact-six question]\label{def:scoped-exact-six-question}
The \emph{scoped exact-six question} over \FATCD{} asks whether, for every
\(D \in \FATCD{}\), there exists at least one decomposition/exhaustion record
\(\Dec{D}\) whose role-channel structure consists of exactly the six typed
closure-role channels \Pone{}, \Ptwo{}, \Pthree{}, \Pfour{}, \Pfive{},
\Psix{}, each assigned exactly one of the five statuses later formalized in
Section~\ref{sec:decomposition-exhaustion}, such that no kind residual of
\(D\) and no element of \(\Residual{D}\) is licensed label (x) or (xi) of
Definition~\ref{def:residual-scheme} relative to \(\Dec{D}\). No uniqueness is built into this
question. For a class \(\mathcal D\subseteq\FATCD{}\), the affirmative answer
for every \(D\in\mathcal D\) is denoted \(\ScopedExactSix{\mathcal D}\).
Section~\ref{sec:scoped-exact-six} establishes \(\ScopedExactSix{\FATCD_{\mathrm{RH}}}\),
where \(\FATCD_{\mathrm{RH}}\subseteq\FATCD{}\) is the class of descriptions
satisfying the recognition hypothesis introduced in
Section~\ref{sec:upper-bound-classification}
(Definition~\ref{def:recognition-hypothesis}); the recognized forms of
probability, causality and agency
(Propositions~\ref{prop:probability-closure}--\ref{prop:agency-closure}) are
classified without that hypothesis. Label (x) is licensed for every in-scope
assertion-bearing record placed by none of (i)--(ix), for instance an unclaimed
predicate record whose closed formula is \(0+0=1\), or the true formula \(0+0=0\), between rational constants (Remark~\ref{rem:recognition-open}); one object record together with such a predicate record already lies in \FATCD{} but not in \(\FATCD_{\mathrm{RH}}\) (it has no claims, so it is audited; the predicate record has a closed assertion field, carries no seed, contributes to no aligned cluster, begins no (R4) chain, is not an annotation and lies in the covered scope, so (R1)--(R5) fail as for \(N\) in Remark~\ref{rem:recognition-open});
a negative answer for \(D\) therefore need not exhibit any candidate role.
\end{definition}

The unrestricted assertion \(\ScopedExactSix{\FATCD{}}\) fails: the finite
audited normalization description of Remark~\ref{rem:recognition-open} lies
in \FATCD{} and has a residual feature labelled (x) in every well-formed decomposition
record (verified in Remark~\ref{rem:recognition-open}; see also
Proposition~\ref{thm:integrated-stress}(d)). Label (x) there records that no
criterion (i)--(ix) places the normalization predicate; it does not exhibit a
seventh role.

\begin{remark}[Channels, not active projections]\label{rem:channel-vs-activation-reminder}
An affirmative answer to the scoped exact-six question asserts the
existence of six role channels per description, not six active derived
projections per description. The distinction introduced in
Remark~\ref{rem:projection-vs-membership} and formalized in
Section~\ref{sec:decomposition-exhaustion} is preserved throughout.
\end{remark}

\begin{remark}[Relation to the Foundations theorem]\label{rem:fatcd-foundations-relation}
Proving \(\ScopedExactSix{\FATCD_{\mathrm{RH}}}\) is not a re-derivation of the
canonical unavoidability result of Foundations~I~\citep{Tsiokos2026SixBirdsFoundations},
and \(\ScopedExactSix{\FATCD_{\mathrm{RH}}}\) does not follow from that result
alone. The Foundations theorem derives the \Pone{}--\Psix{} vocabulary
under minimal hypotheses of composability with limited interface
access; the scoped exact-six claim asserts that, within the finite
audited typed closure regime of \FATCD{} and under the recognition hypothesis of
Section~\ref{sec:upper-bound-classification}, every admissible description
admits a six-channel decomposition/exhaustion record without a seventh
irreducible typed role in the covered scope. The two results share
vocabulary but address different questions.
\end{remark}
 \section{Six-role lower bounds}\label{sec:lower-bounds}



The exact-six question of
\S\ref{subsec:scoped-exact-six-question} has two directions. The
lower-bound direction asserts that each of \Pone{}, \ldots, \Psix{}
is instantiated in \FATCD{}; the upper-bound classification of
Section~\ref{sec:upper-bound-classification} will assert, under a recognition
hypothesis, that no seventh role is needed in the covered scope. This section establishes the
lower-bound direction by constructing one admissible witness for each
role. The content is occurrence, not exhaustion: lower bounds alone
do not answer the exact-six question, but they rule out the
trivialization in which some role has no admissible witness in the
domain.

\subsection{Lower-bound theorem and common witness requirements}\label{subsec:lower-bound-theorem}

\begin{theorem}[Six-role lower bounds]\label{thm:six-role-lower-bounds}
For each role index \(i \in \{1,\ldots,6\}\), there exists a description
\(D_i \in \FATCD{}\) and a derived active projection
\(\pi_i(D_i)\) witnessing the role channel \(P_i\) in the sense of
Definition~\ref{def:role-projection}. In particular, each of
\Pone{}, \Ptwo{}, \Pthree{}, \Pfour{}, \Pfive{}, \Psix{} is instantiated
by at least one admissible description in \FATCD{}. This is an occurrence
theorem, not an exhaustion theorem.
\end{theorem}

A witness for role \(P_i\) consists of the data specified below,
following the role-projection discipline of
\S\ref{subsec:derived-role-projections} and the threshold-attachment
discipline of \S\ref{subsec:threshold-attachment}.

\begin{definition}[Common witness data for role \(P_i\)]\label{def:common-witness-data}
A \emph{witness of role \(P_i\) in \FATCD{}} is a tuple
\[
\bigl(D_i,\; C_i,\; \Theta_i,\; W_i,\; \mathrm{Level}_i,\;
v_i^{\mathrm{fin}},\; v_i^{\mathrm{cl}},\; v_i^{\mathrm{aud}},\;
n_i,\; \pi_i(D_i)\bigr)
\]
where
\begin{itemize}[topsep=2pt,itemsep=1pt]
  \item \(D_i \in \FATCD{}\) is the ambient admissible description;
  \item \(C_i \subseteq \Raw{D_i}\) is the raw feature cluster on which
    the projection is interpreted;
  \item \(\Theta_i\) is the threshold annotation attached to \(C_i\);
  \item \(W_i\) is the witness schema for the role meaning of \(P_i\)
    fixed in the \emph{Witness types} table of
    Section~\ref{sec:primitive-role-architecture};
  \item \(\mathrm{Level}_i\) is the level assignment in the sense of
    Definition~\ref{def:level-trichotomy};
  \item \(v_i^{\mathrm{fin}}, v_i^{\mathrm{cl}}, v_i^{\mathrm{aud}}\) are
    the statements that \(D_i\) satisfies Definitions~\ref{def:finite},
    \ref{def:closed} and~\ref{def:audited} respectively, each with its proof;
  \item \(n_i\) is the explicit nonclaim record carried by the witness;
  \item \(\pi_i(D_i)\) is the derived role projection of
    Definition~\ref{def:role-projection}.
\end{itemize}
The witness is \emph{admissible} when \(v_i^{\mathrm{fin}}\),
\(v_i^{\mathrm{cl}}\) and \(v_i^{\mathrm{aud}}\) show \(D_i\in\FATCD{}\) and
\(\pi_i(D_i,C_i)\) is admissible in the sense of
Definition~\ref{def:role-projection}, with \(\Theta_i\), the \(W_i\)
witness-schema annotation, \(\mathrm{Level}_i\) and \(n_i\) as its threshold,
witness-schema, level and nonclaim annotations; in particular \(W_i(C_i)\)
holds.
\end{definition}

The proof of Theorem~\ref{thm:six-role-lower-bounds} proceeds role by
role: for each \(i \in \{1,\ldots,6\}\) we exhibit an admissible
tuple of the form in Definition~\ref{def:common-witness-data}. The
six constructions are collected in
\S\ref{subsec:per-role-witnesses}.

\subsection{Per-role witnesses}\label{subsec:per-role-witnesses}

Each construction below specifies a raw feature cluster, threshold
annotation, witness schema, level assignment, derived active
projection, and explicit nonclaim record. The finite-description,
closure, and audit/bookkeeping verifications are routine once these
data are placed in a description built from the raw grammar of
\S\ref{subsec:raw-grammar} and checked against the admissibility
conditions of \S\ref{subsec:conditions}. In the \Pone{} construction,
\(\sim\) is a relation on the source records and \(\sim'\) a relation on
the target records; descent asks whether \(x\sim y\) implies
\(F(x)\sim' F(y)\). In the constructions, \(\mathrm{Ann}_\Theta[t]\) is a
threshold annotation whose recorded label is \(t\), and
\(v_i^{\mathrm{fin}}, v_i^{\mathrm{cl}}, v_i^{\mathrm{aud}}\) are the checks
that \(D_i\) satisfies Definitions~\ref{def:finite}, \ref{def:closed}
and~\ref{def:audited}. \par\emph{Annotations.} In each construction, \(D_i\) includes a claim
annotation naming \(C_i\), \(W_i\), the asserted check result, hypotheses and
scope, and an audit annotation referencing it and the finite evidence used by
the check, together with the threshold, level, witness-schema, collapse-case
and nonclaim annotations as declared records; every endpoint or reference used in a
construction is included as a declared record of that \(D_i\), even where the
abbreviated display names only its principal records. In each construction,
the claim explicitly names \(P_i\) and \(C_i\), and the threshold,
witness-schema, collapse-case, level and nonclaim annotations each have a
target field listing exactly \(C_i\). The verification \(v_i^{\mathrm{aud}}\) is the check that
these entries are honest in the sense of Definition~\ref{def:audited}, not
merely present. \par\emph{Fields.} Each \(D_i\) contains only the displayed witness records, their declared
inputs, and the stated annotations. Every content record either contributes a
required field to an aligned subcluster of \(C_i\), carries a covered \(P_i\)
seed, or is role-neutral data without an assertion of its own, such as a record named
in a declared domain, codomain, source or target field of a required-field
contributor, whether or not a check reads it. No content record
makes an assertion independent of the displayed witness. Relation records
other than \(\mathrm{Real}\) in \(D_2\), and map, path, index and state
records, carry only datum and reference fields; every law,
invariant or comparison is a field of the displayed check, comparison or
law-bearing record. Each threshold annotation \(\mathrm{Ann}_\Theta[t]\) carries
only its typing, target and label \(t\), and each level annotation carries only
its typing, target and level; each claim annotation carries, besides its
statement and hypothesis fields, only datum and reference fields; and each
collapse-case, nonclaim, promotion and translation annotation carries only
datum and reference fields. \par\emph{Audits.} In each \(D_i\), the audit of the \(W_i\) claim
lists the row's content certificates carried by members of \(C_i\), each
recomputed true, and a certificate \((W_i(C_i),\mathrm{true})\); its conclusion
is \(W_i(C_i)\). The latter certificate is recomputed after the content checks
and, for \(i=6\), after \(a_\tau\) is established honest in
\(\mathrm{Audit}_0(D_6)\). The conclusion thus follows from a signed premise,
and the audit is honest.

\begin{proposition}[\Pone{} lower bound]\label{prop:p1-witness}
There exists \(D_1 \in \FATCD{}\) with a derived active projection
\(\pi_1(D_1)\) witnessing \Pone{}.
\end{proposition}

\begin{proof}
Let \(D_1\) carry the finite state set \(S=\{x,y,x',y'\}\) with four distinct
state records and a declared object record \(S\) listing \(x,y,x',y'\), the source equivalence relation listing the classes \(\{x,y\}\), \(\{x'\}\), \(\{y'\}\), and the target equality relation on \(S\), listing the four singleton classes, both relation records
of declared use quotient. Let \(F:S\to S\) be a
map record whose declared use is \emph{update}, with \(F(x)=x'\),
\(F(y)=y'\), \(F(x')=x'\), and \(F(y')=y'\). A comparison record checks
\(x\sim y\) but \(F(x)\not\sim'F(y)\); hence the related source pair witnesses
the descent obstruction. Attach a threshold annotation
\(\Theta_1 = \mathrm{Ann}_\Theta[\text{descent obstruction}]\) to the
cluster \(C_1\) comprising these records, and record a witness-schema annotation naming \Pone{} and the \(W_1\) row;
\(W_1(C_1)\) holds by alternative (b), since the comparison record names
\((x,y)\) and carries the certificate
\((x\sim y\wedge((x\in\mathrm{dom}F\leftrightarrow y\notin\mathrm{dom}F)\vee(x\in\mathrm{dom}F\wedge y\in\mathrm{dom}F\wedge F(x)\sim'F(x)\wedge F(y)\sim'F(y)\wedge\neg F(x)\sim'F(y))),\mathrm{true})\),
whose second disjunct recomputes true. Take \(\mathrm{Level}_1\) on the verification
level, and let \(\pi_1(D_1)\) be the derived active projection
interpreting \(C_1\) via \((\Theta_1,W_1,\mathrm{Level}_1)\). The
finiteness, closure, and audit verifications \(v_1^{\mathrm{fin}},
v_1^{\mathrm{cl}}, v_1^{\mathrm{aud}}\) then attach routinely: the claim has
statement \(W_1(C_1)\); its audit references the claim and every member of \(C_1\), lists the comparison certificate and
\((W_1(C_1),\mathrm{true})\), both recomputed true, and concludes \(W_1(C_1)\).
The descriptions \(D_2\)--\(D_6\) follow the same template; for \(D_6\), the
event-or-trace audit is established honest in \(\mathrm{Audit}_0(D_6)\) before
\((W_6(C_6),\mathrm{true})\) is recomputed. The
nonclaim record \(n_1\) excludes any \Pfive{}-packaging
identification (Remark~\ref{rem:repaired-alignments}) and any claim of
an unconditional descended map in the obstruction case.
\end{proof}

In the running example, \(\sim\) and \(\sim'\) are both the grouping by
\(q\), \(x=0\) and \(y=1\): \(q(0)=q(1)\), while \(q(F0)=B\neq A=q(F1)\).

\begin{proposition}[\Ptwo{} lower bound]\label{prop:p2-witness}
There exists \(D_2 \in \FATCD{}\) with a derived active projection
\(\pi_2(D_2)\) witnessing \Ptwo{}.
\end{proposition}

\begin{proof}
Let \(D_2\) carry the state records \(0,1,2,3\), the candidate class of
realizers; an object record \(S_2\) listing \(0,1,2,3\), declared as the domain
of \(q\) and as the candidate class of \(\mathrm{Real}\); a quotient object record \(\{A,B\}\) and a map record \(q\) of
declared use packaging with \(q(0)=q(1)=A\) and \(q(2)=q(3)=B\), whose kernel
\(E\) the comparison uses; a predicate record \(s\) declared as a
macro-specification with formula \(s(r)\equiv(r=0\lor r=2)\); a relation record
\(\mathrm{Real}=\{(0,s),(2,s)\}\), of declared use \emph{realization}, with
candidate class \(K=\{0,1,2,3\}\) and target \(s\),
whose datum fields record the formula's values true, false, true, false at \(0,1,2,3\) and which carries
the certificate \((s(0),\mathrm{true})\) exhibiting the satisfying realizer
\(0\); and a comparison record carrying the certificate
\((q(0)=q(1)\wedge s(0)\wedge\neg s(1),\mathrm{true})\), certifying that the
\(E\)-class \(\{0,1\}\) meets both the realizer set \(\{0,2\}\) and its
complement \(\{1,3\}\) in \(K\), since \(0\in\{0,2\}\) and \(1\notin\{0,2\}\)
although \(q(0)=q(1)\), so that
realizability is not reducible to quotient-class equality alone; the finite quotient-class comparison data are
fields of this comparison record, and any separately declared class
identifiers belong to \(C_2\). Attach
\(\Theta_2 = \mathrm{Ann}_\Theta[\text{nontrivial representability}]\)
to the cluster \(C_2\) of the records \(0,1,2,3\), \(S_2\), \(q\), \(\{A,B\}\), \(s\), \(\mathrm{Real}\) and the comparison record (the comparison reads no class-identifier record; the object records \(A\), \(B\), which occur only as values of \(q\) and are listed by \(\{A,B\}\), are declared records of \(D_2\) outside \(C_2\), have no assertion field and satisfy (R3)), and record a witness-schema annotation naming \Ptwo{} and the \(W_2\) row,
and \(W_2(C_2)\) holds by the satisfying branch: \(s\) has the single free
variable \(r\), \(\mathrm{Real}\) records \(s\) at each member of \(K\) and
exhibits the realizer \(0\), and the comparison certifies that the class
\(\{0,1\}\) of \(\ker q\), whose domain \(S_2\) contains \(K\), meets both
\(\{0,2\}\) and \(\{1,3\}\). Take
\(\mathrm{Level}_2\) on the verification level, and let \(\pi_2(D_2)\)
be the derived active projection determined by
\((C_2,\Theta_2,W_2,\mathrm{Level}_2)\). The finiteness, closure, and
audit checks then follow from the construction. The
nonclaim \(n_2\) excludes collapse to \Pfive{} quotient equality alone
and excludes arbitrary relation promotion.
\end{proof}

The states and the map \(q\) are those of the running example.

\begin{proposition}[\Pthree{} lower bound]\label{prop:p3-witness}
There exists \(D_3 \in \FATCD{}\) with a derived active projection
\(\pi_3(D_3)\) witnessing \Pthree{}.
\end{proposition}

\begin{proof}
Let \(D_3\) carry a state record \(u\) and an object record \(U\) listing \(u\);
generating moves \(a,b\), recorded as map records of declared use
\emph{other} with declared source and target \(U\) and \(a(u)=b(u)=u\);
two path records \(\rho_1=ab\) and \(\rho_2=ba\), each listing the successive
records \((u,u,u)\), both from \(u\) to \(u\) and each referencing the
comparison record; and a comparison record referencing both \(\rho_1\) and \(\rho_2\) and carrying the
rewrite rule \(aa\to a\) as a rule field (Definition~\ref{def:raw-record}) holding the pair of words \((aa,a)\), the finite monoid \(N=\{e,\bar a,\bar b\}\), its elements recorded as the rational constants \(e=0\), \(\bar a=1\), \(\bar b=2\) in a list field and its multiplication as a finite table field, with
\(ey=y\) and \(xy=x\) for \(x\neq e\), the images \(\varphi(a)=\bar a\) and
\(\varphi(b)=\bar b\), a move-list datum field \(M=(a,b)\), and a certificate whose
formula is the full prescribed \(W_3\)(a) condition, reading \(R,N,e,m,\varphi,M\)
from the comparison record and \(w_1,w_2\) from the route records. The
assignments \(R=((aa,a))\), \(M=(a,b)\), \(w_1=ab\), \(w_2=ba\) specify their
recorded values, not literal substitutions in the certificate. It checks that \(e\in N\), that \(N\) is closed under \(m\) and
associative with identity \(e\), that \(\varphi(a),\varphi(b)\in N\), that \(\varphi(\lambda)=\varphi(\mu)\) for each pair \((\lambda,\mu)\) listed in that datum field (here \(\varphi(aa)=\bar a=\varphi(a)\)), and that
\(\varphi^{m,e}(w_{\rho_1})\ne\varphi^{m,e}(w_{\rho_2})\), where \(w_{\rho_j}\) is the recorded word of \(\rho_j\) (here \(ab\) and \(ba\)) and \(m\) is
the recorded multiplication table. On the displayed table,
\(\varphi^{m,e}(ab)=\bar a\) and
\(\varphi^{m,e}(ba)=\bar b\). Let
\(C_3=\{u,U,a,b,\rho_1,\rho_2,\text{comparison}\}\).
Attach \(\Theta_3 = \mathrm{Ann}_\Theta[\text{route mismatch}]\) to the
cluster \(C_3\), and
record a witness-schema annotation naming \Pthree{} and the \(W_3\) row;
\(W_3(C_3)\) holds because the two generated same-source, same-target route
records are distinguished by a finite-monoid non-reducibility
certificate of the \(W_3\) row. Take \(\mathrm{Level}_3\) on the
verification level, and let \(\pi_3(D_3)\) be the derived active
projection determined by
\((C_3,\Theta_3,W_3,\mathrm{Level}_3)\). The routine verifications
\(v_3^{\mathrm{fin}}, v_3^{\mathrm{cl}}, v_3^{\mathrm{aud}}\) then
follow. The nonclaim \(n_3\) excludes conflation with
\Pone{} route undefinedness and with direct cross-level comparison.
Routes are compared in the category presented by the recorded moves and rules,
not by equality of composite maps: here \(a\) and \(b\) have the same graph on
\(\{u\}\), and the witness certifies only that the recorded rules do not
identify \(ab\) with \(ba\).
\end{proof}

The monoid is the first-move monoid of
Section~\ref{sec:primitive-role-architecture}.

\begin{proposition}[\Pfour{} lower bound]\label{prop:p4-witness}
There exists \(D_4 \in \FATCD{}\) with a derived active projection
\(\pi_4(D_4)\) witnessing \Pfour{}.
\end{proposition}

\begin{proof}
Let \(D_4\) carry index records \(0,\tfrac12,1,2\); object records
\(\{0\},\{1\},\{2\}\) listing one index record each, and object records
\(I=\{0,1,2\}\) and \(I'=\{0,\tfrac12,1,2\}\); order records listing the pairs \((0,1),(0,2),(1,2)\)
on \(I\) and all six pairs of \(0<\tfrac12<1<2\) on \(I'\); transition records
\(t_{01}:\{0\}\to\{1\}\) with \(t_{01}(0)=1\) and \(t_{12}:\{1\}\to\{2\}\) with
\(t_{12}(1)=2\); and a refinement record
for the inclusion \(I\hookrightarrow I'\).
Every endpoint and index reference resolves to a declared record. Each
order, transition and refinement record carries its certificate as a law
field: finite certificates check that both orders are irreflexive and
transitive and each relates at least one pair, that each transition is defined on its declared domain with values
in its declared codomain, and that the refinement maps \(I\) into \(I'\) and preserves order, each certificate being the prescribed law formula of Section~\ref{sec:primitive-role-architecture} for its record. Attach
\(\Theta_4=\mathrm{Ann}_\Theta[\text{nontrivial stage/refinement}]\)
to the cluster \(C_4\) comprising these records, and record a
witness-schema annotation naming \Pfour{} and the \(W_4\) row; \(W_4(C_4)\)
holds because the records carry declared endpoints or indices and a law for
them. Take
\(\mathrm{Level}_4\) on the structural-categorical level, and let
\(\pi_4(D_4)\) be the derived active projection determined by
\((C_4,\Theta_4,W_4,\mathrm{Level}_4)\). The routine verifications
\(v_4^{\mathrm{fin}}, v_4^{\mathrm{cl}}, v_4^{\mathrm{aud}}\) then
follow. The nonclaim \(n_4\) excludes bare labels and arbitrary
relations or transitions without laws.
\end{proof}

In the running example the transition record for \(F\) is present, but no
threshold annotation is attached to it, so the \Pfour{} channel of
\(D_{\mathrm{ex}}\) is \(\belowthr\) (Section~\ref{sec:decomposition-exhaustion}).
Only the attachment data are missing there: the cluster formed by the transition
record of \(F\), \(S_0\) and the four state records already satisfies \(W_4\). The description \(D_4\) carries a threshold, a
witness-schema and a level annotation on a larger law-bearing cluster, a stage
tower \(0<1<2\) with transitions and a refinement law.

\begin{proposition}[\Pfive{} lower bound]\label{prop:p5-witness}
There exists \(D_5 \in \FATCD{}\) with a derived active projection
\(\pi_5(D_5)\) witnessing \Pfive{}.
\end{proposition}

\begin{proof}
Take two distinct declared object records \(x,y\) (the micro-elements), an object record
listing them as the declared domain of \(q\), a declared state record \(*\) and a
quotient object record \(Q\) listing it, the relation record \(E\), of declared use
quotient, relating both, and the map record \(q:\{x,y\}\to Q\), with declared use \emph{packaging}, sending both to \(*\). Writing \(X=\{x,y\}\) for the declared domain, the finite check
record carries the certificate, recomputed true,
\begin{multline*}\forall u,v\in X\,(E(u,v)\leftrightarrow q(u)=q(v))\wedge x\ne y\wedge q(x)=q(y)\\ \wedge\forall z\in Q\,\exists u\in X\,(q(u)=z)\wedge\forall u\in|E|\,(u\in X)\wedge\forall u\in X\,(q(u)\in Q).\end{multline*} Thus
\(x\ne y\), \(q(x)=q(y)\), and the complete quotient-validity condition in the
\(W_5\) row holds. Let \(C_5\) comprise these declared records and this check,
and let \(D_5\) carry them. Attach
\(\Theta_5 = \mathrm{Ann}_\Theta[\text{nontrivial quotient packaging}]\)
to the cluster \(C_5\) comprising these records, and record a
witness-schema annotation naming \Pfive{} and the \(W_5\) row; \(W_5(C_5)\)
holds because distinct raw objects are related and packaged together by a
packaging map with quotient-validity evidence. Take
\(\mathrm{Level}_5\) on the verification level, and let \(\pi_5(D_5)\)
be the derived active projection determined by
\((C_5,\Theta_5,W_5,\mathrm{Level}_5)\). The routine verifications
\(v_5^{\mathrm{fin}}, v_5^{\mathrm{cl}}, v_5^{\mathrm{aud}}\) then
follow. The nonclaim \(n_5\) excludes identity packaging and
witnesses whose identified records are not distinct.
\end{proof}

\begin{proposition}[\Psix{} lower bound]\label{prop:p6-witness}
There exists \(D_6 \in \FATCD{}\) with a derived active projection
\(\pi_6(D_6)\) witnessing \Psix{}.
\end{proposition}

\begin{proof}
Let \(D_6\) consist of state records \(s_0,s_1,s_2\); an event record \(e\)
referencing \(s_0\); a trace record \(\tau\) listing \((s_0,s_1,s_2)\) and
referencing \(e\); a provenance record referencing \(e\) and \(\tau\); a map
record \(t:\{s_0,s_1\}\to\{s_0,s_1,s_2\}\) of declared use other, with declared
object records for its source and target, and \(t(s_0)=s_1\), \(t(s_1)=s_2\); and a
check record referencing \(\tau\) and \(t\) that carries the certificate
\(\tau_1=t(\tau_0)\wedge\tau_2=t(\tau_1)\), recomputed true; together with the annotations listed below. Let \(C_6\) consist of \(s_0,s_1,s_2\), \(e\), \(\tau\), the provenance record, \(t\), the source and target object records of \(t\), and the check record. Record a claim annotation \(c_\tau\) referencing \(e\), \(\tau\) and the check record, whose statement is the formula \(\tau_1=t(\tau_0)\wedge\tau_2=t(\tau_1)\) of the check record's certificate (so \(c_\tau\) satisfies the representation condition of Definition~\ref{def:annotations}), with an audit \(a_\tau\) referencing
\(c_\tau\), \(e\), \(\tau\), the provenance record, \(t\) and the check record and certifying
precisely those equations from the check record; both are assigned to the
active \Psix{} projection; then record the claim naming \Psix{} and
\(C_6\) with a second audit \(a_6\) certifying \(W_6(C_6)\), whose
\(\Audit{D}\) conjunct is discharged by \(a_\tau\). Attach
\(\Theta_6 = \mathrm{Ann}_\Theta[\text{nontrivial audit schema}]\) to
the cluster \(C_6\), and record a
witness-schema annotation naming \Psix{} and the \(W_6\) row; \(W_6(C_6)\)
holds because \(C_6\) contains the event \(e\), the trace \(\tau\), which
references \(e\), and the check record, which references \(\tau\) and carries
\(\tau_1=t(\tau_0)\wedge\tau_2=t(\tau_1)\), exactly the conjunction of
successive-entry replay equalities that the \(W_6\) row requires; and
\(a_\tau\in\mathrm{Audit}_0(D_6)\) refers to \(c_\tau\), \(e\), \(\tau\) and the check
record and concludes that formula. Take \(\mathrm{Level}_6\) on the verification
level, and let \(\pi_6(D_6)\) be the derived active projection
determined by \((C_6,\Theta_6,W_6,\mathrm{Level}_6)\). The routine
verifications \(v_6^{\mathrm{fin}}, v_6^{\mathrm{cl}},
v_6^{\mathrm{aud}}\) then follow. The nonclaim \(n_6\) excludes passive
logs, certificate-only content, and partial audit records.
\end{proof}

Propositions~\ref{prop:p1-witness}--\ref{prop:p6-witness} supply, for
each \(i \in \{1,\ldots,6\}\), an admissible witness tuple of the
form in Definition~\ref{def:common-witness-data}. This proves
Theorem~\ref{thm:six-role-lower-bounds}.

The running example \(D_{\mathrm{ex}}\) of Section~\ref{sec:fatcd-domain}
is an explicit description of this kind for three roles: its split pair
witnesses \Pone{}, its quotient \(\{A,B\}\) with packaging map \(q\)
witnesses \Pfive{}, and its trace with replay and check witnesses \Psix{}.

\begin{remark}[What the lower bounds do and do not establish]\label{rem:lower-bounds-scope}
Theorem~\ref{thm:six-role-lower-bounds} asserts that each of the six
roles is instantiated somewhere in \FATCD{}. It does not prove
exact-six, does not rule out candidate seventh-role threats, does not
prove decomposition, and does not prove uniqueness of any later
decomposition/exhaustion record. It also does not assert that any
single description activates all six channels simultaneously;
Proposition~\ref{prop:channel-vs-activation} (case \(k=6\)) constructs one. Those
questions are handled in
Sections~\ref{sec:upper-bound-classification}
and~\ref{sec:decomposition-exhaustion}.
\end{remark}
 \section{Upper-bound classification under the recognition hypothesis}\label{sec:upper-bound-classification}



The lower bounds of Section~\ref{sec:lower-bounds} rule out the
trivialization in which some primitive role lacks an admissible
witness in \FATCD{}. The upper-bound direction addresses the
complementary question: whether some additional, seventh closure role
is needed in the covered scope. This section shows that, under the
recognition hypothesis of Definition~\ref{def:recognition-hypothesis}, every
kind residual in the covered scope, every bookkeeping annotation, and every record residual of every
well-formed decomposition record, is licensed only labels among (i)--(ix) of
Definition~\ref{def:residual-scheme}, and that the recognized forms of
probability, causality and agency receive such labels without it. The argument proceeds in three steps. First, we state the
alignment rules that constrain how raw features may project to
\Pone{}, \ldots, \Psix{}. Second, we close the three high-priority
candidate seventh-role threats: probability and uncertainty, residual
causality, and residual agency. Third, we record the upper-bound
classification theorem that collects these local classifications into
a single covered-scope statement.

\subsection{Role-projection alignment rules}\label{subsec:alignment-rules}

The role meanings of Definition~\ref{def:primitive-roles} carry
alignment rules constraining which raw features may be interpreted as
active projections of each role. These rules refine the
alignments of Remark~\ref{rem:repaired-alignments} and form the
discipline under which the classification theorem is proved.

\begin{proposition}[Role-projection alignment]\label{prop:alignment-rules}
Under honest bookkeeping and the witness discipline of
Definition~\ref{def:role-projection}, each role projection is
constrained as follows.
\begin{itemize}[topsep=2pt,itemsep=1pt]
  \item \Pone{}: admissible only when the candidate cluster carries a map or partial-map record of declared use update, source and target relation records of declared use quotient, and a certificate of alternative (a) or (b) of the \(W_1\) row. In particular, a cluster with no map or partial-map record
    of declared use update, or no relation record of declared use quotient,
    does not satisfy \(W_1\).
  \item \Ptwo{}: admissible only when the candidate carries a macro
    specification, a realization relation, specification-satisfaction
    data, and evidence of certified unrealizability or failure of
    representability over the declared packaging kernel; when a
    satisfying realizer exists, the candidate must also carry a comparison
    with the classes of the kernel of a declared packaging map, whose
    declared domain contains the declared candidate class, certifying
    that some class of that kernel meets both the realizer set and its
    complement in the candidate class.
    Macro-level predicates or objectives count only insofar as they are
    treated as content to be represented. When no realizer exists, the candidate class must be the full declared domain of a map record of the cluster of declared use packaging or update, and the cluster must record a true evaluation of the specification at a declared same-kind record outside that class. Mere quotient equality and
    arbitrary relations without representability structure are not
    admissible as \Ptwo{} projections.
  \item \Pthree{}: admissible only when the candidate carries a route
    grammar (generating moves and all rewrite rules recorded in \(D\)), generated same-source, same-target route terms, and
    route-comparison data with a non-reducibility certificate of either
alternative in the \(W_3\) row. Generic transitions, generic causal arrows, generic
    path records, and route undefinedness arising from a partial update
    map are not admissible as \Pthree{} projections unless they are
    two defined same-source, same-target route records carrying a
    non-reducibility certificate of either alternative in the \(W_3\) row.
  \item \Pfour{}: admissible only when the candidate carries a stage,
    order, transition or refinement record (Definition~\ref{def:raw-record})
    with a law in the sense of the witness table (for a transition record,
    a certificate that its table is defined on its declared domain with
    values in its declared codomain suffices). Bare labels, arbitrary
    relations without order, transition, or refinement law, quotient
    packaging, and audit bookkeeping are not admissible as \Pfour{}
    projections.
  \item \Pfive{}: admissible only when the candidate carries a quotient
    relation, a packaging map, the identification of distinct
    micro-elements, an object record declared as the quotient object, and
    quotient-validity evidence. Generic structural paths, derivations,
    topologies or geometries merely because they have paths, and causal
    chains are not admissible as \Pfive{} projections unless quotient
    or packaging data is present.
  \item \Psix{}: admissible only when the cluster carries a declared event,
    a trace record, a record tied to the event that is that trace record or a
    provenance record also tied to that trace record, a replay or check record tied to the trace carrying a certificate of a replay or provenance property of
    the trace, and \(\mathrm{Audit}_0(D)\) contains an entry \(a\), not necessarily a cluster member, referring to a claim
    annotation, to these event and trace records and to that replay or check
    record, whose conclusion is the formula of that certificate.
    Passive logs alone, arbitrary metadata, and certificate-only
    content without audit/provenance/replay/check structure are not
    admissible as \Psix{} projections.
\end{itemize}
\end{proposition}

\begin{proof}
By Definition~\ref{def:role-projection}, an admissible projection
\(\pi_i(D,C)\) requires \(W_i(C)\), and \(W_i\) is defined by the row for
\(P_i\) of the \emph{Witness types} table in
Section~\ref{sec:primitive-role-architecture}. Each constraint below follows
by reading off the record kinds that row requires; the table is the
definition and this proposition states its consequences for candidate
features. We go through the roles in turn.

For \Pone{}, the fixed meaning is update-vs-quotient compatibility and
descent obstruction. The witness type for that meaning is exactly a
selected update together with source and target relation records of declared use quotient and a
certificate of descent compatibility, descent obstruction, or
update-vs-quotient mismatch; absent such a witness the candidate
exhibits no descent datum to compare, so it cannot be an active \Pone{}
projection. In particular generic semantics, generic predicate meaning,
generic selection criteria, and generic agency carry no descent witness
and are inadmissible, which is the alignment of
Remark~\ref{rem:repaired-alignments}.

For \Ptwo{}, the fixed meaning is macro-specification representability.
Its witness type is a macro specification, a realization relation,
specification-satisfaction data, and a certificate of unrealizability or of
failure of representability over the declared packaging kernel; a macro-level predicate or objective
enters only when treated as content to be represented, since only then
does it instantiate the representability relation. Mere quotient
equality is \Pfive{} data, an arbitrary relation lacking
representability structure supplies no realization relation, and
empirical observations lacking level and threshold data fail the witness
discipline; none is an admissible \Ptwo{} projection.

For \Pthree{}, the fixed meaning is enriched route/coherence mismatch,
and \Pthree{} is not a directionality certificate. Its witness type is
a route grammar (generating moves and all rewrite rules recorded in \(D\)), generated same-source same-target route terms,
and route-comparison data carrying a non-reducibility certificate of either
alternative in the \(W_3\) row. A candidate without such a route
comparison has nothing of the \Pthree{} type to compare. Generic
transitions and generic causal arrows are \Pfour{} data, generic path
records are not route-comparison data, and route undefinedness arising
from a partial update map is \Pone{} content (Table~\ref{tab:boundary-matrix},
row \Pone{}/\Pthree{}); none projects to
\Pthree{} unless the route/coherence threshold is explicitly present and
the comparison is specifically of the \Pthree{} type. This is the
alignment that generic transition is not \Pthree{}.

For \Pfour{}, the fixed meaning is stage, order, transition, and
refinement structure. Its witness type is a stage, order, transition or
refinement record carrying a law beyond bare labels; staged dynamics and
regime ordering qualify when recorded in one of these kinds. Bare labels, arbitrary relations without an order,
transition, or refinement law, quotient packaging (which is \Pfive{}
data), and audit bookkeeping (which is \Psix{} data) supply no such law
and are inadmissible as \Pfour{} projections.

For \Pfive{}, the fixed meaning is quotient packaging. Its witness type
is a quotient relation, a packaging map, the identification of distinct
micro-elements, an object record declared as the quotient object, and
quotient-validity evidence. Generic structural paths, derivations,
topologies or geometries that merely carry paths, and causal chains
supply no quotient relation or packaging map, so they are inadmissible
absent quotient or packaging data; this is the alignment that
generic path or derivation is not \Pfive{}.

For \Psix{}, the fixed meaning is audit, provenance, replay, and
checkability. Its witness type is a declared event, a trace record tied to it directly or through a provenance record, a replay or check record carrying a certified replay or provenance property of the trace, and an entry of \(\mathrm{Audit}_0(D)\) concluding that property. Passive logs without check
structure, arbitrary metadata, and certificate-only content lacking
audit/provenance/replay/check structure supply no checkable audit datum
and are inadmissible as \Psix{} projections.

In each case the admissibility constraint is forced by the fixed role
meaning through the witness discipline: a projection is admissible only
with the witness type that the meaning of the role fixes. Formally, each
bullet is the contrapositive of the row conjunct it names: \(W_1\) requires an
update map and quotient relations; \(W_2\) a macro-specification and a
realization relation; \(W_3\) two defined parallel routes and a certificate;
\(W_4\) a law-bearing record; \(W_5\) a quotient relation, a packaging map and a
quotient object; and \(W_6\) an event, a trace, a certified check and an
\(\mathrm{Audit}_0\) entry. The remaining sentences gloss role meanings and are
not used.
\end{proof}

\begin{remark}[Three excluded alignments, held fixed]\label{rem:three-repaired-alignments}
Proposition~\ref{prop:alignment-rules} fixes three alignments that
are held throughout: generic transition data does not project to
\Pthree{}, generic path or derivation data does not project to
\Pfive{}, and generic semantic or interpretation data does not project
to \Pone{}. These are carried into every step of the classification
theorem below.
\end{remark}

\subsection{Probability, causality and agency}\label{subsec:threat-closures}

Three classes of candidate seventh-role feature require particular
scrutiny: probability and uncertainty, residual causality in the
sense of the interventional and counterfactual
hierarchy~\citep{Pearl2009}, and residual agency in the planning
sense~\citep{Bratman1987}. Each is classified in a proposition. In the titles of
Propositions~\ref{prop:probability-closure}--\ref{prop:agency-closure},
\emph{closure} of a threat means only that each listed form is licensed a label
among (i)--(ix); it is unrelated to closure operations, to the condition
\(\Closed{D}\), or to the closure content of
Definition~\ref{def:recognition-hypothesis}.
In every case the candidate is classified under the alignment rules
of Proposition~\ref{prop:alignment-rules} together with the residual
scheme of Definition~\ref{def:residual-scheme}; no candidate is
dismissed by definition. These three come first because they are the
most natural candidates for a seventh role; the general classification
theorem follows. A \emph{probability or uncertainty feature} is a residual
feature with a field of declared probability, weight, kernel, confidence or
error type, or an empirical-bridge annotation whose identified bridge claim has
such a field in its \(\Audit{D}\) entry; a \emph{residual-causality feature} is one with a field of declared
intervention, counterfactual or causal-dependence type; and a
\emph{residual-agency feature} is one with a field of declared goal,
objective, plan, policy or choice type, or an empirical-bridge annotation
whose identified bridge claim has such a field in its \(\Audit{D}\) entry.
Value-type tags enter no label criterion (Definition~\ref{def:raw-record}), and
the proofs below case only on content; hence each proposition below holds for
every kind residual whose content takes one of its forms, whether or not it
carries such a tag. Each recognized form is specified so that a criterion
among (i)--(ix) applies to it; the propositions list which presentations of
probabilistic, causal and agentive content the witness rows, seed clauses
and annotation kinds already accommodate.

\paragraph{The \(P_i\)-form.}\label{par:pi-form} In these propositions, a feature \(r\) has \emph{\(P_i\)-form} when \(r\) is a
record of a kind occurring in the \(W_i\) row of the witness table,
contributes a required field to a finite cluster \(C\) satisfying \(W_i(C)\),
and every predicate, comparison, invariant, law or claim field of \(r\) is
required for \(W_i\) at \(C\) or is an input to its check
(Section~\ref{subsec:residual-classification}); or when \(r\) carries a covered
\(P_i\) seed and each of its predicate, comparison, invariant, law or claim fields is a named seed field of that clause. Such
a feature satisfies criterion (i) when that criterion holds, and criterion (v)
otherwise. In the bullets below, \(P_i\)-aligned content is
labelled (i) or (v); a bullet does not assert that the decomposition record
names an active projection for it. A seed alone is not an active projection.
\par\emph{Examples.} For instance, a transition record whose fields are its endpoints, a kernel
declared with values in the distributions on its state set, and the law that
each row is such a distribution has \(P_4\)-form: in an aligned cluster in which it is the only law-bearing
record, deleting the row law leaves the cluster without a law, so \(W_4\)
fails. By contrast, the predicate \(N\) in the original \(D_{\mathrm{norm}}\) of
Remark~\ref{rem:recognition-open} has no \(P_i\)-form; neither does a
standalone transition record without a certified law or covered seed. For a
composite form, the feature is carried by a finite set \(R_f\) of records, and
each record of \(R_f\) has the \(P_i\)-form of its own component role (so, by
the second sentence of this paragraph, it satisfies criterion (i) or (v), and
also (ii) when it contributes required fields to two named active clusters and every
predicate, comparison, invariant, law or claim field is required for, or an
input to the check of, either cluster's witness predicate).
\par\emph{Composite forms.} For a composite form, ``classifies as a composite'' is a derived
reading of the finite record set \(R_f\), not an additional residual
label: \(\operatorname{Classify}\) still labels each member of \(R_f\).
Label (ii) applies to a member only under the two-cluster condition just
stated.
\par\emph{How the propositions proceed.} Each proposition takes a kind residual of \(D\) whose content
is of the named kind --- probability or uncertainty; causation in the
interventional and counterfactual sense; agency in the planning sense --- and
passes it through the residual scheme under the alignment rules. Its proof
checks, for each recognized form, that the feature supplies the witness type
of the role it is assigned to, or is an annotation or an outside-domain
structure; the causality and agency propositions also show that one role is
unavailable: \Pfive{} without quotient packaging, and \Pone{} without an
update-vs-quotient descent witness. A stochastic transition record carrying a certified row law, with no assertion
fields beyond the named fields of its covered \(P_4\) seed, has \(P_4\)-form
whether or not attachment data are present. Relative to a fixed well-formed
\(\Dec{D}\), it satisfies criterion (i) if that criterion holds, and criterion
(v) otherwise; an uncertainty record with provenance and replay data
has \(P_6\)-form when each of its predicate, comparison, invariant, law or claim
fields is required for, or an input to, a \(W_6\) check, or is a named field of
its seed; an empirical-bridge annotation attaching a measurement error to a substrate observable is licensed label (vii) when it meets criterion (vii), the error record itself being classified by its own fields; and a
confidence level recorded as the threshold predicate of a threshold annotation
with no other field besides its typing, target and label fields, all of whose
read records satisfy (R1), is an activation-threshold refinement. A normalization predicate with no free variable, such as the
one in \(D_{\mathrm{norm}}\) of Remark~\ref{rem:recognition-open}, is not a
\Ptwo{} seed merely because it states that finitely many weights sum to
one. That remark shows that this
specified predicate receives label (x), and discusses conservation laws and
emergent obstructions separately.
A record of \(P_i\)-form satisfies alternative (R1) of
Definition~\ref{def:recognition-hypothesis}, the threshold and bridge forms below
satisfy (R2), and the outside-domain forms satisfy (R5); the propositions below
are thus the cases, for these three classes, of the pointwise clause of
Theorem~\ref{thm:upper-bound-classification}, and assume nothing about the other
features of \(D\).

\begin{proposition}[Closure of probability and uncertainty]\label{prop:probability-closure}
\textup{(a)} Every kind residual (Definition~\ref{def:residual-scheme}), in particular every probability or uncertainty feature, of a description
\(D \in \FATCD{}\) whose content takes one of the recognized forms below,
where a form naming a \(P_i\) projection or unabsorbed \(P_i\)-aligned content
requires the feature to have \(P_i\)-form (paragraph \emph{The \(P_i\)-form}
before Proposition~\ref{prop:probability-closure}),
is licensed, relative to every well-formed decomposition record \(\Dec{D}\), at least one of the labels of
\begin{itemize}[topsep=2pt,itemsep=1pt]
  \item \(P_2\)-aligned content, label (i) or (v), when the feature has \(P_2\)-form (for
    instance a macro-specification with one free variable and a realization
    relation whose recorded evaluations take both truth values);
  \item \(P_4\)-aligned content, label (i) or (v), when the feature has \(P_4\)-form (for
    instance a stochastic transition record carrying its row law);
  \item \(P_6\)-aligned content, label (i) or (v), when the feature has \(P_6\)-form (for
    instance an audited uncertainty-bookkeeping record);
  \item an activation-threshold refinement, when the feature is a threshold
    annotation targeting a cluster named for an active projection, with a
    confidence level recorded as its threshold predicate
    (Definition~\ref{def:threshold-attachment}) as its only field besides
    typing, target and label, and that predicate reads at least one record, every record whose values it reads is listed in the annotation's target field, and each such record satisfies (R1); relative to a record that names its target cluster for
    no active projection, the annotation satisfies criterion (ix-b);
  \item an empirical-bridge annotation, when the feature is itself an empirical-bridge annotation identifying an admissible bridge claim (Definitions~\ref{prop:honest-bookkeeping} and~\ref{prop:empirical-bridge}) whose \(\Audit{D}\) entry carries the five items of Definition~\ref{prop:empirical-bridge}, with no fields beyond its typing, targets and that identification;
  \item an outside-domain structure, label (viii), when an assertion field of the feature uses a declared predicate whose extension \(D\) does not record or quantifies over a sort whose carrier \(D\) does not record, the asserted property is true in some expansion of \(D\) and false in another, and every other assertion field of that record either also meets the outside-scope condition or is required for \(W_i\) at a cluster aligned with \(P_i\), an input to its check at such a cluster, or named in a seed clause that the record carries (for instance the \(\mathrm{IsDist}\) presentation of Remark~\ref{rem:recognition-open});
  \item label (i), (ii) or (v), as a member of a composite of the foregoing
    projections, when the feature is one of a finite set of records each of which has the \(P_i\)-form of one of the roles above (so, by the paragraph
    \emph{The \(P_i\)-form}, it satisfies criterion (i) or (v), and also (ii)
    when it contributes required fields to two named active clusters and every
    predicate, comparison, invariant, law or claim field is required for, or an
    input to the check of, either cluster's witness predicate);
  \item unabsorbed role-aligned content, label~(v) of
    Definition~\ref{def:residual-scheme}, when it has \(P_i\)-form but
    criterion (i) does not hold for it, for instance because the cluster lacks
    the attachment data of Definition~\ref{def:threshold-attachment}.
\end{itemize}
A content record targeted by such an empirical-bridge annotation is not thereby classified; it is classified by its own fields and is licensed label (ix-c) exactly when either branch of that criterion holds.
Each case licenses a label among (i)--(ix) of
Definition~\ref{def:residual-scheme}, so neither (x) nor (xi) is licensed for
the feature. Probability or uncertainty content of any of these recognized forms therefore
requires no seventh role; content of none of these forms is governed by the
recognition hypothesis (Definition~\ref{def:recognition-hypothesis}).
\textup{(b)} Let a content record of \(D\) have as its only assertion field the formula \(\bigwedge_k v(w_k)\ge0\wedge\sum_k v(w_k)=1\) over value fields of pairwise distinct declared records \(w_1,\dots,w_n\), none of which is the replaced record, suppose the recorded values \(v(w_k)\) are nonnegative and sum to one (so the recomputed row check is true), and let \(D'\) be the description obtained from \(D\) by replacing that
feature, with every incoming reference and affected annotation retargeted in a
type-correct way and with the resulting description \(D'\) in \FATCD{} (in particular, every realization relation still lists exactly the members of its candidate class at which its target's formula is true in \(D'\)), by a transition record \(K:\{*\}\to\mathrm{Dist}(\{w_1,\dots,w_n\})\) (here \(K\) is a transition record, not a candidate class)
whose declared source is an object record \(\{*\}\) listing a declared state
record \(*\), whose declared target is \(\mathrm{Dist}(X)\) for an object record
\(X\) listing \(w_1,\dots,w_n\), recorded by a codomain reference to \(X\) and a
datum field marking the codomain as distribution-valued (paragraph
\emph{Laws}), whose unique row is the list of terms
\(v(w_1),\dots,v(w_n)\), and whose law certificate is
\((\bigwedge_k v(w_k)\ge0\wedge\sum_k v(w_k)=1,\mathrm{true})\); the records
\(*\), \(\{*\}\) and \(X\) are added to \(D'\).
Then \(K\) is a covered \Pfour{} seed of \(D'\) and has \(P_4\)-form, so
relative to every well-formed decomposition record of \(D'\) it is licensed
label (i) or (v), and neither (x) nor (xi) (its only assertion field is the
law certificate, a named seed field of clause (4), so \(K\) has \(P_4\)-form by
the second branch of the paragraph \emph{The \(P_i\)-form}); the added records \(*\),
\(\{*\}\) and \(X\) have no assertion field and satisfy (R1) or (R3). The \(w_i\) and their values are
unchanged, source and target declarations are added, and the claim and audit
annotations of the replaced feature are retargeted to the row-law check of
\(K\); in \(D\) itself the feature is not thereby classified.
\end{proposition}

\begin{proof}
Let \(r\) be a probability or uncertainty feature (a kind residual) of
\(D \in \FATCD{}\), and let it pass through the residual scheme of
Definition~\ref{def:residual-scheme} under the alignment rules of
Proposition~\ref{prop:alignment-rules}. Role-aligned classifications below require the \(P_i\)-form specified in the
proposition. Such a record satisfies criterion (i) when that criterion holds and
criterion (v) otherwise. Composite forms are determined record by record;
criterion (ii) applies under the two-cluster condition of the paragraph
\emph{The \(P_i\)-form}. We case on the content of \(r\).

If the probabilistic content is a macro specification or a
representability or non-representability condition on a description that
contributes a required field to a cluster satisfying \(W_2\), with its
realization relation and certificate, then it supplies the \Ptwo{}
witness, so it has \(P_2\)-form and receives label (i) or (v) (paragraph
\emph{The \(P_i\)-form}). If instead the content is a stochastic transition,
regime, or staged structure carrying a valid stage, order, transition,
or refinement law and has \(P_4\)-form, then it receives label (i) or (v); the
row law of a kernel is such a law, as in the example before this proposition,
and a normalization assertion that is not required for \(W_4\) is carried by a
separate predicate record, classified in its own right
(Remark~\ref{rem:recognition-open}). If the content is an audited record of uncertainty
bookkeeping --- event reference, provenance, replay or check structure
--- it has \(P_6\)-form. A \(P_6\)-form record receives (i) when its named
active cluster carries the certified trace check and honest audit required by
\(W_6\); a covered seed without that full witness receives (v). If the
feature is a threshold annotation of a cluster named for an active projection
whose only field besides typing, target and label is a threshold predicate recording a confidence level
that decides when the projection is attempted, and that predicate reads at least one record, every record whose values it reads is listed in the annotation's target field, and each such record satisfies (R1), then it adds no role content and receives
label (v-b): it is an activation-threshold
refinement. Relative to a record that names its target cluster for no active projection,
the annotation is a structural annotation whose fields are its typing, target
and threshold predicate, with no independent role-theoretic assertion, and
(v), (vi) and (vii) do not apply, so it satisfies criterion (ix-b). If it is an empirical-bridge annotation whose bridge claim is admissible and has an audit
entry with the five items of Definition~\ref{prop:empirical-bridge}, and with
no fields beyond its typing, targets and that identification, it
receives (vii). A content record targeted by such an annotation is not thereby classified; it is classified by its own fields and is licensed label (ix-c) exactly when either branch of that criterion holds. Finally, if several role aspects occur together, take their finite
record set \(R_f\). By the composite-form hypothesis each member has
its component role's \(P_i\)-form, so it satisfies criterion (i) or (v), and
also (ii) under the two-cluster condition stated there;
the set has the composite reading of those classifications. These are the recognized forms of probability and uncertainty content, and
each reduces to the six roles as stated. A probability feature whose content
has none of these forms is not classified by this proposition. In particular,
the bare normalization predicate of Remark~\ref{rem:recognition-open} receives
label (x); other such constraints must be checked against the recognition
alternatives of Definition~\ref{def:recognition-hypothesis}.
A feature whose probabilistic content lies outside the covered scope is classified as an outside-domain structure by Definition~\ref{def:residual-scheme}, not absorbed by default: in the outside-domain case, the asserted property is true in some expansion of \(D\) and false in another, and together with the stated condition on the remaining assertion fields this satisfies criterion (viii). The
closure is therefore scoped, as stated.

For (b), the row of \(K\) is a map from the declared one-point
domain to the declared set \(\mathrm{Dist}(\{w_1,\dots,w_n\})\), and the finite
check of nonnegativity and sum one certifies that its value lies in that
codomain. This is a law in the sense of the witness table, so \(K\) is a
law-bearing transition record, which is a covered \Pfour{} seed by
Definition~\ref{def:channel-status}. Its declared source and target records are fixed by reference fields, and the
value fields of the \(w_k\) are read by that check; none of these records carries
an assertion of its own.
Reference closure and honest bookkeeping of \(D'\) are hypotheses of the
clause; the argument uses only the row of \(K\) and its law certificate.
\end{proof}

\begin{proposition}[Closure of residual causality]\label{prop:causality-closure}
Every kind residual (Definition~\ref{def:residual-scheme}), in particular every residual-causality feature, of a description
\(D \in \FATCD{}\) whose content takes one of the recognized forms below,
where a form naming a \(P_i\) projection or unabsorbed \(P_i\)-aligned content
requires the feature to have \(P_i\)-form (paragraph \emph{The \(P_i\)-form}
before Proposition~\ref{prop:probability-closure}),
is licensed under the residual scheme, relative to every well-formed
decomposition record \(\Dec{D}\), at least one of the labels of
\begin{itemize}[topsep=2pt,itemsep=1pt]
  \item \(P_4\)-aligned content, label (i) or (v), as the primary classification when the
    feature has \(P_4\)-form (for instance a directed transition or
    refinement record with its law);
  \item \(P_3\)-aligned content, label (i) or (v), only when the feature has \(P_3\)-form and is specifically a
    route/coherence comparison between generated same-source,
    same-target routes;
  \item \(P_6\)-aligned content, label (i) or (v), when the feature has \(P_6\)-form;
  \item label (i), (ii) or (v), as a member of a composite of \Pfour{} and
    \Psix{}, when directed dependence and provenance audit are carried by a
    finite set of records containing the feature, each of which has
    \(P_4\)-form or \(P_6\)-form (so it satisfies criterion (i) or (v), and also (ii)
    when it contributes required fields to two named active clusters and every
    predicate, comparison, invariant, law or claim field is required for, or an
    input to the check of, either cluster's witness predicate);
  \item an outside-domain structure, when an assertion field of the feature (for instance a
    metaphysical-causation assertion) uses a declared predicate whose extension \(D\) does not record, the asserted property is true in some expansion of \(D\) and false in another, and every other assertion field of that record either also meets the outside-scope condition or is required for \(W_i\) at a cluster aligned with \(P_i\), an input to its check at such a cluster, or named in a seed clause that the record carries;
  \item unabsorbed role-aligned content, label~(v) of
    Definition~\ref{def:residual-scheme}, when it has \(P_i\)-form but
    criterion (i) does not hold for it, for instance because the cluster lacks
    the attachment data of Definition~\ref{def:threshold-attachment}.
\end{itemize}
Classification as \Pfive{} is inadmissible in the absence of quotient
packaging. Each case licenses a label among (i)--(ix) of
Definition~\ref{def:residual-scheme}, so neither (x) nor (xi) is licensed for
the feature. For these recognized forms, the covered content requires no
seventh role; the outside-domain case receives label (viii) without a claim
about roles beyond the covered scope. None of these classifications uses the
recognition hypothesis (Definition~\ref{def:recognition-hypothesis}); content
of none of these forms is governed by it. An interventional mechanism is covered
when recorded as a law-bearing transition or refinement record of \(P_4\)-form, for instance one whose only assertion field is its law certificate (first bullet); a further assertion field serving no witness can leave it licensed only (x); a
counterfactual assertion is covered only when it takes one of the listed forms,
and otherwise falls under the recognition hypothesis.
\end{proposition}

\begin{proof}
Let \(r\) be a residual-causality feature (a kind residual) of
\(D \in \FATCD{}\), causal content being taken from the hierarchy of~\citet{Pearl2009}, and pass
\(r\) through the residual scheme of Definition~\ref{def:residual-scheme} under
Proposition~\ref{prop:alignment-rules}. The forms of the proposition cover
interventional mechanisms recorded as transitions of \(P_4\)-form; counterfactual assertions
not of these forms fall under the recognition hypothesis. A causal feature presents a
directed dependence, so we case on how that directedness is carried.
Role-aligned and composite forms are labelled as in the paragraph
\emph{The \(P_i\)-form} before Proposition~\ref{prop:probability-closure}: each
record receives (i) or (v), and (ii) under the two-cluster condition stated
there.

If \(r\) has \(P_4\)-form as a directed transition, refinement, or
staged-dependence law, it receives label (i) or (v) by the paragraph
\emph{The \(P_i\)-form}; this is the primary classification. The record
kind and its law alone do not establish the assertion-field guard
defining \(P_4\)-form. One case in which directedness routes elsewhere is when \(r\) is
specifically a route/coherence comparison between generated
same-source, same-target routes carrying a non-reducibility certificate of
either alternative in the \(W_3\) row; only then does it supply the \Pthree{}
witness, have \(P_3\)-form and receive label (i) or (v). We must not assign a
generic causal arrow to \Pthree{}, since by the alignment rules a
generic transition is not a route comparison of the \Pthree{} type. If
\(r\) is instead a provenance or dependency audit --- event reference
with a provenance or trace record --- it has \(P_6\)-form. A \(P_6\)-form
record receives (i) when its named active cluster carries the certified trace
check and honest audit required by \(W_6\); a covered seed without that full
witness receives (v). When directed dependence and provenance audit occur together, take
their finite record set \(R_f\). Its \(P_4\)-form and \(P_6\)-form
members are licensed (i) or (v) individually by their \(P_i\)-form; together they have
the stated composite reading.
Finally, if \(r\) records a metaphysical-causation assertion whose declared
predicate has no recorded extension in \(D\), its truth is not certified by
the finite checks of \(D\). In the outside-domain case, the asserted property is true in some
expansion of \(D\) and false in another. Together with the stated
condition on the remaining assertion fields, this satisfies
criterion (viii).

Classification as a \Pfive{} projection is not available here: by
Proposition~\ref{prop:alignment-rules} a \Pfive{} projection requires a
quotient relation and packaging map, and a causal feature carrying no
quotient packaging supplies none; the directedness of a causal chain is
not packaging data. Hence the \Pfive{}-inadmissibility caveat holds.
The cases above are the recognized forms of residual-causality content; a
causality feature whose content is of none of these forms is not classified here
and falls under the recognition hypothesis
(Definition~\ref{def:recognition-hypothesis}). The closure is scoped and makes no
scope-free claim.
\end{proof}

\begin{proposition}[Closure of residual agency]\label{prop:agency-closure}
Every kind residual (Definition~\ref{def:residual-scheme}), in particular every residual-agency feature, of a description
\(D \in \FATCD{}\) whose content takes one of the recognized forms below,
where a form naming a \(P_i\) projection or unabsorbed \(P_i\)-aligned content
requires the feature to have \(P_i\)-form (paragraph \emph{The \(P_i\)-form}
before Proposition~\ref{prop:probability-closure}),
is licensed under the residual scheme, relative to every well-formed
decomposition record \(\Dec{D}\), at least one of the labels of
\begin{itemize}[topsep=2pt,itemsep=1pt]
  \item label (i), as a constituent of an active \Ptwo{}, \Pfour{} or \Psix{}
    projection, when the feature has the
    corresponding \(P_i\)-form and \(\Dec{D}\) names, for an active \(P_i\) projection, the cluster \(C\) of the paragraph \emph{The \(P_i\)-form} (for the seed branch of that paragraph, any named cluster containing the feature);
  \item label (i), (ii) or (v), as a member of a composite of \Ptwo{},
    \Pfour{}, and \Psix{}, when the feature is one of an objective or
    specification record of \(P_2\)-form, a
    transition, regime, or staged selection record of \(P_4\)-form, and an
    audited decision-provenance record of \(P_6\)-form, each licensed (i) or (v) by its
    \(P_i\)-form;
  \item an empirical-bridge annotation, when the feature is itself an empirical-bridge annotation identifying an admissible bridge claim (Definitions~\ref{prop:honest-bookkeeping} and~\ref{prop:empirical-bridge}) whose \(\Audit{D}\) entry carries the five items of Definition~\ref{prop:empirical-bridge}, with no fields beyond its typing, targets and that identification;
  \item an outside-domain structure, when a proposed irreducible intentional
    or teleological predicate has no recorded extension in \(D\), the asserted property is true in some expansion of \(D\) and false in another, and every other assertion field of that record either also meets the outside-scope condition or is required for \(W_i\) at a cluster aligned with \(P_i\), an input to its check at such a cluster, or named in a seed clause that the record carries;
  \item unabsorbed role-aligned content, label~(v) of
    Definition~\ref{def:residual-scheme}, when it has \(P_i\)-form but
    criterion (i) does not hold for it, for instance because the cluster lacks
    the attachment data of Definition~\ref{def:threshold-attachment}.
\end{itemize}
A content record targeted by such an empirical-bridge annotation is not thereby classified; it is classified by its own fields and is licensed label (ix-c) exactly when either branch of that criterion holds.
An update map receives (i) when it lies in a named active \(W_1\) cluster and
its assertion fields meet criterion (i); otherwise it receives (v) when it
carries a covered \Pone{} seed whose clause names each of its predicate,
comparison, invariant, law or claim fields, or belongs to a failed \(W_1\)
offer each of whose such fields is an input to the recorded check; and (ix-c)
when it supplies no witness field and has no such field; an update map of none of these three forms must be tested against the other recognized forms of this proposition; only when none applies is it not classified here and left to the recognition hypothesis. Each case licenses a label among (i)--(ix) of
Definition~\ref{def:residual-scheme}, so neither (x) nor (xi) is licensed for
the feature. For these recognized forms, the
covered content requires no seventh role; the outside-domain case receives
label (viii) without a claim about roles beyond the covered scope. None of
these classifications uses the recognition hypothesis
(Definition~\ref{def:recognition-hypothesis}); content of none of these forms
is governed by it.
\end{proposition}

\begin{proof}
Let \(r\) be a residual-agency feature (a kind residual) of
\(D \in \FATCD{}\), agency being taken in the planning sense
of~\citet{Bratman1987}, and pass it through the residual scheme of
Definition~\ref{def:residual-scheme} under
Proposition~\ref{prop:alignment-rules}. An agency feature in admissible
form presents an objective pursued through selection and recorded for
later inspection; we case on whether that content reduces to role
witnesses or asserts more. Role-aligned and composite forms are labelled as in the paragraph
\emph{The \(P_i\)-form} before Proposition~\ref{prop:probability-closure}: each
record receives (i) or (v), and (ii) under the two-cluster condition stated
there.

If \(r\) has just one of the stated \(P_2\)-, \(P_4\)- or \(P_6\)-forms and
receives (i), criterion (i) supplies its named active cluster, \(W_i\)
witness and attachments; hence it is the corresponding single-role
projection. In the composite case, take the finite record set \(R_f\) carrying
the objective or specification, staged selection, and decision
provenance. By the stated \(P_2\)-, \(P_4\)- and \(P_6\)-form
hypotheses, each member receives (i), (ii) or (v) individually.
Together those records have the asserted composite reading. If instead \(r\) is an empirical-bridge annotation whose bridge claim is admissible and has an
audit entry with the five items of Definition~\ref{prop:empirical-bridge}, and with
no fields beyond its typing, targets and that identification, it
receives (vii). A content record targeted by such an annotation is not thereby classified; it is classified by its own fields and is licensed label (ix-c) exactly when either branch of that criterion holds. If \(r\) records a
proposed irreducible intentional or teleological predicate whose extension
\(D\) does not record, its truth cannot be checked over \(D\). In the outside-domain case, the asserted property is true in some
expansion of \(D\) and false in another. Together with the stated
condition on the remaining assertion fields, this satisfies
criterion (viii). A finite audited intentional-content claim with a recorded
predicate extension and no recognized six-role disposition remains subject to
the recognition hypothesis.

Classification as a \Pone{} projection is not available here merely
because agency selects an update: by
Proposition~\ref{prop:alignment-rules} a \Pone{} projection requires a
selected update together with source and target relation records of declared use quotient and an
explicit descent-compatibility, descent-obstruction, or
update-vs-quotient mismatch witness. A selection event carrying no such
witness supplies no \Pone{} witness. Its update map receives (i) when it lies
in a named active \(W_1\) cluster and its assertion fields meet criterion (i);
otherwise it receives (v) when it carries a covered \Pone{} seed whose clause
names each of its predicate, comparison, invariant, law or claim fields, or
belongs to a failed \(W_1\) offer each of whose such fields is an input to the
recorded check; and (ix-c) when it supplies no witness field and has no such
field. It has \(P_4\)-form only when recorded as a
law-bearing transition record. This is the stated \Pone{} caveat. The cases
above are the recognized forms of residual-agency content; content of none of
these forms falls under the recognition hypothesis
(Definition~\ref{def:recognition-hypothesis}). The closure is scoped and asserts
nothing beyond it.
\end{proof}

\begin{remark}[Scope of the closures]\label{rem:threat-closures-scope}
Propositions~\ref{prop:probability-closure}--\ref{prop:agency-closure}
close the three threats within the covered scope; they do not make
scope-free claims. A feature that falls outside the covered scope is
classified as an outside-domain structure via
Definition~\ref{def:residual-scheme}, not absorbed into a primitive role
by default. Each named form is checked against its label criterion directly, without
requiring recognition of every feature of \(D\); the propositions give the
applicable label criterion for each named form and state when \Pfive{} and
\Pone{} projections are unavailable. Each recognized form in
Propositions~\ref{prop:probability-closure}--\ref{prop:agency-closure} is
alternative (R1), (R2), (R3) or (R5) of Definition~\ref{def:recognition-hypothesis}:
each individual record in a recognized form satisfies (R1), (R2), (R3) or (R5), and
composite forms are assessed member by member; the propositions verify those alternatives for
the listed forms, and add that a \Pfive{} reading (causality) or a \Pone{}
reading (agency) is unavailable without quotient-packaging or descent data.
\end{remark}

\subsection{The recognition hypothesis and the classification theorem}\label{subsec:classification-theorem}

Each alternative of Definition~\ref{def:recognition-hypothesis} is the
decomposition-independent form of label criteria of
Definition~\ref{def:residual-scheme}: (R1) of (i) or (v); (R2) of (v), (vi),
(vii), (ix-a) and (ix-b); (R3) of (ix-c); (R4) of (iii) and (iv); and (R5) of
(viii). Theorem~\ref{thm:upper-bound-classification} verifies that this
matching does not depend on the chosen decomposition record.

The hypothesis below is imposed on three sets: the kind residuals of Definition~\ref{def:residual-scheme} (content records and structural annotations, whose role reading is not fixed by their kind), all bookkeeping annotations, and the set \(\Residual{D}\) of each well-formed decomposition record (Definition~\ref{def:decomposition-record}). Stated for every record, it
does not depend on a choice among them. Since \(\Residual{D}\) consists of kind
residuals and bookkeeping annotations, the third clause below follows from the first two: \(D\) satisfies the recognition
hypothesis if and only if every kind residual and every bookkeeping annotation of
\(D\) is recognized, a condition that refers to no decomposition
record.

\begin{definition}[Recognition hypothesis]\label{def:recognition-hypothesis}
Terms used below: \emph{aligned}, \emph{contributes a required field},
\emph{required for}, \emph{input to the check}, \emph{declared operation},
\emph{independent assertion} and \emph{outside the covered scope} are defined in
the paragraph \emph{Terms used in the label criteria} after
Definition~\ref{def:residual-scheme}; \emph{covered seed}, \emph{named seed
field} and \emph{failed offer} in Definition~\ref{def:channel-status};
\emph{honestly recorded} (of an offer: its audit is honest,
Definition~\ref{def:audited}(b),(d)).
The \emph{closure content} of a feature \(r\) is the finite induced
subdescription consisting of \(r\), its fields, and all declared records
reached by following its references in \(D\) or read by its formula-valued fields; it is used in
Remark~\ref{rem:recognition-open} (paragraph \emph{Verification}). The alternatives below are evaluated in the
whole of \(D\), and a reference makes data available for checking without by
itself certifying a role witness. A description \(D \in \FATCD{}\)
satisfies the \emph{recognition hypothesis} when every kind residual \(r\)
of \(D\) (Definition~\ref{def:residual-scheme}), every bookkeeping annotation \(r\) of \(D\) (whether or not a decomposition record assigns it to a channel), and every
element \(r\) of \(\Residual{D}\) for every well-formed decomposition record
of \(D\) (Definition~\ref{def:decomposition-record}), is \emph{recognized}, that is,
satisfies at least one of the following:
\begin{enumerate}[label=(R\arabic*),topsep=2pt,itemsep=1pt]
  \item (first branch) \(r\) itself contributes a required field to a finite cluster aligned
    with one of the six witness types of the \emph{Witness types} table of Section~\ref{sec:primitive-role-architecture},
    or itself carries a covered seed for one of them
    (Definition~\ref{def:channel-status}), and every predicate, comparison, invariant, law or claim field \(\varphi\) of \(r\) is, for some index \(i_\varphi\) and some finite cluster \(C_\varphi\ni r\) aligned with \(P_{i_\varphi}\), required for \(W_{i_\varphi}\) at \(C_\varphi\) or an input to its check, or is named in a seed clause that \(r\) carries (so \(r\) carries no assertion beyond these fields; content beyond them must be recorded in a separate declared content record, which is classified on its own); or (failed-offer branch)
    \(r\) is a member of the cluster of an honestly recorded failed offer for
    some \(P_i\), and each of its predicate, comparison, invariant, law or claim
    fields is an input to the recorded check of \(W_i\) on that cluster; merely
    referencing a witness cluster does not suffice;
  \item \emph{Structural annotations.} \(r\) is one of the six structural annotations of
    Definition~\ref{def:annotations}; its fields contain only its declared
    operation, targets, and the finite typing, checking, preserved or
    suppressed data, losses and nonclaims required for that operation, and it
    makes no independent assertion (both terms as defined after
    Definition~\ref{def:residual-scheme}); in particular an instrument-visibility
    annotation whose fields other than its typing and target fields are those
    required by Definition~\ref{prop:visibility}, a lift annotation whose fields are those
    required by Definition~\ref{prop:forgetting-lifts}, and an
    empirical-bridge annotation meeting criterion (vii), whose fields are only
    its typing, targets and the identification of its bridge claim and five
    items, satisfy this
    alternative, provided each of their formula-valued fields meets the
    condition of the next sentence. A formula-valued field of a structural annotation (a threshold predicate, a preserved or lost invariant) is part of its declared operation only when it reads at least one record, every record whose values it reads is listed in the annotation's target field, and each such record satisfies (R1); otherwise it is an independent assertion, and this alternative fails.
    \emph{Bookkeeping annotations.} The alternative is also met when \(r\) is one of the seven bookkeeping annotations (for \(D\in\FATCD{}\) every condition below except the last, on field roles, holds automatically by Definitions~\ref{def:finite}--\ref{def:fatcd}) of
    Definition~\ref{def:annotations}, it is finitely represented, has finite typed fields, and every reference it carries resolves to a declared finite record of \(D\) (Definitions~\ref{def:finite} and~\ref{def:closed}), it is honest if it is an audit annotation, it carries the fields Definition~\ref{def:audited}(g) requires of its kind, and, if it is a claim annotation, it is referred to by an honest audit annotation and satisfies the representation condition (Definition~\ref{def:audited}(e)), and every field of
    \(r\) other than the statement and hypothesis fields of a claim annotation,
    the certificate field of an audit annotation and its conclusion field when its conclusion is true in every well-typed interpretation or is represented, in the sense of Definition~\ref{def:annotations}, by fields of the content records named by the claim it audits, and the
    formula schema of a witness-schema annotation has field role datum or
    reference; a substantive assertion in a claim statement is represented by
    content records (Definition~\ref{def:annotations}), which are assessed
    under (R1)--(R5);
  \item \(r\) is role-neutral data without an assertion of its own: \(r\)
    supplies no witness field and has no predicate, comparison, invariant,
    law or claim field. Typical instances are the inputs that a specified field
    of a feature satisfying (R1) uses in checking a required or seed field, and
    declared objects or states that no role datum uses; also a predicate record supplying no witness field whose only assertion field is a formula with at least one free variable whose truth value is not the same under all assignments respecting the declared sorts of its free variables, using only predicates whose extensions \(D\) records, and quantifying only over sorts whose carriers \(D\) records, each free variable ranging over a record kind or the rational sort, referenced by no claim annotation, and referenced by certificates only through the abbreviation \(s(t)\) of the paragraph \emph{Check language} (a definition, which has no truth value over \(D\));
  \item \(r\) begins a finite chain \(r=r_0,\ldots,r_n\) of distinct
    declared records, with \(n\geq1\), such that, for each \(j<n\), a field of \(r_j\) records a checked
    equality or isomorphism of presentation data with specified fields of
    \(r_{j+1}\), or a checked level translation to \(r_{j+1}\); and, for each \(j<n\), apart from one comparison field recording the equality, isomorphism or level translation with \(r_{j+1}\), and the datum and reference fields it compares, \(r_j\) has no predicate, comparison, invariant, law or claim field. The terminal \(r_n\) satisfies (R1), (R2), (R3) or
    (R5). An equality or isomorphism of presentation data (a value-preserving
    relabelling of datum fields; no structural isomorphism is meant) is a finite
    bijection \(\beta\), recorded in the comparison field, from specified datum fields that are not term-valued
    (or the member list of an object record, a reference field) of \(r_j\) to datum fields that are not term-valued (or the member list of an object record, a reference field) of \(r_{j+1}\), with a
    correct certificate, recorded true, that each recorded value equals that of its \(\beta\)-image;
    its formula is a conjunction of equalities between record identifiers,
    recorded field values and entries of recorded lists (written with indexing,
    as in the running example of Section~\ref{sec:fatcd-domain}). A checked level translation from \(r_j\) to \(r_{j+1}\) is a
    comparison field of \(r_j\) referencing a translation annotation carried
    in \(\Audit{D}\) whose source and target contexts are levels recorded by level annotations whose target fields list \(r_j\) and \(r_{j+1}\), respectively, together with a finite
    bijection \(\beta\) of specified datum fields as above and a correct certificate, recorded true,
    that is such a conjunction of equalities,
    that each value equals its \(\beta\)-image;
  \item \(r\) lies outside the covered scope
    (Definition~\ref{def:residual-scheme}): a field of \(r\) asserts a property whose check is not finite over the records of \(D\): it quantifies over a declared domain for which \(D\) supplies no finite checking range, or uses a declared predicate whose extension \(D\) does not record, and its truth value is not fixed by the recorded values of \(D\) (it is true in some expansion of \(D\) and false in another, in the sense of the paragraph \emph{Certificates, assertion language and laws} of Section~\ref{sec:primitive-role-architecture}; for a formula with free variables: for some assignment \(\alpha\) of its free variables and two expansions of \(D\) whose carriers contain the values \(\alpha\) gives to variables of unlisted sorts, it is true under \(\alpha\) in one and false under \(\alpha\) in the other). In addition, every other predicate, comparison, invariant, law or claim field of \(r\) either has this property or is required for \(W_i\) at a cluster \(C\) aligned with \(P_i\), an input to the check of \(W_i\) at such a cluster, or named in a seed clause that \(r\) carries; \(r\) is then
    recognized as outside-domain structure.
\end{enumerate}
For \(D\in\FATCD{}\), every bookkeeping annotation whose fields other than the
exempt fields listed in the second clause of (R2) are datum or reference fields
satisfies the
second clause of (R2); hence, when all bookkeeping annotations of \(D\) have this
form, \(D\) satisfies the recognition hypothesis if and only if every
kind residual of \(D\) is recognized.
A feature satisfying (R1) is absorbed under label (i) when a decomposition record
places it in an active projection under
Proposition~\ref{prop:alignment-rules}, and is covered under (v) when it
remains unabsorbed. Remark~\ref{rem:recognition-open} shows that some
descriptions in \FATCD{} do not satisfy this hypothesis.
\end{definition}

Alternatives (R1)--(R5) are record-independent forms of label criteria: (R1)
of (i) or (v); (R2) of (v), (vi), (vii), (ix-a) or (ix-b); (R3) of (ix-c);
(R4) of (iii) or (iv); and (R5) of (viii) (the table in the proof of
Theorem~\ref{thm:upper-bound-classification}).

\begin{remark}[The recognition hypothesis is a genuine restriction]\label{rem:recognition-open}
The recognition hypothesis is an explicit assumption on \(D\), not a consequence
of the \FATCD{} definitions, and some descriptions in \FATCD{} do not satisfy
it. The remark establishes five facts, each under the paragraph named:
\(D_{\mathrm{norm}}\in\FATCD{}\setminus\FATCD_{\mathrm{RH}}\), so
\(\ScopedExactSix{\FATCD{}}\) fails (\emph{The counterexample},
\emph{Verification}); a transition-row presentation of the same weights lies
in \(\FATCD_{\mathrm{RH}}\) (\emph{Repair by a transition row}); every closed
check-language formula without atomic \(W_j\) formulas has a covered
\Ptwo{}-seed presentation, given two distinct declared records of one kind (or after adjoining them), whenever the completed description is again in \FATCD{} (\emph{Wrapping a closed formula as a \Ptwo{} seed}; membership can fail when a correct certificate, or a realization relation's target formula, quantifies over all declared records of an adjoined kind);
membership in \(\FATCD_{\mathrm{RH}}\) depends on presentation (\emph{Scope
depends on presentation}); and label (xi) is realized by \(D_{\mathrm{scr}}\)
(\emph{A description passing the screen}). The paragraphs \emph{One record suffices for recognition}, \emph{Other unrecognized assertions} and \emph{Open question} add a one-record failed-offer presentation, further unrecognized forms with the (R2) and (R3) presentation devices, and the open question. A record without an assertion of its own --- one with no witness field and
no predicate, comparison, invariant, law or claim field, such as a declared
object record --- is recognized by (R3). The normalization predicate in the
following description is not recognized.
\par\emph{The counterexample \(D_{\mathrm{norm}}\).} Let \(D_{\mathrm{norm}}\), the description of Remark~\ref{rem:fatcd-minimality}, consist of state records \(w_1,\dots,w_n\)
(\(n\ge1\)), each with a nonnegative rational value field \(v(w_k)\), the values
summing to one (for instance \(n=2\) and \(v(w_1)=v(w_2)=1/2\)); a predicate
record \(N\), declared as a macro-specification, whose formula asserts
\(\bigwedge_k v(w_k)\ge0\wedge\sum_k v(w_k)=1\); a claim annotation naming \(N\) whose statement is the
formula of \(N\); and an honest audit annotation listing the certificate
\((\bigwedge_k v(w_k)\ge0\wedge\sum_k v(w_k)=1,\mathrm{true})\) and certifying that formula.

\emph{Verification.} \(D_{\mathrm{norm}}\) is finite, closed and
audited, so \(D_{\mathrm{norm}}\in\FATCD{}\). It contains no map, relation,
comparison, path, derivation, order, transition, refinement, event,
provenance, trace, replay or check record; in particular the closure content
of \(N\), namely \(N\) with \(w_1,\dots,w_n\), contains none. So no cluster
satisfies any \(W_i\)
and no record is a covered seed (Definition~\ref{def:channel-status}). Hence
\(N\) does not satisfy (R1); being a content record, not an annotation, it does not satisfy (R2); its predicate field contains a closed formula and it is referenced by a claim annotation, so neither branch of (R3) applies;
it records no presentation or level translation, so not (R4); and \(N\) gives
its predicate by the explicit formula \(\bigwedge_k v(w_k)\ge0\wedge\sum_k v(w_k)=1\) over recorded field
values; this formula uses no declared predicate symbol and quantifies over no
sort, so its check is finite over \(D_{\mathrm{norm}}\) and (R5) fails.
Although \(N\) is declared a macro-specification, it is not a covered seed:
\(D_{\mathrm{norm}}\) has no realization relation to a declared nonempty
candidate class, and no realization relation can target \(N\), since
Definition~\ref{def:raw-record} requires its target formula to have exactly one
free variable. For the same reasons no
criterion (i)--(ix) of Definition~\ref{def:residual-scheme} holds for \(N\):
it is not in any cluster or composite, re-presents or restates nothing, is not
an annotation, relates no instrument or substrate observable, is inside the
covered scope, and carries an assertion. With no failed offer and no outside-scope declaration,
every channel of \(D_{\mathrm{norm}}\) is \(\absentrec\), so no active projection absorbs \(N\), and \(N\) receives
label (x) in every well-formed decomposition record. So \(\ScopedExactSix{\FATCD{}}\)
fails, and the six roles are not \emph{recognition-complete} for \FATCD{}: some in-scope
feature of some description in \FATCD{} is placed by none of (i)--(ix).
Label (x) records an unresolved candidate, not a seventh role: (xi) would
further require the recorded witness of category (xi), itself a recorded
screen rather than the identification of a new closure role.

\emph{Repair by a transition row.} For distribution-valued weights the hypothesis constrains presentation, not content. For nonnegative
rational weights summing to one, replace \(N\) by a transition record
\(K:\{*\}\to\mathrm{Dist}(\{w_1,\ldots,w_n\})\) whose unique row is the list
of terms \(v(w_1),\dots,v(w_n)\), so that \(K\) reads the value fields of the
\(w_k\) and records no copies. Declare a state record \(*\), a one-point object record \(\{*\}\) listing it, and an
object record \(X\) listing \(w_1,\dots,w_n\), give \(K\) a codomain reference to
\(X\) and a datum field marking it distribution-valued, and take as the law certificate
of \(K\) the pair \((\bigwedge_k v(w_k)\ge0\wedge\sum_k v(w_k)=1,\mathrm{true})\),
which expresses that the row lies in \(\mathrm{Dist}(X)\). Retarget the claim and its audit to \(K\), replacing every reference
to \(N\) by a reference to \(K\). Set the claim's statement and the
audit's conclusion to the law-certificate formula, and replace the
audit's certificate list by that certificate, recomputed true,
naming \(K\) and its declared inputs. Thus \(K\) represents the
statement (Definition~\ref{def:annotations}), and the audit is honest. Then \(K\) is a covered \Pfour{} seed
satisfying (R1); \(*\), \(\{*\}\), \(X\) and the \(w_k\) carry no assertion field and
satisfy (R1) or (R3); the claim and audit annotations satisfy (R2); and the
resulting description lies in \(\FATCD_{\mathrm{RH}}\).

\emph{Wrapping a closed formula as a \Ptwo{} seed.} In outline: the construction adjoins a predicate record \(s\) with one free variable over two records \(c\ne c'\), a realization relation recording \(s(c)=\mathrm{true}\) and \(s(c')=\mathrm{false}\), and a one-class packaging map \(q\) with a comparison record; it chooses between two formulas for \(s\) according to the value of \(\varphi\) in the completed description \(D'\). More generally, let \(\varphi\) be any closed check-language formula containing
no atomic \(W_j\) formula and
\(c\ne c'\) declared records of one kind. Let \(D'\) be the final description, obtained once every record of this construction (including the records \(X_c,*,Q,q\) and the comparison record introduced under \emph{Activation} below, the \Ptwo{} claim and its audit, all attachment annotations and any fresh candidate records, when used) is adjoined and any removal below is made. The predicate record is \(s(r):\equiv\varphi\wedge r=c\) if \(\varphi\) holds in \(D'\), and \(s(r):\equiv\varphi\vee r=c\) otherwise. \par\emph{Choice of branch.} Fix all additions, removals and retargetings, and all non-formula field values, before choosing the branch. The two branches differ only in the formula field of \(s\), the retargeted claim's statement and its audit's conclusion. A check-language formula never reads a statement or conclusion (formula-valued fields are not terms) and reads a predicate record's formula only through an abbreviation \(s(t)\) naming it; \(s\) is adjoined, so \(\varphi\) contains none. Hence \(\varphi\) has the same value in both candidates for \(D'\). Its value agrees with its value in \(D\) provided every field it reads is unchanged and every quantifier has the same checking range. The construction applies when \(D'\) is again in \FATCD{}. Adjoining records can falsify a correct certificate of \(D\) whose quantifier ranges over all declared records or over the declared records of an adjoined kind (for instance \((\forall x\in\mathrm{State}\,(x=w_1\vee x=w_2),\mathrm{true})\), which the adjoined state record \(*\) falsifies). A sufficient condition is that every retained realization relation still lists exactly its satisfying members and records correct evaluations, every retained certificate and law remains correct, every retained audit remains honest, and all additions, removals and retargetings preserve typing, reference closure, representation and the required annotation fields. \par\emph{Seed and claim.} This predicate record, declared a macro-specification whose variable \(r\) is declared to range over the common kind of \(c\) and \(c'\), with
a realization relation declaring \(K=\{c,c'\}\), listing exactly \(c\) as
satisfying, and recording both evaluations, is a covered \Ptwo{} seed. Let the
realization relation carry the evaluations as certificates
\((s(c),\mathrm{true})\) and \((s(c'),\mathrm{false})\). The record that carried \(\varphi\) as an assertion field, if any (such as
\(N\)), may be removed only if neither \(\varphi\) nor its named inputs
reference it, it is neither \(c\) nor \(c'\), and every incoming reference is a
bookkeeping reference that can be retargeted to \(s\) and the realization
relation while preserving typing, representation and audit honesty; otherwise
that record is retained and classified separately; a claim whose statement
is \(\varphi\) is retargeted to \(s\) and the realization relation, with statement
\(\varphi\) in the first branch and \(\neg\varphi\) in the second; its audit lists
both certificates and concludes that statement. The statement is entailed by
the signed certificates, so the claim satisfies the representation condition
of Definition~\ref{def:annotations} and the audit is honest. Its
predicate field is a named seed field, so it satisfies (R1). It receives (i)
when criterion (i) holds and (v) otherwise. \par\emph{Activation.} Adjoining an object record \(X_c\)
listing \(c,c'\), a state record \(*\) with a quotient object record \(Q\)
listing it, a map record \(q:X_c\to Q\) of declared use packaging with
\(q(c)=q(c')=*\), and a comparison record carrying
\((q(c)=q(c')\wedge s(c)\wedge\neg s(c'),\mathrm{true})\) gives the following. Let \(C\) contain \(c,c',s\), the realization
relation, \(X_c\), \(*\), \(Q\), \(q\) and the comparison record. Then
\(W_2(C)\) holds: the realization relation records both evaluations on
\(K=\{c,c'\}\), and the single kernel class of \(q\) meets the realizer set
\(\{c\}\) and its complement \(\{c'\}\). Record a claim naming \Ptwo{} and \(C\)
with statement \(W_2(C)\), and an honest audit listing the row's content
certificates and \((W_2(C),\mathrm{true})\) and concluding \(W_2(C)\), and
supply all attachments and level and lift conditions of
Definition~\ref{def:role-projection}; fresh candidate records may be used to
avoid pre-existing level, lift or forgetting annotations meeting \(C\). The
resulting projection is active. Since \(q\) is a covered \Pfive{} seed, in \(D'\) channel \Ptwo{} is active and channel \Pfive{} is not \(\absentrec\); the transition-row repair likewise makes \Pfour{} at least \(\belowthr\). Removing \(s\) from this cluster falsifies
\(W_2\), and its sole predicate field is required for that witness, so
criterion (i) applies in a decomposition record naming this projection. \par\emph{What is recorded.} This construction records the truth value of \(\varphi\) in \(D'\) in recognized
seed content: in the first branch the recorded evaluation \(s(c)=\mathrm{true}\)
entails \(\varphi\), while in the second branch \(s(c')=\mathrm{false}\) entails
\(\neg\varphi\). Thus an assertion true in \(D'\) is recorded with its checked content; an assertion false in \(D'\) is represented by its checked negation. Agreement with its truth value in \(D\) requires the preservation conditions above. The transition-row
presentation above likewise preserves the checked normalization content for
nonnegative weights of pairwise distinct records.

\emph{One record suffices for recognition.} Let \(c\) be a declared record, \(s(r):\equiv\varphi\wedge r=c\), and let a realization relation with candidate class \(K=\{c\}\) and target \(s\) carry \((s(c),\mathrm{true})\) if \(\varphi\) holds in the completed description and \((\neg s(c),\mathrm{true})\) otherwise. Fix all additions and fields other than that certificate's formula and the relation's list of satisfying members before selecting the branch. Assume \(\varphi\) reads neither changing field in either completion, including through quantified field lookups or predicate abbreviations. Then \(\varphi\) has the same value in both completions, so the branch choice is well defined. The construction applies when the completed description, with \(s\), the relation, the claim and its audit adjoined, is again in \FATCD{}. Record a claim naming \Ptwo{} and \(C=\{c,s,\mathrm{Real}\}\) with an honest audit listing that certificate and \((W_2(C),\mathrm{false})\). Since \(K=\{c\}\), no kernel class can meet both the realizer set and its complement in \(K\), and no evaluation outside \(K\) is recorded in \(C\); hence \(W_2(C)\) is false in every well-typed interpretation; the formula of \(s\) and the certificate are inputs to the recorded check, so \(s\) and the relation satisfy the failed-offer branch of (R1) and receive (v-c), and the certificate entails \(\varphi\) or \(\neg\varphi\). Channel \Ptwo{} is then \(\collapsedbc\) unless active. The two-record proviso is needed only for the covered-seed and active presentations.

\emph{Scope depends on presentation.} The covered scope also depends on presentation. Replace the formula of \(N\)
by \(\mathrm{IsDist}(v(w_1),\ldots,v(w_n))\), declare \(\mathrm{IsDist}\) by a
predicate record of arity \(n\), each argument place declared rational-sorted, with no formula field and give no relation
record for it, so that \(D\) does not record its extension (the declaring
record has no assertion field and satisfies (R3); had the declaring record
given an explicit formula in its \(n\) rational argument variables, such as \(\bigwedge_k x_k\ge0\wedge x_1+\dots+x_n=1\), \(D\) would record the extension and \(N\) would
again receive (x)), set the claim's statement
to \(\mathrm{IsDist}(v(w_1),\ldots,v(w_n))\), and let the audit list no certificate, conclude the valid formula
\(v(w_1)=v(w_1)\), and list the \(\mathrm{IsDist}\) field of \(N\) in its
non-certification field. This description remains honestly
audited. Its records \(w_k\) and the declaring record of \(\mathrm{IsDist}\)
satisfy (R3), its claim and audit annotations satisfy the second clause of
(R2), and \(N\) satisfies (R5) and receives (viii) rather than (x); so the
description lies in \(\FATCD_{\mathrm{RH}}\) while \(D_{\mathrm{norm}}\) does
not, and no role is certified for \(N\). Hence membership in \(\FATCD_{\mathrm{RH}}\) is not
invariant under replacing an explicit formula by a declared predicate name whose extension \(D\) does not record,
and Theorem~\ref{thm:scoped-exact-six} asserts nothing about content presented
that way. More generally, replace at least one assertion field of an unrecognized
kind residual by a declared predicate name whose extension the resulting
description does not record. The changed record satisfies (R5) provided every
other assertion field, evaluated in the resulting description, also lies
outside the covered scope, is required for or an input to \(W_i\) at an
aligned cluster containing the record, or is named in a seed clause the record
carries. Unchanged records must be recognized separately. The result lies in
\(\FATCD_{\mathrm{RH}}\) if it is again in \FATCD{} (in particular, no replaced field is the formula of a realization relation's target),
every remaining kind residual is recognized, and every bookkeeping annotation
has only datum and reference fields besides the exempt fields of (R2).
Theorem~\ref{thm:scoped-exact-six} certifies no role for the replaced
fields.

\emph{Other unrecognized assertions.} A standalone predicate record stating a conservation or balance law, or an
assembled obstruction or holonomy, by a closed check-language formula over
recorded inputs is not recognized as presented if it carries no covered seed,
contributes no required witness field, is not an input to a prescribed witness
check, and records no checked presentation or level translation. Under these
conditions (R1)--(R5) fail as for \(N\).
An in-scope content record that contributes a required field to an aligned
cluster or carries a covered seed, and whose assertion fields are each required
at, or inputs to the check at, some aligned cluster containing it, possibly
different clusters for different fields, satisfies (R1); one with no assertion
field and no witness field satisfies (R3).
Label (x) is also licensed for an unclaimed predicate record whose closed
formula is true or false in \(D\), such as \(0+0=0\) or \(0+0=1\) between rational constants: label (x) marks an assertion placed
by none of (i)--(ix), which need be neither true nor a candidate role. By contrast, the same record
with formula \(\theta\wedge r=r\) for a fresh variable \(r\), still unclaimed and
referenced by no certificate, does not satisfy the second branch of (R3), since its value does not depend on
\(r\), and still receives (x). Conversely, if \(c\) is a declared record of the kind over which \(r\) ranges and that kind has a second declared record, the unclaimed record with formula \(\theta\wedge r=c\) (for \(\theta\) true in \(D\)) or \(\theta\vee r=c\) (for \(\theta\) false in \(D\)), referenced by no certificate, satisfies the second branch of (R3), since its value at \(c\) differs from its value at the other records of that kind; for unclaimed closed assertions this branch constrains presentation, not content. The same holds for the structural-annotation branch of (R2): for a closed check-language formula \(\theta\) that reads no record values outside a binary relation record \(E\) satisfying (R1), and a fixed declared record identifier \(u\), a threshold annotation targeting \(\{E\}\) with predicate \(\theta\wedge(E(u,u)\vee\neg E(u,u))\) satisfies (R2) and receives (ix-b) (no cluster \(\{E\}\) satisfies any \(W_i\), so (v-b) cannot apply), whether or not \(\Audit{D}\) records the predicate checked true, provided the description with this annotation adjoined is again in \FATCD{}. Since the
\(W_6\) certificate states a replay or provenance property of its trace, an
arbitrary assertion about recorded values does not become a required
\(W_6\) field by adjoining a trace.

\emph{A description passing the screen.} Let \(D_{\mathrm{scr}}\) consist of
index records \(0,1\), an object record \(X\) listing them, an order record \(o\) for
\(0<1\) with its law certificate, a relation record \(r\) (the feature to be screened; not the free variable \(r\) of a macro-specification) of declared use
other carrying a finite table field \(c\) with \(c(0)=c(1)=1\) and a datum field
holding the list \((\mathrm{id}_X)\), with
certificates of extensivity, monotonicity, \(c\circ c=c\), closure of the list
under composition and \(c\ne\mathrm{id}_X\) (the prescribed conditions of criterion (xi)(d)--(e), with \(o\) as the order record and \((\mathrm{id}_X)\) as \(L\)), a witness-schema annotation naming a
candidate outside \(P_1,\dots,P_6\) with \(\psi(C):\equiv r\in C\wedge o\notin
C\), where \(r\in C\) abbreviates the finite disjunction \(\bigvee_{d\in C}r=d\)
over the listed identifiers of \(C\) and \(o\notin C\) the conjunction
\(\bigwedge_{d\in C}o\neq d\), so that each instance \(\psi(C')\) is a closed
check-language formula, a check record carrying \(\psi(\{r\})\), \(\neg W_1(\{r\}),\dots,\neg
W_6(\{r\})\), \(\neg\psi(\{o\})\) and the finite conjunction over all
\(C'\subseteq\Raw{D_{\mathrm{scr}}}\) of \((W_1(C')\vee\dots\vee
W_6(C'))\rightarrow\neg\psi(C')\), a claim referencing \(r\) and that check
record whose statement is these conditions, and an honest audit certifying
it. The \(W_j\) atoms in these certificates are evaluated by the witness rows,
which read no certificate containing a \(W_j\) atom and, for \(W_6\), only
\(\mathrm{Audit}_0(D_{\mathrm{scr}})\), which excludes this audit; so the
certificates are well defined although the check record belongs to some
clusters \(C'\). Only \(o\) is law-bearing and no record supplies data for
\(W_1,W_2,W_3,W_5,W_6\), so \(W_i(C')\) holds only if \(o\in C'\), where
\(\psi(C')\) is false; no map, partial-map, transition or refinement record is
present, so the generated monoid is \(\{\mathrm{id}_X\}\). Recording \(c\) by
a map record would place it in the list of criterion (xi)(e), and \(r\) would
fail the screen. The assertion fields
of \(r\) are required for no witness, \(r\) carries no seed, begins no (R4)
chain and lies in the covered scope, so \(r\) satisfies (x) and (xi); the check record, which serves no witness and carries no seed, also receives (x) but not (xi), since it records no closure operation. Thus
(xi) is realizable in \FATCD{}, and its emptiness in
Theorem~\ref{thm:upper-bound-classification} comes from the hypothesis.

\emph{Open question.} Which assertion-bearing records admit a reduction in which the assertion is
itself a required field of an active cluster (label (i)), without
adjoining wrapper records such as the realization relation and one-class
packaging map above, or a trace reference, is left open and is
recorded with the endogenous-instrument meta-limit of
Sections~\ref{subsec:meta-limit} and~\ref{subsec:future-meta-limit}. The three
threat closures of
Propositions~\ref{prop:probability-closure}--\ref{prop:agency-closure} label
features of their recognized forms without the hypothesis; the classification
theorem below assumes it in general.
\end{remark}

Here and in the entry \emph{Closed unconditionally} of the working vocabulary in Section~\ref{subsec:objects-at-a-glance}, the phrase means that the stated recognized forms of Propositions~\ref{prop:probability-closure}--\ref{prop:agency-closure} receive labels without assuming the recognition hypothesis. It does not assert that hypothesis for every feature in those three threat classes; Theorems~\ref{thm:upper-bound-classification} and~\ref{thm:scoped-exact-six} retain it for general residual exhaustion.

Three ingredients combine in the next theorem. The alignment rules of
Proposition~\ref{prop:alignment-rules} fix which features each role can
claim. The threat closures of
Propositions~\ref{prop:probability-closure}--\ref{prop:agency-closure}
classify the recognized forms of probability, causality and agency. The residual scheme of
Definition~\ref{def:residual-scheme} supplies the labels. Under the
recognition hypothesis, the theorem assigns every residual feature a
label other than an unresolved threat or a seventh role: its proof
matches each recognized disposition to one of labels (i)--(ix). The
three preceding propositions classify their stated recognized forms
without that hypothesis.

The proof sorts a recognized feature by its disposition, following the
label table after Definition~\ref{def:residual-scheme}; a feature of any
raw-grammar or annotation kind receives the label of the disposition its
content has. The alternatives (R1)--(R5) of the recognition hypothesis are
written to correspond to the label criteria, so the proof matches each
alternative to a criterion and uses no substantive property of the particular
witness rows (Remark~\ref{rem:classification-scope}); which features the
hypothesis admits and which it excludes is described in
Remark~\ref{rem:recognition-open}.
Label~(viii) and the channel status \(\outsidescope\) of
Definition~\ref{def:channel-status} both concern content outside the
covered scope: the label sorts a single feature, and the status marks a
channel whose role content lies there.

\begin{theorem}[Upper-bound classification]\label{thm:upper-bound-classification}
Let \(D \in \FATCD{}\) satisfy the recognition hypothesis
(Definition~\ref{def:recognition-hypothesis}), let \(\Dec{D}\) be a
well-formed decomposition record of \(D\)
(Definition~\ref{def:decomposition-record}), and let \(r\) be a kind residual
of \(D\) (Definition~\ref{def:residual-scheme}: any content record or structural
annotation, including members of active clusters), a bookkeeping annotation of \(D\) (whether or not \(\Dec{D}\) assigns it to a channel), or a record residual, an
element of \(\Residual{D}\). Then some criterion among (i)--(ix) holds for \(r\) relative to \(\Dec{D}\); hence neither (x) nor (xi) is licensed for \(r\), and
every label that the label criteria after Definition~\ref{def:residual-scheme}
license for \(r\) relative to \(\Dec{D}\) is one of
\begin{enumerate}[label=(\roman*),topsep=2pt,itemsep=1pt]
  \item a declared constituent of the cluster that \(\Dec{D}\) names for an
    active \Pone{}, \Ptwo{}, \Pthree{}, \Pfour{}, \Pfive{}, or \Psix{}
    projection, in the sense of criterion (i) of Definition~\ref{def:residual-scheme};
  \item a shared constituent of two such projections;
  \item a presentation variant of an already classified feature;
  \item a level shift of an already classified feature;
  \item unabsorbed role-aligned content, covered-seed content with no
    certified witness \(W_i\), content checked in an honestly failed offer, or an activation-threshold refinement;
  \item an instrument-visibility artifact;
  \item an empirical-bridge annotation;
  \item an outside-domain structure;
  \item already-covered bookkeeping or annotation data, or role-neutral data
    without an assertion of its own.
\end{enumerate}
In particular, categories (x) and (xi) of
Definition~\ref{def:residual-scheme}, namely the unresolved-threat category and the recorded seventh-role
screen, are empty: no candidate
seventh-role feature in the covered scope remains as an unresolved threat, and
no feature in the covered scope passes the recorded seventh-role screen of
criterion (xi); here ``no unresolved threat'' means that some criterion among
(i)--(ix) holds for \(r\). Since seed clause (2) of
Definition~\ref{def:channel-status} accepts any one-variable macro-specification with a realization relation
whose recorded evaluations take both truth values, this certifies the form in
which \(r\) is recorded, not that its assertion concerns the role whose seed it
carries (Remark~\ref{rem:classification-scope}). \emph{Record residuals.} If \(r\) is an element of \(\Residual{D}\), every label licensed for it
lies in (iii)--(ix). \emph{Pointwise form.} The conclusion also holds pointwise, without the
recognition hypothesis for the rest of \(D\): for every \(D\in\FATCD{}\), every
well-formed \(\Dec{D}\) and every \(r\) as above that satisfies one of
(R1)--(R5), every label licensed for \(r\) lies in (i)--(ix), and in
(iii)--(ix) if \(r\in\Residual{D}\). For each fixed well-formed decomposition
record, alternatives (R1)--(R5) supply, respectively, a criterion among
(i)/(v), (v)/(vi)/(vii)/(ix-a)/(ix-b), (ix-c), (iii)/(iv), and (viii), as the
proof below establishes.
\end{theorem}

\begin{proof}
Alternatives (R1)--(R5) of Definition~\ref{def:recognition-hypothesis} were
written as record-independent forms of label criteria (i)--(ix) (paragraph
before Definition~\ref{def:recognition-hypothesis}); the proof checks, for a
fixed well-formed record, that each alternative meets one of those criteria.
Let \(D\in\FATCD{}\), fix a well-formed decomposition record
\(\Dec{D}\), and let \(r\) be a kind residual of \(D\)
(Definition~\ref{def:residual-scheme}), a bookkeeping annotation of \(D\), or an element of
\(\Residual{D}\) for that record. Fix any label
\(L\) that the label criteria after Definition~\ref{def:residual-scheme}
license for \(r\) relative to \(\Dec{D}\). The label \(L\) is drawn from
(i)--(xi); we show that, under the recognition hypothesis, it lies in
(i)--(ix), so that (x) and (xi) are empty.

By the label table, \(L\) is (x) only if no criterion (i)--(ix) holds for
\(r\), and \(L\) is (xi) only if \(r\) has the property in (x).

\emph{Conclusion.} Each recognition alternative of
Definition~\ref{def:recognition-hypothesis} supplies a label criterion among
(i)--(ix):
\begin{center}
\small
\begin{tabularx}{\linewidth}{@{}l>{\raggedright\arraybackslash}X@{}}
\toprule
Alternative & Label criterion supplied \\
\midrule
(R1) & (i) if criterion (i) holds; otherwise (v-a) for the first branch of (R1) and (v-c) for its failed-offer branch \\
(R2) & (vi) for visibility; (vii) for an admissible bridge claim (Definition~\ref{prop:honest-bookkeeping}) whose associated audit carries the five items of Definition~\ref{prop:empirical-bridge}; \\
     & (v-b) for a threshold annotation targeting a cluster named for an active projection and carrying only a threshold predicate besides typing, target and label fields, which reads at least one record, every record it reads being listed in its target field and satisfying (R1); \\
     & (ix-b) for remaining structural annotations satisfying (R2); (ix-a) for bookkeeping annotations satisfying the second clause of (R2) \\
(R3) & (ix-c) \\
(R4) & (iii) or (iv), the criterion of the terminal record being supplied by the rows for (R1), (R2), (R3) and (R5); no recursion occurs, since the terminal record satisfies (R1), (R2), (R3) or (R5) \\
(R5) & (viii) \\
\bottomrule
\end{tabularx}
\end{center}
Under the recognition hypothesis
(Definition~\ref{def:recognition-hypothesis}), every kind residual of \(D\), every bookkeeping annotation of \(D\) and every element of \(\Residual{D}\) for every well-formed decomposition record of \(D\), in particular \(r\), satisfies one of (R1)--(R5), so by the table a criterion among (i)--(ix)
holds for it; for (R4) the terminal record of the chain supplies it, and
chains are finite. This uses recognition of \(r\) alone and, for (R4), the
criterion that (R4) requires of the chain's terminal record, so it holds
pointwise in any \(D\in\FATCD{}\). Thus the criterion of (x), which requires that no
criterion (i)--(ix) hold, fails; since (xi) requires (x), it fails as well.
Hence \(L\) lies in (i)--(ix), and categories (x), the
unresolved-threat category, and (xi), the recorded seventh-role screen, are
empty. If \(r\in\Residual{D}\) for a decomposition record, \(L\) is neither (i)
nor (ii): by Definition~\ref{def:decomposition-record}, that set omits the
features already accounted for by an active projection, and (ii) requires
\(r\in C_i\cap C_j\subseteq A(\Dec{D})\); so every label licensed for a record
residual, relative to that record, lies in (iii)--(ix). The
emptiness of (x) and (xi) is the content of the recognition hypothesis, not a
consequence of the \FATCD{} definitions: by Remark~\ref{rem:recognition-open}
a finite audited conservation law or an emergent composite obstruction, if
unrecognized, would fall under (x), and under (xi) only if it also passes the recorded
seventh-role screen; Remark~\ref{rem:recognition-open} separately
exhibits an unrecognized normalization predicate in \FATCD{}.
\end{proof}

\begin{remark}[Scope of the theorem]\label{rem:classification-scope}
Theorem~\ref{thm:upper-bound-classification} is conditional on the recognition
hypothesis (Definition~\ref{def:recognition-hypothesis}) and scoped to \FATCD{}.
Alternative (R1) yields criterion (i) or (v); (R2) yields (v), (vi),
(vii), (ix-a) or (ix-b); (R3) yields (ix-c); (R4) yields (iii) or (iv); and
(R5) yields (viii). The alternatives are stated without reference to a
decomposition record; the theorem records that these
record-independent forms exclude (x). It does not derive that no finite audited feature requires a seventh role; its general conclusion assumes the recognition hypothesis. Propositions~\ref{prop:probability-closure}--\ref{prop:agency-closure} classify only their stated recognized forms without that hypothesis. It does
not claim that no seventh role can arise outside \FATCD{}, that the alignment
rules are exhaustive for richer domains, or that the recognition hypothesis holds
in general. Nor does it prove exact-six on its own: the decomposition/exhaustion
theorem of Section~\ref{sec:decomposition-exhaustion} remains necessary. Which further
classes of assertion-bearing records are recognized, and broader questions, are deferred to
Sections~\ref{subsec:meta-limit} and~\ref{sec:future-work}.
For each fixed well-formed \(\Dec{D}\), the proof matches alternatives
(R1)--(R5) to criteria (i)--(ix). This matching uses no substantive
property of the particular witness rows or seed clauses. It remains
valid for any finite family of check-language witness predicates and
seed clauses, provided the recognition alternatives, label criteria,
named seed fields, alignment and required-field conventions, and
active-cluster assignments are changed consistently. The number six enters
Theorem~\ref{thm:scoped-exact-six} through
Definition~\ref{def:decomposition-record}, and the particular predicates
through Theorem~\ref{thm:six-role-lower-bounds} and
Propositions~\ref{prop:probability-closure}--\ref{prop:agency-closure}. The (R1) and (R4) cases use that record's
active-cluster assignments and checked chains. The theorem proves that recognition excludes (x), using the table in its
proof. The converse does not follow from the label criteria: bookkeeping fields
may receive (ix) while substantive assertions require separate recognition,
and criteria (iii) and (iv) accept a chain whose terminal record lies in \(A(\Dec{D})\cup S(\Dec{D})\), for instance a member of an active or failed-offer cluster with an assertion field serving no witness check, which need not satisfy (R1)--(R5). Criteria (vi) and (ix-b)
impose the (R2) requirement that every record whose values a formula-valued prescribed field reads satisfy (R1), and each implies the structural-annotation branch of
(R2). The upper bound therefore depends on the stated
recognition hypothesis. Recognition is not necessary for the condition of Definition~\ref{def:scoped-exact-six-question}: adding to the nonclaim annotation of \(D_5\) a claim field with formula \(0+0=1\) gives \(D\in\FATCD{}\setminus\FATCD_{\mathrm{RH}}\) in which every well-formed record places that annotation (the only nonclaim annotation targeting \(C_5\), which the active \Pfive{} entry must assign) in \(A(\Dec{D})\), so no kind or record residual is licensed (x) or (xi). The proof uses recognition of kind residuals, bookkeeping annotations and record residuals. Since seed clause (2) accepts any one-variable
macro-specification with a realization relation whose recorded evaluations take both truth values, (R1)
certifies the form of a record, not that its assertion concerns the role
whose seed it carries.
\end{remark}
 \section{Decomposition and exhaustion}\label{sec:decomposition-exhaustion}



Under the recognition hypothesis every well-formed decomposition record meets
the condition of Definition~\ref{def:scoped-exact-six-question}
(Theorem~\ref{thm:upper-bound-classification}), so the exact-six question
reduces to the existence of one well-formed record, which the theorem below
proves for every description. Given \(D \in \FATCD{}\), is there a
well-formed record that lists the six role channels, gives each a
status, and accounts for every remaining raw feature through an
inactive status record or a residual classification? This section shows that there is, and explains why a
channel is not the same as an active role. The lower bounds show that each
of the six channels is materially realizable in the domain; the
upper-bound classification shows that, under the recognition hypothesis, no
kind residual or record residual is licensed label (x) or (xi).

\subsection{Decomposition records and channel statuses}\label{subsec:dec-records}

\begin{definition}[Channel status]\label{def:channel-status}
For \(D\in\FATCD{}\), a decomposition record \(\Dec{D}\)
(Definition~\ref{def:decomposition-record}) and \(i \in \{1,\ldots,6\}\), the
\emph{role channel} \(\Chan{i}{D}\) is the pair of \(i\) and the status entry a
record \(\Dec{D}\) chooses for \(i\); the status value \(\Status{i}{D}\) is fixed
below from \(D\) alone, and the record only chooses the cluster, offer, seed or
declaration it names; the channel exists whether or not the status is
active. A \emph{channel status} for
\(\Chan{i}{D}\) is exactly one of the five values
\begin{itemize}[topsep=2pt,itemsep=1pt]
  \item \(\activeproj\) --- an admissible \(\pi_i(D)\) exists with \(W_i(C)\)
    and all required attachments;
  \item \(\collapsedbc\) --- no admissible \(\pi_i(D)\) exists, and an
    honestly recorded offer of a specified cluster has check result
    \(W_i(C)=\mathrm{false}\);
  \item \(\belowthr\) --- neither preceding case applies, and a covered seed for
    \(P_i\) is present;
  \item \(\outsidescope\) --- none of the preceding cases applies, and an
    honest outside-scope declaration is present;
  \item \(\absentrec\) --- none of the preceding cases applies; the
    decomposition records that absence.
\end{itemize}
In prose these are called active, collapsed, below threshold, outside scope and absent. The status name \(\collapsedbc\) refers to a failed offer; it is unrelated to
the collapse-case annotations of Definition~\ref{def:threshold-attachment} and
to the forbidden collapses of Table~\ref{tab:boundary-matrix}.
The status \(\Status{i}{D}\) depends only on \(D\); \(\Chan{i}{D}\) and
\(\Residual{D}\) are taken relative to a fixed record \(\Dec{D}\), and we write
\(\Chan{i}{\Dec{D}}\) and \(\Residual{\Dec{D}}\) when the record matters. Since
\(\Chan{i}{D}\) contains its status entry, the pair
\((\Chan{i}{D},\Status{i}{D})\) in Definition~\ref{def:decomposition-record}
repeats the status as a separate component. A \emph{covered seed} for \(P_i\) is, respectively:
\begin{enumerate}[label=(\arabic*),topsep=2pt,itemsep=1pt]
  \item a map or partial-map record whose
declared use is \emph{update}, together with a relation record of declared use quotient that it
references as its source relation, or with the relation record
serving as source relation in a cluster satisfying \(W_1\) that contains it;
  \item a predicate record
whose predicate field is declared as a macro-specification whose formula has
exactly one free variable \(r\), occurring in it and ranging over records of
the candidates' kind, and whose evaluations recorded, as in the \(W_2\) row, in a field of the realization relation or of a comparison or check record of \(D\), at declared records of
that kind, are correct (each equals the value of the formula at that record in
\(D\)) and include both truth values, together with a
realization relation (a relation record of declared use realization) whose target field is this predicate record and whose declared candidate class is a nonempty set of declared records of that kind; the seed consists of these two records, and either may be named as carrying it; the named seed fields are the formula field of the predicate record and, on
the realization-relation record, its candidate-class and target fields, its
list of satisfying members, its recorded evaluations of that formula (as datum
values or as certificates \((s(t),b)\) at declared records \(t\) of the
candidates' kind), and its certificate exhibiting a satisfying realizer or
checking that none exists; no other assertion field of either record is a named
seed field;
  \item a certificate, of one of the two types in the \(W_3\) row and naming as inputs two path or derivation records with the same declared endpoints, that these are not related by the rewrite rules recorded in \(D\), carried in a field of a content record of \(D\);
  \item a
law-bearing stage, order, transition or refinement record;
  \item a relation record whose
declared use is \emph{quotient}, or a map record whose declared use is \emph{packaging};
  \item an event, provenance or trace record with its declared event or trace
identity and reference fields, or a replay or check record with its declared
trace reference and a certificate whose formula has exactly one of the two forms prescribed for the \(W_6\) row after the witness table, with that trace as \(\tau\) (annotations are not
seeds).
\end{enumerate}
Unlike the no-realizer branch of the \(W_2\) row, clause (2) does not require
the candidate class to be the full declared domain of a packaging or update
map. Every closed check-language formula without \(W_j\) atoms admits the
seed-(2) presentation of Remark~\ref{rem:recognition-open}, provided the completed wrapped description lies in \FATCD{} and contains two distinct
declared records of one kind (paragraph \emph{Wrapping a closed formula as a \Ptwo{} seed}). Such a seed
supplies recognized \Ptwo{} content but does not by itself certify \(W_2(C)\) or
an active \Ptwo{} projection. In ``covered seed'', ``covered''
names the seed clauses (1)--(6); it is not a reference to the covered scope.
\par The named seed fields are: in (1), the graph of the update map, any
source-relation reference it carries, and the class data of the source
relation; in (3), the
fields of the comparison certificate; in (4), the law certificate and the
endpoint or index fields it checks; in (5), the class data of the relation or
the graph and quotient-object reference of the map; and in (2) and (6), the
fields named there.
\par A record \emph{carries} a seed when it is the map record in
(1), the predicate or realization-relation record in (2), the record holding
the certificate in (3), the law-bearing record in (4), the relation or map
record in (5), or the event, provenance, trace, replay or check record in (6);
a named seed field lying on another record is an input to that record's
seed. A comparison or check record whose recorded evaluations of a clause-(2) formula are used to meet that clause does not carry that seed. This does not exclude recognition through other fields or role data. Its classification must be checked against all the criteria of Definition~\ref{def:residual-scheme}; if none of (i)--(ix) holds, it is licensed (x). By the witness table, every cluster satisfying \(W_i\) contains a
covered seed for \(P_i\) (for \(i=2\) the row, like seed clause (2), requires the recorded evaluations
to be correct): the update map with its source relation data, the
macro-specification with its realization relation, the comparison certificate
of the two routes, the law-bearing record, the quotient relation or packaging
map, and the event or trace record, respectively. A \emph{failed offer} is an
honestly recorded attempt to use a specified cluster \(C\) as a \(P_i\)
witness for which the recorded check gives \(W_i(C)=\mathrm{false}\). An \emph{outside-scope declaration} for \(P_i\) is a claim annotation naming
\(P_i\) and a record \(r\), with no witness-cluster field (so it is not an offer
and its audit records no check of \(W_i\)), whose referenced record \(r\) lies outside the covered scope (the
criterion after Definition~\ref{def:residual-scheme}) and whose kind occurs in
the \(W_i\) row of the witness table; its honest audit lists the
outside-scope field of \(r\) in its non-certification field. The status of \(\Chan{i}{D}\) is the
first that applies in the order: \(\activeproj\) if some cluster supports
an admissible derived projection \(\pi_i(D)\) in the sense of
Definition~\ref{def:role-projection}, including \(W_i(C)\) and every
applicable attachment of Definition~\ref{def:threshold-attachment}; \(\collapsedbc\) if there is a
failed offer; \(\belowthr\) if a covered seed is present;
\(\outsidescope\) if there is an outside-scope declaration; and
\(\absentrec\) otherwise. An inactive status is recorded in \(\Dec{D}\); it
is not a role witness. We write \(\Status{i}{D}\) for the channel status of
\(\Chan{i}{D}\).
\end{definition}

Statuses are assigned only by the ordered rule above. An admissible
projection takes priority over every failed offer or incomplete seed. In its
absence, a failed offer gives \(\collapsedbc\); otherwise a covered seed
lacking a complete admissible projection gives \(\belowthr\). Only then are an
outside-scope declaration and the absent fallback considered.
Proposition~\ref{prop:boundary-distinctions} uses the same two labels for
local attempt verdicts, which need not equal channel statuses. In the running
example, the \Pfour{} channel is \(\belowthr\): its transition record is
present, but no threshold annotation targets a cluster containing it.

\begin{definition}[Decomposition/exhaustion record]\label{def:decomposition-record}
A \emph{decomposition/exhaustion record} for \(D \in \FATCD{}\) is a tuple
\[
\Dec{D} = \bigl(\Raw{D},\; (\Chan{i}{D}, \Status{i}{D})_{i=1}^{6},\;
\Residual{D},\; \mathrm{Classify},\; \mathcal{L},\; \mathcal{V},\;
\Audit{D},\; \mathcal{N}\bigr)
\]
where
\begin{itemize}[topsep=2pt,itemsep=1pt]
  \item the six role channels together with their statuses form the
    six-channel status sequence, each \(\Status{i}{D}\) being the value that the
    ordered rule of Definition~\ref{def:channel-status} assigns. If \(\Status{i}{D} = \activeproj\), the channel
    carries an active projection \(\pi_i(D)\) with its full attachment
    data. If \(\Status{i}{D}\) is inactive, the channel carries the
    corresponding below-threshold, collapsed-boundary,
    absent-with-record, or outside-scope status record;
  \item \(\Residual{D} \subseteq \Raw{D}\) is the residual feature set,
    defined as follows. Let \(A(\Dec{D})\) be the union of the content
    clusters of the active projections and the attachment annotations (the
    threshold, witness-schema, level and collapse-case annotations attached to
    the cluster, the claim annotation naming \(P_i\) and the cluster, the audit
    annotations referring to that claim, the nonclaim annotations targeting it, the translation and lift or
    forgetting annotations that Definition~\ref{def:role-projection} requires for
    a level crossing of the projection, and,
    for \(P_6\), the event-or-trace claim and honest audit used to establish
    the \(W_6\) audit-entry conjunct)
    that the channel entry assigns to its projection (an active entry assigns at least one annotation of each kind that Definition~\ref{def:role-projection} requires for admissibility and may leave further annotations targeting the cluster unassigned). Let \(S(\Dec{D})\) be the union of the
    clusters and offers explicitly named by the inactive status entries: the
    records of the failed offered cluster with the offer's claim annotation and the audits referring to it, for \(\collapsedbc\); all records of the covered seed it names (for clauses (1) and (2), both records the clause names), for
    \(\belowthr\); and the outside-scope declaration with the audits referring to it, for
    \(\outsidescope\); \(\absentrec\) accounts for no feature, and a status
    entry accounts only for records it names. Then
    \(\Residual{D}=\Raw{D}\setminus\bigl(A(\Dec{D})\cup S(\Dec{D})\bigr)\),
    relative to the chosen record, and \(\mathrm{Classify}\) assigns to each
    \(r \in \Residual{D}\), including any bookkeeping annotation not assigned
    to a channel, a label in the scheme of
    Definition~\ref{def:residual-scheme}; already covered bookkeeping may
    receive (ix);
  \item \(\mathcal{L}\) collects the loss and lift records for any level
    crossings involved in projections or classifications;
  \item \(\mathcal{V}\) collects instrument-visibility and
    empirical-bridge annotations where these are asserted;
  \item \(\Audit{D}\) is the honest bookkeeping subrecord;
  \item \(\mathcal{N}\) is the explicit nonclaim record.
\end{itemize}
A decomposition/exhaustion record is \emph{well-formed} when each
\(\Status{i}{D}\) is drawn from the five values of
Definition~\ref{def:channel-status}, each active
\(\pi_i(D)\) is admissible in the sense of
Definitions~\ref{def:role-projection} and~\ref{def:threshold-attachment},
each inactive status record names a failed offer, covered seed or
outside-scope declaration of \(D\) for that channel, as
Definition~\ref{def:channel-status} requires (none for \(\absentrec\)), and every raw feature \(r\) lies in \(A(\Dec{D})\), or in \(S(\Dec{D})\), or in
\(\Residual{D}\) and receives a label, and, for each
\(r\in\Residual{D}\), \(\mathrm{Classify}(r)\) is a label whose criterion in the
table after Definition~\ref{def:residual-scheme} holds for \(r\) relative to
\(\Dec{D}\). The records in
\(\mathcal{L}\) and \(\mathcal{V}\) are auxiliary annotations on this
coverage, not additional role channels. The status entries and the
collections \(\mathcal{L}\), \(\mathcal{V}\) and \(\mathcal{N}\) are data
of \(\Dec{D}\). Each cluster, seed, offer or declaration named by a status is
selected from, or supported by, the records of \(\Raw{D}\) and \(\Audit{D}\)
required above; the selection and the status entries need not themselves be
raw records. Forming \(\Dec{D}\) does not change \(D\).
\end{definition}

By the ordered rule of Definition~\ref{def:channel-status}, the status
\(\Status{i}{D}\) depends only on \(D\), not on the decomposition record; the
cluster, offer, seed or declaration that a record names, and hence
\(A(\Dec{D})\), \(S(\Dec{D})\) and \(\Residual{D}\), may depend on the record.
Because \(\Residual{D}\) leaves out every feature already accounted for by
the active projection named for its channel or by an inactive status record, labels~(i)
and~(ii) of Definition~\ref{def:residual-scheme} do not occur for a residual
feature of \(\Dec{D}\): both require the record to lie in a cluster named for
an active projection. This is why the results below list labels (iii)--(ix).

\paragraph{Two residual sets.} A residual feature in the sense of
Definition~\ref{def:residual-scheme} is also called a \emph{kind residual}, in
statements and elsewhere, and an element of \(\Residual{D}\) a \emph{record
residual}.
Membership in the first depends on the raw-grammar kind of a feature, and
membership in the second on whether the record \(\Dec{D}\) accounts for it
through a channel. In particular, a bookkeeping annotation is not a kind
residual but may be a record residual if the chosen channel assignments do not
account for it. A kind residual absorbed into an active projection is not a
record residual, which is why label~(i) does not occur among record residuals.
The recognition hypothesis and
Theorem~\ref{thm:upper-bound-classification} cover both sets.


\subsection{Existence, six-channel exactness, and residual exhaustion}\label{subsec:dec-theorems}

\begin{theorem}[Decomposition existence]\label{thm:decomposition-existence}
Every description \(D \in \FATCD{}\) admits a well-formed decomposition/exhaustion
record \(\Dec{D}\) in the sense of
Definition~\ref{def:decomposition-record}.
\end{theorem}

\begin{proof}
For each \(i \in \{1,\ldots,6\}\), range over the finitely many clusters
\(C\subseteq\Raw{D}\). If some \(C\) satisfies \(W_i(C)\) together with every
item of Definitions~\ref{def:threshold-attachment}
and~\ref{def:role-projection}, name one such \(C\), record
\(\Status{i}{D} = \activeproj\) and the active projection
\(\pi_i(D)\) it determines; otherwise record the inactive status that the
ordered rule of Definition~\ref{def:channel-status} assigns, and name one
failed offer, covered seed or outside-scope declaration for the first case of
the rule that applies. Six status entries
result, one per role. The raw features absorbed by these channel
records are removed from the residual stage: the statuses, their
supporting clusters and their assigned annotations are selected first, and
\(\Residual{D}\) is then formed by the set difference of
Definition~\ref{def:decomposition-record}. For each \(r\) in it, choose a
licensed label among (i)--(ix) if one applies. If none applies, then \(r\) is
not outside the covered scope (otherwise (viii) would apply), \(D\) is finite,
and every claim concerning \(r\) is referred to by an honest audit annotation; so the criterion of
(x) holds. Thus \(\mathrm{Classify}\) is total, and the record is well-formed
irrespective of which labels occur. An unassigned bookkeeping annotation is
assessed by the same criteria, ordinarily (ix). Relative to the chosen
record, (i) cannot apply to a record residual, since absorbed content was
placed in \(A(\Dec{D})\); under the recognition hypothesis
(Definition~\ref{def:recognition-hypothesis}),
Theorem~\ref{thm:upper-bound-classification} further places every label in
(iii)--(ix), with (x) and (xi) empty. The records \(\mathcal{L}\) and \(\mathcal{V}\) are
read from \(\Raw{D}\) and \(\Audit{D}\); the bookkeeping subrecord is
\(\Audit{D}\); and \(\mathcal{N}\) records the nonclaims carried
forward from the witness discipline. The resulting tuple is a
decomposition/exhaustion record, well-formed by construction.
Non-vacuity of the six channels follows from
Theorem~\ref{thm:six-role-lower-bounds}, which supplies, for each
channel, a description in which that channel is active.
\end{proof}

\begin{theorem}[Six-channel exactness and residual exhaustion]\label{thm:six-channel-exactness}
For every \(D \in \FATCD{}\) and every well-formed
\(\Dec{D}\) in the sense of
Definition~\ref{def:decomposition-record}:
\begin{enumerate}[label=(\roman*),topsep=2pt,itemsep=1pt]
  \item the channel sequence \((\Chan{i}{D})_{i=1}^{6}\) has length
    exactly six, one per primitive role \Pone{}, \ldots, \Psix{};
  \item every raw feature of \(D\) is accounted for at the
    channel/exhaustion level by an active projection \(\pi_i(D)\), a
    channel-status record for an inactive channel, or a residual
    classification under the scheme of
    Definition~\ref{def:residual-scheme}, possibly with label (x) or (xi);
    since label (x) applies to in-scope features for which no criterion
    (i)--(ix) holds, this clause alone does not exclude an unresolved
    candidate (see (iii)); the records in \(\mathcal{L}\)
    and \(\mathcal{V}\) are auxiliary annotations on this coverage;
  \item under the recognition hypothesis
    (Definition~\ref{def:recognition-hypothesis}) and within the covered scope,
    no raw feature is classified as an unresolved seventh-role threat or as
    passing the recorded seventh-role screen
    (Theorem~\ref{thm:upper-bound-classification}); the three principal threat
    classes are closed without the recognition hypothesis, in the
    recognized forms of
    Propositions~\ref{prop:probability-closure}--\ref{prop:agency-closure};
  \item the status sequence \((\Status{i}{D})_{i=1}^{6}\) is the same for every
    well-formed \(\Dec{D}\).
\end{enumerate}
\end{theorem}

\begin{proof}
Exactness of the channel sequence is built into the definition of
\(\Dec{D}\); there are six role channels by
Definition~\ref{def:decomposition-record}. Coverage of raw features is
the well-formedness condition of
Definition~\ref{def:decomposition-record} together with totality of
\(\mathrm{Classify}\) on \(\Residual{D}\), which holds by the case argument in
the proof of Theorem~\ref{thm:decomposition-existence}. Emptiness of categories (x) and (xi),
under the recognition hypothesis (Definition~\ref{def:recognition-hypothesis})
and within the covered scope, is
Theorem~\ref{thm:upper-bound-classification} directly. The last assertion of
(iii) is Propositions~\ref{prop:probability-closure},
\ref{prop:causality-closure} and~\ref{prop:agency-closure}, which hold for
every \(D\in\FATCD{}\) and every well-formed \(\Dec{D}\) without the
recognition hypothesis. For (iv), Definition~\ref{def:decomposition-record}
records each status by the ordered rule of Definition~\ref{def:channel-status},
whose clauses refer only to \(D\).
\end{proof}

\subsection{Channel-versus-activation discipline}\label{subsec:channel-vs-activation}

Theorem~\ref{thm:six-channel-exactness} asserts six channels in every
decomposition/exhaustion record. It does not assert six active
projections in every record, nor does it assert that a well-formed
record is unique. The distinction between channel presence and
channel activation prevents six-channel exactness from being read as
an all-six-active claim.

Membership of \(D_1,\dots,D_6\) in \(\FATCD_{\mathrm{RH}}\) is verified record by
record in the proof of Theorem~\ref{thm:scoped-exact-six}; that verification
uses only Theorem~\ref{thm:six-role-lower-bounds} and the definitions. In brief:
every content record of \(D_i\) contributes a required field to \(C_i\) or
carries a covered seed, each of its assertion fields being required or a named
seed field (R1), or it has no assertion field (R1 or R3); every structural
annotation of \(D_i\) carries only its typing, target and label or level field,
and every bookkeeping annotation only datum and reference fields besides its
exempt fields (R2).

\begin{proposition}[Channel-versus-activation discipline]\label{prop:channel-vs-activation}
Let \(D \in \FATCD{}\) and let \(\Dec{D}\) be a well-formed decomposition/exhaustion
record. The following are distinct statements:
\begin{enumerate}[label=(\roman*),topsep=2pt,itemsep=1pt]
  \item \emph{Six channels}: the channel sequence
    \((\Chan{i}{D})_{i=1}^{6}\) has length six with the fixed role names
    \Pone{}, \ldots, \Psix{}.
  \item \emph{Six active projections}: for each \(i\),
    \(\Status{i}{D} = \activeproj\) and the channel carries an active
    \(\pi_i(D)\).
\end{enumerate}
Theorem~\ref{thm:six-channel-exactness} establishes (i) for every
\(D \in \FATCD{}\). It does not establish (ii). In particular, an
admissible \(D \in \FATCD{}\) may have any number of active channels
between zero and six; inactive channels are recorded with status
\(\belowthr\), \(\collapsedbc\), \(\absentrec\), or \(\outsidescope\)
and are not active witnesses for the corresponding roles. For each
\(k\in\{0,\dots,6\}\) some \(D\in\FATCD_{\mathrm{RH}}\) has exactly \(k\) active
channels, and some single description in \(\FATCD_{\mathrm{RH}}\) has channels of all
five statuses.
\end{proposition}

\begin{proof}
Statements (i) and (ii) differ whenever some channel is inactive, which the
last clause provides for every \(k<6\). \par\emph{Exactly \(k\) active channels.} For \(k=0\), let \(D_0\) consist of a
single object record: it is finite and closed, it has no claims, so it is
audited, and it has no cluster satisfying any \(W_i\), no failed offer, no
covered seed and no outside-scope declaration, so the status rule of
Definition~\ref{def:channel-status} gives every channel \(\absentrec\). Its
object record has no witness field and no predicate, comparison, invariant,
law or claim field, so it satisfies (R3) of
Definition~\ref{def:recognition-hypothesis} and
\(D_0\in\FATCD_{\mathrm{RH}}\). For
\(1\le k\le 6\), choose \(k\) distinct indices \(S\) and let \(D_S\) be the
disjoint union of the descriptions \(D_i\), \(i\in S\), of
Theorem~\ref{thm:six-role-lower-bounds}, with record identifiers made
disjoint by an injective renaming applied recursively to every
record-identifier occurrence in reference, datum, term-valued and
formula-valued fields, including identifiers inside finite lists,
tuples, relation classes and map tables, and the argument lists of
\(W_j\) atoms. Rational constants and other non-identifier values
are unchanged. A numeral in record-sorted position (paragraph \emph{Sorts} of Section~\ref{sec:primitive-role-architecture}), such as \(0\) and \(2\) in the formula of \(s\) in \(D_2\), is a record-identifier occurrence and is renamed; a numeral in rational-sorted position, such as an element of the monoid list of \(D_3\), is not. Each condition of Definition~\ref{def:fatcd} concerns a record, an
annotation or a claim together with the records it refers to, which lie in one component, except that certificate correctness and the evaluation of \(W_j\) atoms may read description-wide data (all declared records or those of one kind, the rule list \(R\), \(\mathrm{Audit}_0(D)\)); the next sentences check that this data leaves every recomputed value unchanged. The certificates' quantifiers are bounded by listed object records or finite field lists, except \(\forall u\in|E|\), which is guarded by \(E(u,u)\); this is false at every record of another component, since renaming keeps the classes of \(E\) inside its component. So the content checks retain their values in \(D_S\). The \(W_1,W_2,W_4,W_5\) checks remain within their components. The
\(W_3\) row also requires the rules of its cluster to be all rewrite rules
recorded in the description; since the indices in \(S\) are distinct and \(D_3\)
is the only component recording rewrite rules, this clause keeps its value in
\(D_S\). For
\(W_6(C_6)\), the honest event-or-trace audit entry of \(D_6\) remains in
\(\mathrm{Audit}_0(D_S)\), so its separate audit conjunct also remains true.
Thus the audits remain honest and \(D_S\in\FATCD{}\). For \(i\in S\), the cluster \(C_i\)
with its attachments still satisfies \(W_i\), so channel \(P_i\) is
\activeproj{}. For \(j\notin S\), no component has a witness-schema
annotation naming \(P_j\), together with a threshold and an honest audit of
\(W_j(C)\) for a cluster \(C\), and making record identifiers disjoint creates
no such annotation or cross-component reference. So no admissible
\(\pi_j(D_S,C)\) in the sense of Definition~\ref{def:role-projection} exists,
and channel \(P_j\) is not active. Hence \(D_S\) has exactly \(k\) active
channels.
Each of these descriptions lies in \(\FATCD_{\mathrm{RH}}\).
The inventory in the proof of the last clause of
Theorem~\ref{thm:scoped-exact-six} (Section~\ref{sec:scoped-exact-six}, p.~\pageref{thm:scoped-exact-six}), which uses
Theorem~\ref{thm:six-role-lower-bounds} and the definitions but not this
proposition, establishes (R1) or (R3) for every content record of each \(D_i\)
using records and checks within that component, and (R2) for its
structural and bookkeeping annotations. Under this renaming each
certificate, claim statement and audit conclusion of \(D_i\) becomes its renamed
image, which recomputes to the same value in \(D_S\). They apply to every
kind residual and every possible record residual of \(D_S\),
independently of the chosen decomposition record. The case \(D_0\)
was checked above.
\par\emph{All five statuses.} For the five statuses, \(D_S\) supplies \(\activeproj\) and \(D_0\) supplies
\(\absentrec\). Let \(o\) be the object record of \(D_0\). Let \(D_{\mathrm c}\) consist of state
records \(x\ne y\), an object record \(X_{\mathrm c}\) listing \(x\) and \(y\), a
relation record \(E\) on \(X_{\mathrm c}\) of declared use quotient with the
single class \(\{x,y\}\), a claim naming \Pfive{} and \(C=\{x,y,E\}\), and an
honest audit listing \((W_5(C),\mathrm{false})\); no packaging map is present, so
this failed offer lies on a cluster carrying a covered \Pfive{} seed, and channel
\Pfive{} is \(\collapsedbc\). Here \(x\), \(y\) and \(X_{\mathrm c}\) satisfy (R3), \(E\) carries a
covered seed whose class data are named seed fields and satisfies (R1), and the
annotations satisfy (R2). Let \(D_{\mathrm{o}}\) be \(D_0\) together with a predicate record \(Q\) (here \(Q\) is a predicate record, unrelated to the quotient objects \(Q\) of the \(W_5\) row and Remark~\ref{rem:recognition-open}, and \(o\) is the object record of \(D_0\), not the order record of \(D_{\mathrm{scr}}\)) of arity
one, its argument place declared to range over object records, with no formula field and no relation record presenting it, a predicate
record \(s\) whose only assertion field is \(Q(o)\), a claim naming \Ptwo{} and
\(s\) with statement \(Q(o)\), and an honest audit that lists no certificate,
records no check of \(W_2\) (so the claim is not an offer), concludes \(o=o\)
and lists the field of \(s\) in its non-certification field. This claim is an
outside-scope declaration, so channel \Ptwo{} of \(D_{\mathrm{o}}\) is
\(\outsidescope\). The declaration is not an admissible claim in the sense of
Definition~\ref{prop:honest-bookkeeping}, since no audit certifies \(Q(o)\);
Definition~\ref{def:channel-status} requires only that its audit be honest, and
the conclusion \(o=o\) makes the audit honest while certifying nothing about
\(Q\). In \(D_{\mathrm{o}}\) \(o\) and \(Q\)
satisfy (R3), \(s\) satisfies (R5), and the annotations satisfy (R2). The
description \(D_{\mathrm{ex}}\) of Section~\ref{sec:fatcd-domain} supplies
\(\belowthr\) for its \Pfour{} channel. It lies in \(\FATCD_{\mathrm{RH}}\). The records \(F\), \(E\), the
split-pair comparison, \(q\), \(\{A,B\}\), the quotient-validity check, the
event, the trace and the check record each contribute a required field to
\(C_1\), \(C_5\) or \(C_6\), and their assertion fields (the split-pair,
validity and replay certificates) are required for \(W_1\), \(W_5\) and \(W_6\)
respectively, so these records satisfy (R1). The provenance record carries a
covered \Psix{} seed, and the transition record carries a covered \Pfour{} seed
whose law field is a named seed field; both satisfy (R1). The state records,
\(S_0\), the object records \(A\), \(B\) and the replay record have no assertion field and satisfy (R1) or
(R3); the index record satisfies (R4) by the one-step chain to \(E\), which satisfies
(R1); and its annotations satisfy (R2).
\par\emph{All five in one description.} Let \(D_{\mathrm{all}}\) be the renamed disjoint union of \(D_6\),
\(D_{\mathrm c}\), \(D_{\mathrm o}\) and \(D_T\), where \(D_T\) has a state record
\(z\), an object record \(Z\) listing it, and a transition record \(T:Z\to Z\),
\(T(z)=z\), whose law field carries the certificate, recomputed true, that \(T\)
is defined on \(Z\) with values in \(Z\). The argument for \(D_S\) gives
\(D_{\mathrm{all}}\in\FATCD{}\); \(T\) carries a covered \Pfour{} seed with named
law field and satisfies (R1), \(z\) and \(Z\) satisfy (R1) or (R3), and the other
components are inventoried above, so \(D_{\mathrm{all}}\in\FATCD_{\mathrm{RH}}\).
No component supplies a \Pone{}, \Ptwo{} or \Pthree{} seed, an active projection
other than that of \(D_6\), or a failed offer other than that of
\(D_{\mathrm c}\); so by the ordered rule \Psix{} is \(\activeproj\), \Pfive{}
\(\collapsedbc\), \Pfour{} \(\belowthr\), \Ptwo{} \(\outsidescope\), and \Pone{}
and \Pthree{} \(\absentrec\).
\end{proof}

\begin{proposition}[Decomposition records are not unique]\label{prop:non-uniqueness}
A description \(D \in \FATCD{}\) may admit more than one well-formed
decomposition/exhaustion record in the sense of
Definition~\ref{def:decomposition-record}. In particular, distinct
well-formed records may name different witness clusters for one active
channel, and then have different sets \(A(\Dec{D})\), \(\Residual{D}\) and
residual labels; Definition~\ref{def:decomposition-record} also leaves the choice among
level annotations attached to the named cluster and the placement of
loss/lift and visibility records to the record. Such a description can be chosen in \(\FATCD_{\mathrm{RH}}\).
Theorems~\ref{thm:decomposition-existence}
and~\ref{thm:six-channel-exactness} assert existence, not uniqueness.
\end{proposition}

\begin{proof}
Let \(D\) be the disjoint union of two copies of the description \(D_5\) of
Theorem~\ref{thm:six-role-lower-bounds} (with the same renaming), as in the proof of
Proposition~\ref{prop:channel-vs-activation}; then \(D\in\FATCD{}\), and each
copy carries a cluster, \(C_5\) and \(C_5'\), that satisfies \(W_5\) with its
attachments. The construction of Theorem~\ref{thm:decomposition-existence}
may name either cluster for the \Pfive{} channel. In the record naming
\(C_5\), the content records of \(C_5'\) that contribute a required field to
\(C_5'\), a six-role-aligned cluster that the chosen projection does not
absorb, receive label (v), and any others, having no predicate, comparison,
invariant, law or claim field and supplying no witness field, receive (ix-c); in the record naming \(C_5'\), the same holds with the
copies exchanged. The \(W_5\) assertion of the unchosen copy is represented by
its check record and other required-field records, labelled (v); its
bookkeeping annotations receive (ix-a) and its threshold and level annotations
(ix-b).
The two records name different clusters, so they differ, and both are
well-formed. The description \(D\) satisfies the recognition hypothesis: each
content record of either copy supplies a witness field to that copy's cluster
satisfying \(W_5\), or is an input without assertion fields, and so satisfies
(R1) or (R3) as in the inventory in the proof of
Theorem~\ref{thm:scoped-exact-six} (p.~\pageref{thm:scoped-exact-six}); its annotations satisfy (R2). Hence
\(D\in\FATCD_{\mathrm{RH}}\).
\end{proof}

\paragraph{The running example, decomposed.} For the description
\(D_{\mathrm{ex}}\) of Section~\ref{sec:fatcd-domain}, a well-formed
decomposition record assigns the six channels the following statuses; the
records named are those of the table in Section~\ref{sec:fatcd-domain}. The
channels \Pone{}, \Pfive{} and \Psix{} have status \activeproj{}: the split pair
\((0,1)\) witnesses the descent obstruction, the quotient object \(\{A,B\}\)
with its packaging map witnesses the packaging, and the replayed and checked
trace witnesses the audit. The channel \Pfour{} has status \belowthr{}: the
transition record of \(F\) is present, but no threshold annotation targets a
cluster containing it. The channels \Ptwo{} and \Pthree{} have status \absentrec{}: the
description records no macro-specification and no path or derivation record. The index
record copying the classes of \(E\) is a residual feature, classified as a
presentation variant of the relation record. So \(D_{\mathrm{ex}}\) carries six
channels and three active projections, as Proposition~\ref{prop:channel-vs-activation}
allows.

\begin{remark}[What the decomposition theorems do and do not prove]\label{rem:dec-scope}
Theorems~\ref{thm:decomposition-existence}
and~\ref{thm:six-channel-exactness} establish the structural side of
the scoped exact-six claim: every \(D \in \FATCD{}\) admits a
six-channel decomposition/exhaustion record, with residual exhaustion in the
covered scope under the recognition hypothesis
(Definition~\ref{def:recognition-hypothesis}). Together with the lower-bound theorem
(Section~\ref{sec:lower-bounds}) and the upper-bound classification
(Section~\ref{sec:upper-bound-classification}), they prepare the
scoped exact-six theorem but do not replace the integrated-stress
check that follows. They do not assert that every description
activates all six roles, that the decomposition is canonical or
unique, or that the scoped exact-six program extends beyond \FATCD{};
those remain nonclaims.
\end{remark}

 \section{Checks against misreadings}\label{sec:integrated-stress}



In this section a misreading is \emph{excluded} when it is not licensed by the
results of
Sections~\ref{sec:fatcd-domain}--\ref{sec:decomposition-exhaustion}.
The section is a checklist. Each of its six remarks lists misreadings of
the results of Sections~\ref{sec:lower-bounds}--\ref{sec:decomposition-exhaustion},
for example asserting exact six without the lower bounds, or reading a
role channel as an active role, and names the earlier result that rules
each misreading out; ``stress'' means a check against such a list, and the
chain \emph{satisfies} a stress family when the result named for each
misreading in its list rules that misreading out.
Proposition~\ref{thm:integrated-stress} checks that the earlier results
block the misreadings listed in the six remarks and combine into the terminal statement.
The section adds no new mathematics; it collects in one place what the
results do not permit.

\begin{center}
\small
\begin{tabularx}{\linewidth}{@{}>{\raggedright\arraybackslash}p{0.34\linewidth}>{\raggedright\arraybackslash}X@{}}
\toprule
Misreadings of & Blocked by \\
\midrule
dependency (Remark~\ref{prop:dependency-stress}) & the \FATCD{} domain and the order of Theorems~\ref{thm:six-role-lower-bounds}, \ref{thm:upper-bound-classification}, \ref{thm:decomposition-existence} and~\ref{thm:six-channel-exactness} \\
circularity (Remark~\ref{prop:circularity-stress}) & Proposition~\ref{prop:non-circularity} and the direct proof of Theorem~\ref{thm:decomposition-existence} \\
channel versus activation (Remark~\ref{prop:channel-activation-stress}) & Proposition~\ref{prop:channel-vs-activation} and Definition~\ref{def:channel-status} \\
role alignment (Remark~\ref{prop:role-alignment-stress}) & Proposition~\ref{prop:alignment-rules} \\
upper-bound residual (Remark~\ref{prop:residual-stress}) & Propositions~\ref{prop:probability-closure}--\ref{prop:agency-closure} and Theorem~\ref{thm:upper-bound-classification} \\
overclaim (Remark~\ref{prop:overclaim-stress}) & Remark~\ref{rem:atlas-global-nonclaims} and Propositions~\ref{prop:nonclaim-closure} and~\ref{prop:non-uniqueness} \\
\bottomrule
\end{tabularx}
\end{center}

\subsection{Dependency stress}\label{subsec:dependency-stress}

\begin{remark}[Dependency stress]\label{prop:dependency-stress}
The scoped exact-six target \(\ScopedExactSix{\FATCD_{\mathrm{RH}}}\)
depends on the \FATCD{} domain of Section~\ref{sec:fatcd-domain}, the
upper-bound classification of Theorem~\ref{thm:upper-bound-classification},
and the decomposition and exhaustion theorems of
Section~\ref{sec:decomposition-exhaustion}; the last clause of
Theorem~\ref{thm:scoped-exact-six} also depends on the lower-bound theorem of
Section~\ref{sec:lower-bounds}. Accordingly, the following are excluded:
\begin{itemize}[topsep=2pt,itemsep=1pt]
  \item asserting the target without the \FATCD{} domain;
  \item asserting the target without the upper-bound classification;
  \item asserting the target without the decomposition and exhaustion
    theorems;
  \item asserting that every channel is active somewhere in
    \(\FATCD_{\mathrm{RH}}\) without the lower-bound theorem.
\end{itemize}
The integrated stress check of Proposition~\ref{thm:integrated-stress} is a
separate consistency check; the proof of Theorem~\ref{thm:scoped-exact-six}
does not use it.
\end{remark}

\subsection{Circularity stress}\label{subsec:circularity-stress}

\begin{remark}[Circularity stress]\label{prop:circularity-stress}
Three circularity patterns are excluded under the theorem chain:
\begin{itemize}[topsep=2pt,itemsep=1pt]
  \item redefining \FATCD{} so that it already contains exactly six role
    slots by construction;
  \item declaring the absence of a seventh role by a membership condition
    of \FATCD{} rather than by an explicitly stated hypothesis;
  \item justifying decomposition existence by assuming scoped
    exact-six.
\end{itemize}
The non-circularity schema of Proposition~\ref{prop:non-circularity} and
the direct proof of Theorem~\ref{thm:decomposition-existence}
block each pattern.
\end{remark}

\subsection{Channel-versus-activation stress}\label{subsec:channel-activation-stress}

\begin{remark}[Channel-versus-activation stress]\label{prop:channel-activation-stress}
The following conflations are excluded:
\begin{itemize}[topsep=2pt,itemsep=1pt]
  \item requiring that every \(D \in \FATCD{}\) activates all six role
    channels;
  \item treating an inactive channel of status \(\belowthr\) as an
    active witness for its role;
  \item treating an inactive channel of status \(\collapsedbc\) as an
    active witness;
  \item treating an inactive channel of status \(\absentrec\) as an
    active witness.
\end{itemize}
Proposition~\ref{prop:channel-vs-activation} enforces each distinction.
The same inactive-status check also applies to
\(\outsidescope\).
\end{remark}

\subsection{Role-alignment stress}\label{subsec:role-alignment-stress}

\begin{remark}[Role-alignment stress]\label{prop:role-alignment-stress}
The three forbidden assignments of
Remark~\ref{rem:three-repaired-alignments} are excluded: generic
transition data does not project to \Pthree{}, generic path or
derivation data does not project to \Pfive{}, and generic semantic or
interpretation data does not project to \Pone{}. The alignment rules
of Proposition~\ref{prop:alignment-rules} enforce each prohibition.
\end{remark}

\subsection{Upper-bound residual stress}\label{subsec:upper-bound-residual-stress}

\begin{remark}[Upper-bound residual stress]\label{prop:residual-stress}
Under Theorem~\ref{thm:upper-bound-classification} and within the covered
scope, the following misreadings are excluded:
\begin{itemize}[topsep=2pt,itemsep=1pt]
  \item probability or uncertainty features of the recognized forms of
    Proposition~\ref{prop:probability-closure} reopening as unresolved
    threats;
  \item causality features of the recognized forms of
    Proposition~\ref{prop:causality-closure} reopening as unresolved threats;
  \item agency features of the recognized forms of
    Proposition~\ref{prop:agency-closure} reopening as unresolved threats;
  \item any residual feature in the covered scope of a description satisfying the
    recognition hypothesis remaining unclassified under
    Definition~\ref{def:residual-scheme};
  \item any residual feature in the covered scope of a description satisfying
    the recognition hypothesis passing the recorded seventh-role screen
    (without that hypothesis such a feature exists: the relation record \(r\) of
    \(D_{\mathrm{scr}}\), Remark~\ref{rem:recognition-open}).
\end{itemize}
Propositions~\ref{prop:probability-closure}--\ref{prop:agency-closure}
settle the three threat closures, and
Theorem~\ref{thm:upper-bound-classification}, under the recognition
hypothesis, delivers the emptiness of categories (x) and (xi).
\end{remark}

\subsection{Overclaim stress}\label{subsec:overclaim-stress}

\begin{remark}[Overclaim stress]\label{prop:overclaim-stress}
The following restatements of the scoped exact-six target are
excluded:
\begin{itemize}[topsep=2pt,itemsep=1pt]
  \item restatement as a full six-primitives algebra;
  \item restatement as a global primitive-independence theorem;
  \item restatement as unique or canonical decomposition;
  \item restatement as empirical realization;
  \item restatement as a claim about all emergence phenomena, or any
    domain broader than \FATCD{}.
\end{itemize}
Remark~\ref{rem:atlas-global-nonclaims},
Proposition~\ref{prop:nonclaim-closure}, and
Proposition~\ref{prop:non-uniqueness} carry the corresponding nonclaims
forward through the theorem chain. In particular, nonclaim closure
survives the exact-six assembly.
\end{remark}

\subsection{The integrated stress certification}\label{subsec:integrated-stress-theorem}

\begin{proposition}[Integrated stress certification]\label{thm:integrated-stress}
Clauses (a)--(d) are statements about descriptions; clause (e) records the
dependency structure of the proof of Theorem~\ref{thm:scoped-exact-six}.
The following hold:
(a) every \(D\in\FATCD{}\), including \(D_{\mathrm{norm}}\notin\FATCD_{\mathrm{RH}}\),
admits a well-formed decomposition record, and every well-formed record of such
\(D\) satisfies clauses (i)--(ii) of Theorem~\ref{thm:six-channel-exactness};
(b) for each \(k\in\{0,\dots,6\}\) some description in \(\FATCD_{\mathrm{RH}}\) has exactly
\(k\) active channels (Proposition~\ref{prop:channel-vs-activation});
(c) some description has two distinct well-formed decomposition records
(Proposition~\ref{prop:non-uniqueness});
(d) the description \(D_{\mathrm{norm}}\) of Remark~\ref{rem:recognition-open}
lies in \FATCD{} but not in \(\FATCD_{\mathrm{RH}}\); and
(e) Theorem~\ref{thm:scoped-exact-six} follows from
Theorems~\ref{thm:six-role-lower-bounds},
\ref{thm:upper-bound-classification}, \ref{thm:decomposition-existence}
and~\ref{thm:six-channel-exactness}, Propositions~\ref{prop:probability-closure}--\ref{prop:agency-closure},
and the inventory of the records of \(D_1,\ldots,D_6\) given in its proof.
\end{proposition}

\begin{proof}
We verify that each of the six stress families is discharged by the earlier
result it rests on, so that none of the listed misreadings can arise in
the combined chain.

\emph{Dependency.} The components named in
Remark~\ref{prop:dependency-stress} are present and ordered: the
\FATCD{} domain is fixed in Section~\ref{sec:fatcd-domain}, the
lower-bound theorem in Section~\ref{sec:lower-bounds}
(Theorem~\ref{thm:six-role-lower-bounds}), the upper-bound
classification in Section~\ref{sec:upper-bound-classification}
(Theorem~\ref{thm:upper-bound-classification}), and the
decomposition and exhaustion theorems in
Section~\ref{sec:decomposition-exhaustion}
(Theorems~\ref{thm:decomposition-existence}
and~\ref{thm:six-channel-exactness}). Theorem~\ref{thm:scoped-exact-six}
uses Theorems~\ref{thm:decomposition-existence}
and~\ref{thm:six-channel-exactness}; clause (iii) of
Theorem~\ref{thm:six-channel-exactness} uses
Theorem~\ref{thm:upper-bound-classification}; and the last clause of
Theorem~\ref{thm:scoped-exact-six} uses
Theorem~\ref{thm:six-role-lower-bounds}. Theorem~\ref{thm:upper-bound-classification} does not use
Theorem~\ref{thm:six-role-lower-bounds}. The existence argument of
Theorem~\ref{thm:decomposition-existence} does not use it either; that proof
cites it only for the separate non-vacuity observation at its end. So the target cannot be asserted
while skipping one of these uses, and no link consumes its own output.

\emph{Circularity.} The three circularity patterns of
Remark~\ref{prop:circularity-stress} are blocked by the
non-circularity schema of Proposition~\ref{prop:non-circularity}, which
defines \FATCD{} without reference to a six-slot count, and by the
direct proof of Theorem~\ref{thm:decomposition-existence}, which
produces a decomposition without assuming scoped exact-six. The absence of
labels (x) and (xi) is not a membership condition of \FATCD{}: it follows,
through the label criteria of Definition~\ref{def:residual-scheme}, from the
recognition hypothesis, which is assumed and fails for the description of
Remark~\ref{rem:recognition-open}.

\emph{Channel versus activation.} The conflations of
Remark~\ref{prop:channel-activation-stress} are excluded by
Proposition~\ref{prop:channel-vs-activation} together with the
five-status taxonomy of Definition~\ref{def:channel-status}: a channel
carrying status \(\belowthr\), \(\collapsedbc\), \(\absentrec\), or
\(\outsidescope\) is an inactive channel and is never read as an active
projection, so no description is required to activate all six channels.

\emph{Role alignment.} The three forbidden assignments of
Remark~\ref{prop:role-alignment-stress} are ruled out by the alignment
rules of Proposition~\ref{prop:alignment-rules}, which refuse generic
transition data as \Pthree{}, generic path or derivation data as
\Pfive{}, and generic semantic or interpretation data as \Pone{}, in
accordance with the alignments of
Remark~\ref{rem:three-repaired-alignments}.

\emph{Upper-bound residual.} The misreadings of
Remark~\ref{prop:residual-stress} are closed within the covered scope by
Propositions~\ref{prop:probability-closure}--\ref{prop:agency-closure},
which settle the probability, causality, and agency threat closures, and
by Theorem~\ref{thm:upper-bound-classification}, which --- under the
recognition hypothesis (Definition~\ref{def:recognition-hypothesis}) ---
classifies every covered residual feature and delivers the emptiness of
categories (x) and (xi); hence, under that hypothesis, no covered residual
reopens as an unresolved threat or passes the recorded seventh-role screen, while the
three principal threat classes are closed without that hypothesis, in the
recognized forms of
Propositions~\ref{prop:probability-closure}--\ref{prop:agency-closure}.

\emph{Overclaim.} The excluded restatements of
Remark~\ref{prop:overclaim-stress} are blocked because the nonclaims of
Remark~\ref{rem:atlas-global-nonclaims} are carried forward by
Proposition~\ref{prop:nonclaim-closure}, and non-uniqueness of the
decomposition is recorded by Proposition~\ref{prop:non-uniqueness}; thus
nonclaim closure survives the exact-six assembly, and no full-algebra,
global-independence, unique-decomposition, empirical-realization, or
broadened-domain reading is admitted.

Thus each misreading listed in the stress remarks is excluded by the result
cited for it; dependencies cohere, the channel-versus-activation distinction
holds and nonclaim closure holds, so the results combine into the terminal
statement of
Section~\ref{sec:scoped-exact-six}.

For the concrete clauses: (a) holds because Theorem~\ref{thm:decomposition-existence}
gives a well-formed decomposition record for every \(D\in\FATCD{}\), and its
existence construction and clauses (i)--(ii) of
Theorem~\ref{thm:six-channel-exactness} use only
Definitions~\ref{def:channel-status} and~\ref{def:decomposition-record} and the
label criteria after Definition~\ref{def:residual-scheme}, and clause (iii)
alone cites Theorem~\ref{thm:upper-bound-classification}; (b) is
Proposition~\ref{prop:channel-vs-activation}; (c) is
Proposition~\ref{prop:non-uniqueness}; (d) is shown in
Remark~\ref{rem:recognition-open}; and (e) holds because the proof of
Theorem~\ref{thm:scoped-exact-six} in Section~\ref{sec:scoped-exact-six} uses exactly
the results listed in (e) and not this proposition.
\end{proof}

In the sense of the checklist of this section, the results named in its table
rule out each listed misreading: dependency, circularity,
channel-versus-activation, role-alignment, upper-bound residual and overclaim
misreadings.
 \section{The scoped exact-six theorem}\label{sec:scoped-exact-six}



The proof assembles Theorems~\ref{thm:six-role-lower-bounds},
\ref{thm:upper-bound-classification},
\ref{thm:decomposition-existence} and~\ref{thm:six-channel-exactness}
over the descriptions in \FATCD{} that satisfy the recognition hypothesis,
together with Propositions~\ref{prop:probability-closure}--\ref{prop:agency-closure}
and an inventory of the records of the lower-bound descriptions.
The theorem statement also records the three threat closures
of Propositions~\ref{prop:probability-closure}--\ref{prop:agency-closure},
and its last clause, from Theorem~\ref{thm:six-role-lower-bounds}, says that
none of the six channels is idle on that class.

In brief: for every \(D\in\FATCD_{\mathrm{RH}}\) a well-formed \(\Dec{D}\) exists; every well-formed \(\Dec{D}\) has the channels \(P_1,\ldots,P_6\), each with the status of Definition~\ref{def:channel-status}; every label licensed for a kind residual or a bookkeeping annotation lies in (i)--(ix), and every label licensed for \(r\in\Residual{D}\) in (iii)--(ix); and \(P_i\) is active for \(D_i\) (\(i=1,\ldots,6\)).

\begin{theorem}[Scoped exact-six]\label{thm:scoped-exact-six}
Let \(\FATCD_{\mathrm{RH}}\) be the class of descriptions \(D \in \FATCD{}\)
that satisfy the recognition hypothesis
(Definition~\ref{def:recognition-hypothesis}), and let
\(D\in\FATCD_{\mathrm{RH}}\). \textup{(a)} Then \(D\) admits a well-formed
decomposition/exhaustion record (Theorem~\ref{thm:decomposition-existence}),
and every well-formed record \(\Dec{D}\) of \(D\) has exactly six typed
closure-role channels:
\begin{itemize}[topsep=2pt,itemsep=1pt]
  \item \Pone{}: update-vs-quotient compatibility / descent obstruction;
  \item \Ptwo{}: macro-specification representability (when active, witnessed by a certified absence
    of realizers in a candidate class that is the full declared domain of a
    packaging or update map, with a recorded true evaluation at a declared
    same-kind record outside that class, or a realizer set such that some class of a declared
    packaging kernel whose domain contains the candidate class meets both
    it and its complement in the candidate class);
  \item \Pthree{}: enriched route/coherence mismatch;
  \item \Pfour{}: stage/order/transition/refinement structure;
  \item \Pfive{}: quotient packaging / quotienting;
  \item \Psix{}: audit/bookkeeping/provenance/replay/checkability.
\end{itemize}

\par\noindent\textup{(b)} Each channel carries exactly one recorded status from the five statuses \activeproj{}, \collapsedbc{}, \belowthr{},
\outsidescope{} and \absentrec{}: the first of these, in the order listed, whose case in
Definition~\ref{def:channel-status} applies. Every active projection on a cluster \(C\) carries
threshold, witness-schema, collapse-case, level and nonclaim annotations
targeting exactly \(C\), and an honest audit certifying \(W_i(C)\)
(Definition~\ref{def:role-projection}).
\par\noindent\textup{(c)} Let \(\Dec{D}\) be any well-formed record of \(D\). Every label that the criteria license for a kind residual of \(D\)
(Definition~\ref{def:residual-scheme}) or a bookkeeping annotation of \(D\) relative to \(\Dec{D}\) lies among
(i)--(ix). Every element of \(\Residual{D}\) for that record is
classified, under the upper-bound classification of
Section~\ref{sec:upper-bound-classification}, as one of:
\begin{itemize}[topsep=2pt,itemsep=1pt]
  \item (iii) presentation variant;
  \item (iv) level shift;
  \item (v) unabsorbed role-aligned content, covered-seed content with no
    certified witness \(W_i\), content checked in an honestly failed offer, or activation-threshold refinement;
  \item (vi) instrument-visibility artifact;
  \item (vii) empirical-bridge annotation;
  \item (viii) outside-domain structure;
  \item (ix) already-covered bookkeeping or annotation data, or role-neutral data
    without an assertion of its own.
\end{itemize}
Under the recognition hypothesis, no kind residual and no element of \(\Residual{D}\) remains as an unresolved candidate
seventh role or passes the recorded seventh-role screen. Since seed clause (2)
of Definition~\ref{def:channel-status} accepts any one-variable macro-specification with a realization relation
whose recorded evaluations take both truth values, this clause certifies the
form in which each feature is recorded, not that its content concerns the role
whose seed or witness it carries (Remark~\ref{rem:classification-scope}).
\par\noindent\textup{(d)} Moreover, for
every \(D'\in\FATCD{}\), whether or not it satisfies the recognition
hypothesis, any kind residual of \(D'\)
whose content takes one of the recognized forms of
Propositions~\ref{prop:probability-closure}--\ref{prop:agency-closure}
is licensed, relative to every well-formed \(\Dec{D'}\), at least one of the
labels those propositions list, and neither (x) nor (xi).

\par\noindent\textup{(e)} The conclusions stated for \(D\in\FATCD_{\mathrm{RH}}\) hold for every
such \(D\); that is, the statement \(\ScopedExactSix{\FATCD_{\mathrm{RH}}}\) holds. Moreover, for each
\(i\in\{1,\dots,6\}\) the description \(D_i\) of
Theorem~\ref{thm:six-role-lower-bounds} lies in \(\FATCD_{\mathrm{RH}}\) and has
a well-formed record, and in every well-formed record of \(D_i\) the channel
\(P_i\) is \activeproj{} (Theorem~\ref{thm:six-channel-exactness}(iv)); so each of the six
channels is active somewhere in \(\FATCD_{\mathrm{RH}}\).
\end{theorem}

The recognition hypothesis restricts how assertions are presented. An
explicit finite normalization formula can fail it; replacing an assertion by a declared predicate
whose extension \(D\) does not record places the feature outside the covered scope when the resulting record meets criterion (viii), while presenting
nonnegative weights as the row of a transition kernel makes them a covered \Pfour{} seed
(Remark~\ref{rem:recognition-open}, paragraphs \emph{Scope depends on presentation} and \emph{Repair by a transition row}). Every closed check-language formula \(\varphi\) containing no atomic \(W_j\)
formula can be wrapped, provided the completed description lies in \FATCD{}, given two
distinct declared records of one kind, in the covered \Ptwo{} seed described in
Remark~\ref{rem:recognition-open} (paragraph \emph{Wrapping a closed formula as a \Ptwo{} seed}). This recognizes the
predicate and realization-relation records and records \(\varphi\)'s truth
value in the resulting description; a false formula is represented by its checked negation. If the wrapped description is again in \FATCD{}, it satisfies the recognition hypothesis exactly when all its other kind residuals, all its bookkeeping annotations and all its record residuals are recognized; for the bookkeeping annotations this holds in particular when each has only datum and reference fields besides its exempt fields (paragraph after Definition~\ref{def:recognition-hypothesis}).

The proof assembles the earlier results. Proposition~\ref{thm:integrated-stress}
checks separately that they combine without the misreadings of
Section~\ref{sec:integrated-stress}; the proof does not use it.

\begin{proof}
Fix \(D \in \FATCD{}\) satisfying the recognition hypothesis. By
Theorem~\ref{thm:decomposition-existence},
\(D\) admits a well-formed decomposition/exhaustion record. Let \(\Dec{D}\) be
any well-formed record of \(D\) in the sense of
Definition~\ref{def:decomposition-record}. By
Theorem~\ref{thm:six-channel-exactness}, the channel sequence of
\(\Dec{D}\) has length six, each channel carries one of the five
statuses of Definition~\ref{def:channel-status}, and every raw feature lies
in \(A(\Dec{D})\), \(S(\Dec{D})\) or \(\Residual{D}\); features in the last
set receive a residual classification, and \(\mathcal L\) and \(\mathcal V\)
annotate this coverage rather than forming a fourth coverage case. By
Theorem~\ref{thm:upper-bound-classification}, which applies because \(D\)
satisfies the recognition hypothesis, every label licensed for a kind residual or a bookkeeping annotation of
\(D\) lies among (i)--(ix) and every label licensed for an element of
\(\Residual{D}\) among (iii)--(ix) (both clauses of Theorem~\ref{thm:upper-bound-classification}), so
categories (x) and (xi) are empty, which is the condition of
Definition~\ref{def:scoped-exact-six-question}. Hence \(\ScopedExactSix{\FATCD_{\mathrm{RH}}}\).
The clause on active projections is the admissibility clause of
Definition~\ref{def:role-projection}, which Definition~\ref{def:decomposition-record}
requires of every active channel of a well-formed record. The clause on \(D'\) is
Propositions~\ref{prop:probability-closure}--\ref{prop:agency-closure}, which hold
for every \(D'\in\FATCD{}\) and every well-formed \(\Dec{D'}\) without the
recognition hypothesis.

For the last clause, fix \(i\) and \(D_i\). Its content records satisfy (R1)
or (R3). In \(D_1\), the update, relation and comparison records supply the
descent-obstruction witness; the four state records and the object record
\(S\) are named inputs or declared-domain records of those checks and carry no assertion field. In \(D_2\), the specification and realization
records supply the representability data, the comparison record supplies the
required comparison in its satisfying-realizer instance, and candidate and \(S_2\) records are direct checking
inputs; the packaging map
\(q\) carries a covered \Pfive{} seed with no assertion field, so it satisfies
(R1), and the object record \(\{A,B\}\), the codomain of \(q\), supplies no
witness field and has no assertion field, so it satisfies (R3). In \(D_3\), the route and
comparison records supply the generated-route and separation data; their
endpoints, and the object record \(U\) listing them, are direct inputs, and the move records \(a,b\) contribute required
fields (their graphs make \(\rho_1,\rho_2\) defined) and, having only datum
fields, satisfy (R1). The rewrite rule \(aa\to a\) is a datum field of the comparison record,
which its monoid certificate reads by quantifying over the listed rules; it
carries no assertion, and the comparison record satisfies (R1). In \(D_4\), each certified order, transition and
refinement record is a covered \Pfour{} seed, while its declared indices and
the object records listing them are inputs to its law check. In \(D_5\), the quotient relation, packaging map and
validity check supply the quotient witness; the micro-objects and quotient
object are declared inputs to those checks. In \(D_6\), the event, trace,
provenance and check records supply required witness data or covered
\Psix{} seeds. Every declared input carries no predicate, comparison,
invariant, law or claim field, so the assertion-field condition of (R1) holds
vacuously for it; it therefore satisfies (R1) if it supplies a witness field,
as the micro-object records \(x,y\) of \(D_5\) do, and (R3) otherwise. The
related-pair comparison in \(D_1\) belongs to its (R1) comparison record. The
structural and bookkeeping annotations satisfy (R2); their recorded assertion
\(W_i(C_i)\) is represented by the content records. These alternatives cover
every kind residual and every possible element of \(\Residual{D_i}\), so
\(D_i\in\FATCD_{\mathrm{RH}}\). Finally \(W_i(C_i)\) and its recorded
attachments give \(P_i\) status \activeproj{} by
Definition~\ref{def:channel-status}.
\end{proof}

For the running example, the record of Section~\ref{sec:decomposition-exhaustion}
illustrates the six-channel and status clauses of the theorem: six channels,
each with one status (three \activeproj{}, one \belowthr{}, two
\absentrec{}), with the index record classified as a presentation variant.

The following table traces the clauses of the theorem to their sources.
Residual exhaustion uses the recognition hypothesis; the integrated-stress check
carries that assumption through its upper-bound residual case.
\begin{center}
\small
\begin{tabularx}{\linewidth}{@{}>{\raggedright\arraybackslash}X>{\raggedright\arraybackslash}p{0.36\linewidth}@{}}
\toprule
Clause & Source \\
\midrule
a well-formed record \(\Dec{D}\) exists & Theorem~\ref{thm:decomposition-existence} \\
six channels, one status each, every raw feature accounted for & Definition~\ref{def:decomposition-record}, Theorem~\ref{thm:six-channel-exactness} \\
each active projection has threshold, witness, level and bookkeeping data & Definition~\ref{def:decomposition-record} \\
record residuals receive labels in (iii)--(ix), kind residuals in (i)--(ix); (x) and (xi) are empty & Theorem~\ref{thm:upper-bound-classification}, under the recognition hypothesis \\
probability, causality and agency, in their recognized forms & Propositions~\ref{prop:probability-closure}--\ref{prop:agency-closure} \\
each channel is active for some description in \(\FATCD_{\mathrm{RH}}\) & Theorem~\ref{thm:six-role-lower-bounds} \\
\bottomrule
\end{tabularx}
\end{center}

\subsection{What scoped exact-six means and does not mean}\label{subsec:exact-six-meaning}

Theorem~\ref{thm:scoped-exact-six} has a reading that should be
preserved throughout what follows.

\begin{itemize}[topsep=2pt,itemsep=1pt]
  \item \emph{Six channels, not six active witnesses.}
    Theorems~\ref{thm:decomposition-existence}
    and~\ref{thm:six-channel-exactness} give six role channels in \(\Dec{D}\)
    for every \(D \in \FATCD{}\), and the theorem restates this for
    \(D\in\FATCD_{\mathrm{RH}}\). It does not assert that every \(D\) activates
    all six roles, that every \(D\) carries six active projections
    \(\pi_1(D), \ldots, \pi_6(D)\), or that a single \(D\) realizes
    all six witness schemas of Section~\ref{sec:lower-bounds}; such a
    description exists in \(\FATCD_{\mathrm{RH}}\) (the case \(k=6\) of
    Proposition~\ref{prop:channel-vs-activation}), but it is not a clause of
    the theorem.
  \item \emph{Existential, not unique, decomposition.} The theorem
    asserts existence of \(\Dec{D}\); it does not assert uniqueness or
    canonicity. Distinct well-formed records of the same \(D\) may
    exist (Proposition~\ref{prop:non-uniqueness}).
  \item \emph{Conditional residual exhaustion.} The residual
    exhaustion clause is conditional on the recognition hypothesis
    (Definition~\ref{def:recognition-hypothesis}) and restricted to the
    covered scope: the probability, causality and agency threats are
    closed without that hypothesis, in the recognized forms of
    Propositions~\ref{prop:probability-closure}--\ref{prop:agency-closure},
    while not every finite audited feature is recognized: the
    normalization description of Remark~\ref{rem:recognition-open} has a
    residual feature labelled (x). Features whose assertions use a predicate whose extension \(D\) does not record, or quantify over a sort whose carrier \(D\) does not record, are classified as outside-domain structures under category~(viii) of Definition~\ref{def:residual-scheme}, not absorbed into primitive roles, when each such assertion is true in some expansion of \(D\) and false in another and the other assertion fields meet that criterion; an assertion of this kind with the same truth value in every expansion, such as \(Q(o)\vee\neg Q(o)\) for an unrecorded \(Q\), lies in the covered scope and is subject to the recognition hypothesis.
  \item \emph{No broader claim.} Theorem~\ref{thm:scoped-exact-six}
    is a statement about \FATCD{}; the claims it does not make are
    listed in Section~\ref{sec:limitations-nonclaims}.
\end{itemize}

 \section{Discussion}\label{sec:discussion}


Theorem~\ref{thm:scoped-exact-six} establishes
\(\ScopedExactSix{\FATCD_{\mathrm{RH}}}\); its six-channel decomposition
clauses hold for all \(D\in\FATCD{}\).
This section discusses what the statement earns and what it leaves
unstated, how the three parts of
Sections~\ref{sec:primitive-role-architecture},
\ref{sec:admissibility-meta-theory}, and
\ref{sec:fatcd-domain}--\ref{sec:integrated-stress} combine into one
theorem stack, and how the result sits next to earlier Six Birds
literature.

\subsection{The shape of the scoped claim}\label{subsec:discussion-shape}

Each of the three scope features --- domain restriction, existence
rather than uniqueness, residual exhaustion within a covered scope
--- is structural rather than cosmetic. An unrestricted exact-six
theorem would require a domain-free argument for which the present
paper does not claim to have the instruments. A uniqueness or
canonicity claim would require a theory of translation records
identifying distinct well-formed decompositions, which lies outside
the present admissibility regime. Scope-free residual exhaustion
would require removing the recognition hypothesis, which fails in
\FATCD{} (Remark~\ref{rem:recognition-open}); the three threat closures of
Section~\ref{sec:upper-bound-classification} are already stated without it
for their recognized forms. Stating
the theorem at the scope the proof supports keeps the
endogenous-instrument meta-limit (Section~\ref{subsec:meta-limit})
separate from the current claim instead of folded into it. The posture
aligns with the broader view that scientific laws form a patchwork of
domain-local regularities rather than a single global
theory~\citep{Cartwright1999}.

\subsection{The three parts as one theorem stack}\label{subsec:discussion-pillars}

The three parts introduced in Section~\ref{subsec:three-pillars} play
complementary roles. The primitive role architecture
(Section~\ref{sec:primitive-role-architecture}) fixes the vocabulary:
six role meanings, the behavioral, verification, and
structural-categorical level pattern, and the local boundary
distinctions. This vocabulary alone does not certify any claim,
since it does not specify which role talk is admissible. The
admissibility and meta-theory regime
(Section~\ref{sec:admissibility-meta-theory}) supplies that
certification through honest bookkeeping, typed non-collapse, the
forgetting and selected-lift architecture, activation thresholds,
instrument-relative visibility, and empirical-bridge admissibility.
Admissibility alone, however, yields no terminal claim; it only
regulates what may be asserted. The exact-six theorem chain
(Sections~\ref{sec:fatcd-domain}--\ref{sec:integrated-stress})
supplies the quantification: the \FATCD{} domain, the six
lower-bound witnesses, the upper-bound classification and the
decomposition/exhaustion theorem together deliver
Theorem~\ref{thm:scoped-exact-six}, and the integrated stress certification
is a separate consistency check that its proof does not use.
Vocabulary without admissibility would be undisciplined; admissibility
without a theorem chain would be inconclusive; a theorem chain without
the first two would lack a subject.

\subsection{Channel versus activation}\label{subsec:discussion-channel-vs-activation}

The channel-versus-activation discipline of
Proposition~\ref{prop:channel-vs-activation}, enforced throughout
Sections~\ref{sec:fatcd-domain}--\ref{sec:scoped-exact-six},
separates two statements that are easily conflated. The structural
statement is that every admissible description carries six typed
closure-role slots. The activation statement is that every admissible
description realizes all six slots with active projections. The
activation statement is too strong: it fails on descriptions that
carry no route grammar and therefore have no active \Pthree{}
projection, and on descriptions whose raw material carries no
\Psix{}-aligned event/provenance/replay/check witness cluster and
therefore have no active \Psix{} projection. The structural
statement, which is what we prove, is compatible with any activation
profile. Inactive channels carry one of the statuses
\(\belowthr{}\), \(\collapsedbc{}\), \(\absentrec{}\), or
\(\outsidescope{}\), describing the way a role fails to activate.
Separating the two statements is what lets the six-channel clauses
(Theorems~\ref{thm:decomposition-existence}
and~\ref{thm:six-channel-exactness}) hold of every \(D \in \FATCD{}\), and
Theorem~\ref{thm:scoped-exact-six} of every \(D\in\FATCD_{\mathrm{RH}}\),
without requiring activation of all six roles in every description; it also lets the lower-bound theorem of
Section~\ref{sec:lower-bounds} be stated existentially over the
domain rather than uniformly. An informal reading that elides the
distinction silently recovers the stronger activation claim, which we
do not prove.

\subsection{Instrument-relative methodology}\label{subsec:discussion-instrument-relative}

A methodological choice underlying
Section~\ref{sec:admissibility-meta-theory} deserves explicit
statement as a contribution. Instruments are treated as mathematical
probes rather than as neutral tools, in the spirit of ontic structural
realism~\citep{LadymanRoss2007} and the sheaf-theoretic treatment of
contextuality~\citep{AbramskyBrandenburger2011}. The technical form
is Definition~\ref{prop:visibility}; the associated nonclaim is that
no result here rests on an instrument-free truth.

The motivation is reflexive. The subject matter --- finite audited
typed closure descriptions and the six typed closure roles that
project them --- is the same descriptional machinery from which any
instrument for detecting closure structure must itself be built.
Operator-presentation, quotient/descent, coherence/route,
constructive/proof-theoretic, and information/regime instruments are
not vantage points outside the domain: each is a closure structure
probing another closure structure. A claim asserted under one such
instrument while silently suppressing what another would have
revealed would privilege that instrument without bookkeeping for the
privilege --- the overread that Definition~\ref{prop:visibility}
blocks.

The consequence is epistemic. Every admissible visibility claim carries an
explicit visibility record: its visible content, its suppressed
content, the instruments under which visibility and suppression are
asserted, and the resulting instrument-relative claim strength
(Definition~\ref{prop:visibility}). Triangulating a role projection
across several instruments is a disciplined way to strengthen an
assertion; reading a projection off a single instrument and declaring
its content canonical is not admissible. The discipline is one of
epistemic hygiene, in the spirit of Bohr's complementarity
doctrine~\citep{Bohr1958}: we do not claim that instruments are
incompatible in principle, only that they are not interchangeable
without explicit visibility bookkeeping and, where contexts shift,
explicit translation records
(Definitions~\ref{prop:visibility} and~\ref{prop:level-profile}).

This posture also explains why the endogenous-instrument meta-limit
(Sections~\ref{subsec:meta-limit} and~\ref{subsec:future-meta-limit})
is deferred to future work. Taking instrument-relativity seriously
from the outset lets the meta-limit be stated as its own question ---
whether the accepted internal instrument family can certify its own
soundness and completeness --- a question made statable, but not
settled, by the stance adopted here.

\subsection{Relation to earlier Six Birds results}\label{subsec:discussion-earlier-results}

Six Birds Foundations~I~\citep{Tsiokos2026SixBirdsFoundations}
introduced the theory-package formalism
\(\mathcal{T} = (Z, f, \Sigma_f, E, \mathcal{A})\) and derived the
\Pone{}--\Psix{} vocabulary under minimal hypotheses of composability
with limited interface access. The present paper does not re-derive
that vocabulary; it takes the six roles as given and asks a
different question: whether, within a specific admissibility regime
and over a specific class of admissible descriptions, the six roles
are exactly enough. The Foundations theorem gives the canonical
unavoidability side of the picture; the terminal theorem here gives
the scoped exhaustion side. The two statements address different
questions and address them at different levels; combining them does
not yield a stronger claim, and neither theorem follows from the
other.

The theorem chain also draws on a nearby
result. \citet{Tsiokos2026SixBirdsProtocolTrap}, complementing the
no-go analysis of \citet{Tsiokos2026SixBirdsNoGo}, established that
protocol holonomy is not a directionality certificate:
noncommutativity of update sequences can be nonzero while
path-reversal asymmetry remains zero, so route mismatch is a
diagnostic, not an arrow of time. That observation is why
Definition~\ref{def:primitive-roles} records that \Pthree{} is not a
directionality certificate, and why Remark~\ref{rem:p6-drive} keeps \Pthree{}
route/coherence mismatch distinct from the drive specialization of
\Psix{}. The next section records the
scope boundaries and nonclaims of the present theorem stack.
 \section{Limitations and nonclaims}\label{sec:limitations-nonclaims}


This section collects the scope restrictions of
Theorem~\ref{thm:scoped-exact-six} and the global nonclaims
preserved throughout.

\subsection{Scope of the terminal theorem}\label{subsec:terminal-scope}

Theorem~\ref{thm:scoped-exact-six} establishes
\(\ScopedExactSix{\FATCD_{\mathrm{RH}}}\); its six-channel decomposition
clauses hold for all \(D\in\FATCD{}\), the finite audited typed closure
description domain of Definition~\ref{def:fatcd}. The theorem asserts:
\begin{itemize}[topsep=2pt,itemsep=1pt]
  \item existence of a decomposition/exhaustion record \(\Dec{D}\)
    with six role channels and one status per channel;
  \item residual-feature exhaustion under the upper-bound classification,
    conditional on the recognition hypothesis
    (Definition~\ref{def:recognition-hypothesis}) and within the covered scope,
    so that no seventh irreducible role is required there (the three principal
    threat classes are closed without the hypothesis, in the recognized
    forms of
    Propositions~\ref{prop:probability-closure}--\ref{prop:agency-closure});
  \item activation of any \(\pi_i(D)\) only under the threshold,
    witness-schema, collapse-case, level, nonclaim and audit data of
    Definitions~\ref{def:threshold-attachment} and~\ref{def:role-projection}.
\end{itemize}

It does not assert:
\begin{itemize}[topsep=2pt,itemsep=1pt]
  \item decomposition uniqueness or canonicity;
  \item six active projections, or six active nontrivial witnesses, in
    every description;
  \item scope-free residual exhaustion, or any claim about features
    outside the covered scope;
  \item that every finite audited feature of a description in \FATCD{}
    satisfies (R1)--(R5); Remark~\ref{rem:recognition-open} exhibits a
    feature with label (x) and a feature with label (xi);
  \item unrestricted exact-six, or coverage of all emergence
    phenomena;
  \item that broader domains require no further work.
\end{itemize}

\subsection{Global nonclaims}\label{subsec:global-nonclaims}

Each nonclaim below is preserved throughout the paper and is enforced
locally by the theorem-chain result that would otherwise overread it.

\begin{itemize}[topsep=2pt,itemsep=1pt]
  \item \emph{No full six-primitives algebra.} Scoped exact-six is
    proved over \FATCD{}; it is not claimed that \Pone{}--\Psix{}
    generate a global algebra of closure operations
    (Remark~\ref{rem:atlas-global-nonclaims},
    Remark~\ref{prop:overclaim-stress}).
  \item \emph{No global primitive-independence theorem.} Typed
    non-collapse (Proposition~\ref{prop:typed-non-collapse}) is the
    only claim made at that level; no global independence notion is defined
    or claimed.
  \item \emph{No empirical realization.} Empirical-bridge admissibility
    (Definition~\ref{prop:empirical-bridge}) admits
    threshold-dependent, level-indexed, instrument-visible substrate
    evidence; it does not assert that any \(D \in \FATCD{}\) is itself
    a substrate observation.
  \item \emph{No unsupported unique selected lift.} The displayed atlas lifts
    are selected and, except where uniqueness is proved and recorded (as for
    a descended map), non-unique; a local uniqueness assertion is admissible
    only with a recorded proof (Definition~\ref{prop:forgetting-lifts}).
  \item \emph{No unsupported lower-to-higher determinacy.} Behavioral or
    verification content is not claimed to determine structural-categorical
    content without a recorded proof, as for a descended map.
  \item \emph{No instrument-free formulation.} Admissible visibility claims carry
    an instrument-relative visibility record
    (Definition~\ref{prop:visibility}).
  \item \emph{No internal instrument-family completeness.} The
    accepted instrument family is not asserted to be sufficient for
    every exact-six judgment formulable in the same regime.
  \item \emph{No unrestricted all-emergence coverage.}
    Theorem~\ref{thm:scoped-exact-six} is scoped to \FATCD{}.
  \item \emph{No uniqueness of primitive presentations.} Distinct
    admissible presentations of the same role may coexist.
  \item \emph{Cross-level comparison requires translation.} Direct
    comparison of claims at different levels is inadmissible without a
    translation record (Definition~\ref{prop:level-profile}).
\end{itemize}

\subsection{The endogenous-instrument meta-limit}\label{subsec:meta-limit}

An adjacent meta-theoretic question, deferred to future work, is
whether the accepted internal instrument family can itself certify
both its soundness and its completeness for exact-six judgments
formulable inside the same regime. The notions of an internal argument and of
the soundness and completeness of an instrument family are not defined in this
paper; what follows is a heuristic target, not a mathematical statement. A
plausible sharp formulation is
that no internal argument should be expected to accomplish all three
of the following at once:
\begin{enumerate}[topsep=2pt,itemsep=1pt]
  \item establish the scoped exact-six claim;
  \item certify the soundness of the accepted instrument family;
  \item certify the completeness of that family for exact-six
    judgments expressible inside the same regime.
\end{enumerate}
Such self-certification limits are of a piece with the classical
G{\"o}del--L{\"o}b tradition on internal consistency and provability
statements~\citep{Godel1931,Lob1955}. This question is not part of
the proof of Theorem~\ref{thm:scoped-exact-six} and does not weaken
or retract the theorem's scoped conclusion. It is taken up as a
follow-on direction in Section~\ref{sec:future-work}.
 \section{Future work}\label{sec:future-work}


Theorem~\ref{thm:scoped-exact-six} and the limitations of
Section~\ref{sec:limitations-nonclaims} suggest several follow-on
research directions. Each subsection below states a direction, not a
commitment.

\subsection{Endogenous-instrument meta-limit}\label{subsec:future-meta-limit}

A formalization of the meta-limit described in
Section~\ref{subsec:meta-limit} would first define internal arguments and the soundness and completeness of
an instrument family, and then target a precise form of the heuristic of
Section~\ref{subsec:meta-limit}: that no internal argument can simultaneously
establish the scoped exact-six claim, certify the soundness of the
accepted instrument family, and certify its completeness for
exact-six judgments formulable inside the same regime.
Theorem~\ref{thm:scoped-exact-six} would remain intact under such a
result; the additional content would be a precise account of
self-certification boundaries. A closely related open question
concerns recognition. Not every finite audited feature of a description in
\FATCD{} is recognized: the bare normalization predicate of
Remark~\ref{rem:recognition-open} receives the unresolved label (x). What remains open is
which classes of assertion-bearing records --- normalization constraints not
presented as the law of a transition record (for instance for signed weights),
finite conservation laws, emergent composite obstructions --- admit a reduction
in which the assertion is itself a required field of an active cluster
(label (i)) without adjoining wrapper records such as a realization relation
and a one-class packaging map (Remark~\ref{rem:recognition-open}), or a trace
reference, and whether any of
them carries the recorded witness of category (xi) of
Definition~\ref{def:residual-scheme}, which is a recorded screen rather than
the identification of a new closure role.

\subsection{Broader-domain exact-six}\label{subsec:future-broader-domain}

Candidate extensions beyond \FATCD{} include admissibility regimes
that weaken finiteness, closure, or audit
(Section~\ref{sec:fatcd-domain}), and regimes that admit annotation
kinds beyond the six \(\mathrm{Ann}_*\) categories of
Definition~\ref{def:annotations}. Any such extension would require
fresh upper-bound and decomposition/exhaustion analysis, since the
alignment rules and threat closures are not stated at scope-free
strength.

\subsection{Uniqueness or canonical decomposition}\label{subsec:future-uniqueness}

Theorem~\ref{thm:scoped-exact-six} is existential. A uniqueness study
would investigate when distinct well-formed decomposition/exhaustion
records of the same \(D \in \FATCD{}\) are equivalent under explicit
translation records, or when a canonical representative can be
selected. The channel-versus-activation discipline
(Proposition~\ref{prop:channel-vs-activation}) is compatible with
such a study; no uniqueness claim is part of the present theorem.

\subsection{Full six-primitives algebra}\label{subsec:future-algebra}

The generator-and-relation structure of \Pone{}--\Psix{} beyond
typed non-collapse remains open. A full algebra would require
scope-free independence data, not just the local boundary
distinctions of Table~\ref{tab:boundary-matrix}, and would likely
require a categorical apparatus richer than the raw grammar of
Section~\ref{subsec:raw-grammar}.

\subsection{Empirical substrate mapping}\label{subsec:future-substrate}

The empirical-bridge admissibility regime
(Definition~\ref{prop:empirical-bridge}) can be connected to
concrete substrate-level observables. The present paper treats
empirical bridges as threshold-dependent, level-indexed, and
instrument-visible annotations; later work can couple this to
substrate-specific measurement theory while preserving the nonclaim
of empirical realization.

\subsection{Mechanization}\label{subsec:future-mechanization}

A Lean~4 development accompanies the present paper;
Appendix~\ref{app:mechanization} records, claim by claim, how far each definition
and theorem is currently mechanized. Strengthening that mechanization toward full
proofs --- in particular a Lean account of the boundary distinctions, the
decomposition-existence construction, and the recognition hypothesis underlying
the upper-bound classification --- is a natural direction, building on the
closure-related infrastructure already mechanized in the Foundations
program~\citep{Tsiokos2026SixBirdsFoundations,DeMouraUllrich2021}.

\subsection{Cross-instrument comparison}\label{subsec:future-cross-instrument}

Decomposition/exhaustion records of the same \(D \in \FATCD{}\) can
be compared under different admissible instrument assignments.
Definition~\ref{prop:visibility} admits visibility records that
suppress different portions of a claim under different instruments; a
cross-instrument study could make the instrument-dependence of
\(\Dec{D}\) explicit and sharpen the notion of instrument-relative
claim strength.
 
\appendix
\section{Provenance and evidence hygiene}\label{app:provenance}


The development of the theorem chain of
Sections~\ref{sec:fatcd-domain}--\ref{sec:integrated-stress} was assisted by
automated drafting and proof-search tooling, the Aristotle
system~\citep{Achim2025Aristotle}. We record this for transparency and to fix the
division of evidential labor precisely. Such tooling produces \emph{candidate}
definitions, statements, and arguments; it carries no evidential weight of its
own, and nothing in this paper rests on it.

The mathematical content is established solely by the proofs given in the body.
The machine-\emph{checked} component of the work --- which declarations are
realized in Lean, and at what fidelity --- is documented separately in
Appendix~\ref{app:mechanization}; the automated tooling named above is distinct
from that mechanization and assisted drafting only. No raw output of that tooling
is offered as mathematical evidence anywhere in the manuscript: where a section
draws on machine assistance, the load-bearing content is the section's own proof,
and the recorded mechanization is exactly as described in
Appendix~\ref{app:mechanization}.

\subsection{Code availability}\label{app:code-availability}

The \LaTeX{} source of this manuscript and the accompanying Lean mechanization
artifacts are available at
\url{https://github.com/ioannist/six-birds-foundations-ii}.
 \section{Mechanization: declarations, coverage, and trust base}\label{app:mechanization}

The theorem chain of this paper is accompanied by a Lean~4 development, the
integrated \texttt{SixBirds} project, available at the code repository of
Appendix~\ref{app:code-availability}. This appendix is the canonical disclosure
venue for that development. For every numbered definition and claim it records
the Lean declaration that realizes it and the \emph{fidelity} of that
realization, and Subsection~\ref{app:result-notes} gives a note on each result.
The mapping is maintained in the repository and reproduced here.

\subsection{Fidelity of the correspondence}\label{app:fidelity-legend}

Each theorem and proposition has a note in Subsection~\ref{app:result-notes}
naming its Lean declaration and its fidelity. We distinguish six fidelities:
\begin{itemize}[topsep=2pt,itemsep=1pt]
  \item \emph{verified} --- a faithful Lean derivation of the statement from
    non-trivial hypotheses;
  \item \emph{narrowed / parametric} --- a genuine Lean proof, but over a single
    concrete witness, a parametric record, or a narrowed input space rather than
    the full generality of the paper statement;
  \item \emph{schema (projection)} --- the Lean declaration reuses
    supplied hypotheses or encoded definitions, or constructs, projects or
    assembles records and their proofs, without formalizing the paper's
    full mathematical content; it is therefore \emph{not} described as a
    verified proof of the corresponding paper statement;
  \item \emph{trivial} --- the Lean statement reduces to \texttt{True} by
    construction, or ranges over a hand-built total classifier; it records a
    bookkeeping fact, not a derivation;
  \item \emph{typed mirror} --- for definitions, a Lean structure or inductive
    mirroring the paper definition;
  \item \emph{none} --- the declaration does not establish the statement for
    any instance.
\end{itemize}
We state plainly that \textbf{no claim in the present development is a faithful
derivation in the first sense}. The Lean track is a typed harness: it fixes the
objects as typed records, exposes the channel, status, and residual machinery,
and checks that malformed and misaligned cases are rejected (the semantic
regression tests of the \texttt{SixBirds.SemanticTests} module); but the
load-bearing classification and closure statements are projections, schemas, or
trivially-true bookkeeping over those records. The notes are worded
accordingly. A reader who wants the mathematics should read the body proofs; the
Lean track certifies typed well-formedness and case rejection, not the theorems.

\subsection{The recognition hypothesis and the upper-bound classifier}\label{app:recognition}

The honest status of the central no-seventh-role result
(Theorem~\ref{thm:upper-bound-classification}) is legible in the Lean. The
classifier \texttt{SixBirds.classifyResidual} is a total function on the finite
raw-record kinds, every branch of which lands in a covered class; the
declarations \texttt{upper\_bound\_classification} and the three closure lemmas
are \texttt{True} by construction over it. This is exactly why the body states
the result conditionally, under the recognition hypothesis
(Definition~\ref{def:recognition-hypothesis}): the Lean classifier imposes
covered outputs by construction; it neither represents nor proves the
paper-side recognition hypothesis or the quantification over every licensed
label. The
paper-side recognition hypothesis has no independent Lean correspondent, and the
failure of recognition-completeness for \FATCD{}
(Remark~\ref{rem:recognition-open}) lies precisely in the gap between the classifier's by-construction coverage and a derived
exhaustiveness.

\subsection{Declaration and coverage table}\label{app:table}

The table maps each result to its Lean declaration in the \texttt{SixBirds}
namespace and its audited coverage. All declarations are \texttt{sorry}-free and
axiom-clean relative to the trust base of \S\ref{app:trust-base}; the coverage
column records fidelity in the sense of \S\ref{app:fidelity-legend}, not mere
presence.

{\footnotesize
\begin{longtable}{@{}c >{\raggedright\arraybackslash}p{0.27\linewidth} >{\raggedright\arraybackslash}p{0.34\linewidth} >{\raggedright\arraybackslash}p{0.20\linewidth}@{}}
\toprule
\S & result & Lean declaration & coverage \\
\midrule
\endfirsthead
\toprule
\S & result & Lean declaration & coverage \\
\midrule
\endhead
\bottomrule
\endfoot
§2 & Definition~\ref{def:primitive-roles} & \texttt{SixBirds.\allowbreak Role} & typed mirror (def.) \\
§2 & Definition~\ref{def:level-trichotomy} & \texttt{SixBirds.\allowbreak Level} & typed mirror (def.) \\
§2 & Definition~\ref{def:selected-lift-atlas} & \texttt{SixBirds.\allowbreak SelectedLiftAtlas} & typed mirror (def.) \\
§2 & Proposition~\ref{prop:boundary-distinctions} & \texttt{SixBirds.\allowbreak boundary\_\allowbreak distinctions} & schema (projection) \\
§3 & Definition~\ref{prop:honest-bookkeeping} & \texttt{SixBirds.\allowbreak honest\_\allowbreak bookkeeping\_\allowbreak admissibility} & schema (projection) \\
§3 & Proposition~\ref{prop:typed-non-collapse} & \texttt{SixBirds.\allowbreak typed\_\allowbreak non\_\allowbreak collapse} & schema (projection) \\
§3 & Definition~\ref{prop:level-profile} & \texttt{SixBirds.\allowbreak level\_\allowbreak profile} & trivial \\
§3 & Definition~\ref{prop:forgetting-lifts} & \texttt{SixBirds.\allowbreak forgetting\_\allowbreak lifts} & schema (projection) \\
§3 & Remark~\ref{prop:activation-thresholds} & \texttt{SixBirds.\allowbreak activation\_\allowbreak thresholds} & schema (projection) \\
§3 & Definition~\ref{prop:visibility} & \texttt{SixBirds.\allowbreak visibility} & schema (projection) \\
§3 & Definition~\ref{prop:empirical-bridge} & \texttt{SixBirds.\allowbreak empirical\_\allowbreak bridge} & schema (projection) \\
§3 & Proposition~\ref{prop:nonclaim-closure} & \texttt{SixBirds.\allowbreak nonclaim\_\allowbreak closure} & schema (projection) \\
§4 & Definition~\ref{def:raw-record} & \texttt{SixBirds.\allowbreak RawRecord} & typed mirror (def.) \\
§4 & Definition~\ref{def:annotations} & \texttt{SixBirds.\allowbreak AnnotationKind} & typed mirror (def.) \\
§4 & Definition~\ref{def:finite} & \texttt{SixBirds.\allowbreak FiniteDescription} & typed mirror (def.) \\
§4 & Definition~\ref{def:closed} & \texttt{SixBirds.\allowbreak ClosedDescription} & typed mirror (def.) \\
§4 & Definition~\ref{def:audited} & \texttt{SixBirds.\allowbreak AuditedDescription} & typed mirror (def.) \\
§4 & Definition~\ref{def:fatcd} & \texttt{SixBirds.\allowbreak FATCD} & typed mirror (def.) \\
§4 & Definition~\ref{def:threshold-attachment} & \texttt{SixBirds.\allowbreak ThresholdAttachment} & typed mirror (def.) \\
§4 & Definition~\ref{def:role-projection} & \texttt{SixBirds.\allowbreak DerivedRoleProjection} & typed mirror (def.) \\
§4 & Definition~\ref{def:residual-scheme} & \texttt{SixBirds.\allowbreak ResidualClassificationScheme} & typed mirror (def.) \\
§4 & Definition~\ref{def:scoped-exact-six-question} & \texttt{SixBirds.\allowbreak ScopedExactSixQuestion} & typed mirror (def.) \\
§4 & Proposition~\ref{prop:non-circularity} & \texttt{SixBirds.\allowbreak non\_\allowbreak circularity} & schema (projection) \\
§4 & Proposition~\ref{prop:non-overbreadth} & \texttt{SixBirds.\allowbreak non\_\allowbreak overbreadth} & schema (projection) \\
§5 & Definition~\ref{def:common-witness-data} & \texttt{SixBirds.\allowbreak CommonWitnessData} & typed mirror (def.) \\
§5 & Theorem~\ref{thm:six-role-lower-bounds} & \texttt{SixBirds.\allowbreak six\_\allowbreak role\_\allowbreak lower\_\allowbreak bounds} & schema (projection) \\
§5 & Proposition~\ref{prop:p1-witness} & \texttt{SixBirds.\allowbreak p1\_\allowbreak witness} & narrowed / parametric \\
§5 & Proposition~\ref{prop:p2-witness} & \texttt{SixBirds.\allowbreak p2\_\allowbreak witness} & narrowed / parametric \\
§5 & Proposition~\ref{prop:p3-witness} & \texttt{SixBirds.\allowbreak p3\_\allowbreak witness} & narrowed / parametric \\
§5 & Proposition~\ref{prop:p4-witness} & \texttt{SixBirds.\allowbreak p4\_\allowbreak witness} & narrowed / parametric \\
§5 & Proposition~\ref{prop:p5-witness} & \texttt{SixBirds.\allowbreak p5\_\allowbreak witness} & narrowed / parametric \\
§5 & Proposition~\ref{prop:p6-witness} & \texttt{SixBirds.\allowbreak p6\_\allowbreak witness} & narrowed / parametric \\
§6 & Proposition~\ref{prop:alignment-rules} & \texttt{SixBirds.\allowbreak alignment\_\allowbreak rules} & schema (projection) \\
§6 & Proposition~\ref{prop:probability-closure} & \texttt{SixBirds.\allowbreak probability\_\allowbreak closure} & trivial \\
§6 & Proposition~\ref{prop:causality-closure} & \texttt{SixBirds.\allowbreak causality\_\allowbreak closure} & trivial \\
§6 & Proposition~\ref{prop:agency-closure} & \texttt{SixBirds.\allowbreak agency\_\allowbreak closure} & trivial \\
§6 & Theorem~\ref{thm:upper-bound-classification} & \texttt{SixBirds.\allowbreak upper\_\allowbreak bound\_\allowbreak classification} & trivial \\
§7 & Definition~\ref{def:channel-status} & \texttt{SixBirds.\allowbreak ChannelStatus} & typed mirror (def.) \\
§7 & Definition~\ref{def:decomposition-record} & \texttt{SixBirds.\allowbreak DecompositionRecord} & typed mirror (def.) \\
§7 & Theorem~\ref{thm:decomposition-existence} & \texttt{SixBirds.\allowbreak decomposition\_\allowbreak existence} & schema (projection) \\
§7 & Theorem~\ref{thm:six-channel-exactness} & \texttt{SixBirds.\allowbreak six\_\allowbreak channel\_\allowbreak exactness} & schema (projection) \\
§7 & Proposition~\ref{prop:channel-vs-activation} & \texttt{SixBirds.\allowbreak channel\_\allowbreak vs\_\allowbreak activation} & schema (projection) \\
§7 & Proposition~\ref{prop:non-uniqueness} & \texttt{SixBirds.\allowbreak non\_\allowbreak uniqueness} & none \\
§8 & Remark~\ref{prop:dependency-stress} & \texttt{SixBirds.\allowbreak dependency\_\allowbreak stress} & schema (projection) \\
§8 & Remark~\ref{prop:circularity-stress} & \texttt{SixBirds.\allowbreak circularity\_\allowbreak stress} & schema (projection) \\
§8 & Remark~\ref{prop:channel-activation-stress} & \texttt{SixBirds.\allowbreak channel\_\allowbreak activation\_\allowbreak stress} & narrowed / parametric \\
§8 & Remark~\ref{prop:role-alignment-stress} & \texttt{SixBirds.\allowbreak role\_\allowbreak alignment\_\allowbreak stress} & narrowed / parametric \\
§8 & Remark~\ref{prop:residual-stress} & \texttt{SixBirds.\allowbreak residual\_\allowbreak stress} & trivial \\
§8 & Remark~\ref{prop:overclaim-stress} & \texttt{SixBirds.\allowbreak overclaim\_\allowbreak stress} & schema (projection) \\
§8 & Proposition~\ref{thm:integrated-stress} & \texttt{SixBirds.\allowbreak integrated\_\allowbreak stress} & schema (projection) \\
§9 & Theorem~\ref{thm:scoped-exact-six} & \texttt{SixBirds.\allowbreak scoped\_\allowbreak exact\_\allowbreak six} & schema (projection) \\
\bottomrule
\end{longtable}
}

\subsection{Trust base and axiom audit}\label{app:trust-base}

The dependency audit accepts only the standard Lean and Mathlib trust base:
\texttt{propext}, \texttt{Classical.choice} and \texttt{Quot.sound}
(function extensionality is a theorem derived from \texttt{Quot.sound}). No project-local axiom is
admitted for any declaration recorded as a theorem. The validation harness
(\texttt{scripts/validate\_lean.sh}) rejects project-local axioms and
\texttt{sorry}, and runs a semantic audit that rejects \texttt{:= True} and
\texttt{:= False} definition stubs and all-purpose residual classifiers. The
trust-base check certifies that the typed harness is axiom-clean relative to that
base; it does not, and is not intended to, upgrade a schema or a trivially-true
statement to a derivation. The fidelity column of \S\ref{app:table} is the
governing record of what each declaration establishes.

\subsection{Notes on individual results}\label{app:result-notes}

For each result, the note below records the Lean declaration that corresponds to it and the extent of that correspondence.

\begin{description}[style=unboxed,leftmargin=0pt,itemsep=3pt]
  \item[Proposition~\ref{prop:boundary-distinctions}.] The Lean declaration \texttt{SixBirds.boundary\_distinctions} has fidelity schema (projection) of the legend above: its encoded hypotheses, definitions or record data supply its conclusion; it does not derive the corresponding paper statement.
  \item[Proposition~\ref{prop:typed-non-collapse}.] The Lean declaration \texttt{SixBirds.typed\_non\_collapse} has fidelity schema (projection) of the legend above: its encoded hypotheses, definitions or record data supply its conclusion; it does not derive the corresponding paper statement.
  \item[Proposition~\ref{prop:nonclaim-closure}.] The Lean declaration \texttt{SixBirds.nonclaim\_closure} has fidelity schema (projection) of the legend above: its encoded hypotheses, definitions or record data supply its conclusion; it does not derive the corresponding paper statement.
  \item[Proposition~\ref{prop:non-circularity}.] The Lean declaration \texttt{SixBirds.non\_circularity} has fidelity schema (projection) of the legend above: its encoded hypotheses, definitions or record data supply its conclusion; it does not derive the corresponding paper statement.
  \item[Proposition~\ref{prop:non-overbreadth}.] The Lean declaration \texttt{SixBirds.non\_overbreadth} has fidelity schema (projection) of the legend above: its encoded hypotheses, definitions or record data supply its conclusion; it does not derive the corresponding paper statement. The Lean declaration is identical in statement and proof to \texttt{SixBirds.non\_circularity} and re-extracts the admissibility fields of the record rather than establishing the structural-restriction content of the proposition.
  \item[Theorem~\ref{thm:six-role-lower-bounds}.] The Lean declaration \texttt{SixBirds.six\_role\_lower\_bounds} has fidelity schema (projection) of the legend above: its encoded hypotheses, definitions or record data supply its conclusion; it does not derive the corresponding paper statement.
  \item[Proposition~\ref{prop:p1-witness}.] The Lean development establishes this for one concrete description \texttt{witnessDescription}\allowbreak\ \texttt{Role.P1} under parametric hypotheses (alignment, threshold attachment, and finite/\allowbreak closed/\allowbreak audited verifications of that fixed witness); the narrowing is recorded in the table of this appendix (\texttt{SixBirds.p1\_witness}).
  \item[Proposition~\ref{prop:p2-witness}.] The Lean development establishes this for one concrete description \texttt{witnessDescription}\allowbreak\ \texttt{Role.P2} under parametric hypotheses (alignment, threshold attachment, and finite/\allowbreak closed/\allowbreak audited verifications of that fixed witness); the narrowing is recorded in the table of this appendix (\texttt{SixBirds.p2\_witness}).
  \item[Proposition~\ref{prop:p3-witness}.] The Lean development establishes this for one concrete description \texttt{witnessDescription}\allowbreak\ \texttt{Role.P3} under parametric hypotheses (alignment, threshold attachment, and finite/\allowbreak closed/\allowbreak audited verifications of that fixed witness); the narrowing is recorded in the table of this appendix (\texttt{SixBirds.p3\_witness}).
  \item[Proposition~\ref{prop:p4-witness}.] The Lean development establishes this for one concrete description \texttt{witnessDescription}\allowbreak\ \texttt{Role.P4} under parametric hypotheses (alignment, threshold attachment, and finite/\allowbreak closed/\allowbreak audited verifications of that fixed witness; the attached level is fixed to the verification level rather than the structural-categorical level of the prose); the narrowing is recorded in the table of this appendix (\texttt{SixBirds.p4\_witness}).
  \item[Proposition~\ref{prop:p5-witness}.] The Lean development establishes this for one concrete description \texttt{witnessDescription}\allowbreak\ \texttt{Role.P5} under parametric hypotheses (alignment over the same raw-kind cluster used for P1, threshold attachment, and finite/\allowbreak closed/\allowbreak audited verifications of that fixed witness); the narrowing is recorded in the table of this appendix (\texttt{SixBirds.p5\_witness}).
  \item[Proposition~\ref{prop:p6-witness}.] The Lean development establishes this for one concrete description \texttt{witnessDescription}\allowbreak\ \texttt{Role.P6} under parametric hypotheses (alignment over event/\allowbreak provenance/\allowbreak replay/\allowbreak check records, threshold attachment, and finite/\allowbreak closed/\allowbreak audited verifications of that fixed witness); the narrowing is recorded in the table of this appendix (\texttt{SixBirds.p6\_witness}).
  \item[Theorem~\ref{thm:upper-bound-classification}.] The Lean declaration is a check of the classifier: every \texttt{RawKind} branch of \texttt{classifyResidual} lands in a class whose \texttt{CoveredResidualClass} branch is \texttt{True}, and the uncovered constructors \texttt{unresolvedThreat}/\allowbreak\texttt{genuineSeventhRole} are never produced; ``no seventh role'' holds by construction rather than by derivation from feature content. (\texttt{SixBirds.upper\_bound\_classification}).
  \item[Proposition~\ref{prop:alignment-rules}.] The Lean declaration \texttt{SixBirds.alignment\_rules} has fidelity schema (projection) of the legend above: its encoded hypotheses, definitions or record data supply its conclusion; it does not derive the corresponding paper statement.
  \item[Proposition~\ref{prop:probability-closure}.] The Lean declaration is a single-record check: \texttt{classifyResidual} maps \texttt{RawKind.probability} to \texttt{composite}, whose \texttt{CoveredResidualClass} branch is \texttt{True}; it does not establish the paper's universal probability/uncertainty closure. (\texttt{SixBirds.probability\_closure}).
  \item[Proposition~\ref{prop:causality-closure}.] The Lean declaration is a single-record check: \texttt{classifyResidual} maps \texttt{RawKind.causality} to \texttt{composite}, whose \texttt{CoveredResidualClass} branch is \texttt{True}; the \Pfive{}-inadmissibility content of the proposition is absent and no universal residual-causality closure is established. (\texttt{SixBirds.causality\_closure}).
  \item[Proposition~\ref{prop:agency-closure}.] The Lean declaration is a single-record check: \texttt{classifyResidual} maps \texttt{RawKind.agency} to \texttt{empiricalBridgeAnnotation}, whose \texttt{CoveredResidualClass} branch is \texttt{True}; the \Pone{}-inadmissibility content is absent and no universal residual-agency closure is established. (\texttt{SixBirds.agency\_closure}).
  \item[Theorem~\ref{thm:decomposition-existence}.] The Lean declaration \texttt{SixBirds.decomposition\_existence} has fidelity schema (projection) of the legend above: its encoded hypotheses, definitions or record data supply its conclusion; it does not derive the corresponding paper statement.
  \item[Theorem~\ref{thm:six-channel-exactness}.] The Lean declaration \texttt{SixBirds.six\_channel\_exactness} has fidelity schema (projection) of the legend above: its encoded hypotheses, definitions or record data supply its conclusion; it does not derive the corresponding paper statement.
  \item[Proposition~\ref{prop:channel-vs-activation}.] The Lean declaration \texttt{SixBirds.channel\_vs\_activation} has fidelity schema (projection) of the legend above: its encoded hypotheses, definitions or record data supply its conclusion; it does not derive the corresponding paper statement.
  \item[Proposition~\ref{prop:non-uniqueness}.] The Lean development exhibits two concrete decomposition records of a single description that differ only in their \texttt{P1} channel status; since the status of a channel is fixed by the description, at most one of these records is well-formed, so the Lean declaration does not witness the proposition, and the further sources of non-uniqueness in the prose (level assignments, choice of witness cluster or representation of the witness predicate, residual labels, loss/lift placement) are not mechanized. The Lean declaration is \texttt{SixBirds.non\_uniqueness}.
  \item[Proposition~\ref{thm:integrated-stress}.] The Lean declaration \texttt{SixBirds.integrated\_stress} has fidelity schema (projection) of the legend above, bundling the six stress conjuncts: its encoded hypotheses, definitions or record data supply its conclusion; it does not derive the corresponding paper statement.
  \item[Theorem~\ref{thm:scoped-exact-six}.] The Lean declaration \texttt{SixBirds.scoped\_exact\_six} has fidelity schema (projection) of the legend above: its encoded hypotheses, definitions or record data supply its conclusion; it does not derive the corresponding paper statement.
\end{description}

\paragraph{The running example.} The module \texttt{SixBirds.Examples.RunningExample}
constructs a typed proxy for \(D_{\mathrm{ex}}\) as a finite, closed and audited
\FATCD{}, proves the split pair, the replay and check of the trace and the
equivalence of the lens classes by computation, constructs the \Pone{}, \Pfive{}
and \Psix{} projections, and gives a well-formed decomposition record with the
statuses stated in Section~\ref{sec:decomposition-exhaustion}, classifying the
index record as a presentation variant. For each active status the record
checks that the raw record kinds its role requires are present in the
description; the three projections are constructed separately and are not
linked to the status field. The absent statuses of \Ptwo{} and \Pthree{} and the
below-threshold status of \Pfour{} are assigned, not derived from the content,
and the record does not track a residual set. The typed layer has no
annotation kind for transitions; the transition is a separate raw record
referring to the map record of \(F\), and no threshold annotation targets it. The paper's expanded \(D_{\mathrm{ex}}\) explicitly includes a \(q\)-map, a
quotient-validity check and claim and audit contents absent from this Lean
raw-record list. Lean checks kind presence and assigned statuses; it does not
verify the honesty predicate of Definition~\ref{def:audited}, the paper's \(W_5\) check, or the
singleton residual calculation.
 \section{Version history}\label{app:version-history}

Version~1 appeared on Zenodo on 20 April 2026, Version~2 on Qeios on 9 June
2026, and Version~3 on 2 October 2026. Each subsection lists the differences
from the preceding version.

\subsection{Version 3}

Version~3 keeps every numbered result of Version~2 with its number, label and
type, in the same order. The titles of Remarks~\ref{rem:repaired-alignments}
and~\ref{rem:three-repaired-alignments} (``corrected'' became ``excluded''
alignments) and of Remark~\ref{rem:recognition-open} changed.

\paragraph{Corrected statements.} Definition~\ref{def:level-trichotomy}
qualifies lossy forgetting, and Definition~\ref{def:selected-lift-atlas}
qualifies non-uniqueness and lossiness.
Definition~\ref{prop:forgetting-lifts} makes the loss-record and
uniqueness-proof requirements explicit; Version~2 already admitted
lossless forgetting and unique lifts when proved and recorded.
Proposition~\ref{prop:boundary-distinctions} is stated as the failure of
\(W_i\) for a cluster built from the bare boundary datum.
Proposition~\ref{prop:non-overbreadth} no longer claims that unaudited
material cannot enter \FATCD{}, which fails for unclaimed records.
Propositions~\ref{prop:probability-closure}--\ref{prop:agency-closure}
concluded that a feature of a recognized form classifies as a \(P_i\)
projection, which needs attachment data; they now license label (i) or (v),
or another listed label, for every well-formed decomposition record.
Definition~\ref{def:channel-status} replaces five conditions that were not
shown to be exclusive and exhaustive by an ordered rule.
Remark~\ref{rem:recognition-open} left open whether the recognition
hypothesis holds on \FATCD{}; Version~3 proves that it does not
(\(D_{\mathrm{norm}}\in\FATCD\setminus\FATCD_{\mathrm{RH}}\)), and
Theorem~\ref{thm:scoped-exact-six}, whose conclusion Version~2 named
\(\operatorname{ScopedExactSix}(\FATCD)\), now concludes
\(\operatorname{ScopedExactSix}(\FATCD_{\mathrm{RH}})\);
Definition~\ref{def:scoped-exact-six-question} records that the unrestricted
statement fails. Proposition~\ref{thm:integrated-stress} replaces the
readiness claim of Version~2 by clauses (a)--(e). The scope statements of
Remarks~\ref{rem:typed-non-collapse-vs-independence},
\ref{prop:activation-thresholds}, \ref{rem:classification-scope},
\ref{prop:dependency-stress} and~\ref{prop:residual-stress} were corrected
in the same way.

\paragraph{Added hypotheses and strengthened results.}
Definition~\ref{def:fatcd} requires realization relations to list exactly the
satisfying candidates, and Definition~\ref{def:threshold-attachment} requires
nonclaim, collapse-case and witness-schema annotations targeting the cluster.
Theorem~\ref{thm:upper-bound-classification} holds for every well-formed
decomposition record and adds a pointwise form;
Theorem~\ref{thm:six-channel-exactness} adds clause (iv), independence of the
status sequence from the decomposition record;
Proposition~\ref{prop:channel-vs-activation} realizes every active-channel
count from \(0\) to \(6\) in \(\FATCD_{\mathrm{RH}}\), and all five
statuses in one description; Proposition~\ref{prop:non-uniqueness}
exhibits distinct well-formed records in \(\FATCD_{\mathrm{RH}}\); Proposition~\ref{prop:typed-non-collapse}
covers every \(W_k\); and Theorem~\ref{thm:scoped-exact-six} adds clauses (d)
and (e).

\paragraph{Definitions and scope.} Version~3 specifies the witness
predicates through the \emph{Witness types} table, audit honesty through
Definition~\ref{def:audited}, and licensed labels through
Definition~\ref{def:residual-scheme}. It distinguishes kind residuals
from record residuals and extends recognition to bookkeeping annotations,
with alternatives (R1)--(R5). Category (xi), called a genuine seventh
irreducible typed role in Version~2, is a recorded screen, not proof
of a new role. Remark~\ref{rem:recognition-open} adds the screen example
and conditional closed-formula wrapping constructions, showing that
recognition depends on presentation.

\paragraph{Additions and removals.} Version~3 adds Subsection~\ref{subsec:objects-at-a-glance}
with its vocabulary, symbols and example index; a running example; figures of
the section dependencies and of the three levels; the label-criteria table;
and collects the Lean-counterpart notes, formerly footnotes to the
results, in Appendix~\ref{app:result-notes}. The
nonclaim ``no surrogate terminal evidence'' and the statement that
\(\mathrm{P6}_{\mathrm{drive}}\) is a face of \(\Psix{}\) were removed.
``Regressions'' and ``guardrails'' are called misreadings and checks, and
Section~\ref{sec:integrated-stress} is titled accordingly.

\subsection{Version 2}

Version~2 made the upper bound conditional on an explicit recognition
hypothesis: it added Definition~\ref{def:recognition-hypothesis} and
Remark~\ref{rem:recognition-open}, so that the upper-bound classification of
Version~1 (Theorem~6.7) became Theorem~\ref{thm:upper-bound-classification} and
its scope remark became Remark~\ref{rem:classification-scope}, and the absence
of a seventh role in the terminal theorem was stated under that hypothesis.
Statements of Section~\ref{sec:admissibility-meta-theory} that fix
admissibility conditions rather than assert consequences were stated as
definitions (Definitions~3.1, 3.5, 3.6, 3.9 and~3.11) or as a remark
(Remark~3.8), the stress checks of Section~\ref{sec:integrated-stress} as
Remarks~8.1--8.6, and the integrated stress theorem as
Proposition~\ref{thm:integrated-stress}. The proof of
Proposition~\ref{prop:boundary-distinctions} was completed over all fifteen
rows of the boundary matrix, corrections from external review were
incorporated, and notes on the Lean counterpart of each result were added,
with an appendix describing the Lean development.
 

\bibliographystyle{plainnat}
\bibliography{bib/references}

\end{document}
